\documentclass[11pt,a4paper]{report}
\usepackage{fontspec,amsmath,amssymb,amsthm,microtype,xcolor}
\setmainfont{lmroman10-regular.otf}[BoldFont=lmroman10-bold.otf,ItalicFont=lmroman10-italic.otf,BoldItalicFont=lmroman10-bolditalic.otf]
\setmonofont{lmmono10-regular.otf}[Scale=MatchLowercase]
\usepackage[margin=25mm]{geometry}
\usepackage[unicode,hidelinks]{hyperref}
\setlength{\emergencystretch}{3em}
\setlength{\parskip}{0.4em}
\setlength{\parindent}{0pt}
\setcounter{tocdepth}{1}
\makeatletter
\renewcommand*\l@section{\@dottedtocline{1}{1.5em}{3.8em}}
\renewcommand*\@pnumwidth{2.5em}
\renewcommand*\@tocrmarg{3.5em}
\makeatother
\setcounter{secnumdepth}{2}
\allowdisplaybreaks[2]
\raggedbottom
\title{\texorpdfstring{Kinetic Langevin Monte Carlo}{Kinetic Langevin Monte Carlo}\\[0.6em]\large Mathematical proofs}
\author{LangevinLean}
\date{}
\begin{document}
\pagenumbering{roman}
\maketitle
\chapter*{About this document}
\addcontentsline{toc}{chapter}{About this document}
This volume collects the theorem statements, mathematical arguments and supporting proofs for this development. Every locally authored proof module required by the displayed results, proof roadmap and linked extensions is included once, in dependency order. Foundational results from Mathlib are used as standard mathematical results.



The actual Brownian kinetic mixing theorem, first-order KLMC sampler theorem, Hessian-corrected KLMC2 theorem, Gibbs stationarity, exact Gaussian and Cholesky constructions, and explicit accuracy corollaries are checked. The source audit records the supplied manuscript corrections and the continuous-L\ensuremath{{}^{2}} Brownian model.



\textbf{Published source; corrected checked proofs}



The audit follows the supplied mainDR\_dec\_18.tex manuscript. Actual canonical Brownian integral-solution mixing, full Gibbs stationarity, both correlated-Gaussian KLMC/\allowbreak{}KLMC2 Wasserstein bounds, literal Cholesky sampling, and explicit accuracy guarantees are proved. The corrected proofs checked here are unpublished: the audit records the continuous prefactor correction, the stationary local reference, and the retained stochastic Hessian-difference terms. KLMC2 keeps the exact fourth-moment factor in every dimension and states p\ensuremath{\ge}2 for the printed coarser simplification; its sharper inverse-step-square tail is proved. Continuous solutions are continuous in L\ensuremath{{}^{2}} on finite horizons; no separate almost-surely continuous sample-path modification theorem is claimed.


\begin{itemize}

\item Arnak S. Dalalyan and Lionel Riou-Durand (2020). On sampling from a log-concave density using kinetic Langevin diffusions. Bernoulli 26(3), 1956–1988. arXiv
\end{itemize}
\tableofcontents
\clearpage
\pagenumbering{arabic}
\chapter{Assumptions and main results}


The statements below concern the constructed Brownian integral solution and actual Gaussian chains. The source audit preserves the printed statements and distinguishes the corrected prefactor, stochastic Hessian comparison, kernel index range, and final dimension simplification.

\section{\texorpdfstring{Kinetic position laws: the complete shifted contraction estimate}{Kinetic position laws: the complete shifted contraction estimate}}


Let \(E\) be a complete real inner product space with its Borel
structure and a countable topological base. Suppose \(f:E\to\mathbb R\) is
differentiable and \(m\)-strongly convex with \(M\)-Lipschitz gradient,
\(m\le M\). Fix \(\gamma\ge0\), \(a<\gamma/2\), and \(0\le t\le T\).
Let \(\mu,\nu\) be probability laws of the initial position with
\(W_2(\mu,\nu)<\infty\). Let \(Q\) be any probability law of the initial
velocity and continuous velocity forcing on \([0,T]\), independent of the
initial position in each process.




The selected flow solves the actual integral equation, with normalized
forcing \(b_t-b_0\),


\[Y_t-Y_0=-\int_0^t(\gamma V_s+\nabla f(X_s),-V_s)ds+(b_t-b_0,0).\]


Denote its position law by \(P_t^Q\mu\). Then


\[\boxed{W_2(P_t^Q\mu,P_t^Q\nu)\le A_a e^{k_at}W_2(\mu,\nu),}\]
\[A_a=\frac{\sqrt{2((\gamma-a)^2+a^2)}}{\gamma-2a},\]
\[k_a=\frac{\max\{a^2-m,M-(\gamma-a)^2\}}{\gamma-2a}.\]


This is the explicit prefactor and exponent of the manuscript's equation
(exp2), proved for the actual independent-product laws and the infimum over
initial couplings. No continuous Hessian assumption is needed. Choosing
\(a=0\) and \(\gamma^2\ge m+M\) gives
\(\sqrt2e^{-mt/\gamma}W_2(\mu,\nu)\).




The proof first cancels the common forcing in two integral solutions.
Strong-smooth gradient interpolation controls the transformed energy
\(\|V-V'+(\gamma-a)(X-X')\|^2+\|-(V-V')-a(X-X')\|^2\).
Gronwall gives the path estimate. Each initial coupling is then extended by
an independent copy of \(Q\); both output marginals and its cost are proved
before taking the Wasserstein infimum. The canonical Brownian source and
its Gibbs-stationary solution are also constructed in the supporting modules;
the separate KLMC theorem below combines them with the numerical analysis.
{\small\color{gray}Supporting proof: Section~\ref{proof-module-KineticLangevin-FlowLaw-lean}.\par}
\section{\texorpdfstring{The actual Brownian diffusion mixes to its proved Gibbs target}{The actual Brownian diffusion mixes to its proved Gibbs target}}


The source is an actual constructed Brownian probability space. With standard Gaussian initial velocity independent of the initial position, the selected continuous L\ensuremath{{}^{2}} solution solves the kinetic integral equation. Its Gibbs stationarity is proved by an exact Hamiltonian–thermostat splitting limit.




For any finite-second-moment initial law \(\nu\), any \(a<\gamma/2\), and \(t\ge0\), its actual position law \(\nu_t\) satisfies


\[W_2(\nu_t,\pi)\le \frac{\sqrt{2((\gamma-a)^2+a^2)}}{\gamma-2a}\,e^{k_at}W_2(\nu,\pi),\]
\[k_a=\frac{\max\{a^2-m,M-(\gamma-a)^2\}}{\gamma-2a}.\]


The theorem takes no stationarity or reference-process premise. The proof constructs the noise, shows that every coupling produces the same output marginals, proves the energy inequality in L\ensuremath{{}^{2}}, and takes the infimum over all initial couplings. At \(a=0\) and \(\gamma^2\ge m+M\), the prefactor becomes \(\sqrt2\) and the exponent is \(-mt/\gamma\).


{\small\color{gray}Supporting proof: Section~\ref{proof-module-KineticLangevin-ContinuousGibbs-lean}.\par}
\section{\texorpdfstring{The actual KLMC sampler: the complete first-\allowbreak{}order bound}{The actual KLMC sampler: the complete first-order bound}}


Let \(f:\mathbb R^p\to\mathbb R\) be differentiable,
\(m\)-strongly convex and have \(M\)-Lipschitz gradient, where
\(0<m\le M\). Write \(\pi(dx)\propto e^{-f(x)}dx\).
Choose \(\gamma>0\) and \(h>0\) with


\[\gamma^2\ge m+M,\qquad 4\gamma Mh\le m.\]


Initialize velocity as a standard Gaussian independent of the initial
position law \(\nu_0\), with \(W_2(\nu_0,\pi)<\infty\).
At every step use the constructed centered Gaussian velocity–position
innovation pair with covariance given by the shared exponential Gram kernels;
pairs are independent across time and spatial coordinates. Let \(\nu_n\)
be the position law of this actual KLMC recursion. Then


\[\boxed{W_2(\nu_n,\pi)\le
\sqrt2\left(1-\frac{3mh}{4\gamma}\right)^nW_2(\nu_0,\pi)
+\frac{Mh\sqrt{2p}}m.}\]


The first term contracts the initial error; the second is the persistent
error caused by freezing the force over a step. The proof constructs the
exact and numerical processes together on every initial coupling extended
by an independent Gaussian velocity and canonical Brownian source. Their
common forcing cancels in the contraction estimate. The local freezing error
is quadratic in \(h\), while geometric accumulation divides by a quantity
of order \(h\), producing the linear bias above.




The stationary reference is derived from the actual kinetic equation by a
Verlet–OU splitting limit. Exact chain-law identities show that changing the
initial coupling does not change the numerical output law. Taking the true
Wasserstein infimum gives the initial \(W_2\) term; an optimal coupling is
not assumed. All these constructions and law identities are internal to the
Lean proof of this theorem.


{\small\color{gray}Supporting proof: Section~\ref{proof-module-KineticLangevin-CanonicalKLMC-lean}.\par}
\section{\texorpdfstring{The actual KLMC2 sampler: the exact-\allowbreak{}moment bound in every dimension}{The actual KLMC2 sampler: the exact-moment bound in every dimension}}


Assume the same strongly convex, smooth potential, and let its actual Hessian be \(M_2\)-Lipschitz in operator norm, with \(M_2>0\). Let the initial position have law \(\nu_0\) with \(W_2(\nu_0,\pi)<\infty\). KLMC2 uses the four correlated Brownian response kernels \(\psi_0,\psi_1,\varphi_2,\varphi_3\), with the Hessian applied to both the deterministic velocity and the corresponding noise corrections. Initialize velocity with the standard Gaussian law \(\mathcal N(0,I_p)\), independently of the initial position and the fresh innovations. Suppose


\[\gamma>0,\qquad\gamma^2\ge m+M,\qquad0<h,\qquad5\gamma Mh\le m,\qquad4\sqrt{5p}M_2h\le m.\]


Define \(q=1-mh/(4\gamma)\), \(\alpha=(m/(4M_2h))^2\), and


\[C_p=M_2\sqrt{p^2+2p}+M\sqrt{Mp}.\]


The actual Gaussian-chain position law satisfies, in every finite dimension,


\[\boxed{\begin{aligned}W_2(\nu_n,\pi)\le{}&\sqrt2q^nW_2(\nu_0,\pi)+\frac{\sqrt2h^2C_p}{m}\\&+\frac{8(M-m)}m h e^{-(\alpha-p)/8}.\end{aligned}}\]


The tail is bounded by both \(e^{-p/2}\) and \(e^{-m^2/(160M_2^2h^2)}\). For \(p\ge2\), simplifying the exact fourth moment yields the supplied manuscript’s displayed bound:


\[\begin{aligned}W_2(\nu_n,\pi)\le{}&\sqrt2q^nW_2(\nu_0,\pi)+\frac{2M_2ph^2}{m}\\&+\frac{Mh^2\sqrt{2Mp}}m+\frac{8M}m h e^{-p/2}.\end{aligned}\]


The dimension restriction belongs to this final simplification: \(p^2+2p\le2p^2\) requires \(p\ge2\). The preceding exact-moment bound provides the valid dimension-one conclusion. The source audit preserves the supplied statement and identifies this correction explicitly.




The proof combines a cubic local Taylor–memory error with a contractive comparison of two actual Hessian-corrected updates. The supplied subtraction omits stochastic Hessian-difference terms. The formal proof retains them and combines the initial Gaussian velocity with the fresh Brownian interval into one stationary free Ornstein–Uhlenbeck process. Its Gaussian cutoff estimate does not assume independence from the random Hessian. The choice of \(\alpha\) absorbs the Hessian contribution into the contraction gap. Exact canonical chain and reference laws then permit the true Wasserstein infimum. No local-error, stationarity, independence, or convergence conclusion is left as a theorem premise.




An explicit accuracy corollary bounds the entire persistent floor by \(B_ph^2\), chooses a positive minimum step satisfying both source restrictions and \(B_ph^2\le0.94\varepsilon\), and rounds the logarithmic transient count to a natural number. It gives a checked error at most \(\varepsilon\) for the actual sampler.


{\small\color{gray}Supporting proof: Section~\ref{proof-module-KineticLangevin-CanonicalKLMC2-lean}.\par}
\section{\texorpdfstring{The supplied simplified KLMC2 bound, with its valid dimension range}{The supplied simplified KLMC2 bound, with its valid dimension range}}


For \(p\ge2\), the exact-moment result above implies the supplied formula with terms \(2M_2ph^2/m\), \(Mh^2\sqrt{2Mp}/m\), and \((8M/m)he^{-p/2}\). The Lean statement makes the dimension restriction explicit. The stronger inverse-step-square tail is available in \texttt{klmc2\_\allowbreak{}wasserstein\_\allowbreak{}bound\_\allowbreak{}sharp\_\allowbreak{}tail} in the same module.


{\small\color{gray}Supporting proof: Section~\ref{proof-module-KineticLangevin-CanonicalKLMC2-lean}.\par}
\section{\texorpdfstring{The refined complexity correction: an explicit bound for the natural count}{The refined complexity correction: an explicit bound for the natural count}}


The source refinement explicitly omits its exponential tail. For \(m>0\), \(M\ge m\), \(M_2>0\), \(p\ge1\), and \(\varepsilon>0\), set \(\gamma=\sqrt{m+M}\). Let \(h_0\) be the minimum of the two source stability limits and \(\sqrt{\varepsilon\sqrt m/(4M_2p+2M\sqrt{2Mp})}\). Put \(\kappa=M/m\), \(\rho=(M_2^{2/3}+Mp^{-1/3})/m\), and \(R=\max\{\kappa,\rho\}\). Take a positive finite initial distance \(D\), and select \(n=\lceil\max\{0,(4\gamma/(mh_0))\log(24D\sqrt m/\varepsilon)\}\rceil\). This natural count obeys

\[n<1+64\left[\sqrt p R^2+\sqrt{p/\varepsilon}R^{5/4}\right]\max\left\{0,\log\frac{24D\sqrt m}{\varepsilon}\right\}.\]

The exact reciprocal-step identity proves the connection to the physical source parameters. The printed form is recovered under the explicit regime \(M(1-p^{-1/3})\le M_2^{2/3}\). A separate Lean theorem proves that no uniform constant repairs the printed expression without that regime or a correction. This is a statement about the source step prescription; the actual sampler accuracy theorem above retains its full tail.


{\small\color{gray}Supporting proof: Section~\ref{proof-module-KineticLangevin-SecondOrderComplexity-lean}.\par}
\chapter{Mathematical argument}
\section{\texorpdfstring{From a density to a position–velocity system}{From a density to a position–velocity system}}


A sampling problem starts with a potential \(f\): lower values should be
visited more often, in proportion to \(e^{-f}\). An overdamped trajectory moves
according to its current force and immediately receives Brownian noise. A
kinetic trajectory carries a velocity. Force changes the velocity, and velocity
changes position:


\[dV_t=-(\gamma V_t+\nabla f(L_t))dt+\sqrt{2\gamma}\,dW_t,
\qquad dL_t=V_tdt.\]


The friction \(\gamma V_t\) removes kinetic energy. The noise adds it back
at exactly the strength needed for a standard Gaussian equilibrium velocity.
The candidate equilibrium has density


\[\Pi(dx,dv)=\pi(dx)\,\varphi(dv)
\propto e^{-f(x)-\|v\|^2/2}dx\,dv.\]


Position and velocity are independent under this law. Integrating out velocity
therefore leaves precisely the desired target. The extra variable changes how
the sampler moves without changing the position distribution at equilibrium.




Strong convexity supplies a quadratic lower bound on \(f\), so the partition
function is finite and all polynomial position moments exist. A proof of
stationarity still needs an actual solution process and a valid invariance
argument; checking that a formal generator integrates to zero is one component
of that argument.



\textbf{Scope.} The actual canonical Brownian kinetic flow and its full Gaussian-times-Gibbs stationarity are proved. The proof uses exact Hamiltonian invariance and a controlled splitting limit.

\section{\texorpdfstring{Inverse mass changes the clock}{Inverse mass changes the clock}}


The general paper equation includes inverse mass \(u>0\):
\(dV=-(\gamma V+u\nabla f(L))dt+\sqrt{2\gamma u}\,dW\).
Its equilibrium velocity has covariance \(uI\). Normalize both the velocity
and the time variable:


\[\bar V_t=u^{-1/2}V_{t/\sqrt u},\qquad
\bar L_t=L_{t/\sqrt u},\qquad \bar\gamma=\gamma/\sqrt u.\]


The position equation becomes \(d\bar L_t=\bar V_tdt\). Brownian scaling
turns \(u^{1/4}W_{t/\sqrt u}\) into a standard Brownian motion; substitution
in the velocity equation gives the same kinetic system with \(u=1\) and
friction \(\bar\gamma\). This is Lemma 1 in the supplied manuscript.




The scaling explains why the rest of the paper fixes \(u=1\): it removes
a redundant clock parameter. When transferring a theorem back to another mass,
the step size, friction, Gaussian velocity law and physical time must all be
transformed together.



\textbf{Scope.} The complete inverse-mass integral equation, transformed noise coefficient, Brownian finite-dimensional laws, and independent changed-clock coordinates are proved.

\section{\texorpdfstring{Drive two trajectories with the same noise}{Drive two trajectories with the same noise}}


To understand mixing, start two copies at different positions and feed them
the same Brownian path. The random forcing cancels in their difference. Put
\(q_t=L_t-L'_t\), \(r_t=V_t-V'_t\). Then


\[q'_t=r_t,\qquad r'_t=-\gamma r_t-H_tq_t,\qquad
H_t=\int_0^1\nabla^2f(L'_t+s(L_t-L'_t))\,ds.\]


The averaged Hessian is symmetric and obeys \(mI\preceq H_t\preceq MI\).
The problem has become a deterministic differential inequality along each pair
of paths. No independent-noise variance is added to this comparison.




The usual Euclidean norm of \((r,q)\) need not decrease monotonically:
energy can move between position and velocity. The paper instead chooses a
coordinate system that follows this transfer. This is the main reason the
mixing proof achieves useful constants.



\textbf{Scope.} Common forcing cancels in the actual integral solutions. The final theorem applies this argument to the constructed Brownian laws and the true Wasserstein coupling infimum.

\section{\texorpdfstring{A transformed energy makes the mixed terms cancel}{A transformed energy makes the mixed terms cancel}}


Choose \(0\le v<\gamma/2\), set \(\lambda_-=v\),
\(\lambda_+=\gamma-v\), and define


\[a_t=r_t+\lambda_+q_t,\qquad b_t=-r_t-\lambda_-q_t,
\qquad \mathcal E_t=\|a_t\|^2+\|b_t\|^2.\]


The change of variables is invertible because \(\lambda_+-\lambda_->0\).
Differentiate \(\mathcal E\), use the equations for \(q,r\), and exploit
symmetry of \(H_t\). The two off-diagonal inner products cancel exactly:


\[\mathcal E'_t=
\frac{2}{\lambda_+-\lambda_-}
\bigl(\langle a_t,(\lambda_-^2I-H_t)a_t\rangle
+\langle b_t,(H_t-\lambda_+^2I)b_t\rangle\bigr).\]


Consequently \(\mathcal E'_t\le2c(v)\mathcal E_t\), where


\[c(v)=\frac{\max\{v^2-m,\ M-(\gamma-v)^2\}}{\gamma-2v}.\]


Gronwall gives \(\sqrt{\mathcal E_t}\le e^{c(v)t}\sqrt{\mathcal E_0}\).
If the initial velocities agree, then
\(\mathcal E_0=((\gamma-v)^2+v^2)\|q_0\|^2\). Since
\(q_t=(a_t+b_t)/(\gamma-2v)\), the position estimate becomes


\[\|q_t\|\le
\frac{\sqrt{2((\gamma-v)^2+v^2)}}{\gamma-2v}
e^{c(v)t}\|q_0\|.\]


This is the explicit form of Proposition 1 and Theorem 1. Squaring,
integrating, and taking the infimum over initial position couplings turns it
into a Wasserstein estimate.



\textbf{Scope.} The transformed energy and general shifted dissipation are proved from the actual strong-convexity and gradient-Lipschitz assumptions.

\section{\texorpdfstring{How to read the friction regimes}{How to read the friction regimes}}


At \(v=0\) the transformed energy is particularly simple:
\(\|r+\gamma q\|^2+\|r\|^2\). If \(\gamma^2\ge m+M\), then
\(c(0)=-m/\gamma\), so


\[W_2(\mu P_t^L,\mu'P_t^L)
\le\sqrt2e^{-mt/\gamma}W_2(\mu,\mu').\]


Here the two initial phase-space laws share the same independent velocity
marginal. The prefactor is \(\sqrt2\); this follows directly by substituting
\(v=0\) in the explicit formula. The supplied source's earlier
\(\sqrt2/\gamma\) prefactor in its equation \texttt{exp1} does not match
that substitution.




Friction that is very large makes the rate \(m/\gamma\) small: velocities
are damped before they transport position very far. The smallest friction
allowed by this simple estimate is \(\sqrt{m+M}\). The full table in the
paper improves the rate in other regimes by optimizing \(v\). For sufficiently
large friction it recovers


\[\beta=\frac{\gamma-\sqrt{\gamma^2-4m}}2,\]


the slow characteristic rate of a quadratic potential with curvature \(m\).
This explains the sharpness example. Singular prefactors at a repeated
characteristic root require separate limiting estimates.



\textbf{Scope.} All three scalar optimality regimes are proved on the admissible nonsingular shift interval. The actual Brownian theorem retains the complete prefactor and exponent.

\section{\texorpdfstring{Integrate the frozen force exactly}{Integrate the frozen force exactly}}


On an interval of length \(h\), replace the changing force
\(\nabla f(L_t)\) by its value at the current position \(x\). The remaining
velocity equation is linear and can be solved exactly. Define


\[\psi_0(h)=e^{-\gamma h},\qquad
\psi_1(h)=\frac{1-e^{-\gamma h}}\gamma,\qquad
\psi_2(h)=\frac h\gamma-\frac{1-e^{-\gamma h}}{\gamma^2}.\]


With the phase ordered as position and velocity, the update is


\[v^+=\psi_0(h)v-\psi_1(h)\nabla f(x)+\sqrt{2\gamma}\,\xi,
\quad x^+=x+\psi_1(h)v-\psi_2(h)\nabla f(x)+\sqrt{2\gamma}\,\xi'.\]


The exponential terms handle friction at its exact continuous-time rate.
The coefficients of the force are its first and second time integrals. For
small \(h\), \(\psi_1(h)\approx h\) and \(\psi_2(h)\approx h^2/2\),
which are the familiar first-order velocity and second-order position changes
under a constant force.




The only approximation here is freezing the nonlinear force. The linear
friction and its associated Gaussian forcing have already been integrated.
This is the KLMC exponential integrator reused by KIPLMC1.



\textbf{Scope.} These are the actual deterministic update definitions reused directly by KIPLMC1.

\section{\texorpdfstring{Position and velocity noise must be correlated}{Position and velocity noise must be correlated}}


The two random increments come from the same Brownian motion:


\[\xi=\int_0^h\psi_0(h-s)dW_s,
\qquad \xi'=\int_0^h\psi_1(h-s)dW_s.\]


For each spatial coordinate their covariance is the two-by-two Gram matrix


\[C=\int_0^h
\begin{pmatrix}\psi_0(t)^2&\psi_0(t)\psi_1(t)\\
\psi_0(t)\psi_1(t)&\psi_1(t)^2\end{pmatrix}dt.\]


The off-diagonal entry is positive: a random impulse to velocity also moves
position during the remainder of the step. Replacing these increments with
two independent Gaussians changes the algorithm. Different spatial coordinates
and different time intervals use independent Gaussian pairs.




For any real \(a,b\),


\[(a,b)C(a,b)^T=\int_0^h(a\psi_0(t)+b\psi_1(t))^2dt\ge0.\]


This proves positive semidefiniteness. When \(h>0\), the two coefficient
functions are linearly independent, which gives strict positivity. One may
therefore sample a pair using a Cholesky factor \(C=HH^T\) and two independent
standard normals. A complete law proof identifies this construction with the
Brownian integrals and establishes independence from the current state.



\textbf{Scope.} The actual Brownian convolutions have the specified joint Gaussian Gram law, with independent spatial and time blocks. No matching-law assumption is supplied to the sampler endpoint. The paper’s explicit positive-diagonal Cholesky factor maps two genuinely independent standard normals to exactly this same innovation law.

\section{\texorpdfstring{Smooth position paths give a smaller local defect}{Smooth position paths give a smaller local defect}}


Compare one frozen-force step with the exact kinetic process started at
the same phase-space state. Their Brownian terms cancel. Gradient Lipschitzness
bounds the velocity defect by


\[\|\Delta v_t\|_{L^2}
\le M\int_0^t\|L_s-L_0\|_{L^2}ds
\le\frac{Mt^2}{2}\sup_{0\le r\le h}\|V_r\|_{L^2}.\]


Position is the integral of velocity, so its defect has an additional time
factor:


\[\|\Delta x_h\|_{L^2}\le\frac{Mh^3}{6}
\sup_{0\le r\le h}\|V_r\|_{L^2}.\]


At stationarity \(\|V_r\|_{L^2}=\sqrt p\). For the auxiliary exact path
started from the numerical state, contraction against the stationary exact
path controls its velocity norm by \(\sqrt p+A_k\), where \(A_k\) is the
current error in transformed coordinates. This prevents the initialization
error from reappearing in every step as a fixed additive penalty.




Under the paper's step restriction, combining the two coordinate defects
gives the transformed local estimate
\(\tfrac34Mh^2(\sqrt p+A_k)\). Constructing the auxiliary process, its
moments and the synchronous Gaussian law is part of the analytic proof behind
this estimate.



\textbf{Scope.} The local integrals are bounded in L\ensuremath{{}^{2}}, and the actual constructed reference supplies the required stationary Gaussian velocity moments.

\section{\texorpdfstring{Contract, add the local defect, and sum}{Contract, add the local defect, and sum}}


Let
\(A_k=\| (\Delta V_k+\gamma\Delta L_k,-\Delta V_k)\|_{L^2}\).
Exact kinetic evolution contracts this transformed norm. A triangle inequality
with the one-step defect gives


\[A_{k+1}\le\left(e^{-mh/\gamma}+\frac34Mh^2\right)A_k
+\frac34Mh^2\sqrt p.\]


For \(0<h\le m/(4\gamma M)\), the quadratic correction is small enough
to retain most of the continuous contraction:


\[A_{k+1}\le qA_k+b,\qquad
q=1-\frac{3mh}{4\gamma},\quad b=\frac34Mh^2\sqrt p.\]


Induction gives the finite geometric sum
\(A_k\le q^kA_0+b\sum_{j=0}^{k-1}q^j\). Because \(0\le q<1\),


\[A_k\le q^kA_0+\frac{Mh\gamma\sqrt p}m.\]


Equal initial velocities give \(A_0=\gamma W_2(\nu_0,\pi)\) for an
optimal initial coupling; approximate couplings and a limiting infimum work as
well. The position projection costs \(\sqrt2/\gamma\), yielding the displayed
KLMC theorem. The recurrence itself is elementary. The substantive stochastic
claim is that the actual chain admits the coupling from which it was derived.



\textbf{Scope.} The actual sampler theorem constructs the reference, derives the numerical local inequality, identifies both output laws, and takes the true initial-coupling infimum internally.

\section{\texorpdfstring{The Hessian-\allowbreak{}corrected method follows the changing force}{The Hessian-corrected method follows the changing force}}


If the Hessian is \(M_2\)-Lipschitz in operator norm, a first-order
Taylor expansion of the gradient leaves a quadratic position remainder:


\[\|\nabla f(y)-\nabla f(x)-\nabla^2f(x)(y-x)\|
\le\frac{M_2}{2}\|y-x\|^2.\]


The paper inserts the leading velocity and Brownian contributions to
\(L_s-L_0\) into this expansion. Define
\(\varphi_{j+1}(h)=\int_0^he^{-\gamma(h-s)}\psi_j(s)ds\). The resulting
KLMC2 update subtracts \(\varphi_2(h)H v\) from velocity and
\(\varphi_3(h)H v\) from position, where \(H=\nabla^2f(x)\). Its random
increments are also corrected by multiplying two additional correlated
Gaussian blocks by \(H\).




The joint covariance is the Gram integral of
\((\psi_0,\psi_1,\varphi_2,\varphi_3)\). This avoids a matrix exponential:
the operations involving the Hessian are matrix–vector products. The second
algorithm is therefore more elaborate than the frozen-force integrator, but
still explicit.




The name KLMC2 here refers to a Hessian method. KIPLMC2 in the particle
paper uses the OBABO splitting and first derivatives; its definition and proof
are separate.



\textbf{Scope.} The four-kernel Hessian update and its exact Brownian Gaussian innovation law are constructed. The separate completion audit tracks the assembled chain endpoint.

\section{\texorpdfstring{Fourth moments and the complete stochastic Hessian correction}{Fourth moments and the complete stochastic Hessian correction}}


The Taylor remainder contains \(\|L_s-L_0\|^2\). Its
\(L^2\) norm therefore uses fourth moments. Under the actual stationary
kinetic law, each \(V_s\) is standard Gaussian and


\[\mathbb E\|V_s\|^2=p,\qquad
\mathbb E\|V_s\|^4=p^2+2p.\]


The actual kinematic equation \(L_s-L_0=\int_0^sV_rdr\) gives
\(\|L_s-L_0\|_{L^4}\le s(p^2+2p)^{1/4}\), without independence
between velocities at different times. The formal proof applies closed-convex
averaging to the existing \(L^2\) time integral. The genuine Taylor forcing
is consequently bounded by \(M_2s^2\sqrt{p^2+2p}/2\).




Its damped integral gives the cubic velocity defect. A further force-memory
term uses the actual stationary score bound
\(\mathbb E_\pi\|\nabla f\|^2\le Mp\). Integrating the velocity error
once more gives the quartic position defect. Proposition 2 is thus checked
for the implemented four-kernel step:


\[\|P^{-1}(V_h-\hat V_h,L_h-\hat L_h)\|_{L^2}
\le\frac{h^3}{4}\bigl(M_2\sqrt{p^2+2p}+M\sqrt{Mp}\bigr),
\qquad \gamma h\le\tfrac15.\]


Proposition 3 requires an additional repair: the source's displayed
subtraction omits the Hessian difference multiplying the fresh
\(\varphi_2\) and \(\varphi_3\) noise. The formal proof keeps both terms.
Write \(D_H=H(L_0)-H(\widetilde L_0)\), and define the actual free-OU velocity
and complete Hessian responses on the same Brownian interval:


\[U_r=\psi_0(r)V_0+\sqrt{2\gamma}\int_0^r\psi_0(r-s)dB_s,\]
\[C_j(t)=\varphi_j(t)V_0+\sqrt{2\gamma}
\int_0^t\varphi_j(t-s)dB_s,\qquad j=2,3.\]


Actual kernel identities, proved by Hilbert-space Fubini and the Brownian
isonormal map, give


\[C_2(t)=\int_0^te^{-\gamma(t-s)}\int_0^sU_rdrds,
\qquad C_2(t)+\gamma C_3(t)=\int_0^t\int_0^sU_rdrds.\]


Each \(U_r\) has the standard Gaussian marginal, proved from the actual
stationary reference velocity and its independent future Brownian interval.
The random operator \(D_H\) can depend on the coupled positions and velocity;
the cutoff argument below requires no independence from \(U_r\). The complete
correction therefore satisfies, for \(a\ge5p\),


\[\|P^{-1}(D_HC_2(h),D_HC_3(h))\|_{L^2}
\le\frac{h^2}{\sqrt2}\left[
M_2\sqrt a\,\|L_0-\widetilde L_0\|_{L^2}
+2(M-m)e^{-(a-p)/8}\right].\]


This retains the source constants while accounting for every stochastic
Hessian term. The frozen-linear part contracts by \(1-mh/(2\gamma)\)
under \(\gamma^2\ge m+M\) and \(5\gamma Mh\le m\); adding the complete
correction gives the repaired actual one-step comparison.



\textbf{Scope.} The actual four-kernel local error and repaired synchronous comparison are proved. Both fresh stochastic Hessian terms are retained through the full free-OU response, whose Gaussian marginals follow from the constructed canonical past/\allowbreak{}future law.

\section{\texorpdfstring{Gaussian truncation and the actual second-\allowbreak{}order sampling theorem}{Gaussian truncation and the actual second-order sampling theorem}}


A Gaussian reference velocity need not be independent of the accumulated
position error. For any actual standard Gaussian vector and \(a\ge5p\),
the checked Gaussian exponential moment gives


\[\|V\|^2\le a+(\|V\|^2-a)_+,\qquad
\mathbb E(\|V\|^2-a)_+\le4e^{-(a-p)/4}.\]


The actual Hessian difference satisfies both
\(\|H(x)-H(y)\|\le M-m\) and
\(\|H(x)-H(y)\|\le M_2\|x-y\|\). On the bounded-velocity part use
the Lipschitz bound; on the tail use the curvature-width bound. Taking the
square root of the integrated estimate gives


\[\|D_HV\|_{L^2}\le M_2\sqrt a\,\|L_0-\widetilde L_0\|_{L^2}
+2(M-m)e^{-(a-p)/8}.\]


Apply this estimate to each actual free-OU velocity in the complete
correction, then combine the repaired comparison with the genuine cubic local
defect. Choose


\[a=\left(\frac{m}{4M_2h}\right)^2,\qquad
h\le\frac{m}{5\gamma M}\wedge\frac{m}{4M_2\sqrt{5p}}.\]


Assume \(m,M_2>0\), \(\gamma^2\ge m+M\), and initial phase law
\(\mathcal N(0,I_p)\otimes\nu_0\), with finite \(W_2(\nu_0,\pi)\).
The second step restriction ensures \(a\ge5p\), and the first enforces the
frozen-linear stability bound. The position-dependent comparison term is
absorbed into \(q=1-mh/(4\gamma)\). Iterating on the constructed common
Brownian source and taking the true infimum over initial couplings proves
for the actual independently initialized, iid four-kernel Gaussian KLMC2 chain


\[W_2(\nu_n^{\rm KLMC2},\pi)\le
\sqrt2q^n W_2(\nu_0,\pi)
+\frac{\sqrt2h^2}{m}\left(M_2\sqrt{p^2+2p}+M\sqrt{Mp}\right)
+\frac{8(M-m)h}{m}e^{-(a-p)/8}.\]


The actual output transport distance is finite; the real-valued expression
does not conceal an infinite distance. No stationary-law, local-error,
independence, Gaussian-innovation or optimal-coupling premise is supplied to
this final sampler theorem.




The cutoff gives both \(e^{-(a-p)/8}\le e^{-p/2}\) and the stronger
\(e^{-(a-p)/8}\le e^{-m^2/(160M_2^2h^2)}\). The printed simplification
of the first bias term uses \(p^2+2p\le2p^2\), which requires
\(p\ge2\). Under that explicit restriction the checked bound implies


\[W_2(\nu_n^{\rm KLMC2},\pi)\le
\sqrt2q^n W_2(\nu_0,\pi)
+\frac{2h^2M_2p}{m}+\frac{h^2M\sqrt{2Mp}}m
+\frac{8Mh}{m}e^{-p/2}.\]


For \(p=1\), the exact fourth-moment expression in the preceding
sampling theorem remains the valid conclusion. The source's silently used
higher-dimensional simplification is not applied.



\textbf{Scope.} The exact-dimensional theorem is proved for the actual KLMC2 Gaussian Markov chain, with all reference, noise, local-error and coupling-infimum obligations discharged internally. The printed coefficient is provided with its necessary dimension restriction p at least two; the one-dimensional exact-moment bound is retained.

\section{\texorpdfstring{Turn an error floor into a computational budget}{Turn an error floor into a computational budget}}


A bound of the form \(q^kD+Bh\) gives a simple design rule. Choose
\(h\) small enough that \(Bh\) receives one part of an error budget
\(\varepsilon\), and choose \(k\) large enough for the decaying initial
term to receive the rest. The admissible \(h\) is the minimum of this accuracy
choice and the stability restriction.




For the first integrator, the bias scales as \(h\sqrt p\), giving
\(h\) of order \(\varepsilon/\sqrt p\), and a step count of order
\(\sqrt p/\varepsilon\) up to logarithmic and curvature factors.
For the Hessian method the leading bias scales as \(h^2p\), giving
\(h\) of order \(\sqrt{\varepsilon/p}\), and a step count of order
\(\sqrt{p/\varepsilon}\). These are asymptotic consequences of actual
algorithmic estimates, with fixed geometry constants.




The paper's comparisons keep track of condition numbers as well. A kinetic
method gains from smooth position paths, but poor conditioning can make its
allowed step size small. This is why neither KLMC nor overdamped LMC uniformly
wins across every dimension, accuracy and curvature regime. A precise
complexity theorem also specifies integer rounding and the treatment of
already-small initialization error.




The source also gives a finer condition-number comparison after explicitly
omitting the Gaussian tail. We can audit that calculation without an asymptotic
shortcut. For target error \(\varepsilon/\sqrt m\), allocate
\(\varepsilon/(2\sqrt m)\) to the printed quadratic bias and set


\[h_0=\min\left\{\frac m{5\gamma M},\;
\frac m{4\sqrt{5p}M_2},\;
\sqrt{\frac{\varepsilon\sqrt m}{4M_2p+2M\sqrt{2Mp}}}\right\}.\]


Every denominator is positive. This step meets both source stability
restrictions and the stated quadratic-bias allocation. Its reciprocal is
exactly


\[Q=\max\left\{\frac{5\gamma M}m,\;
\frac{4\sqrt{5p}M_2}m,\;
\sqrt{\frac{4M_2p+2M\sqrt{2Mp}}{\varepsilon\sqrt m}}\right\}.\]


There is no hidden dimension or curvature factor in this identity. For
initial distance \(D\), the rounded logarithmic iteration prescription satisfies


\[n\le1+\frac{4\gamma}{m}Q\,
\max\left\{0,\log\frac{24D\sqrt m}{\varepsilon}\right\}.\]


This calculation concerns the source comparison with the tail omitted.
The actual sampler accuracy theorem retains the tail and chooses its step
using the complete quadratic envelope proved earlier.




To understand the refinement, write
\(\kappa=M/m\), \(\eta=M_2/m^{3/2}\), and


\[\rho=\kappa_2=\eta^{2/3}+\kappa p^{-1/3}
=\frac{M_2^{2/3}+Mp^{-1/3}}m.\]


With \(\gamma=\sqrt{m+M}\), the source tries to simplify the count
coefficient to \(\sqrt p\,\rho^2+\sqrt{p/\varepsilon}\,\rho^{5/4}\).
That simplification drops a factor. A valid refinement is


\[\sqrt{\frac\kappa\rho}\left[
\sqrt p\,\rho^2+\sqrt{p/\varepsilon}\,\rho^{5/4}\right]\]


up to a universal multiplicative constant and the displayed logarithm.
A convenient checked uniform upper bound uses
\(\widehat\kappa_2=\max\{\kappa,\rho\}\):


\[\mathcal C\le64\left[\sqrt p\,\widehat\kappa_2^2+
\sqrt{p/\varepsilon}\,\widehat\kappa_2^{5/4}\right],\]


where \(\mathcal C=(4\gamma/m)Q\) is the exact count coefficient
above. The natural count is therefore less than one plus this bound times
the visible nonnegative logarithm. If \(\rho\ge\kappa\), the maximum
is simply \(\rho\), so the printed shape follows with the same explicit
constant.




The distinction matters because the modified number discounts gradient
curvature by \(p^{-1/3}\), while the actual stability restriction still
contains \(M/m\). Take a natural \(t\ge2\), dimension \(p=t^6\),
and the valid constants \(m=1\), \(M=t^2\), \(M_2=1\),
\(\varepsilon=1\). A diagonal quadratic potential can have smallest
curvature one and largest curvature \(t^2\); its constant Hessian satisfies
the stated Hessian-Lipschitz bound. Here \(\kappa=t^2\),
\(\eta=1\), and \(\rho=2\) exactly. The normalized iteration
coefficient is at least \(20t^4\), whereas the source expression is at most
\(8t^3\). No universal constant multiplying that expression can cover all
natural \(t\). This diagnoses the source's algebraic budget simplification;
it asserts no lower bound on the best possible behavior of the algorithm.



\textbf{Scope.} Both actual samplers now have explicit epsilon guarantees, positive minimum steps, and rounded natural logarithmic counts. The second-order prescription uses the exact dimension term and the proved inverse-step-square tail envelope. The source paragraph that omits the Gaussian tail is also audited: its exact step and natural count are proved, the corrected condition-number envelope has constant 64, and a formal natural-dimension counterexample rules out the unqualified printed budget.

\chapter{Supporting proofs}
The following mathematical arguments are reproduced in full from the authored exposition of the local proof dependencies. Each module appears after its local imports.

\section{\texorpdfstring{The quadratic upper model}{The quadratic upper model}}
\label{proof-module-UnadjustedLangevin-Analysis-Smooth-lean}
{\small\color{gray}Source module: \texttt{Unadjusted\allowbreak{}Langevin/\allowbreak{}Analysis/\allowbreak{}Smooth.\allowbreak{}lean}.\par}


We next obtain the complementary upper bound from a Lipschitz gradient. This argument uses only first derivatives: restrict the potential to a segment, subtract its affine approximation and the proposed quadratic error, and show that the resulting scalar function decreases.



\texttt{smooth\_\allowbreak{}upper\_\allowbreak{}bound} concerns the gradient of the same differentiable potential occurring in its function values. The argument uses only first derivatives and a Lipschitz gradient; the scalar line restriction and its interval integral provide the quadratic remainder, with no Hessian assumption.

\subsection{\texorpdfstring{The descent lemma}{The descent lemma}}


Let \(H\) be a real Hilbert space and \(L\ge0\). Suppose \(f:H\to\mathbb R\) is differentiable and

\[\|\nabla f(u)-\nabla f(v)\|\le L\|u-v\|\quad(u,v\in H).\]

Then for every \(x,y\),

\[f(y)\le f(x)+\langle\nabla f(x),y-x\rangle+\frac L2\|y-x\|^2.\]

Set \(d=y-x\), \(a=\langle\nabla f(x),d\rangle\), \(b=\|d\|^2\), and introduce the corrected segment function

\[q(t)=f(x+td)-f(x)-ta-\frac L2t^2b.\]

It suffices to show \(q(1)\le q(0)=0\).



The segment has derivative \(d\). The chain rule therefore gives, for every real \(t\),

\[\frac{d}{dt}f(x+td)=\langle\nabla f(x+td),d\rangle.\]

Subtract the derivatives of the affine and quadratic terms:

\[q'(t)=\langle\nabla f(x+td),d\rangle-a-Ltb.\]

For \(t\) in the interior of \([0,1]\), we have \(t\ge0\). Lipschitz continuity gives

\[\|\nabla f(x+td)-\nabla f(x)\|\le L\|td\|=Lt\|d\|.\]

By Cauchy–Schwarz,

\[\langle\nabla f(x+td)-\nabla f(x),d\rangle\le\|\nabla f(x+td)-\nabla f(x)\|\,\|d\|\le Lt\|d\|^2.\]

Substituting \(a\) and \(b\) shows that \(q'(t)\le0\).



The derivative formula also makes \(q\) continuous on \([0,1]\) and differentiable in its interior. The derivative criterion for monotonicity now implies that \(q\) is nonincreasing on this interval.



Compare the two endpoints:

\[q(1)\le q(0),\qquad f(x+d)-f(x)-a-\frac L2b\le0.\]

Finally \(x+d=x+(y-x)=y\). Substitute the definitions of \(a,b,d\) and rearrange to obtain the quadratic upper bound.

\section{\texorpdfstring{Supporting planes and strong monotonicity}{Supporting planes and strong monotonicity}}
\label{proof-module-UnadjustedLangevin-StrongConvexity-lean}
{\small\color{gray}Source module: \texttt{Unadjusted\allowbreak{}Langevin/\allowbreak{}Strong\allowbreak{}Convexity.\allowbreak{}lean}.\par}


We begin with the deterministic geometry behind synchronous Langevin coupling. Throughout this part, \(\langle\cdot,\cdot\rangle\) and \(\|\cdot\|\) are the real inner product and its norm. Strong convexity uses the normalization \(m/2\): the function \(v\mapsto f(v)-\frac m2\|v\|^2\) is convex.



The formal normalization is \texttt{Strong\allowbreak{}Convex\allowbreak{}On Set.\allowbreak{}univ m f}, meaning convexity after subtracting \(m\|x\|^2/2\). \texttt{fderiv \ensuremath{\mathbb R} f x} is the continuous linear derivative; \texttt{gradient f x} is its Hilbert-space representative, so \(Df(x)[v]=\langle\nabla f(x),v\rangle\). The proofs explicitly move between these two representations.

\subsection{\texorpdfstring{The supporting-\allowbreak{}plane inequality}{The supporting-plane inequality}}


Let \(H\) be a real inner-product space and let \(f:H\to\mathbb R\) be differentiable and convex on all of \(H\). For any \(x,y\in H\),

\[f(x)+Df(x)[y-x]\le f(y).\]

Restrict \(f\) to the affine line \(\ell(t)=x+t(y-x)\). Since affine composition preserves convexity, \(\varphi=f\circ\ell\) is convex on \(\mathbb R\), with \(\varphi(0)=f(x)\) and \(\varphi(1)=f(y)\).



The derivative of the line is its constant direction:

\[\ell'(0)=y-x.\]

By the chain rule, the derivative of the restricted function is

\[\varphi'(0)=Df(x)[y-x].\]

For a differentiable convex function, the derivative at the left endpoint is bounded by the secant slope to a later point. Applying this at \(0<1\) gives

\[Df(x)[y-x]\le\frac{\varphi(1)-\varphi(0)}{1-0}.\]

Evaluate the two endpoints and rearrange:

\[Df(x)[y-x]\le f(y)-f(x),\qquad f(x)+Df(x)[y-x]\le f(y).\]\subsection{\texorpdfstring{Strong monotonicity of the gradient}{Strong monotonicity of the gradient}}


Let \(H\) be a real Hilbert space. If \(f:H\to\mathbb R\) is differentiable and \(m\)-strongly convex, where \(m\in\mathbb R\), then for every \(x,y\in H\),

\[m\|x-y\|^2\le\langle\nabla f(x)-\nabla f(y),x-y\rangle.\]

Remove the strong-convexity quadratic. Set

\[g(v)=f(v)-\frac m2\|v\|^2.\]

The definition of strong convexity says precisely that \(g\) is convex.



The squared norm has derivative \(D\|v\|^2[u]=2\langle v,u\rangle\). Consequently \(g\) is differentiable and

\[Dg(v)[u]=Df(v)[u]-m\langle v,u\rangle.\]

Apply the supporting-plane inequality in both directions, and substitute the derivative just computed:

\[g(x)+Dg(x)[y-x]\le g(y),\qquad g(y)+Dg(y)[x-y]\le g(x).\]

Express the derivative of \(f\) using its gradient. The two inequalities become

\[\begin{aligned}f(x)-\tfrac m2\|x\|^2+\langle\nabla f(x),y-x\rangle-m\langle x,y-x\rangle&\le f(y)-\tfrac m2\|y\|^2,\\ f(y)-\tfrac m2\|y\|^2+\langle\nabla f(y),x-y\rangle-m\langle y,x-y\rangle&\le f(x)-\tfrac m2\|x\|^2.\end{aligned}\]

Add these inequalities. The function values and quadratic values cancel. Using \(\langle x-y,x-y\rangle=\|x-y\|^2\), what remains is

\[0\le\langle\nabla f(x)-\nabla f(y),x-y\rangle-m\|x-y\|^2.\]

This is the required monotonicity inequality.

\subsection{\texorpdfstring{The quadratic lower model}{The quadratic lower model}}


Let \(H\) be a real Hilbert space, \(f:H\to\mathbb R\) differentiable, and \(m\in\mathbb R\). If \(f\) is \(m\)-strongly convex on all of \(H\), then for every \(x,y\in H\),

\[f(x)+\langle\nabla f(x),y-x\rangle+\frac m2\|y-x\|^2\le f(y).\]

Set \(g(v)=f(v)-\frac m2\|v\|^2\). Strong convexity makes \(g\) convex.



Differentiate this shifted potential:

\[Dg(v)[u]=Df(v)[u]-m\langle v,u\rangle.\]

In particular, \(g\) is differentiable on \(H\).



Its supporting-plane inequality at \(x\), together with \(Df(x)[u]=\langle\nabla f(x),u\rangle\), yields

\[f(x)-\frac m2\|x\|^2+\langle\nabla f(x),y-x\rangle-m\langle x,y-x\rangle\le f(y)-\frac m2\|y\|^2.\]

Collect the quadratic correction by the identity

\[\|y\|^2-\|x\|^2-2\langle x,y-x\rangle=\|y-x\|^2.\]

Moving terms gives the desired lower model.

\section{\texorpdfstring{A first Euler-\allowbreak{}step estimate}{A first Euler-step estimate}}
\label{proof-module-UnadjustedLangevin-GradientStep-lean}
{\small\color{gray}Source module: \texttt{Unadjusted\allowbreak{}Langevin/\allowbreak{}Gradient\allowbreak{}Step.\allowbreak{}lean}.\par}


Before using the full lower and upper models together, it is useful to record the estimate obtained directly from strong monotonicity and a Lipschitz bound. This estimate applies to arbitrary vectors \(g_x,g_y\); they need not arise as gradients.



The formal inputs \texttt{gx} and \texttt{gy} are arbitrary vectors satisfying the displayed monotonicity and Lipschitz inequalities. Thus \texttt{gradient\_\allowbreak{}step\_\allowbreak{}sq\_\allowbreak{}bound} is an algebraic Hilbert-space result; the potential's actual gradients are substituted later. \texttt{gradient\_\allowbreak{}step\_\allowbreak{}contract} takes a nonnegative square root and does not require a separate sign assumption on the scalar coefficient.

\subsection{\texorpdfstring{The squared Euler-\allowbreak{}step bound}{The squared Euler-step bound}}


Let \(x,y,g_x,g_y\) lie in a real inner-product space, let \(m\in\mathbb R\), \(L\ge0\), and \(h\ge0\). Assume

\[m\|x-y\|^2\le\langle g_x-g_y,x-y\rangle,\qquad\|g_x-g_y\|\le L\|x-y\|.\]

Then

\[\|(x-hg_x)-(y-hg_y)\|^2\le(1-2hm+h^2L^2)\|x-y\|^2.\]

Both sides of the Lipschitz estimate are nonnegative, so squaring preserves the inequality:

\[\|g_x-g_y\|^2\le L^2\|x-y\|^2.\]

Multiply strong monotonicity by \(2h\ge0\), and multiply the squared Lipschitz estimate by \(h^2\ge0\):

\[2hm\|x-y\|^2\le2h\langle g_x-g_y,x-y\rangle,\qquad h^2\|g_x-g_y\|^2\le h^2L^2\|x-y\|^2.\]

Write \(d=x-y\), \(v=g_x-g_y\). Expanding the squared norm and applying the preceding estimates gives

\[\begin{aligned}\|(x-hg_x)-(y-hg_y)\|^2&=\|d-hv\|^2\\&=\|d\|^2-2h\langle v,d\rangle+h^2\|v\|^2\\&\le(1-2hm+h^2L^2)\|d\|^2.\end{aligned}\]\subsection{\texorpdfstring{Taking the square root}{Taking the square root}}


Let \(H\) be a real inner-product space, \(x,y,g_x,g_y\in H\), \(m\in\mathbb R\), \(L\ge0\), and \(h\ge0\). Assume

\[m\|x-y\|^2\le\langle g_x-g_y,x-y\rangle,\qquad\|g_x-g_y\|\le L\|x-y\|.\]

Then

\[\|(x-hg_x)-(y-hg_y)\|\le\sqrt{1-2hm+h^2L^2}\,\|x-y\|.\]

No separate sign assumption on \(1-2hm+h^2L^2\) is required; Lean uses the real square root, which is zero on nonpositive inputs.



Apply the preceding squared estimate. With \(A=\|(x-hg_x)-(y-hg_y)\|\), \(r=\|x-y\|\), and \(c=1-2hm+h^2L^2\), it says \(A^2\le cr^2\).



Since \(A\ge0\), take square roots and use \(r^2\ge0\):

\[A\le\sqrt{cr^2}=\sqrt c\,\sqrt{r^2}=\sqrt c\,r.\]

The square-root product identity is valid here even without a sign hypothesis on \(c\), because the second factor \(r^2\) is nonnegative.

\section{\texorpdfstring{The Langevin update and cancellation of shared noise}{The Langevin update and cancellation of shared noise}}
\label{proof-module-UnadjustedLangevin-ULAStep-lean}
{\small\color{gray}Source module: \texttt{Unadjusted\allowbreak{}Langevin/\allowbreak{}ULAStep.\allowbreak{}lean}.\par}


We now work in finite-dimensional Euclidean space. A ULA step adds the same scaled noise vector to a gradient step. When two starting points use this same noise, it cancels exactly from their difference. The deterministic estimates can therefore be applied pointwise, before introducing any probability law.



\texttt{ula\allowbreak{}Update} is the pointwise Euler map. Its definition accepts arbitrary vector noise; Gaussian laws enter later in realization theorems. \texttt{ula\allowbreak{}Update\_\allowbreak{}sub} uses the same innovation on both sides, so that contribution cancels exactly. The monotonicity of the actual gradient is supplied by the strong-convexity theorem when applying the deterministic step bound.

\subsection{\texorpdfstring{The Euclidean state space}{The Euclidean state space}}


For a natural number \(d\), write \(E_d=\mathbb R^d\) with its Euclidean inner product and norm. Lean represents its vectors as real-valued functions on \(\operatorname{Fin}(d)\), the set of indices \(0,\ldots,d-1\); when \(d=0\), this index set is empty.



Introduce the abbreviation

\[E_d:=\operatorname{EuclideanSpace}_{\mathbb R}(\operatorname{Fin}(d)).\]

This equips the coordinate functions with the Euclidean norm \(\|x\|=(\sum_{i=0}^{d-1}x_i^2)^{1/2}\), so the subsequent gradient and distance estimates use the usual geometry of \(\mathbb R^d\).

\subsection{\texorpdfstring{One unadjusted Langevin step}{One unadjusted Langevin step}}


For a potential \(U:E_d\to\mathbb R\), a real step size \(h\), a state \(x\in E_d\), and a noise vector \(z\in E_d\), define

\[T_h^z(x)=x-h\nabla U(x)+\sqrt{2h}\,z.\]

The definition itself imposes no sign condition on \(h\); later contraction results assume \(h\ge0\).



The update consists of the current state, the drift increment \(-h\nabla U(x)\), and the noise increment \(\sqrt{2h}\,z\).

\subsection{\texorpdfstring{Shared noise cancels}{Shared noise cancels}}


For every potential \(U:E_d\to\mathbb R\), \(h\in\mathbb R\), and vectors \(x,y,z\in E_d\),

\[T_h^z(x)-T_h^z(y)=(x-h\nabla U(x))-(y-h\nabla U(y)).\]

Expand the two updates. The two copies of \(\sqrt{2h}\,z\) subtract to zero:

\[(x-h\nabla U(x)+\sqrt{2h}\,z)-(y-h\nabla U(y)+\sqrt{2h}\,z)=(x-h\nabla U(x))-(y-h\nabla U(y)).\]\subsection{\texorpdfstring{The initial shared-\allowbreak{}noise estimate}{The initial shared-noise estimate}}


Suppose \(U:E_d\to\mathbb R\) is differentiable and \(m\)-strongly convex, \(\nabla U\) is \(L\)-Lipschitz with \(L\ge0\), and \(h\ge0\). For every \(x,y,z\in E_d\),

\[\|T_h^z(x)-T_h^z(y)\|\le\sqrt{1-2hm+h^2L^2}\,\|x-y\|.\]

For \(h>0\) this factor is strictly below one whenever \(hL^2<2m\).



Cancel the common noise. The left-hand side becomes the distance between the two deterministic gradient steps.



Strong convexity gives

\[m\|x-y\|^2\le\langle\nabla U(x)-\nabla U(y),x-y\rangle,\]

and the Lipschitz assumption gives \(\|\nabla U(x)-\nabla U(y)\|\le L\|x-y\|\). Insert these two inequalities into the Euler-step square-root bound.

\section{\texorpdfstring{The sharp deterministic contraction}{The sharp deterministic contraction}}
\label{proof-module-UnadjustedLangevin-SharpContraction-lean}
{\small\color{gray}Source module: \texttt{Unadjusted\allowbreak{}Langevin/\allowbreak{}Sharp\allowbreak{}Contraction.\allowbreak{}lean}.\par}


The preceding square-root factor uses strong monotonicity and Lipschitz continuity separately. The quadratic lower and upper models contain more information. We combine them into the interpolation inequality, then interpret that inequality as a ball containing the gradient difference. This yields the two sharp factors \(1-mh\) and \(Lh-1\), with the change of regime at \(h=2/(m+L)\).



The proof separates the two regimes for the maximum contraction factor \(\max(|1-mh|,|1-Lh|)\). Each scalar bound follows after the strong monotonicity and cocoercivity estimates have supplied the needed inner-product control. The final norm theorem removes the square by nonnegativity, retaining the precise step-regime hypotheses.

\subsection{\texorpdfstring{Cocoercivity from two quadratic models}{Cocoercivity from two quadratic models}}


Let \(H\) be a real inner-product space, \(f:H\to\mathbb R\), \(g:H\to H\), and \(K\ge0\). Suppose that for all \(a,b\in H\),

\[f(a)+\langle g(a),b-a\rangle\le f(b)\le f(a)+\langle g(a),b-a\rangle+\frac K2\|b-a\|^2.\]

Then for all \(x,y\),

\[\|g(x)-g(y)\|^2\le K\langle g(x)-g(y),x-y\rangle.\]

The proof includes \(K=0\), and does not separately assume that \(g\) is a gradient.



Fix \(a,b\), \(t\ge0\), and set \(v=g(a)-g(b)\), \(z=a-tv\). Compare the lower model based at \(b\) with the upper model based at \(a\), both evaluated at \(z\):

\[f(b)+\langle g(b),z-b\rangle\le f(z)\le f(a)+\langle g(a),z-a\rangle+\frac K2\|z-a\|^2.\]

The displacements are \(z-a=-tv\) and \(z-b=(a-b)-tv\). Substitute them into the two models. Since \(t\ge0\), the quadratic term is \(\|z-a\|^2=t^2\|v\|^2\).



The difference of the linear terms satisfies

\[\langle g(a),v\rangle-\langle g(b),v\rangle=\|v\|^2.\]

Eliminating \(f(z)\) therefore gives

\[\left(t-\frac K2t^2\right)\|g(a)-g(b)\|^2\le f(a)-f(b)-\langle g(b),a-b\rangle.\]

Apply this estimate to \((a,b)=(x,y)\) and \((y,x)\), with the same \(t\ge0\). Adding will cancel the function values.



The two gradient norms are equal and the reversed displacement changes the sign of its inner product. Thus, writing \(A=\|g(x)-g(y)\|^2\) and \(B=\langle g(x)-g(y),x-y\rangle\), we have for every \(t\ge0\),

\[(2t-Kt^2)A\le B.\]

First consider \(K=0\). Taking \(t=1\) gives \(2A\le B\). As \(A\ge0\), this also shows \(B\ge0\).



If the desired conclusion \(A\le0\) were false, then \(A>0\). We could choose

\[t=\frac{B+1}{2A}\ge0.\]

The estimate \(2tA\le B\) would then force \(B+1\le B\).



To justify that substitution, \(A>0\) implies \(\|g(x)-g(y)\|\ne0\), so its square can be cancelled from the denominator. Hence indeed

\[2\frac{B+1}{2A}A=B+1.\]

The contradiction \(B+1\le B\) excludes \(A>0\). Thus \(A=0\), proving the claim when \(K=0\).



Now suppose \(K>0\). Choose \(t=K^{-1}\) in the common estimate:

\[\frac AK\le B.\]

Multiplying by \(K\) gives \(A\le KB\), as required.

\subsection{\texorpdfstring{The strong-\allowbreak{}smooth interpolation inequality}{The strong-smooth interpolation inequality}}


Let \(H\) be a real Hilbert space, \(f:H\to\mathbb R\) differentiable and \(m\)-strongly convex, and let \(\nabla f\) be \(L\)-Lipschitz, where \(L\ge0\) and \(m\le L\). Then for every \(x,y\),

\[mL\|x-y\|^2+\|\nabla f(x)-\nabla f(y)\|^2\le(m+L)\langle\nabla f(x)-\nabla f(y),x-y\rangle.\]

The inequality holds also at the boundary \(m=L\).



Subtract the strong-convexity quadratic and its gradient:

\[F(v)=f(v)-\frac m2\|v\|^2,\qquad G(v)=\nabla f(v)-mv.\]

The quadratic lower model for \(f\) becomes an ordinary supporting plane for \(F\):

\[F(a)+\langle G(a),b-a\rangle\le F(b).\]

This follows by expanding the squared distance and collecting the terms involving \(m\).



Apply the smooth upper model to \(f\), and substitute the same definitions of \(F\) and \(G\).



Expanding the quadratic terms shows that the upper curvature decreases from \(L\) to \(L-m\):

\[F(b)\le F(a)+\langle G(a),b-a\rangle+\frac{L-m}{2}\|b-a\|^2.\]

Since \(L-m\ge0\), cocoercivity for these two models gives

\[\|G(x)-G(y)\|^2\le(L-m)\langle G(x)-G(y),x-y\rangle.\]

Write \(d=x-y\) and \(v=\nabla f(x)-\nabla f(y)\). The shifted gradient difference is exactly

\[G(x)-G(y)=v-md.\]

Expand the cocoercivity inequality using this identity:

\[\|v\|^2-2m\langle v,d\rangle+m^2\|d\|^2\le(L-m)\bigl(\langle v,d\rangle-m\langle d,d\rangle\bigr).\]

Use \(\langle d,d\rangle=\|d\|^2\) and collect terms. This gives

\[\|v\|^2+mL\|d\|^2\le(m+L)\langle v,d\rangle.\]

When \(m=L\), the cocoercivity constant above is zero, forcing \(G(x)=G(y)\); consequently \(\nabla f(x)-\nabla f(y)=m(x-y)\). Thus no division by \(L-m\) or limiting argument is needed.

\subsection{\texorpdfstring{The gradient difference lies in a ball}{The gradient difference lies in a ball}}


Let \(H\) be a real inner-product space, \(x,y,g_x,g_y\in H\), and \(m,L\in\mathbb R\) with \(m\le L\). Suppose

\[mL\|x-y\|^2+\|g_x-g_y\|^2\le(m+L)\langle g_x-g_y,x-y\rangle.\]

Then

\[\left\|(g_x-g_y)-\frac{m+L}{2}(x-y)\right\|\le\frac{L-m}{2}\|x-y\|.\]

The proposed radius is nonnegative because \(m\le L\). Both sides of the desired norm bound are therefore nonnegative, so it is enough to prove the corresponding squared inequality.



Set \(c=(m+L)/2\), \(r=(L-m)/2\), \(d=x-y\), and \(v=g_x-g_y\). Completing the square in the interpolation inequality gives

\[\begin{aligned}\|v-cd\|^2&=\|v\|^2-(m+L)\langle v,d\rangle+c^2\|d\|^2\\&\le(c^2-mL)\|d\|^2=r^2\|d\|^2.\end{aligned}\]

Taking nonnegative square roots proves the claim. If \(m=L\), the radius is zero and \(v=md\).

\subsection{\texorpdfstring{The centered Euler-\allowbreak{}step bound}{The centered Euler-step bound}}


Let \(H\) be a real inner-product space, \(m,L,h\in\mathbb R\) with \(m\le L\) and \(h\ge0\), and let \(x,y,g_x,g_y\in H\) satisfy

\[mL\|x-y\|^2+\|g_x-g_y\|^2\le(m+L)\langle g_x-g_y,x-y\rangle.\]

Then

\[\|(x-hg_x)-(y-hg_y)\|\le\left(\left|1-h\frac{m+L}{2}\right|+h\frac{L-m}{2}\right)\|x-y\|.\]

The preceding ball bound gives \(\|v-cd\|\le r\|d\|\), where \(d=x-y\), \(v=g_x-g_y\), \(c=(m+L)/2\), and \(r=(L-m)/2\).



Decompose the Euler difference into its centered part and its deviation:

\[d-hv=(1-hc)d-h(v-cd).\]

The triangle inequality and homogeneity of the norm give

\[\|d-hv\|\le|1-hc|\,\|d\|+h\|v-cd\|,\]

where \(|-h|=h\) because \(h\ge0\).



Multiply the ball bound by \(h\ge0\) and substitute it:

\[\|d-hv\|\le\bigl(|1-hc|+hr\bigr)\|d\|.\]

For \(m=L\), this reduces to the exact scalar factor \(|1-mh|\).

\subsection{\texorpdfstring{The sharp small-\allowbreak{}step factor}{The sharp small-step factor}}


Let \(H\) be a real inner-product space, \(x,y,g_x,g_y\in H\), and \(m,L,h\in\mathbb R\). Suppose \(m\le L\), \(m+L>0\), \(0\le h\le2/(m+L)\), and

\[mL\|x-y\|^2+\|g_x-g_y\|^2\le(m+L)\langle g_x-g_y,x-y\rangle.\]

Then

\[\|(x-hg_x)-(y-hg_y)\|\le(1-mh)\|x-y\|.\]

Apply the centered bound. The step-size assumption and \(m+L>0\) imply \(h(m+L)\le2\), hence \(1-h(m+L)/2\ge0\).



Remove the absolute value with this sign and simplify its coefficient:

\[\left|1-h\frac{m+L}{2}\right|+h\frac{L-m}{2}=1-h\frac{m+L}{2}+h\frac{L-m}{2}=1-mh.\]\subsection{\texorpdfstring{The sharp large-\allowbreak{}step factor}{The sharp large-step factor}}


Let \(H\) be a real inner-product space, \(x,y,g_x,g_y\in H\), and \(m,L,h\in\mathbb R\). Suppose \(m\le L\), \(m+L>0\), \(h\ge0\), \(2/(m+L)\le h\), and

\[mL\|x-y\|^2+\|g_x-g_y\|^2\le(m+L)\langle g_x-g_y,x-y\rangle.\]

Then

\[\|(x-hg_x)-(y-hg_y)\|\le(Lh-1)\|x-y\|.\]

This is a distance bound; it is a strict contraction only when its coefficient is below one.



Apply the centered estimate. Now the step-size assumption gives \(2\le h(m+L)\), so \(1-h(m+L)/2\le0\).



Replace the absolute value by the negative of its argument. The coefficient becomes

\[\left|1-h\frac{m+L}{2}\right|+h\frac{L-m}{2}=h\frac{m+L}{2}-1+h\frac{L-m}{2}=Lh-1.\]\subsection{\texorpdfstring{Small-\allowbreak{}step contraction for a potential}{Small-step contraction for a potential}}


Let \(f:H\to\mathbb R\) be differentiable and \(m\)-strongly convex on a real Hilbert space, with an \(L\)-Lipschitz gradient, \(L\ge0\), \(m\le L\), and \(m+L>0\). If \(0\le h\le2/(m+L)\), then for every \(x,y\),

\[\|(x-h\nabla f(x))-(y-h\nabla f(y))\|\le(1-mh)\|x-y\|.\]

The quadratic models imply the strong-smooth interpolation inequality for \(g_x=\nabla f(x)\) and \(g_y=\nabla f(y)\). Apply the small-step Euler bound to these two vectors. Its hypotheses are exactly the stated curvature and step-size assumptions.

\subsection{\texorpdfstring{Large-\allowbreak{}step contraction for a potential}{Large-step contraction for a potential}}


Let \(H\) be a real Hilbert space and \(f:H\to\mathbb R\) differentiable and \(m\)-strongly convex, with an \(L\)-Lipschitz gradient. Suppose \(L\ge0\), \(m\le L\), \(m+L>0\), \(h\ge0\), and \(2/(m+L)\le h\). Then for every \(x,y\in H\),

\[\|(x-h\nabla f(x))-(y-h\nabla f(y))\|\le(Lh-1)\|x-y\|.\]

Use the potential's interpolation inequality in the large-step Euler bound. This gives the second sharp factor \(Lh-1\).

\subsection{\texorpdfstring{The endpoint-\allowbreak{}factor form}{The endpoint-factor form}}


Let \(f\) be differentiable and \(m\)-strongly convex on a real Hilbert space, with an \(L\)-Lipschitz gradient, \(L\ge0\), \(m\le L\), and \(m+L>0\). For every \(h\ge0\) and every \(x,y\),

\[\|(x-h\nabla f(x))-(y-h\nabla f(y))\|\le\max\{|1-mh|,|1-Lh|\}\,\|x-y\|.\]

If \(0<m\le L\), this factor is below one for \(0<h<2/L\). At \(m=L\), the two endpoint factors coincide.



Split according to whether \(h\le2/(m+L)\). In the first case the small-step theorem bounds the distance by \((1-mh)\|x-y\|\).



Since \(1-mh\le|1-mh|\le\max\{|1-mh|,|1-Lh|\}\), multiplying by \(\|x-y\|\ge0\) gives the desired bound in this case.



In the other case \(h>2/(m+L)\), so the large-step theorem bounds the distance by \((Lh-1)\|x-y\|\).



The elementary absolute-value inequality yields

\[Lh-1=-(1-Lh)\le|1-Lh|.\]

Finally \(|1-Lh|\le\max\{|1-mh|,|1-Lh|\}\). Multiply by the nonnegative distance to finish the second case.

\subsection{\texorpdfstring{Sharp shared-\allowbreak{}noise contraction: small steps}{Sharp shared-noise contraction: small steps}}


Let \(f:E_d\to\mathbb R\) be differentiable and \(m\)-strongly convex, with an \(L\)-Lipschitz gradient, \(L\ge0\), \(m\le L\), and \(m+L>0\). If \(0\le h\le2/(m+L)\), then for all starting states \(x,y\) and every common noise vector \(z\),

\[\|T_h^z(x)-T_h^z(y)\|\le(1-mh)\|x-y\|.\]

The common noise cancels exactly, leaving the distance between the gradient steps from \(x\) and \(y\).



Apply the potential-based small-step estimate. It gives the factor \(1-mh\) independently of the value or distribution of \(z\).

\subsection{\texorpdfstring{Sharp shared-\allowbreak{}noise contraction: large steps}{Sharp shared-noise contraction: large steps}}


Let \(f:E_d\to\mathbb R\) be differentiable and \(m\)-strongly convex, with an \(L\)-Lipschitz gradient, \(L\ge0\), \(m\le L\), and \(m+L>0\). If \(h\ge0\) and \(2/(m+L)\le h\), then for every \(x,y,z\in E_d\),

\[\|T_h^z(x)-T_h^z(y)\|\le(Lh-1)\|x-y\|.\]

Cancel the shared noise, reducing the two ULA updates to their deterministic gradient-step difference.



Apply the potential-based large-step estimate. This yields the factor \(Lh-1\) for every common noise vector.

\section{\texorpdfstring{The local Euler defect}{The local Euler defect}}
\label{proof-module-UnadjustedLangevin-LocalError-lean}
{\small\color{gray}Source module: \texttt{Unadjusted\allowbreak{}Langevin/\allowbreak{}Local\allowbreak{}Error.\allowbreak{}lean}.\par}


Let \(H\) be a real Banach space, with norm \(\|\cdot\|\). All Lipschitz constants \(L\) below are nonnegative. All time integrals are Bochner integrals. For the Langevin application we will take \(H=L^2(P;\mathbb R^d)\); the norms below then measure random variables in mean square. The scalar defect estimate permits a nonnegative real parameter \(d\); its application uses the natural-number dimension, viewed as a real number. Begin with a Lipschitz drift \(g\), a reference curve \(Y\), and its additive noise \(B\). We first derive the error of freezing the drift over one time step directly from the integral equation.



The analytic calculation is conditional on the displayed integral equation and norm bounds. It neither constructs a stationary Langevin diffusion nor proves its score moment. The Gaussian and Brownian noise moments used later are proved separately.



\texttt{local\allowbreak{}Defect} is a vector-valued Bochner interval integral. It exists as a totalized expression before integrability is proved; all norm and linearity steps below explicitly establish the continuous or integrable hypotheses needed for their ordinary mathematical meaning. The interval \(0..h\) is oriented, while the error estimates assume \(h\ge0\).

\subsection{\texorpdfstring{Definition of the local defect}{Definition of the local defect}}


For \(g:H\to H\), \(Y:\mathbb R\to H\), and \(h\in\mathbb R\), define the drift omitted when Euler freezes \(g\) at the initial point:

\[D_h(g,Y):=\int_0^h\bigl(g(Y_t)-g(Y_0)\bigr)\,dt.\]

Integrate the difference between the evolving drift and its frozen value. This defines the defect itself; the estimates below will be consequences of this definition.

\subsection{\texorpdfstring{Separate the exact and frozen drift}{Separate the exact and frozen drift}}


If \(t\mapsto g(Y_t)\) is integrable on the interval between \(0\) and \(h\), then

\[D_h(g,Y)=\int_0^h g(Y_t)\,dt-h\,g(Y_0).\]

Expand the definition. Both the actual drift and the constant function \(g(Y_0)\) are integrable, so linearity separates their difference into two integrals.



The integral of the constant vector is \((h-0)g(Y_0)=h\,g(Y_0)\), yielding the identity.

\subsection{\texorpdfstring{An exact step is Euler minus its defect}{An exact step is Euler minus its defect}}


Assume the preceding drift integrability and the endpoint integral equation

\[Y_h-Y_0=-\int_0^h g(Y_t)\,dt+B_h.\]

Then the exact reference step has the form

\[Y_h=Y_0-h\,g(Y_0)+B_h-D_h(g,Y).\]

Solve the assumed integral equation for \(Y_h\), then replace the defect by the difference between the exact drift integral and its frozen approximation.



Cancel the two frozen-drift terms. The resulting equality is precisely the Euler update with the signed remainder \(-D_h\).

\subsection{\texorpdfstring{Continuity supplies drift integrability}{Continuity supplies drift integrability}}


Let \(H\) be a normed additive group, \(g:H\to H\) be \(L\)-Lipschitz with \(L\ge0\), and \(Y:\mathbb R\to H\) be continuous. For every real \(h\), the function \(t\mapsto g(Y_t)\) is interval integrable from \(0\) to \(h\).



A Lipschitz map is continuous, so \(g\circ Y\) is continuous. A continuous function is integrable on every compact interval, proving the claim.

\subsection{\texorpdfstring{The square-\allowbreak{}root integral}{The square-root integral}}


For every \(h\ge0\),

\[\int_0^h\sqrt t\,dt=\frac23h\sqrt h.\]

Write \(\sqrt t=t^{1/2}\) and use the power-integral formula. The exponent satisfies \(1/2>-1\), so the formula applies at the zero endpoint.



The contribution at zero vanishes. Since \(h\ge0\), the other endpoint is \(h^{1/2+1}=h^{1/2}h=\sqrt h\,h\).



Divide by \(1/2+1=3/2\) and rearrange the factors to obtain \((2/3)h\sqrt h\).

\subsection{\texorpdfstring{Bound a reference increment}{Bound a reference increment}}


Let \(h\ge0\) and suppose that, for every \(0\le s\le h\),

\[\|g(Y_s)\|\le G,\qquad \|B_s\|\le N\sqrt s,\qquad Y_s-Y_0=-\int_0^s g(Y_r)\,dr+B_s.\]

Then, for each \(0\le t\le h\),

\[\|Y_t-Y_0\|\le tG+N\sqrt t.\]

First bound the drift integral on \([0,t]\) by its constant norm bound:

\[\left\|\int_0^t g(Y_s)\,ds\right\|\le G\,|t-0|=Gt,\]

where \(t\ge0\) removes the absolute value.



To justify the constant bound under that integral, note that \(s\in(0,t]\) implies \(0\le s\le t\le h\). Hence the assumed bound on \([0,h]\) applies to every such \(s\).



Insert the integral equation and apply the triangle inequality:

\[\|Y_t-Y_0\|\le\left\|\int_0^t g(Y_s)\,ds\right\|+\|B_t\|.\]

Negation leaves the norm of the drift integral unchanged.



Apply the two bounds to obtain \(Gt+N\sqrt t\), and commute the factors \(G\) and \(t\).

\subsection{\texorpdfstring{Integrate the Lipschitz increment bound}{Integrate the Lipschitz increment bound}}


Assume \(h\ge0\), \(g\) is \(L\)-Lipschitz with \(L\ge0\), and for \(0\le t\le h\) the same integral equation and bounds \(\|g(Y_t)\|\le G\), \(\|B_t\|\le N\sqrt t\) hold. Then

\[\|D_h(g,Y)\|\le L\left(\frac{h^2G}{2}+\frac23h\sqrt h\,N\right).\]

For \(0<t\le h\), Lipschitz continuity gives

\[\|g(Y_t)-g(Y_0)\|\le L\|Y_t-Y_0\|.\]

Multiply the increment estimate by \(L\ge0\) to obtain the pointwise majorant

\[\|g(Y_t)-g(Y_0)\|\le L(tG+N\sqrt t).\]

The scalar majorant \(t\mapsto L(tG+N\sqrt t)\) is continuous, hence integrable on \([0,h]\).



Expand \(D_h\). The norm estimate for a Bochner integral, together with the integrable majorant, gives

\[\|D_h\|\le\int_0^h L(tG+N\sqrt t)\,dt.\]

Pull out constants and split the two integrable summands. Evaluate

\[\int_0^h t\,dt=\frac{h^2}{2},\qquad \int_0^h\sqrt t\,dt=\frac23h\sqrt h.\]

This leaves the two drift and noise contributions in the claimed expression.



Rearrange the scalar factors to place the drift term first and the noise term second.

\subsection{\texorpdfstring{The exact rational constant}{The exact rational constant}}


The numerical factor in the corrected Dalalyan–Karagulyan estimate satisfies

\[\frac76\sqrt2\le\frac{33}{20}=1.65.\]

Use \((\sqrt2)^2=2\) and \(\sqrt2\ge0\). Comparing squares proves the bound, since \(49/18\le1089/400\).

\subsection{\texorpdfstring{The local error with stationary moment scales}{The local error with stationary moment scales}}


Let \(h,d\in\mathbb R\) with \(h,d\ge0\), let \(g\) be \(L\)-Lipschitz with \(L\ge0\), and suppose \(Lh\le2\). Assume that, for every \(0\le t\le h\),

\[Y_t-Y_0=-\int_0^t g(Y_s)\,ds+B_t,\quad \|g(Y_t)\|\le\sqrt{Ld},\quad \|B_t\|\le\sqrt{2d}\sqrt t.\]

Then

\[\boxed{\|D_h(g,Y)\|\le\frac{33}{20}Lh\sqrt{hd}.}\]

In \(L^2\), the drift bound corresponds to a score second moment at most \(Ld\); here that bound remains a hypothesis.



Compare the nonnegative quantities \(h\sqrt{Ld}\) and \(\sqrt2\sqrt{hd}\) by squaring. It suffices to show \(h^2Ld\le2hd\).



Replace each squared square root by its nonnegative argument. Multiply \(Lh\le2\) by \(hd\ge0\), obtaining \(h^2Ld\le2hd\), and thus

\[h\sqrt{Ld}\le\sqrt2\sqrt{hd}.\]

Split products under the square roots and rearrange:

\[h\sqrt h\sqrt{2d}=h\sqrt2\sqrt{hd}.\]

The nonnegativity of \(h\) and \(2\) justifies the square-root product identities.



Apply the two-term local-defect bound with \(G=\sqrt{Ld}\) and \(N=\sqrt{2d}\):

\[\|D_h\|\le L\left(\frac{h^2\sqrt{Ld}}2+\frac23h\sqrt h\sqrt{2d}\right).\]

Bound the first term using \(h\sqrt{Ld}\le\sqrt2\sqrt{hd}\), and rewrite the second with the preceding identity. Since \(1/2+2/3=7/6\),

\[\|D_h\|\le L\frac76\sqrt2\,h\sqrt{hd}.\]

Multiply \((7/6)\sqrt2\le33/20\) by the nonnegative factor \(Lh\sqrt{hd}\). This gives the claimed constant without numerical rounding.

\section{\texorpdfstring{The noise-\allowbreak{}driven iterates}{The noise-driven iterates}}
\label{proof-module-UnadjustedLangevin-Dynamics-lean}
{\small\color{gray}Source module: \texttt{Unadjusted\allowbreak{}Langevin/\allowbreak{}Dynamics.\allowbreak{}lean}.\par}


Write \(T_h(x)=x-h\nabla U(x)\), so one Euler update is \(T_h(x)+\sqrt{2h}\,z\). The noise vectors remain explicit. Common-noise algebra works for any supplied noise sequence; identifying the algorithm with Gaussian ULA also requires independent standard Gaussian noise hypotheses.



\texttt{ula\allowbreak{}Chain} recursively applies \texttt{ula\allowbreak{}Update} to a supplied sequence of vectors. Its zero and successor equations are definitional. \texttt{ula\allowbreak{}Chain\_\allowbreak{}measurable} proves measurable dependence on the sample by induction, combining the current iterate, the Lipschitz gradient's measurability, and the next supplied measurable innovation. No independence or Gaussian-law assumption is needed for this measurability assertion.

\subsection{\texorpdfstring{Definition of the chain}{Definition of the chain}}


Given \(U:\mathbb R^d\to\mathbb R\), \(h\in\mathbb R\), an initial vector \(x_0\), and noise \((\xi_k)\), define

\[X_0=x_0,\qquad X_{k+1}=X_k-h\nabla U(X_k)+\sqrt{2h}\,\xi_k.\]

Start the recursion at the supplied initial vector: \(X_0=x_0\).



Apply the Euler update to \(X_k\) using the next noise vector \(\xi_k\). This determines \(X_{k+1}\).

\subsection{\texorpdfstring{The initial iterate}{The initial iterate}}


For every potential, step, initial vector, and noise sequence, the recursively defined chain satisfies \(X_0=x_0\).



The equality is exactly the initial clause of the recursive definition.

\subsection{\texorpdfstring{The update identity}{The update identity}}


For every \(k\ge0\), the chain satisfies \(X_{k+1}=\operatorname{ulaUpdate}(U,h,X_k,\xi_k)\).



The successor clause of the recursion gives the asserted update formula.

\subsection{\texorpdfstring{Measurability of every iterate}{Measurability of every iterate}}


Let \(\Omega\) be a measurable space. If \(\nabla U\) is \(L\)-Lipschitz, the initial vector \(X_0:\Omega\to\mathbb R^d\) is measurable, and every \(\xi_k:\Omega\to\mathbb R^d\) is measurable, then \(X_n\) is measurable for every \(n\).



A Lipschitz gradient is continuous and therefore measurable.



Induct on \(n\). The initial iterate is measurable by hypothesis. For the successor step, assume \(X_n\) is measurable.



The composition \(\nabla U\circ X_n\) is measurable. Subtracting its scalar multiple from \(X_n\), then adding \(\sqrt{2h}\,\xi_n\), preserves measurability. Hence \(X_{n+1}\) is measurable.

\section{\texorpdfstring{The actual gradient on square-\allowbreak{}integrable random variables}{The actual gradient on square-integrable random variables}}
\label{proof-module-UnadjustedLangevin-LpDrift-lean}
{\small\color{gray}Source module: \texttt{Unadjusted\allowbreak{}Langevin/\allowbreak{}Lp\allowbreak{}Drift.\allowbreak{}lean}.\par}


Let \(P\) be a finite measure on \(\Omega\). The first five results need only a normed additive commutative group \(H\); all Lipschitz constants \(L\) are nonnegative. When \(H\) is a real Banach space, \(L^2(P;H)\) provides the Banach space needed for the local-defect estimate. For a Lipschitz vector field \(g:H\to H\), construct the operator \(\mathcal G[X]=[g\circ X]\), including the case \(g(0)\ne0\). Finiteness of \(P\) ensures the constant offset is square integrable.



The gradient results use a real Hilbert space \(H\). The two score-moment lemmas need only a real inner-product space, without completeness. Whenever \(g=\nabla U\), this is the potential’s actual gradient, and the resulting \(L^2\) recursion is proved to represent the original pointwise ULA chain.



\texttt{lp\allowbreak{}Drift} lifts a Lipschitz pointwise map to \(L^2\), using its linear growth and the finite probability measure to prove integrability. Its representative theorem identifies the lifted map almost everywhere with ordinary composition. All subsequent Hilbert-space estimates concern these quotient classes and are invariant under representative choices.

\subsection{\texorpdfstring{Center a Lipschitz vector field}{Center a Lipschitz vector field}}


In a normed additive group \(H\), if \(g:H\to H\) is \(L\)-Lipschitz, then so is

\[\widetilde g(x):=g(x)-g(0).\]

For arbitrary \(x,y\), the common offset cancels:

\[\|\widetilde g(x)-\widetilde g(y)\|=\|g(x)-g(y)\|\le L\|x-y\|.\]

This is precisely the required Lipschitz inequality.

\subsection{\texorpdfstring{Construct the lifted drift}{Construct the lifted drift}}


Let \(P\) be finite and \(g\) be \(L\)-Lipschitz. For \(X\in L^2(P;H)\), define

\[\mathcal G[X]:=[\widetilde g\circ X]+[g(0)].\]

The centered function satisfies \(\widetilde g(0)=0\) and \(\|\widetilde g(x)\|\le L\|x\|\), so composition produces an \(L^2\) class. Add the constant class \([g(0)]\), which belongs to \(L^2\) because \(P\) is finite.

\subsection{\texorpdfstring{The lift agrees with the original vector field}{The lift agrees with the original vector field}}


For finite \(P\), Lipschitz \(g\), and \(X\in L^2(P;H)\),

\[\mathcal G[X](\omega)=g(X(\omega))\quad\text{for }P\text{-almost every }\omega.\]

Expand \(\mathcal G\) into the centered composition class plus the constant class.



Intersect the full-measure sets on which the representative of the sum, the composition, and the constant class agree with their pointwise formulas.



On this set the value is \((g(X(\omega))-g(0))+g(0)\). Cancel the offset to recover \(g(X(\omega))\).

\subsection{\texorpdfstring{The actual pointwise drift is square integrable}{The actual pointwise drift is square integrable}}


For finite \(P\), Lipschitz \(g\), and \(X\in L^2(P;H)\),

\[g\circ X\in L^2(P;H).\]

The constructed class \(\mathcal G[X]\) has an \(L^2\) representative. Transfer its square integrability to \(g\circ X\) using their almost-everywhere equality.

\subsection{\texorpdfstring{The drift keeps its Lipschitz constant in L\ensuremath{{}^{2}}}{The drift keeps its Lipschitz constant in L²}}


For finite \(P\) and \(L\)-Lipschitz \(g\), the lifted drift is \(L\)-Lipschitz:

\[\|\mathcal G[X]-\mathcal G[Y]\|_{L^2}\le L\|X-Y\|_{L^2}.\]

Fix \(X,Y\). The two constant classes cancel in their drift difference. For the remaining centered compositions, the pointwise Lipschitz estimate passes to the \(L^2\) norm, giving the same constant \(L\).

\subsection{\texorpdfstring{Lift the small-\allowbreak{}step gradient contraction}{Lift the small-step gradient contraction}}


Let \(H\) be a real Hilbert space, \(P\) finite, and \(U:H\to\mathbb R\) differentiable and \(m\)-strongly convex on \(H\), with \(L\)-Lipschitz gradient. Assume \(m\le L\), \(m+L>0\), and \(0\le h\le2/(m+L)\). Write \(\mathcal G[X]=[\nabla U\circ X]\). Then, for all \(X,Y\in L^2(P;H)\),

\[\|(X-h\mathcal G[X])-(Y-h\mathcal G[Y])\|_{L^2}\le(1-mh)\|X-Y\|_{L^2}.\]

It suffices to prove the corresponding norm inequality almost everywhere between representatives: integrating the squared pointwise inequality gives the \(L^2\) bound.



Work on a common full-measure set where the outer difference and both Euler-step differences have their ordinary pointwise subtraction formulas.



Also require the scalar multiples, both lifted gradients, and \(X-Y\) to have their pointwise formulas. The two drift terms are then \(\nabla U(X(\omega))\) and \(\nabla U(Y(\omega))\).



Substitute these representative identities and apply the deterministic sharp contraction to \(x=X(\omega)\), \(y=Y(\omega)\):

\[\|(x-h\nabla U(x))-(y-h\nabla U(y))\|\le(1-mh)\|x-y\|.\]\subsection{\texorpdfstring{Lift the second step-\allowbreak{}size regime}{Lift the second step-size regime}}


Let \(H\) be a real Hilbert space, \(P\) finite, and \(U\) differentiable, \(m\)-strongly convex, with \(L\)-Lipschitz gradient. Assume \(m\le L\), \(m+L>0\), \(h\ge0\), and \(h\ge2/(m+L)\). Then

\[\|(X-h\mathcal G[X])-(Y-h\mathcal G[Y])\|_{L^2}\le(Lh-1)\|X-Y\|_{L^2}.\]

The inequality holds in this whole regime; it is a strict contraction when its factor is less than one.



Reduce the \(L^2\) inequality to its almost-everywhere pointwise counterpart with factor \(Lh-1\).



Choose a full-measure set on which the outer subtraction and each Euler-step subtraction agree with the pointwise vector formulas.



On the same set identify scalar multiplication, both lifted gradient representatives, and the representative of \(X-Y\).



After substitution, invoke the deterministic second-regime estimate at \(X(\omega),Y(\omega)\):

\[\|(x-h\nabla U(x))-(y-h\nabla U(y))\|\le(Lh-1)\|x-y\|.\]\subsection{\texorpdfstring{The local defect for the actual L\ensuremath{{}^{2}} gradient}{The local defect for the actual L² gradient}}


Let \(P\) be finite, \(H\) a real Hilbert space, and \(\nabla U\) be \(L\)-Lipschitz. Let \(Y,B:\mathbb R\to L^2(P;H)\), let \(h,d\in\mathbb R\), and assume \(h,d\ge0\), \(Lh\le2\). For every \(0\le t\le h\), suppose

\[\|\mathcal G[Y_t]\|_{L^2}\le\sqrt{Ld},\quad\|B_t\|_{L^2}\le\sqrt{2d}\sqrt t,\quad Y_t-Y_0=-\int_0^t\mathcal G[Y_s]\,ds+B_t.\]

Then

\[\|D_h(\mathcal G,Y)\|_{L^2}\le\frac{33}{20}Lh\sqrt{hd}.\]

Apply the Banach-space local-defect theorem to \(\mathcal G\). Its required Lipschitz property is the one just proved for the actual lifted gradient; the reference equation and moment-scale bounds are the displayed hypotheses.

\subsection{\texorpdfstring{Identify the drift norm with the score moment}{Identify the drift norm with the score moment}}


For finite \(P\), a real inner-product space \(H\), Lipschitz \(g:H\to H\), and \(X\in L^2(P;H)\),

\[\|\mathcal G[X]\|_{L^2}^2=\int_\Omega\|g(X(\omega))\|^2\,dP(\omega).\]

Write the squared \(L^2\) norm as a self-inner-product and expand the \(L^2\) inner product into the integral of the pointwise inner product.



Almost everywhere, \(\mathcal G[X](\omega)=g(X(\omega))\). Replace the representative under the integral and use \(\langle v,v\rangle=\|v\|^2\).

\subsection{\texorpdfstring{A second-\allowbreak{}moment bound gives the drift bound}{A second-moment bound gives the drift bound}}


Under the preceding finite-measure and inner-product assumptions, if

\[\int_\Omega\|g(X(\omega))\|^2\,dP(\omega)\le C,\]

then

\[\|\mathcal G[X]\|_{L^2}\le\sqrt C.\]

For \(g=\nabla U\) and \(C=Ld\), the moment inequality is an explicit input, not a consequence of this lemma.



The norm-square identity turns the assumed moment bound into \(\|\mathcal G[X]\|^2\le C\). Taking the nonnegative square root yields the conclusion.

\subsection{\texorpdfstring{Define one ULA update in L\ensuremath{{}^{2}}}{Define one ULA update in L²}}


For finite \(P\), a real Hilbert space \(H\), a Lipschitz gradient \(\nabla U\), step \(h\in\mathbb R\), and \(X,Z\in L^2(P;H)\), define

\[\mathsf T_h(X,Z):=X-h\mathcal G[X]+\sqrt{2h}\,Z.\]

Form the drift step and add the scaled noise inside \(L^2\). All terms are already equivalence classes of square-integrable variables, so the vector-space operations produce another such class.

\subsection{\texorpdfstring{The L\ensuremath{{}^{2}} update is the pointwise ULA formula}{The L² update is the pointwise ULA formula}}


For finite \(P\), a real Hilbert space \(H\), Lipschitz \(\nabla U\), and arbitrary \(X,Z\in L^2(P;H)\),

\[\mathsf T_h(X,Z)(\omega)=X(\omega)-h\nabla U(X(\omega))+\sqrt{2h}\,Z(\omega)\quad\text{a.e.}\]

This identity needs no Gaussian or independence assumption on \(Z\).



Expand the definition of \(\mathsf T_h\) into its subtraction, drift, and scaled-noise terms.



Intersect the full-measure sets on which addition, subtraction, both scalar multiples, and the lifted gradient agree with their pointwise formulas.



Substitute these identities. The representative is exactly \(X(\omega)-h\nabla U(X(\omega))+\sqrt{2h}\,Z(\omega)\).

\subsection{\texorpdfstring{Define the genuine L\ensuremath{{}^{2}} ULA chain}{Define the genuine L² ULA chain}}


For a finite measure \(P\), a Lipschitz gradient on \(E_d\), an initial class \(X_0\in L^2(P;E_d)\), and noise classes \(Z_n\in L^2(P;E_d)\), define recursively, with the supplied initial class at index zero,

\[X_{n+1}:=\mathsf T_h(X_n,Z_n).\]

The initial class is the supplied \(X_0\). At the successor index, apply the \(L^2\) ULA update to the preceding iterate and the next noise class \(Z_n\).

\subsection{\texorpdfstring{The initial iterate}{The initial iterate}}


The recursively defined chain satisfies \(X_0=\mathrm{initial}\).



This is the zero branch of the recursion, so both sides reduce to the same class.

\subsection{\texorpdfstring{The successor iterate}{The successor iterate}}


For every \(n\ge0\), the same finite-measure, Lipschitz-gradient, initial-class and noise-class data give

\[X_{n+1}=\mathsf T_h(X_n,Z_n).\]

This is the successor branch of the defining recursion; unfold that branch to obtain the equality.

\subsection{\texorpdfstring{Identify the chain with the original algorithm}{Identify the chain with the original algorithm}}


For finite \(P\), Lipschitz gradient \(\nabla U:E_d\to E_d\), square-integrable initial and noise classes, and every \(n\ge0\),

\[X_n(\omega)=\operatorname{ulaChain}\bigl(U,h,X_0(\omega),(Z_k(\omega))_{k\ge0},n\bigr)\quad\text{for }P\text{-a.e. }\omega.\]

Induct on \(n\). At \(n=0\), both recursions return the same initial representative, so equality holds for every outcome. For the successor case, assume the representatives agree at index \(n\).



On a common full-measure set, combine the one-step representative identity with the induction hypothesis. Replace \(X_n(\omega)\) by the pointwise chain’s preceding iterate and apply the two successor recursion equations. Both next iterates are the same ULA update.

\section{\texorpdfstring{Existence of the additive Langevin integral equation}{Existence of the additive Langevin integral equation}}
\label{proof-module-UnadjustedLangevin-LangevinIntegral-lean}
{\small\color{gray}Source module: \texttt{Unadjusted\allowbreak{}Langevin/\allowbreak{}Langevin\allowbreak{}Integral.\allowbreak{}lean}.\par}


In a complete real normed space, a globally Lipschitz drift and a continuous additive path \(B\), \(B(0)=0\), determine a continuous solution of \(Y_t-x=-\int_0^tg(Y_s)ds+B_t\) on every finite nonnegative horizon. The construction is deterministic and later applies with the Banach space \(L^2(P)\).

\subsection{\texorpdfstring{Solve on a short interval}{Solve on a short interval}}


Remove the additive input by writing \(Y(t)=Z(t)+B(t)\). The new unknown satisfies \(Z(t)=x+\int_0^t f(s,Z(s))ds\), with \(f(s,z)=-g(z+B(s))\). This time-dependent drift is \(L\)-Lipschitz in \(z\) and continuous in time even though \(B\) need not be differentiable.



Bound \(\|B(t)\|\le C'\) on the compact interval and set \(G=\|g(x)\|+LC'\), radius \(a=2hG\), and drift bound \(M=G+La\). The condition \(Lh\le1/2\) proves \(Mh\le a\), giving the ball-invariance condition in \texttt{Is\allowbreak{}Picard\allowbreak{}Lindelof}. The proof checks these bounds and invokes \texttt{ODE.\allowbreak{}Fun\allowbreak{}Space.\allowbreak{}exists\_\allowbreak{}is\allowbreak{}Fixed\allowbreak{}Pt\_\allowbreak{}next}.



The chosen fixed point gives a continuous \(Z\) with its integral identity. Add back \(B\), use \(B(0)=0\), and rewrite the negative drift integral to obtain \(Y(t)-x=-\int_0^tg(Y(s))ds+B(t)\). The solution is represented as a real-time function, with the required identity asserted on \([0,h]\).

\subsection{\texorpdfstring{splice langevin integral}{splice langevin integral}}


Use \(Y(t)\) up to time \(a\), and \(Z(t-a)\) afterward. The equality \(Z(0)=Y(a)\) makes the pasted path continuous at the join. Split \(\int_0^t=\int_0^a+\int_a^t\), translate the second integral, and add the two supplied equations; the terms \(B(a)\) cancel, giving the same equation on \([0,a+b]\).

\subsection{\texorpdfstring{exists langevin integral}{exists langevin integral}}


Choose a positive segment length small enough that \(L\delta\le1/2\). Construct the first Picard segment, then repeatedly solve with initial state at the previous endpoint and shifted input \(B(a+s)-B(a)\). The splicing lemma combines finitely many segments covering the prescribed horizon. The local contraction restriction determines the construction mesh and imposes no restriction on the final horizon.

\subsection{\texorpdfstring{exists lp langevin integral small}{exists lp langevin integral small}}


Apply the Banach-space Picard theorem to \(H=L^2(P;K)\), \(g(X)=\operatorname{lpDrift}(\nabla U)(X)\). The lifted drift has the same Lipschitz constant as \(\nabla U\). Completeness of \(L^2\) and continuity of the supplied \(L^2\)-noise path are the exact fixed-point hypotheses.

\subsection{\texorpdfstring{exists lp langevin integral}{exists lp langevin integral}}


Use the finite-horizon Banach-space existence theorem with the lifted gradient drift. The conclusion is continuity in the \(L^2\) norm and an equality of \(L^2\) classes at every time in the horizon. The Bochner time integral therefore lives in the same Hilbert space; no pathwise differentiability of the random noise is asserted.

\section{\texorpdfstring{Construct the kinetic integral equation by Picard iteration}{Construct the kinetic integral equation by Picard iteration}}
\label{proof-module-KineticLangevin-Dynamics-lean}
{\small\color{gray}Source module: \texttt{Kinetic\allowbreak{}Langevin/\allowbreak{}Dynamics.\allowbreak{}lean}.\par}


For existence, the phase product may use its convenient maximum norm. This norm is equivalent to the Euclidean phase norm but is used here only to establish a global Lipschitz vector field. The input noise is any continuous additive velocity path starting at zero. The theorem therefore constructs paths without postulating a reference kinetic process; a full Brownian solution theorem additionally identifies the stochastic noise and its laws.

\subsection{\texorpdfstring{Write the drift in the shared integral-\allowbreak{}equation sign convention}{Write the drift in the shared integral-equation sign convention}}


The actual kinetic drift is \((-\gamma v-ug(x),v)\). The shared Picard theorem solves an equation with a minus sign in front of its drift integral, so \texttt{kinetic\allowbreak{}Negative\allowbreak{}Drift} stores \((\gamma v+ug(x),-v)\). This sign choice ensures that \(Y_t-Y_0=-\int_0^tD(Y_s)ds+(B_t,0)\) has velocity acceleration \(-\gamma v-ug(x)\) and position derivative \(v\). The second noise coordinate is exactly zero.

\subsection{\texorpdfstring{Prove a global Lipschitz constant for phase space}{Prove a global Lipschitz constant for phase space}}


If \(g\) is \(L\)-Lipschitz and \(\gamma,u,L\ge0\), write \(R=\|(v,x)-(w,y)\|_{\max}\). Each coordinate difference is at most \(R\), and \(\|g(x)-g(y)\|\le LR\). The velocity component of the negative drift is bounded by \((\gamma+uL)R\) using the triangle inequality. The position component is bounded by \(R\). Hence the maximum is bounded by \((\gamma+uL+1)R\). Lean opens the product maximum, proves the two coordinate bounds, rewrites the drift difference into scalar multiples of differences with \texttt{module}, and combines the nonnegative coefficients. The constant is deliberately sufficient for existence rather than optimized for contraction.

\subsection{\texorpdfstring{Invoke the proved continuous-\allowbreak{}path Picard construction}{Invoke the proved continuous-path Picard construction}}


Let the underlying normed space be complete, \(B\) continuous with \(B_0=0\), and \(T\ge0\). The shared theorem \texttt{ULA.\allowbreak{}exists\_\allowbreak{}langevin\_\allowbreak{}integral} constructs a continuous phase path \(Y\), with \(Y_0=\text{initial}\), satisfying the kinetic integral identity for every \(0\le t\le T\). The call supplies the concrete negative drift and its proved Lipschitz constant. It packages the noise as \((B_t,0)\), whose continuity follows from product continuity and whose value at zero follows from \texttt{h\allowbreak{}B0}. This is finite-horizon pathwise existence; Gaussian stationarity, a canonical Brownian law, and full probabilistic chain convergence are not conclusions of this declaration.

\section{\texorpdfstring{Stability and uniqueness for additive-\allowbreak{}noise solutions}{Stability and uniqueness for additive-noise solutions}}
\label{proof-module-UnadjustedLangevin-LangevinUniqueness-lean}
{\small\color{gray}Source module: \texttt{Unadjusted\allowbreak{}Langevin/\allowbreak{}Langevin\allowbreak{}Uniqueness.\allowbreak{}lean}.\par}


Two solutions driven by the same continuous additive noise satisfy \(\|Y_t-Z_t\|\le e^{Lt}\|x-y\|\) on the prescribed interval. Equality of initial values gives uniqueness. Only global Lipschitzness is used here; the sharper contractive estimate uses strong convexity elsewhere.

\subsection{\texorpdfstring{langevin integral stability}{langevin integral stability}}


Subtract the two integral equations so that their common additive input cancels. The integral norm inequality and Lipschitz drift give \(u(t)\le\|x-y\|+L\int_0^t u(s)ds\), where \(u(t)=\|Y(t)-Z(t)\|\) is continuous. Grönwall's inequality yields \(u(t)\le\|x-y\|e^{Lt}\). All time integrals are justified by continuity on the compact horizon.

\subsection{\texorpdfstring{langevin integral unique}{langevin integral unique}}


Set the two initial states equal in the stability estimate. The right side becomes \(0e^{Lt}=0\), forcing \(\|Y(t)-Z(t)\|=0\) for each \(t\in[0,h]\). Norm-zero equivalence gives equality of the vectors, packaged in Lean as \texttt{Eq\allowbreak{}On} on the interval.

\subsection{\texorpdfstring{lp langevin integral stability}{lp langevin integral stability}}


Specialize Banach-space stability to the actual lifted gradient on \(L^2(P;K)\). Its Lipschitz theorem supplies the required constant, so \(\|Y(t)-Z(t)\|_2\le e^{Lt}\|Y(0)-Z(0)\|_2\). This compares solutions with a common \(L^2\)-noise input and possibly different initial random variables.

\subsection{\texorpdfstring{lp langevin integral unique}{lp langevin integral unique}}


Apply the same-input uniqueness theorem in \(L^2(P;K)\). Equality here is equality of \(L^2\) classes for each time, hence almost-everywhere equality of representatives at each fixed time. The conclusion does not by itself select one exceptional null set valid simultaneously at every real time.

\section{\texorpdfstring{The actual kinetic solution depends continuously on its initial state and noise path}{The actual kinetic solution depends continuously on its initial state and noise path}}
\label{proof-module-KineticLangevin-Flow-lean}
{\small\color{gray}Source module: \texttt{Kinetic\allowbreak{}Langevin/\allowbreak{}Flow.\allowbreak{}lean}.\par}


Picard existence alone does not ensure a selected solution can be used as a measurable random process. This module proves stability when both the initial condition and continuous forcing vary, then constructs the actual kinetic flow by selection from the existing existence theorem. The stability estimate proves that the selected endpoint, and indeed the whole compact solution path, is jointly Lipschitz. No measurability, output-law identity or contraction conclusion is postulated.

\subsection{\texorpdfstring{Subtract the rough forcing before differentiating}{Subtract the rough forcing before differentiating}}


Let \(Y_t-x=-\int_0^tg(Y_s)ds+B_t\) and \(Z_t-y=-\int_0^tg(Z_s)ds+C_t\). The inputs \(B,C\) are continuous and may be nowhere differentiable. Set \(F_t=Y_t-Z_t-(B_t-C_t)\). Subtract the integral equations to obtain \(F_t=x-y-\int_0^t(g(Y_s)-g(Z_s))ds\). This expression is differentiable from the right, including at the initial endpoint. If \(\|B_t-C_t\|\le D\), then \(\|Y_t-Z_t\|\le\|F_t\|+D\), and an \(L\)-Lipschitz drift gives \(\|F'_t\|\le(L+1)\|F_t\|+(L+1)D\). Gronwall and the initial equality \(F_0=x-y\) yield \(\|Y_t-Z_t\|\le(\|x-y\|+D)e^{(L+1)t}\). Lean proves equality with the differentiable integral expression on a right neighborhood and uses the within-derivative Gronwall theorem. The harmless extra one makes the coefficient strictly positive and avoids a special zero-Lipschitz case.

\subsection{\texorpdfstring{Use compact paths with a uniform norm}{Use compact paths with a uniform norm}}


\texttt{Velocity\allowbreak{}Forcing T} is the Banach space of continuous velocity-forcing paths on \([0,T]\), equipped with the supremum norm. The extension \(\widehat B_t=B(\operatorname{clip}_{[0,T]}t)-B(0)\) is continuous on the whole real line and starts at zero. This permits direct use of the proved finite-horizon Picard construction. Two normalized extensions satisfy \(\|\widehat B_t-\widehat C_t\|\le2\|B-C\|_\infty\): each of the two evaluation differences is bounded by the uniform norm. Lean uses continuity of the interval projection, evaluates at zero, and applies the norm triangle inequality to the two differences. No derivative of a forcing path is taken.

\subsection{\texorpdfstring{Select an actual solution from proved existence}{Select an actual solution from proved existence}}


\texttt{kinetic\allowbreak{}Flow} uses classical choice on \texttt{exists\_\allowbreak{}kinetic\_\allowbreak{}integral}, supplying the normalized forcing \((\widehat B_t,0)\). Its specification theorem retains all conclusions of that construction: the path is globally continuous, its initial phase is exactly the supplied phase, and it satisfies the concrete kinetic integral equation at every time on the horizon. Classical selection is justified because existence was proved analytically. It does not add a solution axiom. Subsequent stability between any two solutions makes every selected representative agree on the horizon and provides continuous parameter dependence.



\texttt{kinetic\allowbreak{}Flow\_\allowbreak{}eq\_\allowbreak{}of\_\allowbreak{}integral} applies the proved Lipschitz integral-equation uniqueness theorem to identify the selection with any other continuous solution driven by the same normalized forcing and initial phase. Thus a concrete Brownian realization can be identified with this canonical flow without assuming equality of endpoint laws.

\subsection{\texorpdfstring{Prove joint endpoint continuity quantitatively}{Prove joint endpoint continuity quantitatively}}


The kinetic negative drift has Lipschitz constant \(K_0=\gamma+uL+1\) in the product maximum norm. Apply the two-input estimate with forcing difference \(D=2\|B-C\|_\infty\). This gives

\[\|Y_t^{z,B}-Y_t^{w,C}\|\le(\|z-w\|+2\|B-C\|_\infty)e^{(K_0+1)t}.\]

In the product parameter norm both input differences are bounded by the same distance, so the endpoint is Lipschitz with constant \(3e^{(K_0+1)t}\). Lean derives the estimate from the actual selected paths and their specifications, then bounds each coordinate by the product norm. \texttt{kinetic\allowbreak{}Flow\_\allowbreak{}endpoint\_\allowbreak{}continuous} follows from the Lipschitz theorem. Thus the eventual random endpoint map is measurable for a proved reason, even though its original existence definition used classical choice.

\subsection{\texorpdfstring{Control the entire solution as a compact continuous path}{Control the entire solution as a compact continuous path}}


\texttt{kinetic\allowbreak{}Flow\allowbreak{}Path} restricts the selected solution to \([0,T]\), retaining its proved continuity. At each time its difference has the endpoint bound. Since \(K_0+1\ge0\), \(e^{(K_0+1)t}\le e^{(K_0+1)T}\) on the horizon. The supremum-norm characterization then bounds the whole path difference by \(3e^{(K_0+1)T}\) times the initial-state/\allowbreak{}forcing-path distance. Lean applies \texttt{Continuous\allowbreak{}Map.\allowbreak{}norm\_\allowbreak{}le}, proves the pointwise exponential comparison, and uses the endpoint Lipschitz estimate. The resulting path-valued Lipschitz and continuity theorems support pushforwards into a genuine space of continuous trajectories, not merely separate endpoint choices.

\section{\texorpdfstring{Shared-\allowbreak{}noise cancellation and uniqueness of the actual integral solution}{Shared-noise cancellation and uniqueness of the actual integral solution}}
\label{proof-module-KineticLangevin-IntegralCoupling-lean}
{\small\color{gray}Source module: \texttt{Kinetic\allowbreak{}Langevin/\allowbreak{}Integral\allowbreak{}Coupling.\allowbreak{}lean}.\par}


Individual kinetic velocity paths contain continuous noise and need not be differentiable. Their difference under the same additive forcing is differentiable from the right, because the forcing cancels in the integral equations themselves. This module derives that assertion and proves uniqueness by Gronwall, including the initial time.

\subsection{\texorpdfstring{Subtract the two integral equations}{Subtract the two integral equations}}


Suppose \(Y_t-Y_0=-\int_0^t g(Y_s)ds+B_t\) and \(Z_t-Z_0=-\int_0^t g(Z_s)ds+B_t\), with continuous paths and a Lipschitz drift. Subtraction gives

\[Y_t-Z_t=Y_0-Z_0-\int_0^t[g(Y_s)-g(Z_s)]ds.\]

The two drift integrals exist on every compact interval because their integrands are continuous. Lean first establishes interval integrability, distributes the difference through the integral, rewrites both supplied equations as expressions for the paths, and cancels the common forcing with the additive-group tactic. No derivative of the noise is taken.

\subsection{\texorpdfstring{Differentiate after cancellation, including the initial endpoint}{Differentiate after cancellation, including the initial endpoint}}


The integrand \(g(Y_s)-g(Z_s)\) is continuous. The fundamental theorem of calculus therefore differentiates the integral expression, producing the right derivative \(-(g(Y_t)-g(Z_t))\). The identity is known only for \(t\in[0,T]\); consequently the mathematically appropriate conclusion at zero is a derivative within \([t,\infty)\). Lean builds an eventual neighborhood lying inside the valid time interval and transfers the derivative from the integral formula using eventual equality. This avoids imposing an arbitrary differentiable extension of the stochastic path to negative time.

\subsection{\texorpdfstring{Uniqueness from the global Lipschitz bound}{Uniqueness from the global Lipschitz bound}}


For \(D=Y-Z\), the derived derivative obeys \(\|D\prime_t\|\le L\|D_t\|\). Equal initial conditions give \(\|D_0\|=0\). Gronwall then yields \(\|D_t\|\le0\) throughout the finite horizon. Lean passes the right-derivative theorem to its norm version of Gronwall, simplifies the zero initial bound, and uses definiteness of the norm to recover \(Y_t=Z_t\). This result supplies uniqueness for any continuous additive forcing, including the forcing used in the kinetic Picard construction.

\section{\texorpdfstring{The transformed kinetic energy}{The transformed kinetic energy}}
\label{proof-module-KineticLangevin-Geometry-lean}
{\small\color{gray}Source module: \texttt{Kinetic\allowbreak{}Langevin/\allowbreak{}Geometry.\allowbreak{}lean}.\par}


The phase convention is \((v,x)\), velocity first. The paper's transformation is \(T_{\gamma,a}(v,x)=(v+(\gamma-a)x,-v-ax)\). Its squared energy is the sum of the two squared Hilbert norms. This explicit sum matters: Lean's default norm on a product is a maximum norm, not the Euclidean phase norm used in the paper. The identities below also apply with the underlying Hilbert space equal to an L\ensuremath{{}^{2}} space of random vectors.

\subsection{\texorpdfstring{Define the change of variables and its nonnegative energy}{Define the change of variables and its nonnegative energy}}


Set \(y=v+(\gamma-a)x\), \(z=-v-ax\), and \(\mathcal E=\|y\|^2+\|z\|^2\). These are the two coordinates used to expose cancellation in the kinetic drift. The definition \texttt{transform} constructs the pair; \texttt{transform\allowbreak{}Energy} explicitly sums their squares. Unfolding the latter leaves a sum of nonnegative terms, which \texttt{positivity} proves without any sign restriction on \(\gamma\) or \(a\). Thus the energy construction is valid before selecting the admissible friction regime.

\subsection{\texorpdfstring{Recover position and velocity}{Recover position and velocity}}


Adding the new coordinates gives \(y+z=(\gamma-2a)x\). The proof expands the scalar difference with \texttt{sub\_\allowbreak{}smul} and uses \texttt{module} to normalize the vector-space expression. If \(\gamma-2a\ne0\), multiply by its inverse to recover position. A second identity is \(ay+(\gamma-a)z=-(\gamma-2a)v\), giving \(v=-(\gamma-2a)^{-1}(ay+(\gamma-a)z)\). Lean first proves this vector identity, then rewrites nested scalar multiplication and cancels the nonzero denominator. The hypothesis excludes the singular coordinate choice \(a=\gamma/2\).

\subsection{\texorpdfstring{Bound the position projection}{Bound the position projection}}


The triangle inequality and \((\|y\|-\|z\|)^2\ge0\) give \(\|y+z\|^2\le2(\|y\|^2+\|z\|^2)\). Substituting \(y+z=(\gamma-2a)x\) yields \((\gamma-2a)^2\|x\|^2\le2\mathcal E\). In Lean, \texttt{norm\_\allowbreak{}add\_\allowbreak{}le} supplies the triangle bound; \texttt{nlinarith} combines it with nonnegative norms and the square of their difference. The final rewrite uses \texttt{norm\_\allowbreak{}smul} and \(|c|^2=c^2\), so the squared result does not need a sign assumption on the denominator.

\subsection{\texorpdfstring{Initialize a coupling with equal velocities}{Initialize a coupling with equal velocities}}


If the velocity difference is zero, the transformed pair is \(((\gamma-a)x,-ax)\). Its energy is therefore \(((\gamma-a)^2+a^2)\|x\|^2\). The proof unfolds both definitions, removes zero and negation, converts scalar multiplication norms to absolute scalar factors, and uses \texttt{sq\_\allowbreak{}abs} before \texttt{ring}. This is the exact initialization coefficient in the synchronous-position contraction theorem.

\subsection{\texorpdfstring{Read off velocity control and phase energy}{Read off velocity control and phase energy}}


At \(a=0\), the second transformed coordinate is \(-v\), so \(\|v\|^2\le\|v+\gamma x\|^2+\|v\|^2=\mathcal E\). Lean simplifies the zero scalar terms and invokes nonnegativity of the other square. \texttt{phase\allowbreak{}Energy} separately records \(\|v\|^2+\|x\|^2\), the ordinary squared Euclidean phase norm. Both definitions avoid relying on the product's maximum norm; only the transformed energy contains the mixed position–velocity geometry needed for contraction.

\subsection{\texorpdfstring{Use the actual Euclidean transformed norm}{Use the actual Euclidean transformed norm}}


Define \(\|s\|_T\) as the norm of \(T_{\gamma,a}s\) in \texttt{With\allowbreak{}Lp 2 (E \ensuremath{\times} E)}. This mathlib type equips the two components with the Euclidean product norm. Its square is exactly \texttt{transform\allowbreak{}Energy}, by \texttt{With\allowbreak{}Lp.\allowbreak{}prod\_\allowbreak{}norm\_\allowbreak{}sq\_\allowbreak{}eq\_\allowbreak{}of\_\allowbreak{}L2}. Thus the norm, squared-energy calculations, and later triangle inequalities all describe the same geometry. Nonnegativity comes directly from \texttt{norm\_\allowbreak{}nonneg}. The construction applies to a Hilbert space of vectors or a Hilbert space of square-integrable random vectors.

\subsection{\texorpdfstring{Prove the triangle and velocity-\allowbreak{}projection bounds}{Prove the triangle and velocity-projection bounds}}


The transformation is additive: expand the two coordinates and let \texttt{module} normalize the vector identities. After this identity, \texttt{With\allowbreak{}Lp.\allowbreak{}to\allowbreak{}Lp\_\allowbreak{}add} turns the transformed sum into a sum in an actual normed space, so \texttt{norm\_\allowbreak{}add\_\allowbreak{}le} gives its triangle inequality. Rewriting \(s-t=(s-r)+(r-t)\) gives the three-point bound needed to insert an auxiliary exact path. To pass from an energy bound \(\mathcal E(s)\le B^2\) to \(\|s\|_T\le B\), use nonnegativity and \texttt{sq\_\allowbreak{}le\_\allowbreak{}sq\ensuremath{{}_{0}}}, then replace the norm square by energy. Applying the same argument to the already proved velocity-square projection gives \(\|v\|\le\|(v,x)\|_T\) at \(a=0\).

\subsection{\texorpdfstring{Project position only at the final endpoint}{Project position only at the final endpoint}}


The squared position bound is \(\gamma^2\|x\|^2\le2\|s\|_T^2\). Both sides are squares of nonnegative quantities, so \texttt{sq\_\allowbreak{}le\_\allowbreak{}sq\ensuremath{{}_{0}}} gives \(\gamma\|x\|\le\sqrt2\|s\|_T\). Division by positive friction yields the exact projection factor \(\sqrt2/\gamma\). If the velocity difference initially vanishes, the energy identity gives \(\|(0,x)\|_T^2=\gamma^2\|x\|^2\). Comparing squares with nonnegative \(\gamma\) proves \(\|(0,x)\|_T=\gamma\|x\|\). These two identities explain why the eventual position theorem loses only one factor \(\sqrt2\), while friction cancels out of its bias floor.

\section{\texorpdfstring{Spectral cancellation and pathwise exponential decay}{Spectral cancellation and pathwise exponential decay}}
\label{proof-module-KineticLangevin-Contraction-lean}
{\small\color{gray}Source module: \texttt{Kinetic\allowbreak{}Langevin/\allowbreak{}Contraction.\allowbreak{}lean}.\par}


This module proves a deterministic theorem for a kinetic difference equation \(v'=-\gamma v-H_tx\), \(x'=v\). The continuous linear operator \(H_t\) is allowed to vary with time but must be self-adjoint and satisfy quadratic-form bounds \(mI\preceq H_t\preceq MI\). The theorem derives energy decay from this differential equation. Constructing stochastic solutions whose difference has these properties and taking an infimum over probability couplings are separate layers.

\subsection{\texorpdfstring{Write the derivative and prove exact spectral cancellation}{Write the derivative and prove exact spectral cancellation}}


The transformed derivatives are \(y'=-av-Hx\), \(z'=(\gamma-a)v+Hx\). Thus \(\mathcal E'=2\langle y,y'\rangle+2\langle z,z'\rangle\), which \texttt{energy\allowbreak{}Derivative} records. Multiplying by \(d=\gamma-2a\) gives

\[d\mathcal E'=2\bigl(a^2\|y\|^2-\langle y,Hy\rangle+\langle z,Hz\rangle-(\gamma-a)^2\|z\|^2\bigr).\]

The proof expands linear-map applications, scalar multiplication and inner products. Self-adjointness rewrites \(\langle v,Hx\rangle\) as \(\langle x,Hv\rangle\), so cross terms cancel. \texttt{ring} verifies the remaining scalar polynomial identity. Clearing the denominator keeps the identity valid even at singular parameter values.

\subsection{\texorpdfstring{Use lower and upper curvature on different coordinates}{Use lower and upper curvature on different coordinates}}


Apply the lower bound \(\langle y,Hy\rangle\ge m\|y\|^2\) and upper bound \(\langle z,Hz\rangle\le M\|z\|^2\). Let \(k=\max\{a^2-m,M-(\gamma-a)^2\}\). Both coefficients are at most \(k\); multiply these scalar bounds by the nonnegative coordinate squares. The spectral identity then yields \(d\mathcal E'\le2k\mathcal E\). The explicit premise \(d>0\) permits division without reversing the inequality. Lean names each quadratic-form and maximum comparison, combines them with \texttt{nlinarith}, then uses \texttt{le\_\allowbreak{}div\_\allowbreak{}iff\ensuremath{{}_{0}}}. This proves the dissipativity estimate from spectral bounds rather than assuming contraction.

\subsection{\texorpdfstring{Specialize to the simple friction regime}{Specialize to the simple friction regime}}


At \(a=0\), the maximum becomes \(\max\{-m,M-\gamma^2\}\). The assumption \(m+M\le\gamma^2\) gives \(M-\gamma^2\le-m\), so the maximum equals \(-m\). With \(\gamma>0\), the differential bound becomes \(\mathcal E'\le-2m\mathcal E/\gamma\). Lean obtains the general bound, proves the scalar comparison by \texttt{linarith}, and simplifies the maximum. This is the energy exponent; taking a square root later gives the norm exponent \(-m/\gamma\).

\subsection{\texorpdfstring{Differentiate the actual transformed energy along a path}{Differentiate the actual transformed energy along a path}}


Assume \texttt{Has\allowbreak{}Deriv\allowbreak{}At} witnesses the two kinetic equations at time \(t\). The linear combinations \(v+(\gamma-a)x\) and \(-v-ax\) inherit derivatives by addition, negation and constant scalar multiplication. The vector normalization tactic \texttt{module} reduces them to \(-av-Hx\) and \((\gamma-a)v+Hx\). The Hilbert-space squared-norm chain rule gives derivative \(2\langle y,y'\rangle\), and adding the two rules produces \texttt{energy\allowbreak{}Derivative}. No derivative of \(H_t\) is needed: only the value at the current time enters the first-order kinetic equation.

\subsection{\texorpdfstring{Apply Gronwall with an explicitly proved derivative bound}{Apply Gronwall with an explicitly proved derivative bound}}


Define scalar functions \(F(t)=\mathcal E(v_t,x_t)\) and \(F'(t)=\texttt{energyDerivative}(H_t,v_t,x_t)\), with \(K=2k/(\gamma-2a)\). Continuous paths make \(F\) continuous on the closed horizon, and the preceding derivative theorem gives its derivative on the half-open interval. Spectral dissipativity proves \(F'\le KF+0\). Lean feeds these facts to Mathlib's right-slope Gronwall theorem and simplifies its zero-forcing formula, obtaining \(F(t)\le F(0)e^{Kt}\). The zero additive term is significant: this is exact synchronous-flow contraction rather than a numerical defect estimate.

\subsection{\texorpdfstring{Project the energy estimate to the position difference}{Project the energy estimate to the position difference}}


Suppose \(v_0=0\). Substitute the exact initial energy \(((\gamma-a)^2+a^2)\|x_0\|^2\) into the Gronwall estimate. The position projection theorem gives \((\gamma-2a)^2\|x_t\|^2\le2F(t)\). Combining the two inequalities yields the stated squared version of Proposition 1. The final \texttt{nlinarith} only assembles already proved bounds. Its square root, under \(\gamma-2a>0\), is the paper's prefactor \(\sqrt{2((\gamma-a)^2+a^2)}/(\gamma-2a)\). This theorem concerns paths satisfying the explicit ODE; it is not yet a probability-law mixing result.

\subsection{\texorpdfstring{Include the initial time through right derivatives}{Include the initial time through right derivatives}}


Finite-horizon integral solutions naturally provide a derivative within \([t,\infty)\), including at \(t=0\). The coordinate differentiation proof works unchanged with \texttt{Has\allowbreak{}Deriv\allowbreak{}Within\allowbreak{}At}: add the two right derivatives, normalize their vector coefficients, apply the squared-norm derivative rule, and sum. The resulting derivative is still \texttt{energy\allowbreak{}Derivative}. The right-slope Grönwall theorem therefore gives the same exponential energy bound under the spectral hypotheses, without imposing a differentiable extension of the stochastic paths to negative time.

\subsection{\texorpdfstring{Recover the full arbitrary-\allowbreak{}shift prefactor}{Recover the full arbitrary-shift prefactor}}


Write \(d=\gamma-2a>0\) and \(k=\max\{a^2-m,M-(\gamma-a)^2\}/d\). The energy estimate and equal initial velocities give \(d^2\|x_t\|^2\le2((\gamma-a)^2+a^2)e^{2kt}\|x_0\|^2\). Lean proves \((e^{kt})^2=e^{2kt}\) using \texttt{Real.\allowbreak{}exp\_\allowbreak{}add}, proves the square-root square identity from nonnegativity of its argument, and compares the squares of two nonnegative quantities. This first yields \(d\|x_t\|\le\sqrt{2((\gamma-a)^2+a^2)}e^{kt}\|x_0\|\). Dividing by the positive \(d\) gives precisely Proposition 1's prefactor and exponent.

\section{\texorpdfstring{Nonlinear kinetic contraction directly from strong convexity and smoothness}{Nonlinear kinetic contraction directly from strong convexity and smoothness}}
\label{proof-module-KineticLangevin-PotentialContraction-lean}
{\small\color{gray}Source module: \texttt{Kinetic\allowbreak{}Langevin/\allowbreak{}Potential\allowbreak{}Contraction.\allowbreak{}lean}.\par}


The spectral proof can be replaced by a sector inequality for an actual gradient difference. This avoids selecting a time-varying averaged Hessian and only needs a differentiable strongly convex potential with Lipschitz gradient. The module proves the simple \(\sqrt2e^{-mt/\gamma}\) pathwise position estimate from the concrete noise-cancelled equations. Stochastic solution construction and Wasserstein law comparison are separate steps.

\subsection{\texorpdfstring{Write the concrete energy derivative and the gradient sector inequality}{Write the concrete energy derivative and the gradient sector inequality}}


For velocity difference \(v\), position difference \(x\), and force difference \(g\), define \(D=2\langle v+\gamma x,-g\rangle+2\langle-v,\gamma v+g\rangle\). The strong-smooth interpolation inequality is \(mM\|x\|^2+\|g\|^2\le(m+M)\langle g,x\rangle\). Subtract the linear part \(mx\), expand its squared norm, and rearrange to obtain \(\|g-mx\|^2\le(M-m)\langle x,g-mx\rangle\). Lean performs this expansion with the Hilbert norm-square and inner-product identities, then uses \texttt{nlinarith}. This sector estimate contains the joint lower–upper curvature information needed for contraction.

\subsection{\texorpdfstring{Complete a square to prove dissipativity}{Complete a square to prove dissipativity}}


Put \(k=g-mx\). Expanding the transformed energy yields

\[\gamma D+2m\mathcal E
=2\bigl[-(\gamma^2-2m)\|v\|^2-2\gamma\langle v,k\rangle-\gamma^2\langle x,k\rangle\bigr].\]

When \(M=m\), the sector inequality forces \(k=0\); friction gives \(\gamma^2\ge2m\), proving nonpositivity. When \(M>m\), combine the sector estimate with \(\|\gamma k+(M-m)v\|^2\ge0\) and \((M-m)(\gamma^2-m-M)\|v\|^2\ge0\). These imply \(\gamma D+2m\mathcal E\le0\). Lean handles the equal-curvature case separately, expands the square in the strict case, and combines the scalar inequalities. Positive \(\gamma\) then permits division, giving \(D\le-2m\mathcal E/\gamma\).

\subsection{\texorpdfstring{Obtain the sector condition from the actual potential}{Obtain the sector condition from the actual potential}}


For \(g=\nabla f(x)-\nabla f(y)\) and displacement \(x-y\), the shared ULA library has already proved the strong-smooth interpolation inequality from differentiability, strong convexity and gradient Lipschitzness. \texttt{potential\_\allowbreak{}force\_\allowbreak{}dissipativity} applies that theorem, converts it to the sector condition, and feeds it into the completed-square estimate. Consequently no Hessian, sector bound, or contraction conclusion is assumed in this potential theorem. The explicit hypotheses retain \(m\le M\), \(\gamma>0\), and \(\gamma^2\ge m+M\).

\subsection{\texorpdfstring{Differentiate the noise-\allowbreak{}cancelled path energy}{Differentiate the noise-cancelled path energy}}


Assume \(v'=-\gamma v-g\) and \(x'=v\). Then \((v+\gamma x)'=-g\), while \((-v)'=\gamma v+g\). Lean obtains these two derivative formulas by addition, negation and constant scalar multiplication, normalizes them with \texttt{module}, and applies the squared-norm derivative rule to each coordinate. Their sum is precisely \texttt{force\allowbreak{}Energy\allowbreak{}Derivative}. The theorem assumes differentiability of the difference paths; it does not incorrectly differentiate an individual Brownian velocity path.

\subsection{\texorpdfstring{Integrate energy decay and take the position square root}{Integrate energy decay and take the position square root}}


For two continuous position paths \(x,y\), set \(F(t)=\mathcal E(v_t,x_t-y_t)\). The derivative theorem and potential dissipativity give \(F'\le-2mF/\gamma\). Closed-interval continuity and the explicit derivative witnesses allow the same Gronwall theorem as in the spectral proof, giving \(F(t)\le F(0)e^{-2mt/\gamma}\). If \(v_0=0\), then \(F(0)=\gamma^2\|x_0-y_0\|^2\). Combine with the position projection bound, cancel the positive \(\gamma^2\), and compare nonnegative squares. The exponential square identity \((e^{-mt/\gamma})^2=e^{-2mt/\gamma}\) yields \(\|x_t-y_t\|\le\sqrt2e^{-mt/\gamma}\|x_0-y_0\|\). This is a proved nonlinear pathwise estimate with the paper's correct prefactor.

\subsection{\texorpdfstring{Use right derivatives at the initial time}{Use right derivatives at the initial time}}


A finite-horizon integral solution is constrained only for nonnegative time. Its extension to negative time need not have the same derivative. The correct initial derivative is therefore the derivative within the right half-line. The product and squared-norm derivative rules remain valid in this form, so differentiating the transformed energy gives the same force-energy expression. Lean applies those within-derivative rules to both transformed coordinates, cancels the opposing friction terms, and transfers the resulting derivative to the energy expression. The right-slope form of Gronwall requires exactly these hypotheses, and gives the same exponential energy decay throughout the closed interval. This is the version used by the actual integral-flow theorem.

\section{\texorpdfstring{Arbitrary-\allowbreak{}friction contraction directly from the potential}{Arbitrary-friction contraction directly from the potential}}
\label{proof-module-KineticLangevin-GeneralContraction-lean}
{\small\color{gray}Source module: \texttt{Kinetic\allowbreak{}Langevin/\allowbreak{}General\allowbreak{}Contraction.\allowbreak{}lean}.\par}


The paper's full shifted contraction can be derived from differentiability, strong convexity and a Lipschitz gradient alone. This proof centers velocity at \(v+\gamma x/2\), uses the strong-smooth sector inequality, and completes a scalar square. It removes any need to assume that the Hessian is continuous when constructing a mean Hessian, while recovering the exact maximum of the two endpoint curvature exponents.

\subsection{\texorpdfstring{Center the phase energy and its derivative}{Center the phase energy and its derivative}}


The concrete force-difference derivative is \(2\langle v+(\gamma-a)x,-av-F\rangle+2\langle-v-ax,(\gamma-a)v+F\rangle\). Put \(y=v+\gamma x/2\), \(d=\gamma-2a\), and \(c=((\gamma-a)^2+a^2)/2\). Expanding the inner products gives the two exact identities

\[\mathcal E=2\|y\|^2+\tfrac12d^2\|x\|^2,\qquad \mathcal E'=-\gamma\mathcal E+4\langle y,cx-F\rangle.\]

The Lean proof rewrites squared norms as self-inner-products, distributes inner products over sums, differences and scalar multiplication, uses symmetry, and finishes with \texttt{ring}. These identities expose a scalar residual rather than a matrix diagonalization.

\subsection{\texorpdfstring{Bound the residual around an arbitrary scalar curvature}{Bound the residual around an arbitrary scalar curvature}}


Strong-smooth interpolation places \(F\) in the ball centered at \(\mu x\), where \(\mu=(m+M)/2\), with radius \(R\|x\|\), where \(R=(M-m)/2\). Decompose \(cx-F=(c-\mu)x-(F-\mu x)\). The triangle inequality gives \(\|cx-F\|\le(|c-\mu|+R)\|x\|\). The coefficient is exactly \(\max\{c-m,M-c\}\): Lean splits on \(c\ge\mu\), simplifies the absolute value and the corresponding maximum, and checks the remaining arithmetic. No gradient or Hessian operator is invented in this argument.

\subsection{\texorpdfstring{Complete the centered square to obtain the exact exponent}{Complete the centered square to obtain the exact exponent}}


Let \(K=\max\{c-m,M-c\}\ge0\). Cauchy–Schwarz and the residual bound give \(\langle y,cx-F\rangle\le K\|y\|\|x\|\). The nonnegative square \((2\|y\|-d\|x\|)^2\) implies \(4d\|y\|\|x\|\le2\mathcal E\). Multiplying by \(K\) and inserting the exact energy-derivative identity yields \(d\mathcal E'\le(-\gamma d+2K)\mathcal E\). Finally \(K=\max\{a^2-m,M-(\gamma-a)^2\}+\gamma d/2\), so the \(\gamma d\) terms cancel. Lean records the multiplied identities explicitly before \texttt{nlinarith}, then divides by the positive \(d\). This proves the precise spectral differential estimate even in a friction regime with a positive growth exponent.

\subsection{\texorpdfstring{Supply the sector inequality from the actual potential}{Supply the sector inequality from the actual potential}}


For \(x-y\) and \(F=\nabla f(x)-\nabla f(y)\), the project's proved \texttt{strong\_\allowbreak{}smooth\_\allowbreak{}gradient\_\allowbreak{}interpolation} gives the strong-smooth interpolation inequality from \texttt{Differentiable}, \texttt{Strong\allowbreak{}Convex\allowbreak{}On} and \texttt{Lipschitz\allowbreak{}With}. Its centered-ball consequence is exactly the residual theorem's premise. The proof composes these already proved results, yielding arbitrary-shift dissipativity for the actual nonlinear gradient. It assumes neither a Hessian representation nor the desired contraction conclusion.

\subsection{\texorpdfstring{Integrate the right-\allowbreak{}derivative inequality}{Integrate the right-derivative inequality}}


The two kinetic difference equations imply the transformed-coordinate right derivatives \(-av-F\) and \((\gamma-a)v+F\). Their squared-norm derivative rules add to \texttt{general\allowbreak{}Force\allowbreak{}Energy\allowbreak{}Derivative}. For paths on \([0,T]\), define their transformed energy \(E(t)\). Continuity follows from continuity of the paths, while the preceding nonlinear result gives \(E'(t)\le2kE(t)\), with \(k=\max\{a^2-m,M-(\gamma-a)^2\}/(\gamma-2a)\). The right-slope form of Grönwall then proves \(E(t)\le E(0)e^{2kt}\), including the initial endpoint. The hypotheses are the actual force-difference differential equations, not an exponential bound supplied by the caller.

\section{\texorpdfstring{Contraction of actual shared-\allowbreak{}input kinetic integral solutions}{Contraction of actual shared-input kinetic integral solutions}}
\label{proof-module-KineticLangevin-SynchronousFlow-lean}
{\small\color{gray}Source module: \texttt{Kinetic\allowbreak{}Langevin/\allowbreak{}Synchronous\allowbreak{}Flow.\allowbreak{}lean}.\par}


This joins the Picard integral construction to the transformed-energy argument. Two actual solutions with the same continuous forcing have a noise-free difference equation as a consequence of their integral identities. The energy estimate and the paper's position decay then follow without assuming either a contraction result or a differentiable Brownian path.

\subsection{\texorpdfstring{Derive the phase difference equation and contract its energy}{Derive the phase difference equation and contract its energy}}


Apply the shared-input derivative theorem to the concrete kinetic drift \(G(v,x)=(\gamma v+\nabla f(x),-v)\). Project its right derivative to velocity and position using the two continuous linear projection maps. The velocity derivative becomes \(-\gamma(V-V\prime)-[\nabla f(X)-\nabla f(X\prime)]\); the position derivative is \(V-V\prime\). Scalar multiplication and additive rearrangement justify both identities. The right-derivative energy theorem, which is valid at the initial endpoint, yields

\[\mathcal E_t\le e^{-2mt/\gamma}\mathcal E_0,\quad\mathcal E_t=\|V_t-V\prime_t+\gamma(X_t-X\prime_t)\|^2+\|V_t-V\prime_t\|^2.\]

In Lean the required Lipschitz constant for \(G\) is derived from the actual gradient bound; it is not an extra Hessian or solution assumption. The argument then supplies the original potential's strong convexity and smoothness hypotheses to the proved energy inequality.

\subsection{\texorpdfstring{Equal initial velocities give the sharp position prefactor}{Equal initial velocities give the sharp position prefactor}}


If the initial velocities agree, \(\mathcal E_0=\gamma^2\|X_0-X\prime_0\|^2\). The projection estimate \(\gamma^2\|X_t-X\prime_t\|^2\le2\mathcal E_t\) cancels the factor \(\gamma^2\), giving the squared position bound. Positivity of \(\gamma\), the norm, the exponential and the square root permits taking square roots. Lean explicitly proves \(e^{-mt/\gamma\,\cdot2}=(e^{-mt/\gamma})^2\), rewrites \((\sqrt2)^2=2\), and applies the nonnegative square comparison theorem. The resulting prefactor is \(\sqrt2\), consistent with the explicit general prefactor in the supplied source.

\subsection{\texorpdfstring{Contract the full transformed norm}{Contract the full transformed norm}}


The actual integral-flow energy theorem gives \(\mathcal E_t\le e^{-2mt/\gamma}\mathcal E_0\). Replacing energy by the square of \texttt{transform\allowbreak{}Norm} and using \((e^{-mt/\gamma})^2=e^{-2mt/\gamma}\) gives \(\|Y_t-Z_t\|_T\le e^{-mt/\gamma}\|Y_0-Z_0\|_T\). Lean proves the exponential square identity by \texttt{Real.\allowbreak{}exp\_\allowbreak{}add}, notes that both norms and the exponential are nonnegative, and invokes the energy-to-norm lemma. This stronger phase estimate controls velocity differences and accumulates discretization errors without a separate projection loss at each step.

\subsection{\texorpdfstring{Apply the full shifted theorem to actual integral flows}{Apply the full shifted theorem to actual integral flows}}


Take two actual kinetic integral solutions with the same additive input. Their noise-cancelled phase difference has a right derivative by \texttt{integral\_\allowbreak{}difference\_\allowbreak{}has\allowbreak{}Deriv\allowbreak{}Within\allowbreak{}At}. Compose this derivative with the continuous linear velocity and position projections; \texttt{module} normalizes their drift formulas. The centered-sector theorem then supplies the full arbitrary-shift energy exponent from the original potential assumptions. Equal initial velocities identify the exact initial energy, and the position-square projection supplies the final prefactor. Squaring the desired nonnegative quantities, using the exponential-addition identity, and dividing by \(\gamma-2a>0\) proves

\[\|X_t-Y_t\|\le\frac{\sqrt{2((\gamma-a)^2+a^2)}}{\gamma-2a}\exp\!\left(\frac{\max\{a^2-m,M-(\gamma-a)^2\}}{\gamma-2a}t\right)\|X_0-Y_0\|.\]

No continuous-Hessian assumption or difference-ODE assumption remains; both the ODE and the coefficient are derived from the concrete integral paths.

\section{\texorpdfstring{A noise-\allowbreak{}driven recursion has a law independent of its realization}{A noise-driven recursion has a law independent of its realization}}
\label{proof-module-UnadjustedLangevin-ChainLaw-lean}
{\small\color{gray}Source module: \texttt{Unadjusted\allowbreak{}Langevin/\allowbreak{}Chain\allowbreak{}Law.\allowbreak{}lean}.\par}


For measurable \(F(x,z)\), the recursion \(X_{n+1}=F(X_n,Z_n)\) has a law determined by the initial law and the full noise law when the initialization and noise are independent.

\subsection{\texorpdfstring{noise\allowbreak{}Driven\allowbreak{}Chain}{noiseDrivenChain}}


Define \(X_0=x\), \(X_{n+1}=F(X_n,w_n)\) for any update and noise path. No distribution assumptions are part of the recursion; they enter only when a law is assigned to its inputs.

\subsection{\texorpdfstring{noise\allowbreak{}Driven\allowbreak{}Chain zero}{noiseDrivenChain zero}}


The zero clause of the recursion evaluates to the supplied initial state. This definitional identity is the base case for all law and measurability inductions.

\subsection{\texorpdfstring{noise\allowbreak{}Driven\allowbreak{}Chain succ}{noiseDrivenChain succ}}


The successor clause unfolds to \(X_{n+1}=F(X_n,w_n)\). This exact recursion lets the generic construction later be identified with ULA by induction.

\subsection{\texorpdfstring{noise\allowbreak{}Driven\allowbreak{}Chain measurable}{noiseDrivenChain measurable}}


Induct on \(n\). The base map is the initial-state projection. At the successor, pair the measurable current state with the measurable coordinate evaluation \(w\mapsto w_n\), then compose with the measurable uncurried update. Thus each iterate is a measurable function of initialization and the full noise path.

\subsection{\texorpdfstring{chain\allowbreak{}Law}{chainLaw}}


Push \(\mu\otimes Q\) through the measurable recursion output at index \(n\). The product gives independent initialization and noise path. This canonical law depends only on the initial law, noise-path law, and update, not on a particular sample-space implementation.

\subsection{\texorpdfstring{chain\allowbreak{}Law zero}{chainLaw zero}}


At index zero the output map is first-coordinate projection. Since \(Q\) has mass one, the first marginal of \(\mu\otimes Q\) is exactly \(\mu\). This proves that the canonical law starts from the supplied initial law.

\subsection{\texorpdfstring{noise\allowbreak{}Driven\allowbreak{}Chain law of indep}{noiseDrivenChain law of indep}}


Independence gives the joint input law \(P_{(X,W)}=P_X\otimes P_W\). Push that identity through the measurable recursion output and use pushforward composition. The resulting output law is exactly \texttt{chain\allowbreak{}Law F (P.\allowbreak{}map X) (P.\allowbreak{}map W) n}.

\subsection{\texorpdfstring{noise\allowbreak{}Driven\allowbreak{}Chain law of iid}{noiseDrivenChain law of iid}}


Joint independence of the innovations and their common law \(\eta\) identify the entire noise-path law with \(\bigotimes_n\eta\). Substitute this identity into the independent-initialization law theorem. Countability of the time index supplies path-valued measurability.

\subsection{\texorpdfstring{noise\allowbreak{}Driven\allowbreak{}Chain law unique}{noiseDrivenChain law unique}}


Apply the canonical iid-law formula on both probability spaces. Equality of their initial laws and the shared innovation law makes the right sides identical, so their iterate laws agree. The theorem compares distributions without coupling the two implementations.

\subsection{\texorpdfstring{noise\allowbreak{}Driven\allowbreak{}Chain eq ula\allowbreak{}Chain}{noiseDrivenChain eq ulaChain}}


Induct on the index with \(F=\operatorname{ulaUpdate}\). The zero values agree, and the successor equations apply the same update to equal previous states and the same noise vector. This proves pointwise equality of the generic recursion and the repository's ULA chain.

\subsection{\texorpdfstring{ula\allowbreak{}Update measurable}{ulaUpdate measurable}}


Lipschitzness makes \(\nabla U\) continuous. The joint ULA update is obtained from coordinate projections and that gradient by scalar multiplication and addition, hence is measurable. This discharges the generic recursion's update-measurability premise.

\subsection{\texorpdfstring{ula\allowbreak{}Chain law of iid}{ulaChain law of iid}}


Use measurability of \texttt{ula\allowbreak{}Update} in the generic iid law theorem, then rewrite the generic recursion using \texttt{noise\allowbreak{}Driven\allowbreak{}Chain\_\allowbreak{}eq\_\allowbreak{}ula\allowbreak{}Chain}. The result identifies the actual ULA implementation with its canonical product-noise law.

\subsection{\texorpdfstring{lp\allowbreak{}Ula\allowbreak{}Chain law of iid}{lpUlaChain law of iid}}


The \(L^2\) chain representative agrees almost everywhere with the pointwise ULA recursion of the input representatives. Apply the pointwise law theorem and transfer pushforwards across that almost-everywhere equality. Quotient representative choices therefore do not change the algorithm law.

\section{\texorpdfstring{Couplings and quadratic transport}{Couplings and quadratic transport}}
\label{proof-module-UnadjustedLangevin-Coupling-lean}
{\small\color{gray}Source module: \texttt{Unadjusted\allowbreak{}Langevin/\allowbreak{}Coupling.\allowbreak{}lean}.\par}


Let \(H\) be a normed additive group with its Borel measurable structure and a second countable topology. Let \(P\) be a measure on a measurable space \(\Omega\), and write \(\mathcal L_P(X)=X_\#P\). The symbol \(\|Z\|_{2,P}\) denotes the extended \(L^2\) seminorm, which may be infinite. In the probability setting the following definitions give the usual quadratic transport cost.



\texttt{Coupling} stores a joint measure with exact first and second pushforward marginals. \texttt{coupling\allowbreak{}Cost} is its extended \(L^2\) displacement norm, and \texttt{wasserstein2} is the infimum over all such couplings. Finiteness guards precede \texttt{to\allowbreak{}Real} conversions; no existence of an optimal coupling is assumed.

\subsection{\texorpdfstring{Definition of a coupling}{Definition of a coupling}}


A coupling of measures \(\mu,\nu\) on \(H\) is a joint measure \(\gamma\) on \(H\times H\) with those two marginals. Denote the family by \(\Gamma(\mu,\nu)\).



Choose the joint measure \(\gamma\) on pairs \((x,y)\).



Require \((\operatorname{pr}_1)_\#\gamma=\mu\) and \((\operatorname{pr}_2)_\#\gamma=\nu\).

\subsection{\texorpdfstring{Definition of quadratic coupling cost}{Definition of quadratic coupling cost}}


For a joint measure \(\gamma\), define

\[c(\gamma)=\left(\int\|x-y\|^2\,d\gamma(x,y)\right)^{1/2}\in[0,+\infty].\]

Take the extended \(L^2(\gamma)\) norm of the displacement map \((x,y)\mapsto x-y\).

\subsection{\texorpdfstring{Definition of the transport infimum}{Definition of the transport infimum}}


Define the extended quadratic transport distance by

\[W_2(\mu,\nu)=\inf_{\gamma\in\Gamma(\mu,\nu)}c(\gamma).\]

Take the infimum over all couplings without selecting a minimizer.

\subsection{\texorpdfstring{Real conversion of transport distance}{Real conversion of transport distance}}


Define \(W_2^{\mathbb R}(\mu,\nu)=\operatorname{toReal}(W_2(\mu,\nu))\). The comparison lemmas below establish finite upper bounds before using this conversion.



Apply real conversion to the extended distance; for a finite distance this returns its ordinary real value.

\subsection{\texorpdfstring{Every coupling supplies an upper bound}{Every coupling supplies an upper bound}}


For every \(\gamma\in\Gamma(\mu,\nu)\), \(W_2(\mu,\nu)\le c(\gamma)\).



An infimum is at most each member of the family. Apply this to the chosen coupling.

\subsection{\texorpdfstring{The joint law of two random vectors}{The joint law of two random vectors}}


If \(X,Y:\Omega\to H\) are measurable, then \((X,Y)_\#P\) is a coupling of \(X_\#P\) and \(Y_\#P\).



Set \(\gamma=(\omega\mapsto(X(\omega),Y(\omega)))_\#P\).



Compose the pushforward with the first projection. Since \(\operatorname{pr}_1(X(\omega),Y(\omega))=X(\omega)\), the first marginal is \(X_\#P\).



The same calculation with the second projection gives the second marginal \(Y_\#P\).

\subsection{\texorpdfstring{Joint laws from almost everywhere measurable vectors}{Joint laws from almost everywhere measurable vectors}}


If \(X,Y\) are measurable almost everywhere with respect to \(P\), their joint law is still a coupling of their two laws.



Again define \(\gamma=(X,Y)_\#P\).



Apply pushforward composition for almost everywhere measurable maps to identify the first marginal with \(X_\#P\).



Apply its second-projection version to identify the second marginal with \(Y_\#P\).

\subsection{\texorpdfstring{Cost of the joint-\allowbreak{}law coupling}{Cost of the joint-law coupling}}


For measurable \(X,Y:\Omega\to H\),

\[c((X,Y)_\#P)=\|X-Y\|_{2,P}.\]

Expand the two definitions. The cost is the \(L^2\) norm of \((x,y)\mapsto x-y\) under the joint law.



Change variables through the pushforward. The displacement composed with the paired random vector is \(\omega\mapsto X(\omega)-Y(\omega)\), giving the identity.

\subsection{\texorpdfstring{A coupled pair bounds transport}{A coupled pair bounds transport}}


For measurable \(X,Y\),

\[W_2(X_\#P,Y_\#P)\le\|X-Y\|_{2,P}.\]

Use the joint law as an admissible coupling, then replace its cost by the \(L^2\) distance just computed.

\subsection{\texorpdfstring{The almost everywhere measurable transport bound}{The almost everywhere measurable transport bound}}


The same transport comparison holds for almost everywhere measurable \(X,Y\):

\[W_2(X_\#P,Y_\#P)\le\|X-Y\|_{2,P}.\]

The almost everywhere measurable joint-law construction supplies an admissible coupling, whose cost bounds the infimum.



The pushforward change of variables identifies that cost with \(\|X-Y\|_{2,P}\).

\subsection{\texorpdfstring{Finite coupled cost gives a real bound}{Finite coupled cost gives a real bound}}


If \(X,Y\) are measurable and \(X-Y\in L^2(P;H)\), then

\[W_2^{\mathbb R}(X_\#P,Y_\#P)\le\operatorname{toReal}(\|X-Y\|_{2,P}).\]

The extended comparison has a finite upper bound because \(X-Y\in L^2\). Apply monotonicity of real conversion under a finite upper bound.

\subsection{\texorpdfstring{The transport bound for L\ensuremath{{}^{2}} classes}{The transport bound for L² classes}}


For \(X,Y\in L^2(P;H)\), their representatives define laws satisfying

\[W_2(X_\#P,Y_\#P)\le\operatorname{ofReal}(\|X-Y\|_{L^2}).\]

The representative of \(X-Y\) agrees almost everywhere with the difference of the representatives. Their extended \(L^2\) norms agree.



The L\ensuremath{{}^{2}} norm is the real value of the finite extended seminorm. Embedding it back into extended reals recovers that seminorm.



Apply the almost everywhere measurable transport bound to the two representatives and use the equality of displacement seminorms.

\subsection{\texorpdfstring{The real L\ensuremath{{}^{2}}-\allowbreak{}to-\allowbreak{}transport bridge}{The real L²-to-transport bridge}}


For \(X,Y\in L^2(P;H)\),

\[W_2^{\mathbb R}(X_\#P,Y_\#P)\le\|X-Y\|_{L^2}.\]

The preceding extended bound is dominated by the finite number \(\operatorname{ofReal}(\|X-Y\|_{L^2})\), so convert monotonically to real numbers.



Real conversion of the embedding of a nonnegative norm returns the norm itself.

\subsection{\texorpdfstring{Pointwise domination passes to L\ensuremath{{}^{2}}}{Pointwise domination passes to L²}}


Let \(q\ge0\), let \(A\) be almost everywhere strongly measurable, and suppose \(\|A(\omega)\|\le q\|D(\omega)\|\) almost everywhere. Then

\[\|A\|_{2,P}\le q\|D\|_{2,P}.\]

Apply monotonicity and positive homogeneity of the extended \(L^2\) seminorm to the almost everywhere domination. No independence is required.

\subsection{\texorpdfstring{One-\allowbreak{}step error from contraction and a defect}{One-step error from contraction and a defect}}


Suppose \(q\ge0\), \(A\) is almost everywhere strongly measurable, and almost everywhere

\[X'-Y'=A-V,\qquad\|A\|\le q\|X-Y\|.\]

Then

\[\|X'-Y'\|_{2,P}\le q\|X-Y\|_{2,P}+\|V\|_{2,P}.\]

Replace \(X'-Y'\) by \(A-V\) in the seminorm using the almost everywhere error identity.



Minkowski gives \(\|A-V\|_2\le\|A\|_2+\|V\|_2\). Pointwise contraction bounds the first term by \(q\|X-Y\|_2\).

\subsection{\texorpdfstring{One-\allowbreak{}step transport estimate}{One-step transport estimate}}


Under the preceding decomposition and contraction hypotheses, if \(X',Y'\) are measurable, then

\[W_2(X'_\#P,Y'_\#P)\le q\|X-Y\|_{2,P}+\|V\|_{2,P}.\]

The joint-law coupling bounds transport by \(\|X'-Y'\|_2\). Compose this with the synchronous \(L^2\) estimate.

\subsection{\texorpdfstring{The finite-\allowbreak{}moment one-\allowbreak{}step inequality}{The finite-moment one-step inequality}}


Assume \(q\ge0\), \(A\) is almost everywhere strongly measurable, \(X'-Y'=A-V\) and \(\|A\|\le q\|X-Y\|\) almost everywhere. If \(X-Y,V\in L^2(P;H)\), then

\[\operatorname{toReal}(\|X'-Y'\|_{2,P})\le q\operatorname{toReal}(\|X-Y\|_{2,P})+\operatorname{toReal}(\|V\|_{2,P}).\]

The product \(q\|X-Y\|_2\) is finite because both factors are finite.



Its sum with the finite defect norm is finite. Use this finite upper bound to convert the synchronous extended inequality monotonically to real numbers.



Distribute real conversion through the finite sum and the product, giving the stated real inequality.

\section{\texorpdfstring{Passing from a chosen initial coupling to an infimum}{Passing from a chosen initial coupling to an infimum}}
\label{proof-module-UnadjustedLangevin-InitialCoupling-lean}
{\small\color{gray}Source module: \texttt{Unadjusted\allowbreak{}Langevin/\allowbreak{}Initial\allowbreak{}Coupling.\allowbreak{}lean}.\par}


The analytic estimate first uses the cost of one initial coupled pair. The infimum calculation here replaces that cost by \(W_2(\mu,\nu)\) when all finite-cost couplings give the same output laws. \texttt{Product\allowbreak{}Reference.\allowbreak{}lean} proves this common-law property on concrete product extensions; \texttt{Canonical\allowbreak{}ULA.\allowbreak{}lean} supplies their stationary references.



The mathematical infimum here is the infimum in \texttt{wasserstein2}. The proof works with arbitrary finite-cost couplings and approximates that infimum when necessary. This gives the true initial transport term without postulating an optimal transport plan.

\subsection{\texorpdfstring{A uniform affine estimate passes to the infimum}{A uniform affine estimate passes to the infimum}}


Let \(a\ge0\) be finite, \(b\in[0,+\infty]\), and \(\Gamma(\mu,\nu)\ne\varnothing\). If the fixed output measures satisfy

\[W_2(\mu',\nu')\le ac(\gamma)+b\quad\text{for every }\gamma\in\Gamma(\mu,\nu),\]

then \(W_2(\mu',\nu')\le aW_2(\mu,\nu)+b\).



Use the nonempty coupling family and expand transport distances as infima.



Multiplication by finite \(a\ge0\) and addition of \(b\) commute with the nonempty extended-real infimum, including \(a=0\). Thus \(aW_2(\mu,\nu)+b=\inf_\gamma(ac(\gamma)+b)\).



The output distance lies below every member of this last family, hence below its infimum.

\subsection{\texorpdfstring{A finite infimum supplies a finite-\allowbreak{}cost coupling}{A finite infimum supplies a finite-cost coupling}}


If \(W_2(\mu,\nu)<+\infty\), then some \(\gamma\in\Gamma(\mu,\nu)\) has \(c(\gamma)<+\infty\). Attainment of the infimum is not asserted.



The coupling-cost infimum is strictly below \(+\infty\) by the finiteness hypothesis.



An infimum strictly below a threshold has a member strictly below that threshold. Apply this with threshold \(+\infty\).

\subsection{\texorpdfstring{The real uniform-\allowbreak{}coupling estimate}{The real uniform-coupling estimate}}


Let \(a,b\ge0\) be finite and assume \(W_2(\mu,\nu)<+\infty\). If \(W_2(\mu',\nu')\le ac(\gamma)+b\) for every initial coupling, then

\[W_2^{\mathbb R}(\mu',\nu')\le aW_2^{\mathbb R}(\mu,\nu)+b.\]

A finite initial distance supplies a coupling. Apply the extended uniform-coupling lemma to obtain the desired inequality in extended reals.



The product \(aW_2(\mu,\nu)\) is finite.



Its sum with finite \(b\) is finite, so the extended output inequality permits monotone real conversion.



Distribute real conversion through the finite sum and product to obtain the real inequality.

\subsection{\texorpdfstring{Finite-\allowbreak{}cost initial couplings suffice}{Finite-cost initial couplings suffice}}


Let \(a,b\ge0\) be finite. Assume both the initial and output transport distances are finite and that every finite-cost \(\gamma\in\Gamma(\mu,\nu)\) satisfies

\[W_2^{\mathbb R}(\mu',\nu')\le a\operatorname{toReal}(c(\gamma))+b.\]

Then \(W_2^{\mathbb R}(\mu',\nu')\le aW_2^{\mathbb R}(\mu,\nu)+b\).



For each finite-cost coupling, embed the real estimate in extended nonnegative reals. Finiteness of the output distance and coupling cost recovers their original extended values, giving \(W_2(\mu',\nu')\le ac(\gamma)+b\).



Apply the all-couplings theorem. For a finite-cost coupling, the needed extended estimate has just been proved.



For an infinite-cost coupling, split into the cases \(a=0\) and \(a\ne0\).



If \(a=0\), choose one finite-cost coupling supplied by the finite initial distance. Its estimate gives \(W_2(\mu',\nu')\le b\), which also supplies the required zero-coefficient estimate for the infinite-cost coupling.



If \(a\ne0\), then \(a>0\), so \(a\cdot(+\infty)+b=+\infty\) and the required inequality is automatic.

\section{\texorpdfstring{Finite quadratic transport transfers second moments}{Finite quadratic transport transfers second moments}}
\label{proof-module-UnadjustedLangevin-TransportMoments-lean}
{\small\color{gray}Source module: \texttt{Unadjusted\allowbreak{}Langevin/\allowbreak{}Transport\allowbreak{}Moments.\allowbreak{}lean}.\par}


Finite \(W_2\) distance to a finite-second-moment law supplies the other second moment automatically. This removes a redundant integrability premise when passing an initial law to the L\ensuremath{{}^{2}} formulation.

\subsection{\texorpdfstring{mem\allowbreak{}Lp id of wasserstein2 ne top}{memLp id of wasserstein2 ne top}}


Finite Wasserstein distance supplies a finite-cost coupling. On that coupling space the second coordinate is in \(L^2\) by the reference moment, and the coordinate difference is in \(L^2\) by finite coupling cost. Their sum is the first coordinate, so its law \(\mu\) has a finite second moment. No optimal coupling is selected.

\subsection{\texorpdfstring{wasserstein2 ne top of mem\allowbreak{}Lp}{wasserstein2 ne top of memLp}}


Use the explicit independent coupling \(\mu\otimes\pi\). Each coordinate is in \(L^2\), hence their difference is in \(L^2\) and has finite cost. Wasserstein distance is at most this coupling cost, proving finiteness.

\subsection{\texorpdfstring{mem\allowbreak{}Lp of law of wasserstein2 ne top}{memLp of law of wasserstein2 ne top}}


The finite-transport theorem first proves \(\mathrm{id}\in L^2(\mu)\). Replace \(\mu\) by the supplied pushforward law of \(X\), then use the change-of-law \(L^2\) criterion to pull this moment back to \(X\) on its original sample space.

\subsection{\texorpdfstring{mem\allowbreak{}Lp of wasserstein2 map ne top}{memLp of wasserstein2 map ne top}}


Specialize the preceding theorem to \(\mu=P\circ X^{-1}\), whose law identity is reflexive. Finite transport of the actual pushforward to a finite-moment reference therefore gives \(X\in L^2(P)\).

\section{\texorpdfstring{Realize every initial coupling independently of the noise}{Realize every initial coupling independently of the noise}}
\label{proof-module-UnadjustedLangevin-InitialExtension-lean}
{\small\color{gray}Source module: \texttt{Unadjusted\allowbreak{}Langevin/\allowbreak{}Initial\allowbreak{}Extension.\allowbreak{}lean}.\par}


Given any coupling \(\gamma\) of \(\mu\) and \(\pi\), extend it to \(\gamma\otimes Q\). This preserves its quadratic cost and supplies independent noise without requiring an optimal coupling.

\subsection{\texorpdfstring{Coupling.\allowbreak{}is\allowbreak{}Probability\allowbreak{}Measure}{Coupling.isProbabilityMeasure}}


The coupling's first marginal is the initial probability. Pushforward preserves total mass, so \(c(H\times H)=\mu(H)=1\). This constructs normalization of the joint measure using a marginal identity already stored in \texttt{Coupling}.

\subsection{\texorpdfstring{initial\allowbreak{}Extension}{initialExtension}}


Extend an initial coupling \(c(dx,dy)\) by an independent noise sample \(s\sim Q\), using \(c\otimes Q\) on \((H\times H)\times S\). Both algorithms retain the prescribed initial joint law while sharing fresh noise independent of the pair.

\subsection{\texorpdfstring{initial\allowbreak{}Extension is\allowbreak{}Probability\allowbreak{}Measure}{initialExtension isProbabilityMeasure}}


The coupling is a probability by its first marginal, and \(Q\) is a probability by hypothesis. Their product therefore has total mass one. This instance supports all subsequent \(L^2\) constructions on the extension.

\subsection{\texorpdfstring{initial\allowbreak{}Extension joint}{initialExtension joint}}


Projecting \(c\otimes Q\) to the first factor integrates out the probability noise coordinate. Since \(Q(S)=1\), the resulting measure is exactly \(c\), preserving the entire initial dependence structure.

\subsection{\texorpdfstring{initial\allowbreak{}Extension initial}{initialExtension initial}}


Compose the joint projection with its first state projection. The extension's joint marginal is \(c\), whose first state marginal is \(\mu\). Pushforward composition therefore yields the specified algorithm initialization law.

\subsection{\texorpdfstring{initial\allowbreak{}Extension reference}{initialExtension reference}}


Compose projection to the initial pair with its second state projection. The coupling's second marginal is \(\pi\), so the reference initialization on the extension has exactly law \(\pi\).

\subsection{\texorpdfstring{initial\allowbreak{}Extension noise}{initialExtension noise}}


The coupling measure has mass one. Projecting its product with \(Q\) to the noise factor therefore returns \(Q\). The \texttt{SFinite} hypothesis supplies the product marginal theorem's measure-theoretic condition.

\subsection{\texorpdfstring{initial\allowbreak{}Extension independent}{initialExtension independent}}


The two coordinate projections of a product probability measure are independent. Here those coordinates are the full coupled initial pair and the fresh noise sample. Independence of either individual initial state from any noise-derived random object follows by composition.

\subsection{\texorpdfstring{initial\allowbreak{}Extension initial noise}{initialExtension initial noise}}


Map \(((x,y),s)\) to \((x,s)\). Product pushforward and the first marginal identity of \(c\) give \(\mu\otimes Q\). Thus removing the reference coordinate leaves the same algorithm input law for every initial coupling.

\subsection{\texorpdfstring{initial\allowbreak{}Extension map}{initialExtension map}}


Factor the map through \(((x,y),s)\mapsto(x,s)\), whose law is \(\mu\otimes Q\). Pushforward composition proves that every measurable output depending only on algorithm initialization and noise has the coupling-independent law \(f_*(\mu\otimes Q)\).

\subsection{\texorpdfstring{initial\allowbreak{}Extension chain\allowbreak{}Law}{initialExtension chainLaw}}


Apply \texttt{initial\allowbreak{}Extension\_\allowbreak{}map} to the measurable recursion output. Its canonical pushforward is exactly \texttt{chain\allowbreak{}Law}, so the same output law is realized on every initial-coupling extension. This equality is derived rather than made a field of the coupling.

\subsection{\texorpdfstring{initial\allowbreak{}Extension iid noise\allowbreak{}Law}{initialExtension iid noiseLaw}}


First project the extension to its product-path noise law, then evaluate coordinate \(n\). The marginal of \(\bigotimes_k\eta\) at that coordinate is \(\eta\), giving the requested noise law on the extension.

\subsection{\texorpdfstring{initial\allowbreak{}Extension iid independent}{initialExtension iid independent}}


The iid path has independent coordinate evaluations. Pull this independence back along the extension's measure-preserving noise projection. The resulting innovation sequence remains jointly independent on the coupled initialization space.

\subsection{\texorpdfstring{initial\allowbreak{}Extension initial independent}{initialExtension initial independent}}


The algorithm's initial state is a projection of the initial pair. Compose the pair/\allowbreak{}noise independence theorem with this state projection. This proves independence from the entire noise sample, not just from each individual noise vector.

\subsection{\texorpdfstring{initial\allowbreak{}Extension cost}{initialExtension cost}}


The difference \(x-y\) depends only on the initial pair. Its \(L^2\) extended norm under \(c\otimes Q\) equals the same norm under \(c\), because the noise factor has mass one. The latter is exactly \texttt{coupling\allowbreak{}Cost c.\allowbreak{}joint}.

\subsection{\texorpdfstring{initial\allowbreak{}Extension mem\allowbreak{}Lp difference}{initialExtension memLp difference}}


The state-difference function is measurable. Its extended \(L^2\) norm equals the finite initial coupling cost by the previous theorem. These two facts prove \texttt{Mem\allowbreak{}Lp} for the difference on the extension.

\subsection{\texorpdfstring{initial\allowbreak{}Extension mem\allowbreak{}Lp reference}{initialExtension memLp reference}}


The reference coordinate has law \(\pi\). Pull its supplied \(L^2(\pi)\) identity moment back through this law to get square integrability on the extension. This establishes the reference initialization without choosing a separate realization.

\subsection{\texorpdfstring{initial\allowbreak{}Extension mem\allowbreak{}Lp initial}{initialExtension memLp initial}}


Write \(x=(x-y)+y\). The difference is in \(L^2\) by finite coupling cost and the reference coordinate is in \(L^2\) by its law. Closure of \(L^2\) under addition gives square integrability of the initial algorithm state; no independent moment premise on \(\mu\) is introduced.

\subsection{\texorpdfstring{initial\allowbreak{}Extension lp\allowbreak{}Initial}{initialExtension lpInitial}}


Use \texttt{initial\allowbreak{}Extension\_\allowbreak{}mem\allowbreak{}Lp\_\allowbreak{}initial} to convert the first state coordinate into an actual \(L^2\) class. The class is represented almost everywhere by \(((x,y),s)\mapsto x\), preserving its concrete relation to the coupling coordinates.

\subsection{\texorpdfstring{initial\allowbreak{}Extension lp\allowbreak{}Reference}{initialExtension lpReference}}


Convert the square-integrable second state coordinate into an \(L^2\) class. Its representative equals \(((x,y),s)\mapsto y\) almost everywhere, with law \(\pi\), and it shares the same product probability space as the algorithm initialization.

\subsection{\texorpdfstring{initial\allowbreak{}Extension lp\allowbreak{}Initial law}{initialExtension lpInitial law}}


The representative chosen by \texttt{to\allowbreak{}Lp} agrees almost everywhere with the first state coordinate. Pushforwards are unchanged by such an equality. Combine this with \texttt{initial\allowbreak{}Extension\_\allowbreak{}initial} to obtain the initial law \(\mu\).

\subsection{\texorpdfstring{initial\allowbreak{}Extension lp\allowbreak{}Reference law}{initialExtension lpReference law}}


Replace the selected \(L^2\) representative by the second state projection almost everywhere. Its already proved marginal is \(\pi\), hence the class used as reference initialization has exactly that law.

\subsection{\texorpdfstring{initial\allowbreak{}Extension lp cost}{initialExtension lp cost}}


The representative of the class difference equals \(x-y\) almost everywhere. The norm of that \(L^2\) difference is the real value of its extended norm; \texttt{initial\allowbreak{}Extension\_\allowbreak{}cost} identifies it with \((\operatorname{couplingCost}c)_{\mathbb R}\). The finite-cost premise justifies this real conversion.

\section{\texorpdfstring{Actual independent-\allowbreak{}source kinetic laws and their Wasserstein infimum}{Actual independent-source kinetic laws and their Wasserstein infimum}}
\label{proof-module-KineticLangevin-FlowLaw-lean}
{\small\color{gray}Source module: \texttt{Kinetic\allowbreak{}Langevin/\allowbreak{}Flow\allowbreak{}Law.\allowbreak{}lean}.\par}


The position law is an explicit pushforward of an independent product of the initial position law and a fixed source law for common initial velocity and continuous forcing. Every initial coupling is extended by that same source. The proof computes the two output marginals, preserves the initial cost, derives a coupling estimate and takes the infimum. The kinetic pointwise contraction is separately proved from the actual integral equation and the potential hypotheses; it is not an assumed convergence field.

\subsection{\texorpdfstring{Define the independent random-\allowbreak{}map law and compute both output marginals}{Define the independent random-map law and compute both output marginals}}


For a measurable endpoint \(F(x,s)\), set \(\mathcal T_Q\mu=F_\#(\mu\otimes Q)\). Given any coupling \(c\) of \(\mu,\nu\), extend it to \(c\otimes Q\). Projecting to \((y,s)\) produces \(\nu\otimes Q\), because the second marginal of \(c\) is \(\nu\). Pushing further by \(F\) therefore gives \(\mathcal T_Q\nu\). Lean proves this with the product-map identity, the coupling's marginal equality, and pushforward composition. The corresponding first marginal uses the existing \texttt{initial\allowbreak{}Extension\_\allowbreak{}map}. Both output-law identities are derived from the concrete product construction and are independent of which coupling was chosen.

\subsection{\texorpdfstring{Lift pathwise contraction through L\ensuremath{{}^{2}} and take all initial couplings}{Lift pathwise contraction through L² and take all initial couplings}}


Suppose the actual endpoint obeys \(\|F(x,s)-F(y,s)\|\le q\|x-y\|\) for each common source. On \(c\otimes Q\), the two outputs form a measurable coupling of \(\mathcal T_Q\mu,\mathcal T_Q\nu\). Pointwise contraction passes to the extended L\ensuremath{{}^{2}} norm. The initial difference depends only on \(c\), so its norm on the product is exactly the original coupling cost. Thus \(W_2(\mathcal T_Q\mu,\mathcal T_Q\nu)\le q\,\operatorname{cost}(c)\) for every \(c\). Lean then applies the infimum theorem, using the independent product \(\mu\otimes\nu\) as a witness that couplings exist. No optimal-coupling attainment and no supplied output-law independence are required.

\subsection{\texorpdfstring{Retain probability and justify the real-\allowbreak{}valued distance}{Retain probability and justify the real-valued distance}}


A measurable pushforward of \(\mu\otimes Q\) has total mass one when both inputs are probability measures; \texttt{random\allowbreak{}Map\allowbreak{}Law\_\allowbreak{}probability} records this. The primary distance bound lives in the extended nonnegative reals, so infinite distances are preserved. If the initial Wasserstein distance is finite, the finite coefficient \(q\) makes its upper bound finite, and monotonicity of \texttt{ENNReal.\allowbreak{}to\allowbreak{}Real} gives the real-valued estimate. Lean proves that finiteness before converting the inequality, then rewrites the real value of the product. This avoids accidentally converting an infinite output distance to zero.

\subsection{\texorpdfstring{Make the concrete kinetic endpoint a measurable random map}{Make the concrete kinetic endpoint a measurable random map}}


The compact forcing space receives its Borel measurable structure. The source \(s=(v,B)\) contains initial velocity and a continuous forcing path; the separate input is initial position \(x\). \texttt{kinetic\allowbreak{}Position\allowbreak{}Map} evaluates the actual selected flow at phase \((v,x)\) and takes its position coordinate. The coordinate rearrangement \((x,(v,B))\mapsto((v,x),B)\) is continuous. Composing it with the proved joint endpoint continuity and projecting position gives continuity, hence measurability, of the actual map. Independence of source and position is implemented later by the explicit product measure rather than encoded as an unspecified output-law premise.

\subsection{\texorpdfstring{Derive pointwise contraction from actual integral solutions}{Derive pointwise contraction from actual integral solutions}}


Take two values of the selected kinetic flow with phases \((v,x)\) and \((v,y)\) and the same forcing path. Their specification gives the actual integral equations and equal initial velocities. The proved integral-synchronous theorem cancels their common forcing, derives the right-derivative difference equation, and applies the nonlinear potential energy estimate. Under differentiability, \(m\)-strong convexity, \(M\)-Lipschitz gradient and \(\gamma^2\ge m+M\), this yields \(\|F_t(x,s)-F_t(y,s)\|\le\sqrt2e^{-mt/\gamma}\|x-y\|\). Lean obtains the integral premises by rewriting the flow's initial-value specification, then simplifies the initial position and velocity projections. No differential equation or contraction conclusion is assumed of the chosen flow.

\subsection{\texorpdfstring{Conclude the actual law-\allowbreak{}to-\allowbreak{}law kinetic theorem}{Conclude the actual law-to-law kinetic theorem}}


\texttt{kinetic\allowbreak{}Position\allowbreak{}Law} is the explicit independent-product pushforward for this concrete endpoint. Its probability theorem follows from endpoint measurability. Substituting the derived coefficient \(q=\sqrt2e^{-mt/\gamma}\) in the general random-map transport theorem gives

\[W_2(\mu_t,\nu_t)\le\sqrt2e^{-mt/\gamma}W_2(\mu_0,\nu_0).\]

The source law may be any fixed probability law of initial velocity and continuous forcing, independent of initial position. The statement therefore applies to any such stochastic realization; a stationary Gibbs law is unnecessary for comparing two evolving laws. The extended-valued theorem retains infinite initial costs, and the real-valued theorem explicitly requires finite initial Wasserstein distance. The full all-friction spectral prefactor and a canonical Brownian-source identification are separate refinements beyond this simple-friction law theorem.

\subsection{\texorpdfstring{The full shifted prefactor and exponent pass to actual output laws}{The full shifted prefactor and exponent pass to actual output laws}}


The actual-integral general contraction theorem uses the centered gradient sector estimate, so it requires differentiability, strong convexity and a Lipschitz gradient without adding Hessian continuity. For \(d=\gamma-2a>0\), its coefficient is

\[q_a(t)=\frac{\sqrt{2((\gamma-a)^2+a^2)}}{\gamma-2a}
\exp\!\left(\frac{\max\{a^2-m,M-(\gamma-a)^2\}}{\gamma-2a}t\right).\]

The proof instantiates that actual-integral theorem with the same two selected flows and their proved specifications. The coefficient is nonnegative, so it can be bundled as an \(\mathbb R_{\ge0}\) scalar in the already proved random-map transport theorem. This yields \(W_2(\mu_t,\nu_t)\le q_a(t)W_2(\mu_0,\nu_0)\), with the genuine coupling infimum and internally derived output marginals. The real-valued version uses the same finite-initial-cost conversion. The theorem holds for every independent continuous-forcing source law; tying a constructed Brownian source and its required independence to the paper’s named diffusion remains separate, and no Gibbs invariant law is assumed.

\section{\texorpdfstring{Exponential response kernels and valid Gaussian covariance matrices}{Exponential response kernels and valid Gaussian covariance matrices}}
\label{proof-module-KineticLangevin-Coefficients-lean}
{\small\color{gray}Source module: \texttt{Kinetic\allowbreak{}Langevin/\allowbreak{}Coefficients.\allowbreak{}lean}.\par}


The functions \(\psi_0(t)=e^{-\gamma t}\), \(\psi_{k+1}(t)=\int_0^t\psi_k(s)ds\) encode exactly integrated friction. The Hessian method also uses \(\varphi_{k+1}(t)=\int_0^te^{-\gamma(t-s)}\psi_k(s)ds\). Integral definitions avoid division by friction, so zero friction is handled exactly. The module proves their calculus and constructs positive semidefinite Gram covariances; it does not yet construct the corresponding Gaussian random variables.

\subsection{\texorpdfstring{Define iterated response coefficients and prove their calculus}{Define iterated response coefficients and prove their calculus}}


The recursive definition chooses an exponential at index zero and an interval integral at a successor index; the two \texttt{rfl} lemmas expose these cases to rewriting. Continuity follows by induction: exponential continuity handles zero, while continuity of the primitive handles the successor. The fundamental theorem of calculus then proves \(\psi'_{k+1}(t)=\psi_k(t)\), using interval integrability and continuity of the integrand. The separate zero-index derivative is \(\psi'_0=-\gamma\psi_0\), obtained from the exponential chain rule; \texttt{convert} and simplification reconcile its scalar expression with the recursive notation.

\subsection{\texorpdfstring{Derive friction balance before dividing by friction}{Derive friction balance before dividing by friction}}


Differentiate \(1-\psi_0(t)\): its derivative is \(\gamma\psi_0(t)\), and its value at zero is zero. Integrating gives \(\gamma\psi_1(t)=1-\psi_0(t)\). Integrating this identity once more gives \(\gamma\psi_2(t)=t-\psi_1(t)\). Lean first applies the integral form of the fundamental theorem, then commutes a scalar through the integral. The second proof integrates equality of two functions, distributes the integral over subtraction with its integrability proofs, and evaluates the constant integral. Only after establishing these division-free identities do \texttt{psi\_\allowbreak{}one\_\allowbreak{}closed} and \texttt{psi\_\allowbreak{}two\_\allowbreak{}closed} use \(\gamma\ne0\) to derive the familiar quotient formulas.

\subsection{\texorpdfstring{Prove exact time scaling of every primitive}{Prove exact time scaling of every primitive}}


For \(c\ne0\), substitution \(r=cs\) gives \(\psi_{c\gamma,k}(t)=c^{-k}\psi_{\gamma,k}(ct)\). At index zero this is the scalar identity \(-(c\gamma)t=-\gamma(ct)\). At a successor, use the induction hypothesis inside the integral, pull out \(c^{-k}\), and change variables in the remaining primitive; one additional inverse factor appears. Lean uses \texttt{integral\_\allowbreak{}comp\_\allowbreak{}mul\_\allowbreak{}left}, so oriented interval integrals also cover negative \(c\). Substituting \(t=h/c\) yields \(\psi_{c\gamma,k}(h/c)=\psi_{\gamma,k}(h)/c^k\). This is the exact coefficient identity used in mass and particle rescalings.

\subsection{\texorpdfstring{Establish initial values, signs and polynomial upper bounds}{Establish initial values, signs and polynomial upper bounds}}


At time zero the exponential coefficient is one and every primitive is zero. For \(t\ge0\), induction and nonnegative integration show every coefficient is nonnegative, even without a sign restriction on \(\gamma\). If also \(\gamma\ge0\), then \(e^{-\gamma t}\le1\). Integrating once gives \(\psi_1(t)\le t\), and integrating again gives \(\psi_2(t)\le t^2/2\). Lean's comparison theorem receives continuous, hence interval-integrable, functions and the pointwise inequality on the interval. These elementary estimates control the local time-step error without approximating the exponential by an unproved series.

\subsection{\texorpdfstring{Check the zero-\allowbreak{}friction limit exactly}{Check the zero-friction limit exactly}}


Putting \(\gamma=0\) in the integral definition yields \(\psi_0=1\), \(\psi_1=t\), and \(\psi_2=t^2/2\). The first two proofs simplify the exponential and constant integral. The third uses the integral of the first power. These are identities at zero, not limits inferred from quotient formulas, so no removable-singularity argument is needed. They later turn the kinetic exponential step into the constant-acceleration position–velocity update.

\subsection{\texorpdfstring{Define the Hessian response and its useful factorization}{Define the Hessian response and its useful factorization}}


\texttt{phi} assigns zero to index zero as an explicit convention and defines positive indices by exponential convolution. Its rewrite lemmas expose those cases. Nonnegativity follows from positivity of the exponential and nonnegativity of \(\psi_k\). Factoring \(e^{-\gamma t}\) out of the integral gives \(\varphi_{k+1}(t)=e^{-\gamma t}\int_0^te^{\gamma s}\psi_k(s)ds\). Lean proves the integrands agree by the exponential addition identity and scalar arithmetic. This factorization expresses the convolution as a product of a continuous exponential and a continuous primitive, proving continuity even though the upper limit and the kernel both depend on time. Every \(\varphi_k(0)\) is zero.

\subsection{\texorpdfstring{Compare convolution coefficients in the correct direction}{Compare convolution coefficients in the correct direction}}


For \(\gamma,t\ge0\), the factor \(e^{-\gamma(t-s)}\le1\) on \(0\le s\le t\), and \(\psi_k(s)\ge0\). Thus \(\varphi_{k+1}(t)\le\int_0^t\psi_k(s)ds=\psi_{k+1}(t)\). The proof explicitly establishes integrability of the convolution and primitive before applying integral comparison. At zero friction the exponential factor is one, so equality holds. In particular \(\varphi_2\le\psi_2\), correcting the opposite ordering printed in the supplied Proposition 3; both coefficients can still be bounded above by \(t^2/2\).

\subsection{\texorpdfstring{A Gram quadratic form is an integral of a square}{A Gram quadratic form is an integral of a square}}


Given continuous scalar kernels \(k_i\), define \(C_{ij}=\int_0^hk_i(t)k_j(t)dt\). Commutativity proves symmetry. For a finite coefficient vector \(z\), expand the quadratic form and interchange its finite sums with integrals:

\[\sum_{i,j}z_iC_{ij}z_j=\int_0^h\left(\sum_i z_ik_i(t)\right)^2dt.\]

Lean supplies interval-integrability witnesses for each product, pulls scalar factors inside, applies \texttt{integral\_\allowbreak{}finset\allowbreak{}Sum} twice, and normalizes the finite-sum expansion. For \(h\ge0\) the integral is nonnegative. Combining symmetry with this quadratic-form bound proves positive semidefiniteness. This is a matrix property, not yet an identification with Itô covariance.

\subsection{\texorpdfstring{Instantiate the first-\allowbreak{} and second-\allowbreak{}order noise covariances}{Instantiate the first- and second-order noise covariances}}


The first integrator uses kernels \((\psi_0,\psi_1)\); the Hessian integrator uses \((\psi_0,\psi_1,\varphi_2,\varphi_3)\). \texttt{klmc\allowbreak{}Kernels} and \texttt{klmc2\allowbreak{}Kernels} list them on \(\mathrm{Fin}(2)\) and \(\mathrm{Fin}(4)\). Their covariance definitions apply the same Gram construction. Finite case splits reduce kernel continuity to the proved coefficient continuity. Each final positive-semidefinite theorem is then a direct application of the general Gram result. In the stochastic algorithm the noise is additionally required to be centered jointly Gaussian with these covariances, independent across time and coordinates; those law properties are not encoded by matrix positivity alone.

\subsection{\texorpdfstring{Strict positivity of every positive-\allowbreak{}time response}{Strict positivity of every positive-time response}}


For \(t>0\), every response coefficient \(\psi_k(t)\) is strictly positive. The exponential base case is positive at every time. At a successor, the integrand is continuous and nonnegative on the interval, and the induction hypothesis makes it strictly positive at the right endpoint. Mathlib's positive-integral theorem then gives a positive primitive. The proof explicitly supplies both the pointwise nonnegative bound and an endpoint witnessing positivity. No positive-friction hypothesis is needed because an exponential remains positive even when its exponent grows.

\subsection{\texorpdfstring{All higher responses are bounded by their frictionless Taylor coefficient}{All higher responses are bounded by their frictionless Taylor coefficient}}


At zero friction induction gives \(\psi_{0,k}(t)=t^k/k!\). Integrating the induction hypothesis uses \(\int_0^ts^kds=t^{k+1}/(k+1)\), and the factorial recursion \((k+1)!=(k+1)k!\) matches the denominator. For nonnegative friction and time, compare each \(\psi_{\gamma,k}\) with its zero-friction counterpart. The exponential comparison starts the induction, and integral monotonicity preserves it at each successor. Substitution of the exact zero-friction identity yields \(0\le\psi_{\gamma,k}(t)\le t^k/k!\). These estimates provide every higher-order coefficient bound without a formal power-series approximation.

\subsection{\texorpdfstring{Evaluate the first Hessian response without division}{Evaluate the first Hessian response without division}}


The convolution of the two exponential responses has a constant integrand: \(e^{-\gamma(t-s)}\psi_0(s)=e^{-\gamma t}\). Hence \(\varphi_1(t)=t\psi_0(t)\). Reversing the integration variable also gives \(\int_0^te^{-\gamma(t-s)}ds=\psi_1(t)\). To compute the next balance, multiply the \(\varphi_2\) integrand by \(\gamma\), substitute \(\gamma\psi_1(s)=1-\psi_0(s)\), and simplify the exponential product. Integrating yields \(\gamma\varphi_2(t)=\psi_1(t)-t\psi_0(t)\). Lean proves equality of the integrand functions, commutes constants through integrals, and evaluates the constant integral. These identities remain valid at zero friction.

\subsection{\texorpdfstring{Compute covariance entries and the exact friction–noise balance}{Compute covariance entries and the exact friction–noise balance}}


The velocity variance coefficient is \(C_{00}=\int_0^he^{-2\gamma t}dt=\psi_{2\gamma,1}(h)\), proved by combining exponentials. Since \((\psi_1^2/2)'=\psi_0\psi_1\) and \(\psi_1(0)=0\), the fundamental theorem gives \(C_{01}=\psi_1(h)^2/2\). Finally the friction identity at rate \(2\gamma\) yields \(2\gamma C_{00}=1-e^{-2\gamma h}\). Therefore \(\psi_0(h)^2+2\gamma C_{00}=1\), the exact stationary velocity variance balance. Lean proves the exponential square identity and reduces the final expression with \texttt{ring}. An independent unit-variance input velocity and the specified Gaussian noise consequently retain unit variance after the linear thermostat part.

\subsection{\texorpdfstring{Relate the second-\allowbreak{}order coefficients with an exact integrated balance}{Relate the second-order coefficients with an exact integrated balance}}


The derivative of \((t-s)\psi_1(s)+\psi_2(s)\), with \(t\) fixed, is \((t-s)\psi_0(s)\). Integrating and reversing the variable yields \(\int_0^te^{-\gamma(t-s)}s\,ds=\psi_2(t)\). Next put the defining integrals for \(\varphi_2\) and \(\varphi_3\) under one integral. The identity \(\psi_1(s)+\gamma\psi_2(s)=s\) simplifies their integrand, giving \(\varphi_2(t)+\gamma\varphi_3(t)=\psi_2(t)\). Lean explicitly supplies interval-integrability before combining the integrals and uses \texttt{linear\_\allowbreak{}combination} for the scalar integrand equality. For nonnegative friction and time, \(\gamma\varphi_3(t)\ge0\), hence \(\varphi_2(t)\le\psi_2(t)\). This provides a direct corrected coefficient relation, rather than the reversed intermediate inequality printed in the supplied manuscript.

\section{\texorpdfstring{From the actual integral equation to the exponential mild equation}{From the actual integral equation to the exponential mild equation}}
\label{proof-module-KineticLangevin-MildEquation-lean}
{\small\color{gray}Source module: \texttt{Kinetic\allowbreak{}Langevin/\allowbreak{}Mild\allowbreak{}Equation.\allowbreak{}lean}.\par}


The KLMC update freezes the force inside an exponentially weighted integral. To compare it with an exact path, that weighted formula must follow from the actual integral equation driven by the same Brownian input. The theorem below proves this bridge for any complete real normed space, including the L\ensuremath{{}^{2}} space of random Euclidean vectors. No time derivative of Brownian motion is used.

\subsection{\texorpdfstring{Subtract the filtered noise and apply an integrating factor}{Subtract the filtered noise and apply an integrating factor}}


Assume continuous velocity \(V\), continuous force path \(F\), and a continuous filtered input \(B\), satisfying \(V_t-V_0=-\int_0^t(\gamma V_s+F_s)ds+W_t\) and \(B_t=W_t-\gamma\int_0^tB_sds\). The conclusion is the exact identity \(V_T=e^{-\gamma T}V_0-\int_0^Te^{-\gamma(T-s)}F_sds+B_T\). In the proof, evaluation at zero first gives \(B_0=0\). Subtracting the two equations shows \(D=V-B\) solves \(D_t=V_0-\int_0^t(\gamma D_s+F_s)ds\). Continuous integrands allow the fundamental theorem of calculus on the interval interior. Lean derives \(D\!\prime=-\gamma D-F\), differentiates \(e^{\gamma t}D_t\), and obtains derivative \(-e^{\gamma t}F_t\). A second application of the fundamental theorem integrates this identity over the closed interval; the proof handles the endpoints by continuity. Finally it multiplies by \(e^{-\gamma T}\), commutes scalar multiplication with the Bochner integral, and simplifies the exponent. The hypotheses are the genuine integral equations, and the resulting mild equation is proved rather than postulated. The argument works for arbitrary real friction and includes \(T=0\).

\section{\texorpdfstring{One shared exponential integrator and the Hessian correction}{One shared exponential integrator and the Hessian correction}}
\label{proof-module-KineticLangevin-Integrator-lean}
{\small\color{gray}Source module: \texttt{Kinetic\allowbreak{}Langevin/\allowbreak{}Integrator.\allowbreak{}lean}.\par}


The deterministic maps accept a force vector field and supplied noise vectors. Keeping the force explicit permits particle and inverse-mass reuse. Throughout the module, the first phase coordinate is velocity and the second is position. The coefficient library specifies the covariance matrices required by the probabilistic interpretation; these definitions by themselves make no assertion that arbitrary noise inputs have those laws.

\subsection{\texorpdfstring{Define phase space and the frozen-\allowbreak{}force step}{Define phase space and the frozen-force step}}


\texttt{Phase\allowbreak{}Space E} is the pair \(E\times E\). For supplied additive innovation \((n_v,n_x)\), define

\[v^+=\psi_0v-\psi_1g(x)+n_v,\qquad x^+=x+\psi_1v-\psi_2g(x)+n_x.\]

\texttt{exponential\allowbreak{}Step} implements these two expressions. \texttt{klmc\allowbreak{}Step} passes innovation \(\sqrt{2\gamma}(\xi,\xi')\), matching the paper's noise normalization. The velocity and position projection theorems are definitional equalities: \texttt{rfl} exposes the corresponding pair coordinate. This separation lets KIPLMC1 use exactly this implementation with its own temperature amplitude rather than copying the formula.

\subsection{\texorpdfstring{Check zero time and cancel shared additive noise}{Check zero time and cancel shared additive noise}}


At zero time, \(\psi_0=1\), \(\psi_1=\psi_2=0\); with zero innovation the step is the identity. Lean proves equality of pairs by coordinate extensionality and simplification. For two starting phases driven by the same innovation, subtraction cancels each additive noise term before any estimate:

\[\Delta v^+=\psi_0\Delta v-\psi_1\Delta g,\quad
\Delta x^+=\Delta x+\psi_1\Delta v-\psi_2\Delta g.\]

The proof expands pair subtraction, distributes scalar multiplication across differences, and uses \texttt{abel} for the additive group rearrangement. The KLMC cancellation theorem reuses this exact identity with the scaled noise argument. It does not need Gaussianity.

\subsection{\texorpdfstring{Include every Hessian correction, including noise multipliers}{Include every Hessian correction, including noise multipliers}}


With \(H=\nabla^2f(x)\), the second-order map subtracts \(\varphi_2Hv\) from velocity and \(\varphi_3Hv\) from position. Its innovations are \(\sqrt{2\gamma}(\xi_0-H\xi_2)\) and \(\sqrt{2\gamma}(\xi_1-H\xi_3)\). The definition \texttt{klmc2\allowbreak{}Step} includes all four terms. With zero Hessian, the deterministic and random corrections vanish and the map becomes the original KLMC map using \((\xi_0,\xi_1)\). Pair extensionality and simplification prove this reduction. When comparing two nonzero Hessians, the noise terms differ too; this explicit map prevents silently discarding that contribution.

\subsection{\texorpdfstring{The frictionless map is constant acceleration}{The frictionless map is constant acceleration}}


At \(\gamma=0\), the coefficient identities give \(\psi_0=1\), \(\psi_1=h\), \(\psi_2=h^2/2\), while \(\sqrt{2\gamma}=0\). Therefore \(v^+=v-hg(x)\) and \(x^+=x+hv-h^2g(x)/2\). The proof simplifies the actual map with the exact zero-friction coefficient lemmas. It asserts exact integration of a frozen force, not exact integration of a spatially varying Hamiltonian force along the whole interval.

\subsection{\texorpdfstring{Iterate a prescribed stream of innovations}{Iterate a prescribed stream of innovations}}


\texttt{klmc\allowbreak{}Chain} starts at the supplied phase and recursively applies \texttt{klmc\allowbreak{}Step} with innovation number \(k\). Its zero and successor equations are definitional equalities, providing reliable rewrite points for later induction and law identification. \texttt{klmc2\allowbreak{}Chain} uses the same recursion pattern with the four-block innovation required by \texttt{klmc2\allowbreak{}Step}. A prescribed noise stream is a pathwise object: assigning independent Gaussian coordinate laws and proving finite moments of this recursion is the next probabilistic task, not a hidden assumption in the definitions.

\subsection{\texorpdfstring{Verify that the frozen deterministic response solves its differential equations}{Verify that the frozen deterministic response solves its differential equations}}


For constant force \(g\), define \(v_t=\psi_0(t)v_0-\psi_1(t)g\) and \(x_t=x_0+\psi_1(t)v_0-\psi_2(t)g\). The coefficient derivative identities immediately give \(x'_t=v_t\). For velocity they give \(v'_t=-\gamma\psi_0(t)v_0-\psi_0(t)g\). Substituting \(\gamma\psi_1=1-\psi_0\) proves this equals \(-\gamma v_t-g\). Lean builds the vector derivatives from the scalar coefficient derivatives and rewrites the friction balance before normalizing scalar actions. \texttt{exponential\allowbreak{}Step\_\allowbreak{}eq\_\allowbreak{}frozen} is then a definitional equality: the deterministic step evaluates these exact frozen responses at \(h\) and adds the supplied noise. Thus the term “exponential integrator” is supported by an actual ODE identity.

\section{\texorpdfstring{Constructing the actual reference and auxiliary coupling segments}{Constructing the actual reference and auxiliary coupling segments}}
\label{proof-module-KineticLangevin-CoupledSegments-lean}
{\small\color{gray}Source module: \texttt{Kinetic\allowbreak{}Langevin/\allowbreak{}Coupled\allowbreak{}Segments.\allowbreak{}lean}.\par}


The local KLMC analysis compares a numerical state with an exact reference state. On each step it inserts an auxiliary exact path started at the numerical state and driven by precisely the same additive input as the reference. This module constructs those paths from the proved kinetic flow. The differential, kinematic and mild identities used in the local-error theorem follow from their actual integral equations.

\subsection{\texorpdfstring{Restart an actual integral solution at any time}{Restart an actual integral solution at any time}}


If \(Y_t-Y_0=-\int_0^tg(Y_s)ds+W_t\), then \(Y_{a+t}-Y_a=-\int_0^tg(Y_{a+s})ds+(W_{a+t}-W_a)\) whenever the shifted interval lies in the original horizon. Lean subtracts the two actual equations at \(a+t\) and \(a\), splits the continuous drift integral at \(a\), and changes variables in the remaining integral. Additive group algebra cancels the shared starting point. This supplies the exact time-shift identity required when segmenting a continuous diffusion.

\subsection{\texorpdfstring{Use a raw continuous input with the selected kinetic flow}{Use a raw continuous input with the selected kinetic flow}}


A continuous input \(W\) is restricted to the compact solution horizon, producing the existing \texttt{Velocity\allowbreak{}Forcing} object. If \(W_0=0\), the flow's normalized clipped extension equals \(W_t\) at every time in that horizon. Lean uses the exact projection identity \texttt{proj\allowbreak{}Icc\_\allowbreak{}of\_\allowbreak{}mem} to prove this. The selected segment is then the existing Picard-constructed \texttt{kinetic\allowbreak{}Flow} with that compact input. Its specification gives continuity, its true initial condition, and the exact integral equation with the original raw input.

\subsection{\texorpdfstring{Derive velocity, kinematic and mild equations by projection}{Derive velocity, kinematic and mild equations by projection}}


Applying the velocity continuous linear projection to the kinetic integral equation gives \(V_t-V_0=-\int_0^t(\gamma V_s+u g(X_s))ds+W_t\). Applying the position projection gives \(X_t-X_0=\int_0^tV_sds\), because the positional input is exactly zero. In both cases Lean commutes the projection with the Bochner integral. If the actual filtered input satisfies \(B_t=W_t-\gamma\int_0^tB_sds\), the proved Banach-space integrating-factor theorem gives \(V_T=\psi_0(T)V_0-\int_0^Te^{-\gamma(T-s)}u g(X_s)ds+B_T\). Thus the exponential mild equation is a consequence of the kinetic equation, rather than a new stochastic hypothesis.

\subsection{\texorpdfstring{Construct exact reference segments and the implemented numerical chain}{Construct exact reference segments and the implemented numerical chain}}


The reference chain starts from its supplied actual phase vector. At each iteration, its next state is the endpoint of the selected exact kinetic segment driven by that iteration's input. The corresponding reference segment is the same selected path, so its endpoint matches the next reference state by definition. The numerical chain instead performs the actual frozen-force exponential update with innovation \((B_n(h),\int_0^hB_n(s)ds)\). This uses the constructed filtered input and its genuine positional primitive.

\subsection{\texorpdfstring{Construct the auxiliary exact path from the current numerical state}{Construct the auxiliary exact path from the current numerical state}}


The auxiliary segment is an actual selected exact kinetic solution whose initial condition is the current exponential-integrator state. It receives exactly the same raw Brownian input as the reference segment on that iteration. The two specification theorems prove continuity, exact starting values, reference endpoint matching, and their shared-input integral equations. Every field comes from the kinetic flow construction; no contraction or local-error inequality is inserted into the definitions.

\subsection{\texorpdfstring{Discharge the local-\allowbreak{}error proof's kinematic and mild requirements}{Discharge the local-error proof's kinematic and mild requirements}}


The constructed auxiliary path has its positional primitive identity by the general kinetic projection theorem. Its exponentially weighted force formula follows by restricting its actual integral equation to any smaller time \(t\le h\), restricting the filtered-input balance to the same interval, and applying the generic mild theorem. Lean explicitly checks the interval inclusions and simplifies the inverse mass \(u=1\). The remaining stochastic work concerns the law and stationarity of the reference; the coupling segments and their analytic identities are concrete constructions.

\section{\texorpdfstring{Actual kinetic flows under measure-\allowbreak{}preserving source changes}{Actual kinetic flows under measure-preserving source changes}}
\label{proof-module-KineticLangevin-FlowPullback-lean}
{\small\color{gray}Source module: \texttt{Kinetic\allowbreak{}Langevin/\allowbreak{}Flow\allowbreak{}Pullback.\allowbreak{}lean}.\par}


Changing a probability source by a measure-preserving projection preserves the actual random variables and their output laws. This module proves that the genuine lifted nonlinear drift and the selected kinetic integral solution commute with that change of source. It supplies coupling-independent continuous-flow laws without assuming samplewise continuity of an arbitrary L\ensuremath{{}^{2}} representative.

\subsection{\texorpdfstring{Define continuous linear source and phase pullbacks}{Define continuous linear source and phase pullbacks}}


Measure-preserving composition is a linear isometry on L\ensuremath{{}^{2}}. The source pullback is this isometry viewed as a continuous linear map, and the phase pullback applies it separately to velocity and position. Their representative theorems identify the resulting random variables with actual composition. Pushforward composition proves preservation of every output law. A further representative argument proves that two successive pullbacks equal pullback along the composed projection.

\subsection{\texorpdfstring{Commute the actual nonlinear drift with pullback}{Commute the actual nonlinear drift with pullback}}


The lifted vector field is pointwise application of the original Lipschitz field. Lean compares its actual representatives before and after source projection, obtaining exact equality in L\ensuremath{{}^{2}}. Combining this with linearity proves the kinetic drift commutes with the phase pullback, including both friction and force terms and the position equation.

\subsection{\texorpdfstring{Identify the selected actual integral solutions}{Identify the selected actual integral solutions}}


Apply the phase pullback to the original kinetic integral equation. Continuity and linearity commute it with the Bochner time integral, and the drift commutation identity produces the target-space kinetic equation. The transformed initial condition and additive input are exact. The proved integral-solution uniqueness theorem therefore identifies the selected target-space solution with the pulled original solution at every time in the horizon. This equality transfers laws between an independent initial-state/\allowbreak{}noise product and any larger initial-coupling extension.



The reusable integral-equation theorem proves this transformation for any genuine continuous kinetic solution before specializing to selected flows.

\section{\texorpdfstring{The actual nonlinear kinetic coupling in L\ensuremath{{}^{2}}}{The actual nonlinear kinetic coupling in L²}}
\label{proof-module-KineticLangevin-LpContraction-lean}
{\small\color{gray}Source module: \texttt{Kinetic\allowbreak{}Langevin/\allowbreak{}Lp\allowbreak{}Contraction.\allowbreak{}lean}.\par}


Random vectors form a Hilbert space under the L\ensuremath{{}^{2}} norm. The gradient acts on these vectors through the project's pointwise \texttt{lp\allowbreak{}Drift} construction. This module lifts the centered strong-smooth inequality to that actual vector field, then proves contraction of actual kinetic integral paths in L\ensuremath{{}^{2}}. It does not introduce a hypothetical differentiable potential on the infinite-dimensional space of random variables.

\subsection{\texorpdfstring{Lift the centered gradient inequality pointwise}{Lift the centered gradient inequality pointwise}}


For deterministic points, strong convexity and Lipschitz gradient imply \(\|\nabla f(x)-\nabla f(y)-\frac{m+M}{2}(x-y)\|\le\frac{M-m}{2}\|x-y\|\). The lifted gradient agrees almost everywhere with pointwise application of the original gradient. Lean collects the almost-everywhere identities for subtraction, scalar multiplication and \texttt{lp\allowbreak{}Drift\_\allowbreak{}coe}, rewrites the random-variable representatives, and applies the deterministic centered inequality at each sample. \texttt{Lp.\allowbreak{}norm\_\allowbreak{}le\_\allowbreak{}mul\_\allowbreak{}norm\_\allowbreak{}of\_\allowbreak{}ae\_\allowbreak{}le\_\allowbreak{}mul} integrates this norm comparison, yielding the same constant in L\ensuremath{{}^{2}}. No independence or moment estimate beyond membership in L\ensuremath{{}^{2}} is needed for this step.

\subsection{\texorpdfstring{Derive the L\ensuremath{{}^{2}} energy inequality from actual integral paths}{Derive the L² energy inequality from actual integral paths}}


Two phase paths satisfy the concrete kinetic integral equations with the actual lifted gradient and common additive input. The noise cancels in those equations, so their difference has a right derivative. Continuous linear projections extract the velocity and position equations. The general centered dissipativity theorem applies in the Hilbert space L\ensuremath{{}^{2}} because the preceding pointwise lift supplies its sector premise. Under \(\gamma^2\ge m+M\), the maximum curvature coefficient equals \(-m\), giving \(E'(t)\le-2mE(t)/\gamma\). Continuity and the right-slope Grönwall theorem prove \(E(t)\le E(0)e^{-2mt/\gamma}\). Taking nonnegative square roots, using the exponential-addition identity, gives the full transformed L\ensuremath{{}^{2}} norm contraction.

\section{\texorpdfstring{The arbitrary-\allowbreak{}shift kinetic contraction for actual L\ensuremath{{}^{2}} random vectors}{The arbitrary-shift kinetic contraction for actual L² random vectors}}
\label{proof-module-KineticLangevin-LpGeneralContraction-lean}
{\small\color{gray}Source module: \texttt{Kinetic\allowbreak{}Langevin/\allowbreak{}Lp\allowbreak{}General\allowbreak{}Contraction.\allowbreak{}lean}.\par}


The continuous mixing theorem permits an arbitrary shifted kinetic metric, not only the convenient zero-shift case used by KLMC. The pointwise strong-smooth gradient sector bound already lifts to L\ensuremath{{}^{2}}. Here it gives the complete shifted energy and position estimates for actual shared-input L\ensuremath{{}^{2}} kinetic solutions.

\subsection{\texorpdfstring{Prove the full shifted L\ensuremath{{}^{2}} energy decay}{Prove the full shifted L² energy decay}}


Let \(d=\gamma-2a>0\). The exact common-noise integral equations make the difference differentiable, with the original gradient difference lifted pointwise. Lean differentiates the explicit shifted energy and applies the general dissipativity inequality using the genuine lifted sector bound. The resulting rate is \(2\max(a^2-m,M-(\gamma-a)^2)/d\). A right-derivative Grönwall argument integrates this differential inequality, including the starting endpoint. No potential on the L\ensuremath{{}^{2}} space is introduced.

\subsection{\texorpdfstring{Project the shifted energy to the position error}{Project the shifted energy to the position error}}


Equal initial velocities make the initial energy equal to \(((\gamma-a)^2+a^2)\|X_0-Y_0\|_2^2\). The invertible transformed coordinates bound the position error by the energy divided by \((\gamma-2a)^2\). Lean combines these identities with the energy decay, proves the exact squared-exponential identity, and takes the nonnegative square-root inequality. The conclusion is the full source-paper prefactor \(\sqrt{2((\gamma-a)^2+a^2)}/(\gamma-2a)\) multiplied by the explicit exponential rate and the actual initial L\ensuremath{{}^{2}} position cost.

\section{\texorpdfstring{The L\ensuremath{{}^{2}} recursion is the original implemented KLMC chain}{The L² recursion is the original implemented KLMC chain}}
\label{proof-module-KineticLangevin-LpDynamics-lean}
{\small\color{gray}Source module: \texttt{Kinetic\allowbreak{}Langevin/\allowbreak{}Lp\allowbreak{}Dynamics.\allowbreak{}lean}.\par}


Analysis on equivalence classes of random vectors is useful only if it represents the intended algorithm. The lemmas here prove that the actual lifted exponential update, KLMC noise scaling and every recursively defined iterate agree almost everywhere with the original pointwise vector formulas. No Gaussian or independence assumptions are required for this algebraic correspondence.

\subsection{\texorpdfstring{Read L\ensuremath{{}^{2}} phase points and identify the exponential update}{Read L² phase points and identify the exponential update}}


\texttt{lp\allowbreak{}Phase\allowbreak{}Coe} reads the velocity and position representatives of a pair of L\ensuremath{{}^{2}} classes at the same sample. A private affine-operation lemma combines the almost-everywhere identities for scalar multiplication, subtraction and addition. Apply it to the velocity formula \(\psi_0v-\psi_1G(x)+\eta_v\). For position, rewrite its affine expression as \(\psi_1v-\psi_2G(x)+(x+\eta_x)\), then use the same lemma. Finally \texttt{lp\allowbreak{}Drift\_\allowbreak{}coe} identifies \(G(x)(\omega)\) with \(\nabla f(x(\omega))\). The two component equalities give equality of the actual phase update and the original pointwise \texttt{exponential\allowbreak{}Step}.

\subsection{\texorpdfstring{Expand affine operations on L\ensuremath{{}^{2}} representatives}{Expand affine operations on L² representatives}}


The representative of a sum or scalar multiple of Lp classes agrees with the corresponding operation almost everywhere. The private helper gathers the four rules needed for \(aX-bY+Z\) into one common full-measure set. On that set, it rewrites the representative of the outer sum, the inner difference, and the two scalar multiples. The resulting equality is precisely the pointwise affine formula. This step accounts for the fact that Lp elements are equivalence classes; pointwise equalities of their arbitrarily selected representatives cannot be assumed everywhere.

\subsection{\texorpdfstring{Retain the exact KLMC noise scaling through every iterate}{Retain the exact KLMC noise scaling through every iterate}}


The two phase components of a scaled L\ensuremath{{}^{2}} innovation agree almost everywhere with the same scalar times their representatives. Combining this with the exponential-step correspondence proves the KLMC correspondence, including its \(\sqrt{2\gamma}\) fluctuation-dissipation factor. Induction then proves the entire chain identity: at zero both chains use the supplied initial point; at a successor, intersect the almost-everywhere previous-chain identity with the current-step identity, and rewrite the recursive definitions. This shows the L\ensuremath{{}^{2}} error estimates concern the original numerical sampler for the supplied innovation stream, rather than a different abstract sequence.

\section{\texorpdfstring{Iteration of the scalar error estimate}{Iteration of the scalar error estimate}}
\label{proof-module-UnadjustedLangevin-Recurrence-lean}
{\small\color{gray}Source module: \texttt{Unadjusted\allowbreak{}Langevin/\allowbreak{}Recurrence.\allowbreak{}lean}.\par}


Once coupling gives \(e_{k+1}\le qe_k+b\), the remaining calculation is deterministic. We retain the exact geometric sum before obtaining the simpler bias-floor bound. The upper envelope approaches \(b/(1-q)\); this does not assert that the actual error converges to that value.



The scalar sequence need not have an algorithmic interpretation for these lemmas. The proofs separate the exact geometric sum from its uniform bias-floor envelope and explicitly require the contraction factor in \([0,1)\). Later the coupling construction proves that the actual ULA error satisfies this recurrence.

\subsection{\texorpdfstring{Iteration gives a geometric sum}{Iteration gives a geometric sum}}


If \(q\ge0\) and \(e_{k+1}\le qe_k+b\) for every \(k\), then

\[e_n\le q^ne_0+b\sum_{j=0}^{n-1}q^j.\]

Induct on \(n\). At zero the sum is empty and both sides equal \(e_0\). Assume the bound at \(n\).



Multiply the induction inequality by \(q\ge0\), preserving its direction.



Write \(S_n=\sum_{j=0}^{n-1}q^j\). The geometric identity gives \(qS_n+1=S_n+q^n\); multiply this equality by \(b\).



Combine the one-step bound with the multiplied induction bound. Since \(S_{n+1}=S_n+q^n\), obtain \(e_{n+1}\le q^{n+1}e_0+bS_{n+1}\).

\subsection{\texorpdfstring{The exact finite-\allowbreak{}horizon envelope}{The exact finite-horizon envelope}}


If \(0\le q<1\) and \(e_{k+1}\le qe_k+b\), then

\[e_n\le q^ne_0+\frac{b(1-q^n)}{1-q}.\]

Apply the geometric-sum bound and evaluate the sum as \((q^n-1)/(q-1)\). Its denominator is nonzero because \(q<1\).



Change the signs of numerator and denominator simultaneously and put \(b\) in the numerator to obtain \(b(1-q^n)/(1-q)\).

\subsection{\texorpdfstring{A uniform bias-\allowbreak{}floor bound}{A uniform bias-floor bound}}


If \(0\le q<1\), \(b\ge0\), and \(e_{k+1}\le qe_k+b\), then

\[e_n\le q^ne_0+\frac b{1-q}.\]

Start from the exact finite-horizon bound. Its denominator \(1-q\) is positive.



Since \(q^n\ge0\) and \(b\ge0\), the numerator satisfies \(b(1-q^n)\le b\).



Divide by the positive denominator, add \(q^ne_0\), and compose with the exact bound.

\subsection{\texorpdfstring{Limit of the upper envelope}{Limit of the upper envelope}}


For \(0\le q<1\) and arbitrary real \(b,e_0\),

\[q^ne_0+\frac b{1-q}\longrightarrow\frac b{1-q}\quad(n\to\infty).\]

The geometric powers satisfy \(q^n\to0\).



Multiply this limit by \(e_0\) and add the constant \(b/(1-q)\).

\subsection{\texorpdfstring{Every accuracy above the floor is eventually attained}{Every accuracy above the floor is eventually attained}}


Assume \(0\le q<1\), \(b\ge0\), and \(e_{k+1}\le qe_k+b\). For every \(\varepsilon>b/(1-q)\), the errors satisfy \(e_n<\varepsilon\) for every sufficiently large \(n\).



The upper envelope tends to a value strictly below \(\varepsilon\), so it is eventually below \(\varepsilon\).



At each such index, place \(e_n\) below the envelope by the recurrence bound and then below \(\varepsilon\).

\section{\texorpdfstring{The shared Hilbert-\allowbreak{}space proof and both step regimes}{The shared Hilbert-space proof and both step regimes}}
\label{proof-module-UnadjustedLangevin-HilbertDalalyan-lean}
{\small\color{gray}Source module: \texttt{Unadjusted\allowbreak{}Langevin/\allowbreak{}Hilbert\allowbreak{}Dalalyan.\allowbreak{}lean}.\par}


Work on a probability space with \(\mathcal H=L^2(P;H)\), and write \(g=\widehat{\nabla U}\). This is the single analytic proof used by both Euclidean ULA and the lifted IPLA system. The reference record exposes the integral equation, stationary laws, and moment bounds used at this stage. \texttt{Canonical\allowbreak{}ULA.\allowbreak{}lean} constructs these inputs from strong convexity, a Lipschitz gradient, and independent Gaussian innovations.



\texttt{L2\allowbreak{}Reference} stores a stationary reference path in the same \(L^2\) space and with the same innovations as the numerical chain. The analytic theorem derives its recurrence and final transport estimate. The reference is a supplied argument at this layer; the canonical Brownian construction later proves its fields from the potential assumptions.

\subsection{\texorpdfstring{The common L\ensuremath{{}^{2}} algorithm}{The common L² algorithm}}


For an initial \(X_0\in L^2(P;H)\) and innovations \(\xi_k\in L^2(P;H)\), set \(X_{k+1}=X_k-h\widehat{\nabla U}(X_k)+\sqrt{2h}\,\xi_k\).



The initial iterate is the supplied random vector.



Each successor uses the actual-gradient Euler update.

\subsection{\texorpdfstring{Unfolding one step}{Unfolding one step}}


The next iterate is exactly one L\ensuremath{{}^{2}} ULA update of the current iterate.



Evaluate the successor clause of the definition.

\subsection{\texorpdfstring{The continuous reference assumptions}{The continuous reference assumptions}}


For an \(L\)-Lipschitz gradient, step \(h\), noise vectors \(\xi_k\in\mathcal H\), and a measure \(\pi\), a reference consists of segments \(Y_{k,t}\in\mathcal H\) and increments \(B_{k,t}\in\mathcal H\) satisfying the following conditions.



Supply the two families \((Y_{k,t})\) and \((B_{k,t})\) in the Hilbert space \(\mathcal H\).



Require continuity of \(t\mapsto Y_{k,t}\) in the \(L^2\) norm for every \(k\).



Require compatible endpoints, common endpoint noise, and stationary grid laws:

\[Y_{k,h}=Y_{k+1,0},\qquad B_{k,h}=\sqrt{2h}\,\xi_k,\qquad(Y_{k,0})_\#P=\pi.\]

For \(0\le t\le h\), require

\[Y_{k,t}-Y_{k,0}=-\int_0^t g(Y_{k,s})\,ds+B_{k,t}.\]

Require the score moment bound \(\|g(Y_{k,t})\|_2\le\sqrt{Ld}\) throughout each closed step interval.



Require the increment moment bound \(\|B_{k,t}\|_2\le\sqrt{2d}\sqrt t\) on the same interval.

\subsection{\texorpdfstring{Full convergence calculation from the integral reference}{Full convergence calculation from the integral reference}}


Suppose a reference with the preceding properties is supplied, \(h\ge0\), and \(Lh\le2\). Let \(0\le q<1\), and assume the lifted Euler map contracts:

\[\|(X-hg(X))-(Y-hg(Y))\|_2\le q\|X-Y\|_2.\]

For any initial vector \(X_0\in\mathcal H\), the \(L^2\) chain driven by the reference endpoint noise satisfies

\[W_2^{\mathbb R}((X_n)_\#P,\pi)\le q^n\|X_0-Y_{0,0}\|_2+\frac{(33/20)Lh\sqrt{hd}}{1-q}.\]

The statement is conditional on the supplied reference, and the initial term is the cost of the chosen coupling.



Set \(g=\widehat{\nabla U}\), name the actual iterates \(X_k\), and define

\[b=\frac{33}{20}Lh\sqrt{hd},\qquad D_k=\int_0^h\bigl(g(Y_{k,t})-g(Y_{k,0})\bigr)\,dt.\]

The bound \(b\) is nonnegative. Apply the local-defect theorem to the actual lifted gradient and the supplied equation and moment bounds, deriving \(\|D_k\|_2\le b\).



We claim the reference endpoints obey

\[Y_{k+1,0}=Y_{k,0}-hg(Y_{k,0})+\sqrt{2h}\,\xi_k-D_k.\]

The lifted gradient is Lipschitz and the reference segment is continuous, so their composition is continuous and Bochner integrable on \([0,h]\).



Apply the Euler-defect identity at \(t=h\), then substitute the segment join and endpoint noise identities. This proves the claimed reference update. Its defect has a minus sign.



Fix \(k\) and compare the actual update with the reference endpoint equation, both using the same \(\xi_k\).



Cancel their identical noise terms. Because the reference subtracts \(D_k\),

\[X_{k+1}-Y_{k+1,0}=\bigl(X_k-hg(X_k)\bigr)-\bigl(Y_{k,0}-hg(Y_{k,0})\bigr)+D_k.\]

Apply the triangle inequality, Euler contraction, and the defect bound. Setting \(e_k=\|X_k-Y_{k,0}\|_2\) gives \(e_{k+1}\le qe_k+b\).



The actual pair supplies a finite-cost coupling, so its transport distance is at most \(e_n\). The second marginal is \(\pi\) by the stationary grid-law assumption.



Apply the scalar recurrence bound to the derived inequality and compose it with the transport comparison.

\subsection{\texorpdfstring{The small-\allowbreak{}step bound}{The small-step bound}}


Suppose \(U\) is differentiable and \(m\)-strongly convex on \(H\), \(\nabla U\) is \(L\)-Lipschitz, and \(0<m\le L\). Let \(0<h\le2/(m+L)\), and supply a reference with the preceding properties. For any initial vector \(X_0\in\mathcal H\),

\[W_2^{\mathbb R}((X_n)_\#P,\pi)\le(1-mh)^n\|X_0-Y_{0,0}\|_2+\frac{33}{20}\frac Lm\sqrt{hd}.\]

This uses a chosen initial coupling and a supplied stationary reference. Replacing the initial cost by \(W_2(\mathcal L(X_0),\pi)\) needs the additional uniform family of reference couplings explained in the final section.



Since \(m+L>0\), clear the positive denominator in the step hypothesis to obtain \(h(m+L)\le2\).



Use \(m\le L\) to get \(2mh\le h(m+L)\le2\), hence \(mh\le1\).



The same threshold gives \(Lh\le2\). Thus \(q=1-mh\ge0\); since \(mh>0\), also \(q<1\).



The sharp small-step contraction of the lifted Euler map gives the contraction hypothesis with \(q=1-mh\). Apply the integral-reference theorem.



Because \(m,h>0\), simplify the bias term:

\[\frac{(33/20)Lh\sqrt{hd}}{1-(1-mh)}=\frac{(33/20)Lh\sqrt{hd}}{mh}=\frac{33}{20}\frac Lm\sqrt{hd}.\]

Substitute the bias identity into the integral-reference bound.

\subsection{\texorpdfstring{The large-\allowbreak{}step bound}{The large-step bound}}


Suppose \(U\) is differentiable and \(m\)-strongly convex on \(H\), \(\nabla U\) is \(L\)-Lipschitz, and \(0<m\le L\). Let \(h>0\), \(h\ge2/(m+L)\), and \(Lh<2\), and supply a reference with the stated properties. For any initial vector \(X_0\in\mathcal H\),

\[W_2^{\mathbb R}((X_n)_\#P,\pi)\le(Lh-1)^n\|X_0-Y_{0,0}\|_2+\frac{33}{20}\frac{Lh}{2-Lh}\sqrt{hd}.\]

The denominator is positive by strict stability. The reference remains an input, and the initial term remains the chosen coupling cost.



Since \(m+L>0\), clear the denominator in the lower step threshold to obtain \(2\le h(m+L)\).



Use \(m\le L\) and \(h\ge0\): \(2\le h(m+L)\le2Lh\). Thus \(Lh\ge1\) and \(q=Lh-1\ge0\).



Strict stability \(Lh<2\) gives \(q=Lh-1<1\).



Apply the sharp large-step Euler contraction with \(q=Lh-1\), then invoke the integral-reference theorem.



Use \(1-(Lh-1)=2-Lh\) and factor out \(\sqrt{hd}\):

\[\frac{(33/20)Lh\sqrt{hd}}{1-(Lh-1)}=\left(\frac{(33/20)Lh}{2-Lh}\right)\sqrt{hd}.\]

Substitute this identity into the integral-reference bound to obtain the large-step formula.

\section{\texorpdfstring{A reference for every initial coupling with one output law}{A reference for every initial coupling with one output law}}
\label{proof-module-UnadjustedLangevin-ProductReference-lean}
{\small\color{gray}Source module: \texttt{Unadjusted\allowbreak{}Langevin/\allowbreak{}Product\allowbreak{}Reference.\allowbreak{}lean}.\par}


This construction turns the analytic reference estimate into an exact initial Wasserstein term by realizing each finite initial coupling on a product with one fixed noise source.

\subsection{\texorpdfstring{pointwise\allowbreak{}Update}{pointwiseUpdate}}


Define the ordinary Euler update \(F_h(x,z)=x-h\nabla U(x)+\sqrt{2h}z\) in an arbitrary Hilbert-space presentation. It is the pointwise function whose lifted composition gives the \(L^2\) numerical chain.

\subsection{\texorpdfstring{pointwise\allowbreak{}Update measurable}{pointwiseUpdate measurable}}


The Lipschitz gradient is continuous. Compose it with the first coordinate projection and combine the remaining continuous scalar and vector operations. The uncurried map \((x,z)\mapsto F_h(x,z)\) is therefore measurable.

\subsection{\texorpdfstring{product\allowbreak{}Noise mem\allowbreak{}Lp}{productNoise memLp}}


The noise projection from \(c\otimes Q\) to \(Q\) is measure preserving because the coupling is a probability. Pull back the supplied \(L^2(Q)\) noise function along this projection. Its square integrability is unchanged.

\subsection{\texorpdfstring{product\allowbreak{}Noise}{productNoise}}


Use the preceding square-integrability proof to form the class of \(((x,y),s)\mapsto W_n(s)\) on the concrete extension. Both the numerical chain and reference use this same \(L^2\) innovation.

\subsection{\texorpdfstring{product\allowbreak{}Noise coe}{productNoise coe}}


The representative selected by \texttt{to\allowbreak{}Lp} agrees almost everywhere with the pulled-back original noise \(W_n(s)\). This identity connects quotient-space noise arithmetic to the actual source random functions.

\subsection{\texorpdfstring{product\allowbreak{}Output\allowbreak{}Law}{productOutputLaw}}


Push \(\mu\otimes Q\) through the pointwise noise-driven chain starting at \(x\) with source path \(W(s)\). This definition omits the reference coupling entirely, making explicit the common algorithm output law for all coupling realizations.

\subsection{\texorpdfstring{product chain coe}{product chain coe}}


Induct through the lifted chain and use the representative identities for initial state, product noise, lifted drift, addition, and scalar multiplication. At each fixed index the result agrees almost everywhere with the ordinary recursion driven by \(x\) and \(W(s)\).

\subsection{\texorpdfstring{product chain output\allowbreak{}Law}{product chain outputLaw}}


Replace the \(L^2\) chain by its almost-everywhere pointwise recursion. That recursion depends only on the algorithm's initial coordinate and noise sample, so \texttt{initial\allowbreak{}Extension\_\allowbreak{}map} gives its coupling-independent law. The result is exactly \texttt{product\allowbreak{}Output\allowbreak{}Law}.

\subsection{\texorpdfstring{product\allowbreak{}Probability}{productProbability}}


A coupling of the initial probability with the reference measure is itself a probability, and its product with the fixed probability noise source is a probability. This local instance supplies normalization on every extension separately.

\subsection{\texorpdfstring{Product\allowbreak{}Reference}{ProductReference}}


For every finite-cost initial coupling, the structure supplies a concrete \texttt{L2\allowbreak{}Reference} on that coupling's product extension, with its initial state matching the second coordinate. Algorithm output laws and initial costs are proved outside the structure; no convergence estimate is a field.

\subsection{\texorpdfstring{Product\allowbreak{}Reference.\allowbreak{}output finite}{ProductReference.output finite}}


Finite initial Wasserstein distance supplies at least one finite-cost initial coupling. On its concrete extension, the numerical and reference variables belong to \(L^2\), giving a finite-cost output coupling. Their established laws identify its marginals with the common output law and \(\pi\), hence the output Wasserstein distance is finite.

\subsection{\texorpdfstring{dalalyan small step of product reference}{dalalyan small step of product reference}}


For each finite-cost initial coupling \(c\), apply the analytic \(L^2\) Dalalyan theorem on \(c\otimes Q\). The initial class-distance theorem identifies its initial error exactly with the real coupling cost, and the common output-law theorem identifies its numerical law.



The resulting bound holds for every finite \(c\), with coefficient \((1-mh)^n\ge0\). Infimizing their costs gives the genuine \(W_2(\mu,\pi)\), without selecting an optimal coupling. Output finiteness is proved first so the final real-valued inequality cannot conceal infinity.

\subsection{\texorpdfstring{dalalyan large step of product reference}{dalalyan large step of product reference}}


Repeat the couplingwise analytic argument for the stable large-step branch, whose contraction factor is \(Lh-1\in[0,1)\). The geometric error sum has denominator \(2-Lh>0\). The finite-coupling infimum theorem replaces the initial class distance by the actual initial Wasserstein distance and yields the displayed large-step bias term.

\section{\texorpdfstring{Actual KLMC output laws on concrete initial-\allowbreak{}coupling extensions}{Actual KLMC output laws on concrete initial-coupling extensions}}
\label{proof-module-KineticLangevin-ProductLaw-lean}
{\small\color{gray}Source module: \texttt{Kinetic\allowbreak{}Langevin/\allowbreak{}Product\allowbreak{}Law.\allowbreak{}lean}.\par}


An estimate for one coupling must eventually be optimized over all initial couplings. Its numerical output law must therefore be independent of that choice. This module proves that independence for the actual KLMC chain. For every coupling \(c\) of the initial position law \(\mu\) and reference law \(\pi\), the sample measure is explicitly \(c\otimes Q\), where \(Q\) supplies the common initial velocity and all phase innovations. The velocity and position components of each innovation stay together, so their correlation is preserved. No convergence or stationary-reference assumption enters these law identities.

\subsection{\texorpdfstring{Build velocity and innovations on each actual product extension}{Build velocity and innovations on each actual product extension}}


The initial coupling has a probability first marginal, hence its joint measure is a probability measure. Its product with the fixed probability source \(Q\) is therefore a probability measure; the local instance records this derived fact. Pull the source velocity \(V\) back through the second projection and take its canonical L\ensuremath{{}^{2}} class. Do the same separately for the two components of each supplied phase innovation \(W_n\). Both use the same source point. This gives the pair of L\ensuremath{{}^{2}} classes needed by the kinetic analysis without assuming that velocity and position noise are independent. The implementation reuses the existing checked product-noise construction.

\subsection{\texorpdfstring{Identify the chosen L\ensuremath{{}^{2}} representatives almost everywhere}{Identify the chosen L² representatives almost everywhere}}


The canonical Lp representative equals the original pulled-back random variable almost everywhere. For velocity this is the existing product-noise representative theorem. For a phase innovation, intersect the two full-measure sets for its components and apply pair extensionality. Thus the L\ensuremath{{}^{2}} pair represents precisely \(W_n(s)\), including all its correlation. This is an almost-everywhere equality of actual random variables, not a comparison of moments or covariances.

\subsection{\texorpdfstring{Define one numerical output law from initialization and fixed source}{Define one numerical output law from initialization and fixed source}}


Given a position \(x\) and source point \(s\), start the original pointwise KLMC recursion at \((V(s),x)\), use its phase innovations \(W_n(s)\), and project the final phase state to position. This is \texttt{product\allowbreak{}KLMCMap}. Its law is the pushforward of \(\mu\otimes Q\), recorded by \texttt{product\allowbreak{}KLMCLaw}. The definition contains no reference coupling. In particular, it is the same candidate output law for every coupling used to estimate its transport error.

\subsection{\texorpdfstring{Prove the L\ensuremath{{}^{2}} numerical chain is exactly the pointwise chain}{Prove the L² numerical chain is exactly the pointwise chain}}


On \(c\otimes Q\), the L\ensuremath{{}^{2}} numerical start consists of the pulled-back common velocity and first initial-position coordinate. The reference position's finite second moment together with finite coupling cost supplies a square-integrable first coordinate. Apply the shared L\ensuremath{{}^{2}}-to-pointwise KLMC chain theorem. A countable intersection supplies agreement of every phase innovation on one full-measure set; intersect it with the initial velocity and position agreements. Pair extensionality identifies the exact initial phase state, function extensionality identifies the entire innovation sequence, and projection to the position coordinate finishes the equality. Equal initial velocities arise by using the same constructed L\ensuremath{{}^{2}} variable; no separate velocity coupling cost is inserted.

\subsection{\texorpdfstring{Derive law independence with almost-\allowbreak{}everywhere measurable data}{Derive law independence with almost-everywhere measurable data}}


The projection \(((x,y),s)\mapsto(x,s)\) sends \(c\otimes Q\) exactly to \(\mu\otimes Q\). Consequently every almost-everywhere measurable function of \((x,s)\) has the same law on both spaces. Lean proves this using the actual projection marginal and the almost-everywhere version of the composition rule for pushforwards. This avoids an unnecessary assumption that arbitrarily selected Lp representatives are everywhere measurable.



The output map is indeed almost-everywhere measurable: the source L\ensuremath{{}^{2}} assumptions give that property for velocity and both innovation components. The force is continuous by Lipschitz regularity, so the KLMC update is measurable. Induction propagates measurability through every phase iterate, and projection gives the position output. The final output-law theorem rewrites the L\ensuremath{{}^{2}} chain with its proved representative equality, then applies the projection-law identity. Both its initial marginal and its numerical law are derived rather than supplied as fields of a reference interface.

\subsection{\texorpdfstring{Check that the common output has total mass one}{Check that the common output has total mass one}}


The independent product of the initial probability law and source law is a probability measure. The preceding almost-everywhere measurability theorem ensures its pushforward remains a probability measure. Lean applies \texttt{Measure.\allowbreak{}is\allowbreak{}Probability\allowbreak{}Measure\_\allowbreak{}map\_\allowbreak{}iff}. This establishes an actual numerical output law before any Wasserstein estimate is attached to it. Canonical Gaussian source construction and the stationary continuous reference are additional layers needed for the final unconditional sampler theorem.

\section{\texorpdfstring{Actual continuous kinetic laws and the true Wasserstein infimum}{Actual continuous kinetic laws and the true Wasserstein infimum}}
\label{proof-module-KineticLangevin-ProductContinuousFlow-lean}
{\small\color{gray}Source module: \texttt{Kinetic\allowbreak{}Langevin/\allowbreak{}Product\allowbreak{}Continuous\allowbreak{}Flow.\allowbreak{}lean}.\par}


The initial position is independent of a common source carrying the initial velocity and continuous forcing. This construction makes the output law independent of any chosen coupling. We then couple two initial position laws, pull their genuine kinetic solutions onto the same source, and apply the proved arbitrary-shift kinetic energy contraction. Taking the infimum over all initial couplings gives the actual Wasserstein estimate.

\subsection{\texorpdfstring{Construct the actual independent-\allowbreak{}source diffusion}{Construct the actual independent-source diffusion}}


The probability space is the product of the initial position law and the common source. The first L\ensuremath{{}^{2}} coordinate is the pulled-back velocity; the second is the original position coordinate. Pulling back the continuous additive forcing gives the genuine kinetic input. The selected Picard solution is then pushed forward at the terminal time to define the position law. The solution specification proves continuity, the exact initial condition and both integral equations. Measurability makes the output a probability measure, and its L\ensuremath{{}^{2}} construction proves a finite second moment.

\subsection{\texorpdfstring{Put any initial coupling on the same source}{Put any initial coupling on the same source}}


Adjoin the same independent source to an arbitrary coupling of the two initial position distributions. The two projections recover the original independent-source spaces exactly, so they preserve measure. Composition of L\ensuremath{{}^{2}} pullbacks shows that both projected solutions receive exactly the same velocity and forcing. The representative formula identifies their initial positions with the two original coupling coordinates.

\subsection{\texorpdfstring{Prove the actual bound before taking an infimum}{Prove the actual bound before taking an infimum}}


The source-pullback theorem transforms each genuine solution equation to the coupling extension. Their common velocity and noise cancel, so the arbitrary-shift L\ensuremath{{}^{2}} contraction theorem applies. The initial position L\ensuremath{{}^{2}} norm equals the actual root-mean-square transport cost of the supplied coupling: the independent extra source contributes a factor one. At the endpoint, measure-preserving projection recovers each independently defined output law. The L\ensuremath{{}^{2}} transport inequality therefore gives the stated bound for this particular initial coupling.

\subsection{\texorpdfstring{Pass to the genuine Wasserstein distance}{Pass to the genuine Wasserstein distance}}


Both initial laws and both output laws have finite second moments. The prefactor is nonnegative because the shifted damping denominator is positive. The proved affine bound over every finite-cost initial coupling therefore passes to the actual infimum defining Wasserstein distance. The result has the paper’s full arbitrary-shift exponential factor and holds for any common square-integrable velocity and any genuine continuous common input, so Brownian forcing is a direct specialization.

\section{\texorpdfstring{The OBABO splitting and its exact pathwise identities}{The OBABO splitting and its exact pathwise identities}}
\label{proof-module-KineticLangevin-OBABO-lean}
{\small\color{gray}Source module: \texttt{Kinetic\allowbreak{}Langevin/\allowbreak{}OBABO.\allowbreak{}lean}.\par}


OBABO brackets a kick–drift–kick Hamiltonian step between two Ornstein–Uhlenbeck half steps. The module proves the composition and its reversal and cancellation identities for arbitrary supplied vectors. The Gaussian-law and convergence analysis is separate. Its thermostat coefficient is chosen so that the damping and fresh noise exactly preserve a target scalar velocity variance.

\subsection{\texorpdfstring{Define each subflow and compose O–B–A–B–O}{Define each subflow and compose O–B–A–B–O}}


In velocity–position ordering, \(O(v,x)=(\delta v+\sigma\xi,x)\), \(B_h(v,x)=(v-hg(x),x)\), and \(A_h(v,x)=(v,x+hv)\). \texttt{obabo\allowbreak{}Step} composes \(O,B_{h/2},A_h,B_{h/2},O\), using distinct noise arguments in the two thermostat steps. Expanding the nested definitions gives \(v_{1/2}=\delta v+\sigma\xi-hg(x)/2\), \(x^+=x+hv_{1/2}\), and \(v^+=\delta(v_{1/2}-hg(x^+)/2)+\sigma\xi'\). \texttt{obabo\allowbreak{}Step\_\allowbreak{}expanded} is proved by \texttt{rfl} because this formula follows directly from the composition; the new-position force is structurally enforced.

\subsection{\texorpdfstring{Reverse the deterministic kick and drift}{Reverse the deterministic kick and drift}}


The kick leaves position fixed, so subtracting \(hg(x)\) and then adding it cancels exactly. The drift leaves velocity fixed, so adding \(hv\) and then subtracting it also cancels. Their inverse formulas hold for arbitrary real \(h\) and arbitrary force, with no smoothness assumption. \texttt{reverse\allowbreak{}Velocity} negates velocity and keeps position; applying it twice is the identity. Conjugating a kick or drift by this reversal changes \(h\) to \(-h\). The proofs use pair extensionality and simplification of scalar signs, making the time-reversal algebra explicit.

\subsection{\texorpdfstring{The palindromic Hamiltonian composition has an exact inverse}{The palindromic Hamiltonian composition has an exact inverse}}


\texttt{verlet\allowbreak{}Step} is \(B_{h/2}A_hB_{h/2}\). To invert it, reverse the order and invert each factor. Since the composition is palindromic, this is precisely \(B_{-h/2}A_{-h}B_{-h/2}\). Lean expands the two compositions, rewrites \((-h)/2=-(h/2)\), and cancels the three adjacent inverse pairs with the preceding lemmas. The OBABO zero-time theorem separately checks that \(h=0,\delta=1,\sigma=0\) makes every substep the identity. Reversal of the deterministic map alone does not assert reversibility of a stochastic transition with a particular probability density.

\subsection{\texorpdfstring{Synchronous noises cancel while the new forces remain}{Synchronous noises cancel while the new forces remain}}


Apply two OBABO steps with the same \(\xi,\xi'\). The difference of half velocities is \(a=\delta\Delta v-h(g(x)-g(y))/2\). Thus \(\Delta x^+=\Delta x+ha\). The final shared refresh cancels too, leaving \(\Delta v^+=\delta(a-h(g(x^+)-g(y^+))/2)\). Lean unfolds all subflows, proves the velocity and position identities as a conjunction, and normalizes vector expressions with \texttt{module}. Although additive noises disappear explicitly from the difference, the new positions still depend on the common first noise. The identity preserves their force evaluations instead of treating them as old-position gradients.

\subsection{\texorpdfstring{Choose damping and prove the thermal variance identity}{Choose damping and prove the thermal variance identity}}


The OU half-step damping is \(\delta=e^{-\gamma h/2}>0\). The exponential addition identity gives \(\delta^2=e^{-\gamma h}\). With \(\gamma,h\ge0\), this is at most one, so \(s^2=1-\delta^2\ge0\). For a target velocity variance \(u\ge0\), set \(\sigma=\sqrt{u(1-\delta^2)}\). Then \(\delta^2u+\sigma^2=u\). Lean proves positivity directly from \texttt{exp\_\allowbreak{}pos}, rewrites the square with \texttt{exp\_\allowbreak{}add}, uses the nonpositive exponent to bound the exponential, and finally applies \texttt{Real.\allowbreak{}sq\_\allowbreak{}sqrt} followed by \texttt{ring}. For independent centered input velocity and fresh Gaussian noise this is the scalar covariance balance; Gaussian-law preservation additionally uses those distributional facts.

\section{\texorpdfstring{Exact Hamiltonian energy and volume-\allowbreak{}preserving Verlet geometry}{Exact Hamiltonian energy and volume-preserving Verlet geometry}}
\label{proof-module-KineticLangevin-Hamiltonian-lean}
{\small\color{gray}Source module: \texttt{Kinetic\allowbreak{}Langevin/\allowbreak{}Hamiltonian.\allowbreak{}lean}.\par}


The Hamiltonian subflow is the kinetic equation with zero friction and zero noise. Its total energy is conserved by an exact cancellation between force work and kinetic-energy change. This module proves conservation for the actual Picard-constructed integral flow and proves that the explicit kick/\allowbreak{}drift/\allowbreak{}Verlet maps preserve product Haar volume. These are genuine prerequisites for Gibbs stationarity; energy conservation alone and volume preservation of a numerical map do not yet establish invariance of the full kinetic diffusion.

\subsection{\texorpdfstring{Differentiate total energy along the actual Hamiltonian equations}{Differentiate total energy along the actual Hamiltonian equations}}


Define \(H(v,x)=f(x)+\|v\|^2/2\). If \(x'=v\) and \(v'=-\nabla f(x)\), the chain rule gives \((f(x))'=\langle\nabla f(x),v\rangle\), while the squared-norm derivative gives \((\|v\|^2/2)'=-\langle v,\nabla f(x)\rangle\). Symmetry of the real inner product makes their sum zero. Lean composes the actual Fréchet derivative of \(f\) with the position derivative, uses the gradient representation of that derivative, and combines it with the Hilbert squared-norm derivative. The statement uses right derivatives within \([t,\infty)\), allowing the initial endpoint of an integral solution.

\subsection{\texorpdfstring{Turn vanishing right derivatives into exact conservation}{Turn vanishing right derivatives into exact conservation}}


A continuous scalar function with zero right derivative on \([0,T)\) is constant on \([0,T]\), including the final endpoint. The proof applies the scalar Gronwall comparison with coefficient and forcing both zero, first to \(F\) and then to \(-F\). These give \(F(t)\le F(0)\) and \(-F(t)\le-F(0)\). Their combination proves equality. This avoids assuming a derivative of the chosen extension beyond the terminal time or on negative time.

\subsection{\texorpdfstring{Conserve energy for the constructed Hamiltonian integral flow}{Conserve energy for the constructed Hamiltonian integral flow}}


\texttt{hamiltonian\allowbreak{}Flow} is the already constructed \texttt{kinetic\allowbreak{}Flow} with friction zero, inverse mass one, and the zero compact forcing path. Its specification reduces to \(Y_t=Y_0-\int_0^tD(Y_s)ds\), with \(D(v,x)=(\nabla f(x),-v)\). Continuity of the Lipschitz drift gives the integral derivative. Lean identifies the selected path with its differentiable primitive on a right neighborhood, projects its derivative onto velocity and position, and applies the energy derivative cancellation. The scalar constancy theorem then yields \(H(Y_t)=H(Y_0)\). This is an exact statement about an existing solution, not an energy-preservation hypothesis on an abstract flow.

\subsection{\texorpdfstring{Build the reversible proposal and preserve the Hamiltonian under reversal}{Build the reversible proposal and preserve the Hamiltonian under reversal}}


Velocity reversal \(R(v,x)=(-v,x)\) conjugates every kick and drift to its negative-time inverse. Applying this to the palindromic kick–drift–kick composition gives \(R\,V_h\,R=V_{-h}\). Since \(V_{-h}V_h=\mathrm{id}\), the proposal \(J_h=R\,V_h\) satisfies \(J_h^2=\mathrm{id}\). Lean rewrites the three component conjugacies and then uses the proved inverse Verlet identity. Negating velocity leaves its squared norm unchanged, so \(H(Rz)=H(z)\). The exact involution will give symmetric old/\allowbreak{}new proposal graph measures in the Metropolis argument.

\subsection{\texorpdfstring{Expose the exact numerical energy defect}{Expose the exact numerical energy defect}}


Let \(w=v-hg(x)/2\), \(y=x+hw\), and \(v^+=w-hg(y)/2\). Expanding both kinetic-energy squares gives

\[H(v^+,y)-H(v,x)=f(y)-f(x)-\frac h2\langle w,g(x)+g(y)\rangle
+\frac{h^2}{8}\bigl(\|g(y)\|^2-\|g(x)\|^2\bigr).\]

The first part is a trapezoidal force-work remainder. The second is an explicit force-norm correction. Lean first proves the two-half-kick norm identity for arbitrary vectors using squared-norm and inner-product expansions, then substitutes the actual intermediate Verlet velocity and position. This identity preserves the error terms that must later be bounded to control rejection; no second-order error estimate is assumed by definition.

\subsection{\texorpdfstring{Prove the elementary subflows preserve Haar volume}{Prove the elementary subflows preserve Haar volume}}


A drift \((v,x)\mapsto(v,x+hv)\) is a skew product whose second component is a translation for each fixed velocity. Translation invariance of Haar measure therefore proves preservation of the product measure. A kick \((v,x)\mapsto(v-hg(x),x)\) is the same construction after swapping the two phase coordinates; measurability of \(g\) suffices. Lean applies the measure-preserving skew-product theorem to translations, then conjugates the kick by the measure-preserving swap. Composition proves the complete Verlet map preserves product Haar volume. This derives the Jacobian-one property without presupposing differentiability of the force.

\subsection{\texorpdfstring{The reversed Verlet proposal is a measure-\allowbreak{}preserving involution}{The reversed Verlet proposal is a measure-preserving involution}}


For a Haar factor also invariant under negation, velocity reversal preserves product measure: apply the negation-preserving map in the velocity coordinate and identity in the position coordinate. Composing this with the proved volume-preserving Verlet map gives volume preservation of the explicit proposal \(J_h\). Together with \(J_h^2=\mathrm{id}\), this supplies the two exact properties used by deterministic Metropolis detailed balance. Neither property asserts that the uncorrected finite-step Verlet map itself preserves the Gibbs weight.

\section{\texorpdfstring{Euler approximations converge to the integral solution}{Euler approximations converge to the integral solution}}
\label{proof-module-UnadjustedLangevin-EulerLimit-lean}
{\small\color{gray}Source module: \texttt{Unadjusted\allowbreak{}Langevin/\allowbreak{}Euler\allowbreak{}Limit.\allowbreak{}lean}.\par}


This module proves a deterministic convergence theorem for additive Euler discretization. With mesh \(h=T/n\), its terminal value converges to the continuous integral solution, including when the state space itself is \(L^2(P)\).

\subsection{\texorpdfstring{additive\allowbreak{}Euler\allowbreak{}Chain}{additiveEulerChain}}


Define \(X_0=x\) and \(X_{k+1}=X_k-hg(X_k)+B((k+1)h)-B(kh)\). The additive input is one fixed function evaluated at mesh endpoints; consequently the numerical chain and the integral solution can be compared with exactly matching noise increments.

\subsection{\texorpdfstring{additive\allowbreak{}Euler\allowbreak{}Chain zero}{additiveEulerChain zero}}


Evaluating the recursive definition at index zero gives \(X_0=x\). Lean proves this by definitional equality; no continuity of the input or positivity of the mesh is required for this initialization identity.

\subsection{\texorpdfstring{additive\allowbreak{}Euler\allowbreak{}Chain succ}{additiveEulerChain succ}}


Unfold the successor clause once to expose \(X_{k+1}=X_k-hg(X_k)+\Delta B_k\). This definitional identity is the equation used when subtracting the exact solution's increment and deriving the error recurrence.

\subsection{\texorpdfstring{langevin integral increment}{langevin integral increment}}


Subtract the integral equations at \(b\) and \(a\). Continuity of \(g\circ Y\) makes it interval integrable, so interval-integral additivity gives \(\int_0^b-\int_0^a=\int_a^b\). Thus \(Y(b)-Y(a)=-\int_a^b g(Y(s))ds+B(b)-B(a)\), including oriented intervals when \(b<a\).

\subsection{\texorpdfstring{finite euler recurrence bound}{finite euler recurrence bound}}


From \(e_{k+1}\le qe_k+b\), \(e_0=0\), \(q\ge1\), and \(b\ge0\), induction proves \(e_n\le nbq^n\). In the successor step bound \(q\cdot kbq^k+b\) by \((k+1)bq^{k+1}\), using \(q^{k+1}\ge1\). This coarse bound avoids dividing by \(q-1\), so it also covers \(q=1\).

\subsection{\texorpdfstring{explicit\allowbreak{}Euler norm sub le}{explicitEuler norm sub le}}


Write the difference as \((x-y)-h(g(x)-g(y))\). The triangle inequality and Lipschitzness give \(\|\cdot\|\le\|x-y\|+hL\|x-y\|=(1+Lh)\|x-y\|\). The hypothesis \(h\ge0\) removes the absolute value in the scalar norm.

\subsection{\texorpdfstring{additive\allowbreak{}Euler\allowbreak{}Chain error le of drift modulus}{additiveEulerChain error le of drift modulus}}


At one step the common noise cancels. The Euler map contributes \((1+Lh)e_k\), and the exact integral remainder contributes at most \(h\varepsilon\), because the drift changes by at most \(\varepsilon\) within one mesh interval. Apply the finite recurrence bound and \((1+Lh)^n\le e^{Lnh}\), then use \(nh=T\) to obtain \(e_n\le T\varepsilon e^{LT}\).

\subsection{\texorpdfstring{additive\allowbreak{}Euler\allowbreak{}Chain tendsto}{additiveEulerChain tendsto}}


The continuous drift path \(g\circ Y\) is uniformly continuous on \([0,T]\). Given any positive error tolerance, choose a drift modulus \(\varepsilon\) making \(T\varepsilon e^{LT}\) small, then take the mesh \(T/(n+1)\) below its continuity radius. The previous error bound gives endpoint convergence. The case \(T=0\) is separated so no division by the horizon is needed.

\subsection{\texorpdfstring{lp langevin integral euler tendsto}{lp langevin integral euler tendsto}}


Apply the deterministic matching-input Euler convergence theorem in \(L^2(P;H)\), using the Lipschitz lifted gradient. The numerical endpoint converges in the \(L^2\) norm to the constructed solution at \(T\). This is the topology later used to pass smooth bounded tests and identify laws.

\section{\texorpdfstring{Prove Verlet converges to the actual Hamiltonian flow}{Prove Verlet converges to the actual Hamiltonian flow}}
\label{proof-module-KineticLangevin-HamiltonianConvergence-lean}
{\small\color{gray}Source module: \texttt{Kinetic\allowbreak{}Langevin/\allowbreak{}Hamiltonian\allowbreak{}Convergence.\allowbreak{}lean}.\par}


Exact reversibility and volume preservation of the Verlet map do not by themselves identify the continuous Hamiltonian flow. This module supplies that analytic limit. For a globally Lipschitz force, the actual Verlet iterates on a mesh of size \(T/(n+1)\) converge to every continuous Hamiltonian integral solution at time \(T\). Specializing to the gradient force identifies the limit with the Picard-selected \texttt{hamiltonian\allowbreak{}Flow}. The proof uses a genuine local state-error calculation, uniform growth of an Euler comparison chain, a discrete Grönwall estimate, and the existing Euler convergence theorem. It needs no Hessian assumption.

\subsection{\texorpdfstring{Control the elementary shears and their symmetric composition}{Control the elementary shears and their symmetric composition}}


Use the standard maximum norm on phase pairs \((v,x)\). A kick of duration \(h\ge0\) is \((1+Mh)\)-Lipschitz when the force is \(M\)-Lipschitz, and a drift is \((1+h)\)-Lipschitz. The two half kicks around one drift therefore have factor \((1+Mh/2)^2(1+h)\). For \(0\le h\le1\), expansion and \(h^2,h^3\le h\) bound it by

\[\|\operatorname{Verlet}_h(z)-\operatorname{Verlet}_h(w)\|\le\bigl(1+(M+1)^2h\bigr)\|z-w\|.\]

Lean first bounds the velocity and position projections by the phase norm. It rewrites each difference so the shared force difference is explicit, applies the triangle inequality and the force Lipschitz estimate, and handles the product norm using \texttt{max\_\allowbreak{}le}. The Verlet estimate composes the three bounds; a scalar polynomial inequality controls the resulting product. This is a stability estimate for the actual symmetric map, not an assumed numerical property.

\subsection{\texorpdfstring{Compute and bound the actual local Verlet state defect}{Compute and bound the actual local Verlet state defect}}


Let \(v_{1/2}=v-hg(x)/2\) and \(x_+=x+hv_{1/2}\). Compare Verlet with the explicit Hamiltonian Euler step \((v-hg(x),x+hv)\). Their exact coordinate differences are

\[\Delta v=-\frac h2\bigl(g(x_+)-g(x)\bigr),\qquad\Delta x=-\frac{h^2}{2}g(x).\]

The Lipschitz bound gives \(\|g(x)\|\le(M+1)(\|z\|+\|g(0)\|)\) and \(\|v_{1/2}\|\le(M+2)(\|z\|+\|g(0)\|)\) for \(h\le1\). Consequently,

\[\|\operatorname{Verlet}_h(z)-(z-hD(z))\|\le h^2(M+1)(M+2)(\|z\|+\|g(0)\|),\quad D(v,x)=(g(x),-v).\]

Lean proves the two displayed coordinate identities by module algebra. It then applies the force Lipschitz estimate to \(x_+-x=hv_{1/2}\), rewrites scalar-multiple norms, and bounds each coordinate before taking their maximum. This deliberately sufficient quadratic local estimate proves convergence under only a Lipschitz force; it does not claim a second-order global rate under those minimal assumptions.

\subsection{\texorpdfstring{Bound the entire Euler comparison chain on a fixed horizon}{Bound the entire Euler comparison chain on a fixed horizon}}


Write \(X_{k+1}=X_k-hd(X_k)\), with \(d\) globally \(L\)-Lipschitz and \(nh=T\). The distance from the actual initial point satisfies

\[e_{k+1}\le(1+Lh)e_k+h\|d(X_0)\|,\qquad e_0=0.\]

The finite recurrence estimate and \((1+Lh)^k\le e^{Lkh}\le e^{LT}\) yield, for every \(k\le n\),

\[\|X_k\|\le\|X_0\|+T\|d(X_0)\|e^{LT}.\]

The bound is uniform over the number of mesh steps. Lean derives the recurrence by comparing one Euler update with the update of the fixed initial point, invoking the previously proved Lipschitz bound for explicit Euler. \texttt{finite\_\allowbreak{}euler\_\allowbreak{}recurrence\_\allowbreak{}bound} sums the recurrence. The inequality \(1+r\le e^r\), the exponential power identity, and the partition equation \(nh=T\) control the amplification. No boundedness of the numerical trajectory is added as a hypothesis.

\subsection{\texorpdfstring{Propagate the local difference between Verlet and Euler}{Propagate the local difference between Verlet and Euler}}


Let \(V_k\) be the actual Verlet iterates and \(X_k\) the actual Euler iterates from the same initial phase. Split their next difference into a Verlet stability term and the local Verlet-versus-Euler defect evaluated at \(X_k\). The proved Euler growth bound makes the local coefficient uniform in \(k\). With \(A=(M+1)^2\) and

\[C=(M+1)(M+2)\bigl(\|z\|+T\|D(z)\|e^{(M+1)T}+\|g(0)\|\bigr),\]

the actual error satisfies \(e_{k+1}\le(1+Ah)e_k+h^2C\), hence

\[\|V_n-X_n\|\le ThC e^{AT}.\]

Lean constructs the two concrete chains and proves their recursive update formulas. The triangle inequality combines the established stability and local-defect theorems. The finite recurrence bound gives \(nh^2C(1+Ah)^n\); the partition identity converts \(nh^2\) into \(Th\), and the exponential bound removes dependence on \(n\). This is the full finite-horizon numerical comparison required by the limit theorem.

\subsection{\texorpdfstring{Send the mesh size to zero}{Send the mesh size to zero}}


Choose \(h_n=T/(n+1)\). This mesh tends to zero, and eventually \(h_n\le1\), so the finite comparison applies. Its right side is the vanishing mesh times a fixed, finite coefficient depending only on the horizon, force, and initial phase. Therefore the norm of the difference between the two actual endpoints tends to zero.



Lean proves the mesh limit from \texttt{tendsto\_\allowbreak{}const\_\allowbreak{}div\_\allowbreak{}at\allowbreak{}Top\_\allowbreak{}nhds\_\allowbreak{}zero\_\allowbreak{}nat}, obtains the eventual small-step condition from the neighborhood of zero, and applies \texttt{squeeze\_\allowbreak{}zero'} to the nonnegative error norm. \texttt{tendsto\_\allowbreak{}iff\_\allowbreak{}norm\_\allowbreak{}sub\_\allowbreak{}tendsto\_\allowbreak{}zero} upgrades this scalar estimate to convergence of the phase-vector difference itself.

\subsection{\texorpdfstring{Identify the selected Hamiltonian flow as the numerical limit}{Identify the selected Hamiltonian flow as the numerical limit}}


The existing Euler theorem states that a globally Lipschitz drift's matching-input Euler approximations converge to every continuous solution of its actual integral equation. Here the additive input is zero and the negative Hamiltonian drift is \(D(v,x)=(g(x),-v)\), which has Lipschitz constant at most \(M+1\). Adding the vanishing Verlet-minus-Euler difference to that convergent Euler endpoint proves convergence of Verlet itself.



For \(g=\nabla f\), the selected \texttt{hamiltonian\allowbreak{}Flow} is already constructed by Picard from this precise integral equation. \texttt{kinetic\allowbreak{}Flow\_\allowbreak{}spec} supplies its continuity, initialization, and integral identity. The final theorem therefore proves

\[\operatorname{Verlet}_{T/(n+1)}^{\,n+1}(z)\longrightarrow\operatorname{hamiltonianFlow}(f,T,z)(T).\]

Lean invokes \texttt{additive\allowbreak{}Euler\allowbreak{}Chain\_\allowbreak{}tendsto} with zero input, combines the two limits using \texttt{Tendsto.\allowbreak{}add}, and simplifies the exact vector identity \((V_n-X_n)+X_n=V_n\). The specialization unfolds only the zero-friction, unit-inverse-mass flow definition and its normalized zero forcing. Neither the desired numerical convergence nor a stationarity conclusion occurs among the hypotheses.

\section{\texorpdfstring{Concrete Verlet–OU splitting and genuine local defects}{Concrete Verlet–OU splitting and genuine local defects}}
\label{proof-module-KineticLangevin-KineticSplitting-lean}
{\small\color{gray}Source module: \texttt{Kinetic\allowbreak{}Langevin/\allowbreak{}Kinetic\allowbreak{}Splitting.\allowbreak{}lean}.\par}


The splitting is an explicit numerical map: velocity Verlet for the Hamiltonian drift, then exponential velocity damping and a supplied filtered additive increment. Its stability and local kinetic Euler defect are proved from the map algebra. The filtered increment is written by deterministic integration by parts, matching the canonical Brownian segment identity proved separately.

\subsection{\texorpdfstring{Damp velocity while retaining position}{Damp velocity while retaining position}}


The map is \(D_\delta(v,x)=(\delta v,x)\). For \(0\le\delta\le1\), its velocity difference has norm \(\delta\|v-w\|\) and its position difference is unchanged. Each is bounded by the original maximum product norm, hence the full map is nonexpansive. Lean proves the two coordinate bounds and combines them using \texttt{max\_\allowbreak{}le}. This uses the actual usual product norm rather than an undeclared Euclidean phase norm.

\subsection{\texorpdfstring{Construct and stabilize the actual Verlet–OU step}{Construct and stabilize the actual Verlet–OU step}}


The explicit step is \(S_h(z,N)=O_{e^{-\gamma h},N}(V_hz)\). Two synchronously driven steps have exactly the same additive increment, which cancels. Exponential damping lies in \((0,1]\) for nonnegative friction and step, so it is nonexpansive. The already proved Verlet stability bound gives \[\|S_h(z,N)-S_h(w,N)\|\le[1+(M+1)^2h]\|z-w\|\] for \(h\le1\). Lean first checks the exact noise-cancellation identity and then composes these two genuine bounds.

\subsection{\texorpdfstring{Derive exponential damping errors without dividing by friction}{Derive exponential damping errors without dividing by friction}}


The response balances \(\gamma\psi_1=1-e^{-\gamma h}\) and \(\gamma\psi_2=h-\psi_1\) imply \(e^{-\gamma h}-1+\gamma h=\gamma^2\psi_2\). Since \(0\le\psi_2\le h^2/2\), the absolute Euler remainder is bounded by \(\gamma^2h^2/2\). Also the exponential tangent inequality gives \(0\le1-e^{-\gamma h}\le\gamma h\). Lean obtains both estimates from exact identities and checks all sign conditions. The formulas remain valid at zero friction.

\subsection{\texorpdfstring{Compare the damped Hamiltonian Euler stage with full kinetic Euler}{Compare the damped Hamiltonian Euler stage with full kinetic Euler}}


Damp the Hamiltonian Euler stage and subtract full kinetic Euler. Their position components agree exactly, while their velocity difference is \((e^{-\gamma h}-1+\gamma h)v+(1-e^{-\gamma h})h\,g(x)\). The preceding scalar estimates bound this by a quadratic expression. Lipschitz force growth \(\|g(x)\|\le M\|z\|+\|g(0)\|\) yields the displayed \(h^2[\gamma^2/2+\gamma(M+1)](\|z\|+\|g(0)\|)\) bound. Lean proves the exact vector identity, bounds its two terms, and combines nonnegative scalar factors.

\subsection{\texorpdfstring{Bound the complete deterministic kinetic Euler defect}{Bound the complete deterministic kinetic Euler defect}}


Insert the damped Hamiltonian Euler stage between the actual split step and full kinetic Euler. Nonexpansiveness of damping preserves the proved Verlet-versus-Hamiltonian-Euler defect. Add the damped-stage defect just established. The noise increment is retained exactly on both sides and cancels, giving \[\|S_h(z,N)-[z-hd(z)+(N,0)]\|\le h^2C(\|z\|+\|g(0)\|),\quad C=(M+1)(M+2)+\gamma^2/2+\gamma(M+1).\] Lean applies the triangle inequality through this explicit intermediate state and normalizes the coefficient algebra.

\subsection{\texorpdfstring{Use the actual filtered additive increment}{Use the actual filtered additive increment}}


For a continuous additive velocity input \(B\), define \[N_{a,h}=B_{a+h}-B_a-\gamma\int_0^h e^{-\gamma(h-s)}(B_{a+s}-B_a)\,ds.\] This is the deterministic integration-by-parts representation of the exact OU response. If the input oscillation over the interval is at most \(\eta\), exponential damping is at most one, so \(\|N_{a,h}-(B_{a+h}-B_a)\|\le\gamma h\eta\). Lean subtracts the matching increment exactly, moves the scalar outside the norm, and bounds the actual integral by its interval length and pointwise integrand bound. The Brownian segment module independently proves precisely this identity for the constructed Gaussian source.

\section{\texorpdfstring{The actual Verlet–OU chain converges to the kinetic integral solution}{The actual Verlet–OU chain converges to the kinetic integral solution}}
\label{proof-module-KineticLangevin-KineticSplittingConvergence-lean}
{\small\color{gray}Source module: \texttt{Kinetic\allowbreak{}Langevin/\allowbreak{}Kinetic\allowbreak{}Splitting\allowbreak{}Convergence.\allowbreak{}lean}.\par}


This module proves the missing numerical limit directly in a Banach space, including the intended space of square-integrable random variables. The actual scheme uses explicit Verlet stages, exact exponential damping, and the filtered increments of the same continuous additive input as the actual kinetic equation. Compact continuity supplies every local modulus used in the proof. No stationary solution, local-error contract, or numerical-convergence hypothesis is assumed.

\subsection{\texorpdfstring{Construct the splitting on the actual grid input}{Construct the splitting on the actual grid input}}


Starting from \(X_0=z_0\), the next state is \(X_{k+1}=S_h(X_k,N_{kh,h})\), where \(S_h\) is the explicit Verlet–OU step and \(N_{a,h}\) is the actual filtered increment of the supplied continuous input \(B\). Lean defines this recursion directly. Each stage uses the correct grid time \(kh\); the forcing is not replaced by an unrelated Gaussian with the same marginal variance.

\subsection{\texorpdfstring{Compare one numerical increment with the exact kinetic path}{Compare one numerical increment with the exact kinetic path}}


Let the actual path satisfy \(Y_t-z_0=-\int_0^t d(Y_s)\,ds+(B_t,0)\), where \(d(v,x)=(\gamma v+g(x),-v)\). Subtract its equations at \(a\) and \(a+h\). Freezing its drift gives an exact Euler identity with remainder \(D_{a,h}=\int_0^h[d(Y_{a+s})-d(Y_a)]\,ds\). A drift oscillation bound \(\eta\) gives \(\|D_{a,h}\|\le h\eta\). The matching OU input differs from \(B_{a+h}-B_a\) by at most \(\gamma h\eta\) when the input oscillation is also at most \(\eta\). Finally, the proved explicit splitting defect contributes at most \(h^2C(R+\|g(0)\|)\), where \(\|Y_t\|\le R\) and \(C=(M+1)(M+2)+\gamma^2/2+\gamma(M+1)\). Lean constructs the grid and interval-membership facts, derives the exact increment equation through adjacent-interval integration, proves the local-defect identity, and bounds each of these three genuine terms.



For \(e_k=\|X_k-Y_{kh}\|\), synchronous numerical stability then gives \[e_{k+1}\le[1+(M+1)^2h]e_k+h\{hC(R+\|g(0)\|)+(1+\gamma)\eta\}.\] The actual initial equation and \(B_0=0\) prove \(e_0=0\). The checked finite recurrence estimate and \((1+ah)^n\le e^{aT}\), with \(nh=T\), yield the displayed finite-horizon bound. Every modulus premise belongs to this intermediate estimate and is derived from continuity in the next theorem.

\subsection{\texorpdfstring{Derive all moduli and prove convergence}{Derive all moduli and prove convergence}}


The continuous exact path is bounded on \([0,T]\). Its drift path is continuous by global Lipschitz continuity of \(g\), and the additive input \(B\) is continuous. Compactness makes both uniformly continuous. Given any target error \(\varepsilon>0\), choose a common oscillation tolerance \(\eta=\varepsilon/(D+1)\), where the explicit finite horizon constant \(D\) contains the state bound, the splitting defect coefficient, friction, and the stability exponential. Since \(h_n=T/(n+1)\to0\), the mesh is eventually at most one, at most \(\eta\), and below both continuity radii. The finite error estimate is then at most \(\eta D<\varepsilon\). Lean obtains the two compact uniform-continuity radii, combines the four eventual mesh bounds, instantiates the actual finite estimate, and checks this final scalar comparison. Thus the concrete scheme converges to every actual continuous kinetic integral solution driven by the same input.

\subsection{\texorpdfstring{Identify the limit with the constructed kinetic flow}{Identify the limit with the constructed kinetic flow}}


The earlier Picard construction supplies a continuous solution for every initial state and every compact continuous forcing path, with its actual kinetic integral identity already proved. Its normalized forcing is globally continuous and vanishes at zero. Substituting these checked facts into the generic Banach convergence theorem proves that the splitting limit is exactly \texttt{kinetic\allowbreak{}Flow} at the requested horizon. This ties the approximation to the existing jointly continuous solution map, not to an unspecified solution chosen by a new contract.

\section{\texorpdfstring{The L\ensuremath{{}^{2}} splitting is the actual vector algorithm almost everywhere}{The L² splitting is the actual vector algorithm almost everywhere}}
\label{proof-module-KineticLangevin-KineticSplittingLp-lean}
{\small\color{gray}Source module: \texttt{Kinetic\allowbreak{}Langevin/\allowbreak{}Kinetic\allowbreak{}Splitting\allowbreak{}Lp.\allowbreak{}lean}.\par}


The deterministic convergence theorem applies to a Banach space of square-integrable random variables. To connect that analysis to the intended stochastic method, this module checks every elementary stage on actual representatives. The lifted force is the original force almost everywhere, and the supplied noise is retained unchanged. Induction then identifies the entire L\ensuremath{{}^{2}} recursion with the vector recursion.

\subsection{\texorpdfstring{Expose the chain driven by explicit filtered innovations}{Expose the chain driven by explicit filtered innovations}}


The innovation chain starts at its supplied phase state and recursively applies the explicit Verlet–OU step with innovation \(N_k\). The earlier continuous-input chain is exactly this chain with \(N_k=N_{kh,h}(B)\). Lean proves this equality by induction, unfolding the two actual recursions at the successor step. This representation lets the probability layer identify the innovation array law independently of the deterministic input analysis.

\subsection{\texorpdfstring{Check force kicks on actual L\ensuremath{{}^{2}} representatives}{Check force kicks on actual L² representatives}}


The lifted force \(G\) is constructed from the actual Lipschitz force \(g\), and the shared lifting theorem gives \(G(X)(\omega)=g(X(\omega))\) almost everywhere. The representative identities for subtraction and scalar multiplication therefore give \([V-hG(X)](\omega)=V(\omega)-hg(X(\omega))\). Position is unchanged. Lean intersects the three almost-everywhere events and proves equality of the two phase coordinates. No new vector-field or output-law assumption is introduced.

\subsection{\texorpdfstring{Check free transport and the full supplied thermostat noise}{Check free transport and the full supplied thermostat noise}}


The representative of \(X+hV\) is \(X(\omega)+hV(\omega)\), proving the free drift identity. Likewise the representative of \(\delta V+\sigma\Xi\) is \(\delta V(\omega)+\sigma\Xi(\omega)\), proving the OU stage identity. Lean explicitly intersects the addition and scalar-multiplication representative events and retains the exact supplied innovation \(\Xi\). These are pathwise stage identities and do not presume any Gaussian distribution.

\subsection{\texorpdfstring{Compose the checked stages into the actual Verlet–OU map}{Compose the checked stages into the actual Verlet–OU map}}


The Verlet step is the kick–drift–kick composition. Apply the checked representative identity at each actual intermediate L\ensuremath{{}^{2}} state, intersect the resulting almost-everywhere events, and substitute them in order. The complete splitting then composes this checked Verlet map with the checked OU stage, using exactly the damping \(e^{-\gamma h}\) and supplied filtered innovation. Lean performs these successive substitutions, so the state at which the second force is evaluated remains the updated position.

\subsection{\texorpdfstring{Identify every numerical iterate with the original vector chain}{Identify every numerical iterate with the original vector chain}}


The initial phase representatives agree by definition. Assuming agreement at step \(n\), the checked full-step identity and the induction hypothesis give agreement at step \(n+1\), driven by the same innovation representative \(N_n(\omega)\). Lean intersects these two almost-everywhere events and rewrites the recursive definitions. Hence every L\ensuremath{{}^{2}} iterate analyzed by the convergence theorem is the actual vector splitting algorithm on the same probability source.

\section{\texorpdfstring{Construct correlated Gaussian blocks and an independent noise stream}{Construct correlated Gaussian blocks and an independent noise stream}}
\label{proof-module-KineticLangevin-Gaussian-lean}
{\small\color{gray}Source module: \texttt{Kinetic\allowbreak{}Langevin/\allowbreak{}Gaussian.\allowbreak{}lean}.\par}


A covariance matrix alone is insufficient to define a random algorithm. This module constructs actual centered finite-dimensional Gaussian probability measures with the Gram covariance and an infinite product giving one independent block per time and spatial coordinate. Correlation is retained within each block. The exact covariance theorem requires a nonnegative step and continuous kernels; the general Gaussian constructor remains a probability measure even outside that parameter regime.

\subsection{\texorpdfstring{Construct the centered Gaussian probability measure and its vector moments}{Construct the centered Gaussian probability measure and its vector moments}}


For a finite kernel index set, \texttt{kernel\allowbreak{}Gaussian} is Mathlib's multivariate Gaussian with mean zero and the defined Gram matrix. The named instances establish Gaussianity and total probability mass using the distribution constructor. The mean theorem computes its vector expectation as zero, and \texttt{kernel\allowbreak{}Gaussian\_\allowbreak{}mem\allowbreak{}Lp} invokes the finite Gaussian second-moment theorem. These are derived properties of a specified probability measure, rather than moment fields of an abstract noise source. They hold even if the covariance argument is degenerate; identifying its entries with the requested matrix uses the positive-semidefinite theorem below.

\subsection{\texorpdfstring{Evaluate coordinate means, moments and exact covariances}{Evaluate coordinate means, moments and exact covariances}}


Each coordinate projection is a continuous linear functional. Commuting it with the vector integral proves zero coordinate mean, and composing the finite-second-moment identity with that functional proves coordinate L\ensuremath{{}^{2}} integrability. For continuous kernels and \(h\ge0\), Gram positivity permits the multivariate-Gaussian covariance formula:

\[\operatorname{Cov}(X_i,X_j)=\int_0^hk_i(t)k_j(t)dt,\quad
\operatorname{Var}(X_i)=\int_0^hk_i(t)^2dt.\]

The two distribution theorems receive the already proved positive-semidefinite matrix witness. The variance proof only rewrites a product as a square. This identifies the exact correlated coordinate law required by the integrators, while its representation as actual Brownian stochastic integrals is a separate identity.

\subsection{\texorpdfstring{Instantiate the two-\allowbreak{}block and four-\allowbreak{}block kinetic laws}{Instantiate the two-block and four-block kinetic laws}}


\texttt{klmc\allowbreak{}Gaussian} uses the kernel vector \((\psi_0,\psi_1)\); \texttt{klmc2\allowbreak{}Gaussian} uses \((\psi_0,\psi_1,\varphi_2,\varphi_3)\). They are therefore concrete measures on two- and four-dimensional Euclidean space. Each probability instance unfolds the specialization and reuses the general Gaussian probability instance. The first law correlates a coordinate's velocity and position increments; the second also couples the Hessian-correction integrals. Separate spatial coordinates will use separate copies of these entire blocks, rather than separate copies of their individual components.

\subsection{\texorpdfstring{Build one independent Gaussian block for each time and coordinate}{Build one independent Gaussian block for each time and coordinate}}


The path space is \(\Omega=(\mathbb R^{\iota})^{\mathbb N\times\mathrm{Fin}(d)}\). Define its law as the infinite product of the kernel Gaussian. Evaluation at \((n,i)\) is measurable and has exactly the original Gaussian law, by the infinite-product marginal theorem. The coordinate-evaluation family is jointly independent by \texttt{i\allowbreak{}Indep\allowbreak{}Fun\_\allowbreak{}infinite\allowbreak{}Pi}. Thus a time step retains all within-block covariance while distinct iteration–coordinate pairs are independent. The product's probability instance establishes total mass one; the \texttt{Has\allowbreak{}Law} proof records both measurability and the exact pushforward equality.

\subsection{\texorpdfstring{Transfer block moments and assemble Euclidean innovations}{Transfer block moments and assemble Euclidean innovations}}


Evaluation on the infinite product is measure preserving, so composing the finite-block L\ensuremath{{}^{2}} identity with evaluation gives a finite second moment for each block. Its expectation is zero by the established \texttt{Has\allowbreak{}Law} integral transformation. \texttt{klmc\allowbreak{}Innovation} collects component zero across spatial coordinates into velocity and component one into position. \texttt{klmc2\allowbreak{}Innovation} similarly collects each of the four components into a Euclidean vector. \texttt{With\allowbreak{}Lp.\allowbreak{}to\allowbreak{}Lp} equips the coordinate functions with the Euclidean norm. Their measurable-map theorems unfold these finite coordinate constructions and apply compositional measurability. These maps supply concrete innovation streams for pathwise integrator recursions.

\subsection{\texorpdfstring{Turn covariance into actual second and cross moments}{Turn covariance into actual second and cross moments}}


The Gaussian coordinates have zero mean. Thus their variance equals \(\int X_i^2\), and their covariance equals \(\int X_iX_j\). The second-moment proof starts from the exact variance formula, rewrites variance as the integral of the centered square, and substitutes the proved zero mean. The cross-moment proof unfolds covariance and substitutes both zero means. Each conclusion is an equality for the actual constructed Gaussian measure: \(\mathbb E X_i^2=\int_0^hk_i^2\) and \(\mathbb E X_iX_j=\int_0^hk_ik_j\). These formulas can be used in numerical-error calculations without assuming moments of an unspecified noise source.

\subsection{\texorpdfstring{The scalar KLMC transition is an actual Gaussian pushforward}{The scalar KLMC transition is an actual Gaussian pushforward}}


\texttt{klmc\allowbreak{}Scalar\allowbreak{}Transition} maps the constructed correlated two-dimensional Gaussian through the scalar KLMC step at a fixed phase state. The map is continuous in its innovation, so the pushforward remains a probability measure. Both output coordinates are a constant frozen-force response plus \(\sqrt{2\gamma}\) times a Gaussian coordinate. Their finite second moments follow by adding a constant L\ensuremath{{}^{2}} function to a scalar multiple of the proved Gaussian L\ensuremath{{}^{2}} coordinate. To compute the transition mean, change variables through the pushforward, integrate the two coordinates separately using their established integrability, and substitute zero Gaussian means. The result is exactly the frozen deterministic position–velocity pair. This is a genuine one-step law and moment result, with no assumption that an abstract transition has the desired mean.

\subsection{\texorpdfstring{Correlated transition covariance and independent initialization}{Correlated transition covariance and independent initialization}}


The actual scalar phase transition retains covariance \(2\gamma C_{01}\) between its velocity and position coordinates. Lean pulls covariance through the transition pushforward, removes the deterministic means, moves both factors \(\sqrt{2\gamma}\) outside covariance, and uses the Gaussian Gram-covariance theorem; their product is \(2\gamma\) when friction is nonnegative. This checks the correlation in the update rather than replacing the two noises by independent draws.



To initialize independently, the canonical probability space is the product \(\mu_0\otimes Q\), where \(Q\) is the countable independent Gaussian-block path measure. The first projection has law \(\mu_0\), the second projection has law \(Q\), and the two projections are independent. Lean constructs this product explicitly, obtains probability from the two probability instances, and applies the product-marginal and independent-projection theorems. No later-time independence or convergence inequality is assumed by this construction.

\section{\texorpdfstring{Gaussian and Brownian noise moments}{Gaussian and Brownian noise moments}}
\label{proof-module-UnadjustedLangevin-GaussianMoments-lean}
{\small\color{gray}Source module: \texttt{Unadjusted\allowbreak{}Langevin/\allowbreak{}Gaussian\allowbreak{}Moments.\allowbreak{}lean}.\par}


Fix a natural-number dimension \(d\in\mathbb N\), write \(E_d=\mathbb R^d\) with its Euclidean norm, and let \(P\) be a measure on \(\Omega\). A law statement \(X_i\sim\mathcal N(0,v)\) below means that the coordinate is almost everywhere measurable and has the stated pushforward measure. We prove the required \(L^2\) norms from these laws. Coordinate independence is unnecessary for second moments because \(\|X\|^2=\sum_iX_i^2\) has no cross terms.



For Brownian processes the input is mathlib’s actual scalar pre-Brownian law predicate, whose marginal and increment laws supply the Gaussian hypotheses. These facts establish the noise moments; existence and stationarity of a coupled Langevin reference require separate inputs.



The \(L^2\) norms and expectations below concern the actual measure \(P\), with \texttt{Mem\allowbreak{}Lp} giving square integrability before passing to a quotient class. Laws are pushforward identities; independence is used only for cross-term cancellation, while marginal Gaussian second moments supply the dimension factor.

\subsection{\texorpdfstring{Second moment of a centered Gaussian}{Second moment of a centered Gaussian}}


For nonnegative variance \(v\), the centered Gaussian law has second moment

\[\int_{\mathbb R}x^2\,d\mathcal N(0,v)(x)=v.\]

The known Gaussian variance is \(v\). Expand variance as \(\int(x-\mathbb E[x])^2\), and insert the Gaussian mean \(\mathbb E[x]=0\).



Simplify \(x-0=x\). The variance integral is now exactly the required second moment.

\subsection{\texorpdfstring{Gaussian coordinates give a square-\allowbreak{}integrable vector}{Gaussian coordinates give a square-integrable vector}}


Let \(X:\Omega\to E_d\), and suppose every coordinate satisfies \(X_i\sim\mathcal N(0,v)\) under \(P\), with \(v\ge0\). Then \(X\in L^2(P;E_d)\).



In finite-dimensional Euclidean space, square integrability is equivalent to square integrability of each coordinate. Each coordinate belongs to \(L^2\) because its Gaussian law has a finite second moment; transfer that membership along the equality of laws.

\subsection{\texorpdfstring{Sum the coordinate moments}{Sum the coordinate moments}}


Under the same coordinate Gaussian laws,

\[\int_\Omega\|X(\omega)\|^2\,dP(\omega)=dv.\]

For each coordinate \(i\), retain the proved square integrability \(X_i\in L^2(P)\). It will justify integrating the finite sum of squared coordinates.



Change variables using the law of \(X_i\), then apply the scalar Gaussian moment:

\[\int_\Omega X_i^2\,dP=\int_{\mathbb R}x^2\,d\mathcal N(0,v)=v.\]

Expand the squared Euclidean norm and interchange the finite sum and integral:

\[\int_\Omega\|X\|^2\,dP=\sum_{i=0}^{d-1}\int_\Omega X_i^2\,dP.\]

Every summand is integrable by its \(L^2\) membership.



Each of the \(d\) summands equals \(v\); hence their sum is \(dv\).

\subsection{\texorpdfstring{The Gaussian vector as an L\ensuremath{{}^{2}} class}{The Gaussian vector as an L² class}}


Let \([X]\) be the \(L^2\) equivalence class constructed from \(X\), whose coordinates have laws \(\mathcal N(0,v)\). Then

\[\|[X]\|_{L^2}^2=dv.\]

Express the squared \(L^2\) norm as the self-inner-product, then use the definition of the \(L^2\) inner product as an integral.



The chosen representative of \([X]\) equals \(X\) almost everywhere. Therefore the integrand \(\langle[X](\omega),[X](\omega)\rangle\) agrees almost everywhere with \(\|X(\omega)\|^2\), so the two integrals agree.



Substitute the integral identity and the vector second-moment calculation \(\int\|X\|^2\,dP=dv\).

\subsection{\texorpdfstring{The square-\allowbreak{}root Gaussian norm}{The square-root Gaussian norm}}


For the same Gaussian coordinate laws,

\[\|[X]\|_{L^2}=\sqrt{dv}.\]

This includes \(d=0\) and \(v=0\).



Replace \(dv\) by \(\|[X]\|^2\). Since norms are nonnegative, \(\sqrt{\|[X]\|^2}=\|[X]\|\).

\subsection{\texorpdfstring{Include the Langevin diffusion factor}{Include the Langevin diffusion factor}}


If every coordinate of \(X\) has law \(\mathcal N(0,v)\), then the Langevin-scaled vector satisfies

\[\|\sqrt2[X]\|_{L^2}=\sqrt{2d}\sqrt v.\]

Norm homogeneity gives \(\|\sqrt2[X]\|=\sqrt2\|[X]\|\), because \(\sqrt2\ge0\). Insert \(\|[X]\|=\sqrt{dv}\), then split the square roots of the nonnegative products.



Rearrange \(\sqrt2\sqrt d\sqrt v\) into \(\sqrt{2d}\sqrt v\).

\subsection{\texorpdfstring{Assemble Brownian coordinates}{Assemble Brownian coordinates}}


For scalar processes \(W_i:[0,\infty)\times\Omega\to\mathbb R\), define the Euclidean vector

\[W_t(\omega):=(W_0(t,\omega),\ldots,W_{d-1}(t,\omega))\in E_d.\]

Collect the coordinates at the fixed time \(t\) and outcome \(\omega\), using the finite-dimensional \(\ell^2\) norm.

\subsection{\texorpdfstring{Brownian marginals are Gaussian}{Brownian marginals are Gaussian}}


Assume each scalar \(W_i\) satisfies \(\mathrm{IsPreBrownianReal}\) under \(P\). For each \(t\ge0\) and coordinate \(i\),

\[W_t^{(i)}\sim\mathcal N(0,t).\]

Apply the scalar Brownian marginal-law theorem at time \(t\). The \(i\)-th coordinate of the assembled vector is definitionally the original scalar process.

\subsection{\texorpdfstring{Brownian vectors belong to L\ensuremath{{}^{2}}}{Brownian vectors belong to L²}}


Under the scalar pre-Brownian assumptions, for every \(t\ge0\),

\[W_t\in L^2(P;E_d).\]

Apply Gaussian vector square integrability with variance \(v=t\), using the just-proved marginal laws of all coordinates.

\subsection{\texorpdfstring{The exact Brownian noise norm}{The exact Brownian noise norm}}


For scalar pre-Brownian coordinates and \(t\ge0\),

\[\|\sqrt2[W_t]\|_{L^2}=\sqrt{2d}\sqrt t.\]

Specialize the scaled Gaussian vector norm to the Brownian vector, whose coordinate variance is \(t\). This derives the square-root noise scale from the Brownian laws.

\subsection{\texorpdfstring{Define an increment from any starting time}{Define an increment from any starting time}}


For \(a,t\ge0\), define the vector increment over the interval \([a,a+t]\) by

\[\Delta_aW_t(\omega):=W_{a+t}(\omega)-W_a(\omega).\]

Subtract the two Euclidean vectors at the interval endpoints. Each coordinate is the corresponding scalar Brownian increment.

\subsection{\texorpdfstring{Increment variance is elapsed time}{Increment variance is elapsed time}}


Under the scalar pre-Brownian assumptions, for every \(a,t\ge0\) and every coordinate \(i\),

\[W_i(a+t)-W_i(a)\sim\mathcal N(0,t).\]

Expand the vector increment to its scalar coordinate. Apply the Brownian difference-law theorem at times \(a+t\) and \(a\); its variance is the nonnegative distance between those times.



It remains to identify that variance. As real numbers,

\[|(a+t)-a|=|t|=t,\]

because \(t\ge0\). Thus the increment law is exactly \(\mathcal N(0,t)\).

\subsection{\texorpdfstring{Increments are square integrable}{Increments are square integrable}}


For pre-Brownian coordinates and any \(a,t\ge0\),

\[\Delta_aW_t\in L^2(P;E_d).\]

Each increment coordinate has the centered Gaussian law of variance \(t\), so the Gaussian vector integrability theorem applies.

\subsection{\texorpdfstring{The same noise norm on every reference segment}{The same noise norm on every reference segment}}


For scalar pre-Brownian coordinates and every \(a,t\ge0\),

\[\boxed{\|\sqrt2[W_{a+t}-W_a]\|_{L^2}=\sqrt{2d}\sqrt t.}\]

Apply the scaled Gaussian vector norm to the increment coordinate laws. The answer depends only on elapsed time \(t\), so every reference interval has the same moment bound.

\section{\texorpdfstring{Characteristic functions are continuous in L\ensuremath{{}^{2}}}{Characteristic functions are continuous in L²}}
\label{proof-module-UnadjustedLangevin-LpCharacteristic-lean}
{\small\color{gray}Source module: \texttt{Unadjusted\allowbreak{}Langevin/\allowbreak{}Lp\allowbreak{}Characteristic.\allowbreak{}lean}.\par}


At each fixed frequency, the characteristic function of the law of a real \(L^2\) variable depends continuously on that variable. Consequently centered Gaussian laws are closed under \(L^2\) convergence, with variance given by the squared norm.

\subsection{\texorpdfstring{centered\allowbreak{}Characteristic}{centeredCharacteristic}}


Set \(c_t(x)=e^{itx}-1\), represented by the unit-radius circle map centered at \(-1\). Subtracting one makes \(c_t(0)=0\), which is the condition needed by the Lipschitz composition construction on \(L^2\).

\subsection{\texorpdfstring{centered\allowbreak{}Characteristic lipschitz}{centeredCharacteristic lipschitz}}


Multiplication by \(t\) is \(|t|\)-Lipschitz, and the unit circle parametrization is \(1\)-Lipschitz. Composing gives \(|c_t(x)-c_t(y)|\le|t||x-y|\). The proof identifies the nonnegative Lipschitz constants before applying the composition theorem.

\subsection{\texorpdfstring{centered\allowbreak{}Characteristic zero}{centeredCharacteristic zero}}


At zero the exponential is one, hence \(c_t(0)=e^0-1=0\). Unfolding the circle-map definition makes this a scalar identity; this vanishing proof permits \texttt{comp\allowbreak{}Lp} to lift the map without an integrability assumption on a nonzero constant.

\subsection{\texorpdfstring{characteristic\allowbreak{}Lp}{characteristicLp}}


Lift \(c_t\) to \(L^2\) by Lipschitz composition, then add the constant class one. Thus \(\operatorname{characteristicLp}(t,X)\) represents \(e^{itX}\). The probability assumption makes the constant function square integrable.

\subsection{\texorpdfstring{characteristic\allowbreak{}Lp coe}{characteristicLp coe}}


Combine the almost-everywhere representative identities for lifted composition, \(L^2\) addition, and the constant class. Their intersection has full measure; there the representative is \((e^{itX}-1)+1=e^{itX}\). This connects quotient-space operations with the random function used in the characteristic integral.

\subsection{\texorpdfstring{char\allowbreak{}Fun map eq inner characteristic\allowbreak{}Lp}{charFun map eq inner characteristicLp}}


The pushforward integral theorem gives \(\widehat{P_X}(t)=\int e^{itX}dP\). In complex \(L^2\), inner product with the constant one has exactly this integral because its conjugate is one. The representative identity above therefore yields \(\widehat{P_X}(t)=\langle1,\operatorname{characteristicLp}(t,X)\rangle\).

\subsection{\texorpdfstring{continuous char\allowbreak{}Fun map Lp}{continuous charFun map Lp}}


Lipschitz composition is continuous on \(L^2\); adding the constant one preserves continuity. Pairing with one is a continuous linear functional. The preceding integral identity rewrites their composition as \(X\mapsto\widehat{P_X}(t)\), proving continuity of the actual-law characteristic function at every fixed frequency.

\subsection{\texorpdfstring{continuous gaussian\allowbreak{}Real char\allowbreak{}Fun Lp}{continuous gaussianReal charFun Lp}}


The centered Gaussian characteristic function is \(e^{-t^2\|X\|_2^2/2}\). Norm, squaring, conversion of a nonnegative real to \texttt{\ensuremath{\mathbb R}\ensuremath{\ge}0}, and the exponential are continuous. Lean unfolds the Gaussian formula and composes these continuity facts; the variance is the actual squared \(L^2\) norm.

\subsection{\texorpdfstring{has\allowbreak{}Law gaussian\allowbreak{}Real of tendsto Lp eventually}{hasLaw gaussianReal of tendsto Lp eventually}}


At each frequency, \(L^2\) convergence makes the actual characteristic functions converge. It also makes the centered Gaussian characteristic functions with variances \(\|X_i\|_2^2\) converge to variance \(\|Y\|_2^2\). Eventual Gaussianity identifies the two nets eventually, so uniqueness of limits identifies their limits. \texttt{Measure.\allowbreak{}ext\_\allowbreak{}of\_\allowbreak{}char\allowbreak{}Fun} then gives the claimed Gaussian law of \(Y\); a nontrivial filter is required for uniqueness.

\subsection{\texorpdfstring{has\allowbreak{}Law gaussian\allowbreak{}Real of tendsto Lp}{hasLaw gaussianReal of tendsto Lp}}


If every approximant has the stated Gaussian law, then the law identity holds eventually for any filter. Feed \texttt{Eventually.\allowbreak{}of\_\allowbreak{}forall} to the preceding closure theorem. This version is suited to finite orthogonal sums converging to an infinite Gaussian series.

\section{\texorpdfstring{Construct an isonormal Gaussian family from product coordinates}{Construct an isonormal Gaussian family from product coordinates}}
\label{proof-module-UnadjustedLangevin-GaussianHilbert-lean}
{\small\color{gray}Source module: \texttt{Unadjusted\allowbreak{}Langevin/\allowbreak{}Gaussian\allowbreak{}Hilbert.\allowbreak{}lean}.\par}


An explicit product of standard normal coordinates realizes a linear isometry \(I:H\to L^2\) with \(I(x)\sim\mathcal N(0,\|x\|^2)\) and covariance \(\mathbb E[I(x)I(y)]=\langle x,y\rangle\).

\subsection{\texorpdfstring{source}{source}}


Use the product probability \(Q=\bigotimes_{i\in\iota}\mathcal N(0,1)\) on real coordinate families. The construction uses \texttt{Measure.\allowbreak{}infinite\allowbreak{}Pi}; independence and coordinate marginals are properties of this concrete source rather than extra assumptions on an unspecified Gaussian process.

\subsection{\texorpdfstring{coordinate has\allowbreak{}Law}{coordinate hasLaw}}


The coordinate projection \(\omega\mapsto\omega_i\) pushes the product measure to its \(i\)-th factor \(\mathcal N(0,1)\). The product marginal theorem proves both almost-everywhere measurability and the pushforward identity, packaged as \texttt{Has\allowbreak{}Law}.

\subsection{\texorpdfstring{coordinate\allowbreak{}Lp}{coordinateLp}}


A standard Gaussian has finite second moment, so the coordinate projection belongs to \(L^2(Q)\). Apply its \texttt{Mem\allowbreak{}Lp.\allowbreak{}to\allowbreak{}Lp} construction to obtain the quotient class \(Z_i\); the definition stores the original coordinate function together with its integrability proof.

\subsection{\texorpdfstring{coordinate\allowbreak{}Lp coe}{coordinateLp coe}}


The representative selected for \(Z_i\) agrees with \(\omega_i\) almost everywhere by the defining property of \texttt{to\allowbreak{}Lp}. This equality is the bridge used to transfer coordinate laws, means, and independence calculations to the \(L^2\) classes.

\subsection{\texorpdfstring{coordinates independent}{coordinates independent}}


The coordinate projections of an infinite product measure are jointly independent. Invoke that product theorem for the scalar Gaussian factors. This is independence of every finite set of coordinates, which is precisely what finite Gaussian linear combinations require.

\subsection{\texorpdfstring{coordinate mean}{coordinate mean}}


Change variables through the coordinate law and use the centered Gaussian mean: \(\int\omega_i\,dQ=\int x\,d\mathcal N(0,1)=0\). Gaussian integrability makes this an ordinary expectation identity.

\subsection{\texorpdfstring{coordinate\allowbreak{}Lp inner}{coordinateLp inner}}


For \(i=j\), the \(L^2\) inner product is \(\mathbb EZ_i^2=1\). For \(i\ne j\), independence factors \(\mathbb EZ_iZ_j=(\mathbb EZ_i)(\mathbb EZ_j)=0\). Almost-everywhere representative identities move between \(Z_i\) and the product coordinates. Thus \(\langle Z_i,Z_j\rangle=\mathbf1_{i=j}\).

\subsection{\texorpdfstring{coordinates orthonormal}{coordinates orthonormal}}


The preceding inner-product formula is the orthonormality criterion for a family in a real Hilbert space. Lean converts the diagonal value one and off-diagonal value zero into \texttt{Orthonormal \ensuremath{\mathbb R} coordinate\allowbreak{}Lp}, the input for the orthogonal-series extension.

\subsection{\texorpdfstring{source is\allowbreak{}Probability\allowbreak{}Measure}{source isProbabilityMeasure}}


Each scalar Gaussian factor is a probability measure. The product-probability instance therefore gives \(Q(\Omega)=1\). This instance is what permits constant functions and integrable functions to be treated in the usual probability-space \(L^2\) theory below.

\subsection{\texorpdfstring{synthesis}{synthesis}}


Extend the orthonormal family \(Z_i\) to a linear isometry \(S:\ell^2(\iota)\to L^2(Q)\). The orthogonal-family extension theorem constructs the limit of finite sums and preserves squared norm, \(\|S(a)\|_2^2=\sum_i a_i^2\). It constructs an \(L^2\) limit rather than assuming pointwise convergence of a random series.

\subsection{\texorpdfstring{synthesis has\allowbreak{}Sum}{synthesis hasSum}}


The synthesis extension theorem states \(\sum_i a_iZ_i=S(a)\), with convergence of finite partial sums in \(L^2\). In Lean \texttt{Has\allowbreak{}Sum} is convergence along finite subsets, so the identity does not require choosing an enumeration of the basis index set.

\subsection{\texorpdfstring{isonormal}{isonormal}}


For a Hilbert basis \(b\), compose its coordinate isometry \(b.\mathrm{repr}:H\to\ell^2(\iota)\) with \texttt{synthesis}. The resulting linear isometry \(I_b\) sends \(x\) to the Gaussian series with coefficients \(\langle b_i,x\rangle\), and preserves all inner products.

\subsection{\texorpdfstring{isonormal has\allowbreak{}Sum}{isonormal hasSum}}


Apply \texttt{synthesis\_\allowbreak{}has\allowbreak{}Sum} to the coefficient family \(b.\mathrm{repr}(x)\). Unfolding the composition gives \(I_b(x)=\sum_i\langle b_i,x\rangle Z_i\) in \(L^2\). This is the exact series represented by the bundled isonormal isometry.

\subsection{\texorpdfstring{integral eq inner one}{integral eq inner one}}


On a probability space, the constant one is in \(L^2\). The definition of the real \(L^2\) inner product gives \(\mathbb EX=\langle1,X\rangle\), after replacing the selected constant representative by one almost everywhere. This exhibits expectation as a continuous functional.

\subsection{\texorpdfstring{coordinate\allowbreak{}Lp mean}{coordinateLp mean}}


Use \(Z_i=\omega_i\) almost everywhere to replace the \(L^2\) representative inside the integral. The already established centered coordinate expectation then gives \(\mathbb EZ_i=0\).

\subsection{\texorpdfstring{synthesis mean}{synthesis mean}}


Each finite partial sum \(\sum_{i\in F}a_iZ_i\) has mean zero by linearity and the coordinate means. Pair the \(L^2\)-convergent series with the constant one; continuity of the inner product passes the zero expectations to the limit, giving \(\mathbb ES(a)=0\).

\subsection{\texorpdfstring{isonormal mean}{isonormal mean}}


The isonormal variable is \texttt{synthesis} applied to the Hilbert-basis coefficients of \(x\). Specialize the synthesis mean theorem to these coefficients to get \(\mathbb EI_b(x)=0\), without redoing the infinite-series limit.

\subsection{\texorpdfstring{variance eq norm sq}{variance eq norm sq}}


For an \(L^2\) random variable with mean zero, the variance formula reduces to \(\operatorname{Var}(X)=\int X^2dQ\). The \(L^2\) norm identity identifies this with \(\|X\|_2^2\). The explicit mean premise is used to remove the centering term.

\subsection{\texorpdfstring{partial sum coe}{partial sum coe}}


Induct over the finite set. The zero sum has zero representative almost everywhere; addition and scalar multiplication of \(L^2\) classes agree almost everywhere with the corresponding pointwise operations. Intersect the finitely many full-measure sets to obtain \(\sum_{i\in F}a_iZ_i=\sum_{i\in F}a_i\omega_i\) almost everywhere.

\subsection{\texorpdfstring{partial sum has\allowbreak{}Gaussian\allowbreak{}Law}{partial sum hasGaussianLaw}}


Independent scalar standard Gaussians remain Gaussian after a finite linear combination. Apply that finite-sum Gaussian theorem to \(a_i\omega_i\), and transfer the law across \texttt{partial\_\allowbreak{}sum\_\allowbreak{}coe}. This proves Gaussianity of the actual \(L^2\) partial sum, not merely of a formally similar expression.

\subsection{\texorpdfstring{partial sum mean}{partial sum mean}}


The integral of a finite \(L^2\) sum is the sum of its integrals, because every summand is integrable. Each coordinate mean is zero, hence \(\mathbb E\sum_{i\in F}a_iZ_i=0\). This supplies the centering parameter needed to identify its Gaussian law.

\subsection{\texorpdfstring{partial sum has\allowbreak{}Law}{partial sum hasLaw}}


A scalar Gaussian law is determined by its mean and variance. The finite sum is Gaussian and centered; its variance is its squared \(L^2\) norm. Therefore its law is \(\mathcal N(0,\|\sum_{i\in F}a_iZ_i\|_2^2)\), with the variance converted to a nonnegative real.

\subsection{\texorpdfstring{synthesis has\allowbreak{}Law}{synthesis hasLaw}}


Finite partial sums converge in \(L^2\) to \(S(a)\), and each has the centered Gaussian law with its own squared norm as variance. Apply \texttt{has\allowbreak{}Law\_\allowbreak{}gaussian\allowbreak{}Real\_\allowbreak{}of\_\allowbreak{}tendsto\_\allowbreak{}Lp}; the isometry then rewrites the limiting variance as \(\|a\|_{\ell^2}^2\). This proves Gaussianity of the completed infinite series.

\subsection{\texorpdfstring{isonormal has\allowbreak{}Law}{isonormal hasLaw}}


Apply the synthesis law theorem to \(b.\mathrm{repr}(x)\). The Hilbert-basis coordinate map preserves norm, so the variance becomes \(\|x\|^2\). Hence \(I_b(x)\sim\mathcal N(0,\|x\|^2)\), including the degenerate law for \(x=0\).

\subsection{\texorpdfstring{isonormal has\allowbreak{}Gaussian\allowbreak{}Law}{isonormal hasGaussianLaw}}


The exact law theorem identifies \(I_b(x)\) with a scalar Gaussian probability. Forgetting its explicit mean and variance parameters yields the bundled \texttt{Has\allowbreak{}Gaussian\allowbreak{}Law} predicate used by the finite-dimensional Gaussian-process construction.

\subsection{\texorpdfstring{isonormal covariance}{isonormal covariance}}


Both variables are centered, so their covariance is \(\mathbb E[I_b(x)I_b(y)]\), the \(L^2\) inner product. The linear isometry preserves that inner product, giving \(\operatorname{Cov}(I_b(x),I_b(y))=\langle x,y\rangle_H\). Lean checks integrability through their \(L^2\) membership.

\section{\texorpdfstring{Scalar Gaussian linear maps give jointly Gaussian processes}{Scalar Gaussian linear maps give jointly Gaussian processes}}
\label{proof-module-UnadjustedLangevin-GaussianLinearProcess-lean}
{\small\color{gray}Source module: \texttt{Unadjusted\allowbreak{}Langevin/\allowbreak{}Gaussian\allowbreak{}Linear\allowbreak{}Process.\allowbreak{}lean}.\par}


If a linear map \(T:H\to L^2(P;\mathbb R)\) sends every input to a Gaussian variable, all finite collections of its evaluations are jointly Gaussian.

\subsection{\texorpdfstring{has\allowbreak{}Gaussian\allowbreak{}Law pi of linear\allowbreak{}Map Lp}{hasGaussianLaw pi of linearMap Lp}}


Use the finite-dimensional Gaussian criterion: every scalar linear functional of \((T(f_i))_i\) must be Gaussian. A finite linear combination equals \(T(\sum_i a_if_i)\) as an \(L^2\) class by linearity, hence agrees almost everywhere as a random variable. The assumed Gaussianity of every \(T(x)\) proves the criterion.

\subsection{\texorpdfstring{is\allowbreak{}Gaussian\allowbreak{}Process of linear\allowbreak{}Map Lp}{isGaussianProcess of linearMap Lp}}


For each finite set of times, apply the preceding joint-Gaussian theorem to the corresponding finite family of vectors \(f(t)\). This verifies the finite-dimensional-law definition of \texttt{Is\allowbreak{}Gaussian\allowbreak{}Process}. The index set may be uncountable; this statement alone asserts no sample-path continuity.

\section{\texorpdfstring{An explicit Gaussian construction of pre-\allowbreak{}Brownian coordinates}{An explicit Gaussian construction of pre-Brownian coordinates}}
\label{proof-module-UnadjustedLangevin-CanonicalBrownian-lean}
{\small\color{gray}Source module: \texttt{Unadjusted\allowbreak{}Langevin/\allowbreak{}Canonical\allowbreak{}Brownian.\allowbreak{}lean}.\par}


The isonormal family on the time Hilbert space constructs a Gaussian process with covariance \(\min(s,t)\). Independent product copies supply the Brownian coordinates required by the canonical ULA theorem. The construction proves pre-Brownian laws; continuous sample paths are not asserted or needed, because continuity is established in \(L^2\).

\subsection{\texorpdfstring{Time\allowbreak{}Hilbert}{TimeHilbert}}


Take the deterministic Hilbert space \(L^2(\mathbb R,dt)\). Its vectors will encode time intervals before an isonormal map turns their inner products into covariances. It is distinct from the random-variable \(L^2\) space on the Gaussian source.

\subsection{\texorpdfstring{time\allowbreak{}Indicator}{timeIndicator}}


Let \(e_t=\mathbf1_{(0,t]}\) as a class in \(L^2(dt)\), for \(t\ge0\). Lean constructs it with \texttt{indicator\allowbreak{}Const\allowbreak{}Lp} and the finite Lebesgue measure of \texttt{Ioc 0 t}. Endpoints do not affect this Lebesgue \(L^2\) class, but this half-open convention makes interval intersection identities direct.

\subsection{\texorpdfstring{time\allowbreak{}Indicator inner}{timeIndicator inner}}


The product of the two indicators is the indicator of \([0,\min(s,t)]\) almost everywhere. Integrating it gives \(\langle e_s,e_t\rangle=\min(s,t)\). The proof evaluates the Lebesgue measure of that finite interval, producing the Brownian covariance before introducing randomness.

\subsection{\texorpdfstring{basis\allowbreak{}Index}{basisIndex}}


Choose the index type supplied by a Hilbert basis of the time Hilbert space. It is a set of Hilbert-space vectors used as a type; no explicit countable enumeration is required by the orthogonal-family construction.

\subsection{\texorpdfstring{basis}{basis}}


Select a Hilbert basis for \texttt{Time\allowbreak{}Hilbert} by Hilbert-space basis existence. This basis supplies the coordinate isometry for the explicit isonormal Gaussian family; all subsequent covariance identities follow from isometry and do not depend on which basis was selected.

\subsection{\texorpdfstring{Sample}{Sample}}


A sample is a real family indexed by the chosen time-Hilbert basis. This concrete product space stores the independent standard Gaussian coefficients used to synthesize all Brownian times.

\subsection{\texorpdfstring{probability}{probability}}


Use \texttt{Gaussian\allowbreak{}Hilbert.\allowbreak{}source} on the basis index type, namely the product of standard scalar Gaussians. Its established probability instance provides normalization and independent coordinate laws for the Brownian construction.

\subsection{\texorpdfstring{instance at 53}{instance at 53}}


The scalar source is the Gaussian coordinate product probability. Reuse \texttt{Gaussian\allowbreak{}Hilbert.\allowbreak{}source\_\allowbreak{}is\allowbreak{}Probability\allowbreak{}Measure} to provide normalization for the chosen basis index. This instance enables the \(L^2\) expectation and covariance theorems for the constructed Brownian process.

\subsection{\texorpdfstring{value\allowbreak{}Lp}{valueLp}}


Apply the isonormal isometry to \(e_t\), obtaining \(B_t=I(e_t)\) as an \(L^2\) random variable. The construction immediately gives a centered Gaussian with variance \(\|e_t\|^2=t\); joint Gaussianity follows from linearity.

\subsection{\texorpdfstring{process}{process}}


At each time select the strongly measurable representative of \texttt{value\allowbreak{}Lp}. This defines an actual function \(B_t(\omega)\) for use in \texttt{Has\allowbreak{}Law} and finite-dimensional distributions. Equalities of \(L^2\) classes still translate to almost-everywhere equalities of these representatives.

\subsection{\texorpdfstring{process measurable}{process measurable}}


The representative selected for an \(L^2\) class is measurable in this setting. Apply the representative measurability theorem to \texttt{value\allowbreak{}Lp t}, obtaining measurability of the actual time-\(t\) random variable.

\subsection{\texorpdfstring{process mean}{process mean}}


The isonormal mean theorem gives \(\mathbb EI(e_t)=0\). Since \texttt{process t} is precisely the selected representative of this class, this is the expectation of the displayed random function \(B_t\).

\subsection{\texorpdfstring{process covariance}{process covariance}}


The isonormal covariance formula gives \(\operatorname{Cov}(I(e_s),I(e_t))=\langle e_s,e_t\rangle\). Substitute \texttt{time\allowbreak{}Indicator\_\allowbreak{}inner} to obtain \(\min(s,t)\). This calculation establishes the Brownian covariance for the concrete source and representatives.

\subsection{\texorpdfstring{is\allowbreak{}Gaussian\allowbreak{}Process}{isGaussianProcess}}


The isonormal isometry is linear and every output is scalar Gaussian. Apply \texttt{is\allowbreak{}Gaussian\allowbreak{}Process\_\allowbreak{}of\_\allowbreak{}linear\allowbreak{}Map\_\allowbreak{}Lp} to the family \(t\mapsto e_t\). Every finite collection of the selected random variables is therefore jointly Gaussian.

\subsection{\texorpdfstring{is\allowbreak{}Pre\allowbreak{}Brownian}{isPreBrownian}}


Combine centered Gaussian finite-dimensional laws with covariance \(\min(s,t)\). The Gaussian-process Brownian characterization yields \texttt{Is\allowbreak{}Pre\allowbreak{}Brownian\allowbreak{}Real}. This proves Brownian finite-dimensional laws; the prefix “pre” records that sample-path continuity is not part of this predicate.

\subsection{\texorpdfstring{Vector\allowbreak{}Sample}{VectorSample}}


For a coordinate index type \(I\), take one scalar Brownian sample for each \(i\in I\). This is the product sample type \(I\to\mathrm{Sample}\), allowing independent entire coordinate processes rather than only independent values at one time.

\subsection{\texorpdfstring{vector\allowbreak{}Probability}{vectorProbability}}


Equip the vector sample space with the infinite product of the scalar Brownian source probabilities. The coordinate processes inherit independence from this actual product construction.

\subsection{\texorpdfstring{coordinate}{coordinate}}


Define the \(i\)-th coordinate process by \(B_i(t,\omega)=B(t,\omega_i)\). It first projects to the \(i\)-th scalar sample, then evaluates the already constructed scalar process at \(t\).

\subsection{\texorpdfstring{coordinate is\allowbreak{}Pre\allowbreak{}Brownian}{coordinate isPreBrownian}}


The product coordinate projection has the scalar Brownian source as its law. Pull the scalar pre-Brownian process back along this projection; its finite-dimensional distributions, mean, and covariance are preserved. Thus each vector coordinate has the established scalar Brownian laws.

\subsection{\texorpdfstring{coordinate paths independent}{coordinate paths independent}}


The product coordinate projections are independent as sample-valued maps. Compose each projection with the map sending a scalar sample to its entire time-indexed process. Independence is preserved under these coordinatewise functions, giving independence of the whole coordinate paths.

\subsection{\texorpdfstring{coordinate grid paths independent}{coordinate grid paths independent}}


Restrict each independent coordinate path to the deterministic grid \(nh\). This is a coordinatewise function of the path, so independence is preserved. The result supplies exactly the countable grid-path hypothesis used in the innovation theorems.

\subsection{\texorpdfstring{instance at 89}{instance at 89}}


Each scalar Brownian sample source has mass one. The infinite product over vector coordinates therefore also has mass one. This instance concerns the actual vector sample measure, not an assumed normalization of an abstract process law.

\section{\texorpdfstring{Construct the correlated innovations from the same Brownian source}{Construct the correlated innovations from the same Brownian source}}
\label{proof-module-KineticLangevin-BrownianConvolution-lean}
{\small\color{gray}Source module: \texttt{Kinetic\allowbreak{}Langevin/\allowbreak{}Brownian\allowbreak{}Convolution.\allowbreak{}lean}.\par}


The exponential schemes use correlated Brownian integrals, for example \(\xi_v=\int_0^t e^{-\gamma(t-s)}\,dB_s\) and \(\xi_x=\int_0^t\psi_1(t-s)\,dB_s\). Drawing two independent Gaussians would give the wrong algorithm. This module constructs these random variables on the explicit Gaussian probability space already used for the canonical Brownian motion. Their joint Gaussian law, all covariances, and continuity in L\ensuremath{{}^{2}} are proved from the deterministic time kernels. The construction uses the existing isonormal isometry, which is the deterministic-kernel Wiener integral, rather than adding a hypothesis asserting that an unspecified noise has the required law.

\subsection{\texorpdfstring{Reuse the existing Brownian probability space}{Reuse the existing Brownian probability space}}


The time Hilbert space is \(H=L^2(\mathbb R,ds)\). The existing canonical Brownian construction selects a Hilbert basis of this space and places independent standard Gaussian coordinates on an explicit product probability space. Its isonormal map \(I:H\to L^2(\Omega,\mathbb P)\) is linear and preserves inner products. Brownian motion is represented there by \(B_t=I(1_{(0,t]})\). The three abbreviations reuse exactly this Hilbert space, sample space, and probability measure. No second independent Brownian source is introduced.

\subsection{\texorpdfstring{Turn a continuous response into a genuine L\ensuremath{{}^{2}} time kernel}{Turn a continuous response into a genuine L² time kernel}}


For a continuous scalar response \(k\), define \(K_{k,t}(s)=1_{(0,t]}(s)k(t-s)\). The time reversal is essential: a Brownian fluctuation arriving at time \(s\) has \(t-s\) units of damping before the endpoint. Continuity of \(k\) gives a finite bound on the compact interval \([0,t]\). Since \((0,t]\) has finite Lebesgue measure, this bound proves square integrability.



Lean first uses \texttt{mem\allowbreak{}Lp\_\allowbreak{}indicator\_\allowbreak{}iff\_\allowbreak{}restrict} to work on the finite interval measure. \texttt{Is\allowbreak{}Compact.\allowbreak{}exists\_\allowbreak{}bound\_\allowbreak{}of\_\allowbreak{}continuous\allowbreak{}On} provides an actual bound, and \texttt{Mem\allowbreak{}Lp.\allowbreak{}of\_\allowbreak{}bound} proves membership in L\ensuremath{{}^{2}}. \texttt{time\allowbreak{}Kernel\allowbreak{}Lp} then uses \texttt{Mem\allowbreak{}Lp.\allowbreak{}to\allowbreak{}Lp} to form the equivalence class. Its selected representative agrees almost everywhere with the explicit indicator function, recorded by \texttt{time\allowbreak{}Kernel\allowbreak{}Lp\_\allowbreak{}coe}. Subsequent pointwise calculations therefore apply to the actual L\ensuremath{{}^{2}} object.

\subsection{\texorpdfstring{Compute the time-\allowbreak{}kernel inner product exactly}{Compute the time-kernel inner product exactly}}


The fundamental identity is

\[\langle K_{k,t},K_{\ell,t}\rangle_H=\int_0^t k(s)\ell(s)\,ds.\]

Start with the definition of the L\ensuremath{{}^{2}} inner product as an integral. Replace both chosen representatives using their almost-everywhere formulas. Their product is zero outside \((0,t]\), and inside it equals \(k(t-s)\ell(t-s)\). Integrating the indicator restricts the integral to that interval. The substitution \(u=t-s\) reverses the endpoints and gives the displayed expression. Lean implements exactly these steps with \texttt{L2.\allowbreak{}inner\_\allowbreak{}def}, \texttt{integral\_\allowbreak{}congr\_\allowbreak{}ae}, \texttt{integral\_\allowbreak{}indicator}, and \texttt{interval\allowbreak{}Integral.\allowbreak{}integral\_\allowbreak{}comp\_\allowbreak{}sub\_\allowbreak{}left}. This identity is the source of the covariance matrix, including its off-diagonal entries.

\subsection{\texorpdfstring{Prove time continuity in L\ensuremath{{}^{2}}, including the moving endpoint}{Prove time continuity in L², including the moving endpoint}}


A moving indicator has a discontinuity at its endpoint, so pointwise continuity at every time coordinate is unavailable. Fix an endpoint \(t\). For every \(r\ne t\), the membership of \(r\) in \((0,s]\) is constant for all \(s\) sufficiently close to \(t\). If \(r<t\), ordinary continuity of \(k(s-r)\) applies; if \(r>t\), the kernel stays zero. The single exceptional time coordinate has Lebesgue measure zero.



The exact norm formula is

\[\|K_{k,s}-K_{k,t}\|_H^2=\int_{\mathbb R}|K_{k,s}(r)-K_{k,t}(r)|^2\,dr.\]

For \(s\) near \(t\), both kernels are supported in the fixed finite interval \((0,t+1]\). Compactness gives a common bound \(C\) for their values. The squared difference is therefore bounded by the integrable function \(4C^2 1_{(0,t+1]}\). Dominated convergence makes the squared L\ensuremath{{}^{2}} distance tend to zero; continuity of the square root gives convergence of the distance itself.



Lean proves the pointwise limit by splitting the cases \(r<t\), \(r>t\), and \(r\le0\), with eventual inequalities supplied by the neighborhood filter. It obtains the norm formula from the L\ensuremath{{}^{2}} inner product and the almost-everywhere subtraction formula. The continuity proof then applies \texttt{tendsto\_\allowbreak{}integral\_\allowbreak{}filter\_\allowbreak{}of\_\allowbreak{}dominated\_\allowbreak{}convergence}, explicitly supplies the compact bound and the null singleton exception, and finishes with \texttt{Real.\allowbreak{}continuous\_\allowbreak{}sqrt}. This proves continuity as an L\ensuremath{{}^{2}}-valued map; it does not assert that the selected representatives have continuous sample paths.

\subsection{\texorpdfstring{Apply the isonormal map and derive Gaussian moments}{Apply the isonormal map and derive Gaussian moments}}


Define \(Z_k(t)=I(K_{k,t})\). This is the Wiener integral \(\int_0^t k(t-s)\,dB_s\) for a deterministic continuous kernel. The existing construction proves that every output of \(I\) is centered Gaussian and that \(\operatorname{Cov}(Iu,Iv)=\langle u,v\rangle_H\). Consequently,

\[\mathbb E Z_k(t)=0,\qquad \operatorname{Cov}(Z_k(t),Z_\ell(t))=\int_0^t k(s)\ell(s)\,ds.\]

The same linear map proves joint Gaussianity of every finite collection. This last statement is stronger than proving Gaussianity of each coordinate separately: all linear combinations must be Gaussian. Lean invokes the previously proved \texttt{has\allowbreak{}Gaussian\allowbreak{}Law\_\allowbreak{}pi\_\allowbreak{}of\_\allowbreak{}linear\allowbreak{}Map\_\allowbreak{}Lp}, whose proof uses finite linear combinations and the linearity of the isonormal map. Covariance follows by rewriting the isonormal covariance theorem and the time-kernel inner-product identity. Time continuity follows by composing the continuous isometry with \texttt{continuous\_\allowbreak{}time\allowbreak{}Kernel\allowbreak{}Lp}.

\subsection{\texorpdfstring{Identify the exact multivariate innovation law}{Identify the exact multivariate innovation law}}


For finitely many kernels \((k_i)\), assemble \(Z=(Z_{k_i}(t))_i\) as a Euclidean random vector. Its law is proved equal to \(\mathcal N(0,C)\), where \(C_{ij}=\int_0^t k_i(s)k_j(s)\,ds\). The conclusion compares the actual pushforward probability measure with \texttt{kernel\allowbreak{}Gaussian}; equality of laws is not a premise.



Lean first transports the jointly Gaussian coordinate family through the continuous linear equivalence from a finite function space to Euclidean space. Coordinate L\ensuremath{{}^{2}} membership gives vector L\ensuremath{{}^{2}} membership, hence integrability. To identify the mean, the continuous coordinate projections commute with the Bochner integral and reduce every coordinate to the already proved zero mean. To identify covariance, \texttt{covariance\allowbreak{}Bilin\_\allowbreak{}apply\_\allowbreak{}pi} expands the covariance bilinear form as a finite double sum of coordinate covariances. The Brownian covariance identity turns that double sum into the matrix expression for the Gram covariance. \texttt{Is\allowbreak{}Gaussian.\allowbreak{}ext} then concludes equality of the two Gaussian probability measures from equality of their means and covariance bilinear forms.

\subsection{\texorpdfstring{Specialize to both kinetic schemes and check zero friction}{Specialize to both kinetic schemes and check zero friction}}


The first-order velocity and position noises use \(k=\psi_0\) and \(k=\psi_1\), respectively. Their covariance is exactly \texttt{klmc\allowbreak{}Covariance \ensuremath{\gamma} t 0 1}, which generally is nonzero. The two-dimensional law theorem supplies the full first-order KLMC innovation, and the four-dimensional theorem uses \((\psi_0,\psi_1,\varphi_2,\varphi_3)\) to supply every correlated innovation needed by the Hessian-corrected second-order scheme. Both are direct specializations of the finite-family joint-law theorem, so each uses one common Brownian motion.



At zero friction, \(\psi_0\equiv1\), so the velocity convolution must reduce to the original Brownian variable \(B_t\). Lean verifies this as an equality in L\ensuremath{{}^{2}}: the time kernel is precisely \(1_{(0,t]}\), the same indicator used in the canonical Brownian construction. The proof compares the actual almost-everywhere representatives with \texttt{Lp.\allowbreak{}ext} and \texttt{indicator\allowbreak{}Const\allowbreak{}Lp\_\allowbreak{}coe\allowbreak{}Fn}. The position and velocity L\ensuremath{{}^{2}}-continuity statements follow from the general continuous-kernel theorem.

\subsection{\texorpdfstring{Justify stochastic Fubini in the deterministic Hilbert space}{Justify stochastic Fubini in the deterministic Hilbert space}}


The essential primitive identity is

\[K_{\psi_1,t}=\int_0^t K_{\psi_0,s}\,ds\quad\text{in }L^2(\mathbb R,dr).\]

Point evaluation on L\ensuremath{{}^{2}} is not continuous, so one cannot simply evaluate a Bochner integral at a time coordinate and interchange the two integrals. The proof instead tests both sides against an arbitrary L\ensuremath{{}^{2}} function \(f\). The inner product with a kernel is the ordinary integral \(\int_0^s f(r)e^{-\gamma(s-r)}\,dr\), a fact obtained from the actual almost-everywhere kernel representatives.



On the finite square \((0,t]^2\), consider \(1_{\{r\le s\}}e^{-\gamma(s-r)}f(r)\). For nonnegative friction the triangular exponential weight is bounded by one. The restriction of an L\ensuremath{{}^{2}} function to a finite-measure interval is integrable, so the product is integrable on the square. This establishes the hypothesis of Fubini's theorem; the interchange is proved rather than assumed. After swapping, the inner integral is

\[\int_r^t e^{-\gamma(s-r)}\,ds=\psi_1(t-r).\]

Thus the two inner products agree for every \(f\), and Hilbert-space extensionality gives the vector identity.



Lean proves the two triangular-domain reductions by converting indicator integrals to restricted measures and explicitly identifying the intersections \((0,t]\cap(-\infty,s]=(0,s]\) and \((0,t]\cap[r,\infty)=[r,t]\). It removes the endpoint distinction using the null singleton property of Lebesgue measure. \texttt{Integrable.\allowbreak{}comp\_\allowbreak{}snd} and \texttt{Integrable.\allowbreak{}bdd\_\allowbreak{}mul} establish product integrability. \texttt{integral\_\allowbreak{}integral\_\allowbreak{}swap} applies Fubini, and \texttt{ext\_\allowbreak{}inner\_\allowbreak{}left} finishes after commuting the continuous inner-product functional with the Bochner integral.

\subsection{\texorpdfstring{Obtain the actual kinetic Brownian integral equation}{Obtain the actual kinetic Brownian integral equation}}


The scalar coefficient identity \(\psi_0(u)+\gamma\psi_1(u)=1\) holds pointwise, including at zero friction. Multiplying it by the time indicator and applying linearity of the isonormal map gives the exact relation

\[V_t^{\mathrm{noise}}+\gamma X_t^{\mathrm{noise}}=B_t.\]

The deterministic Hilbert primitive identity just proved transfers through the continuous linear isonormal map. Consequently,

\[X_t^{\mathrm{noise}}=\int_0^tV_s^{\mathrm{noise}}\,ds,\qquad V_t^{\mathrm{noise}}=B_t-\gamma\int_0^tV_s^{\mathrm{noise}}\,ds.\]

These are equations for the actual L\ensuremath{{}^{2}} random variables whose joint Gaussian innovation law has already been identified. They connect the numerical exponential-integrator noise to the same additive Brownian forcing used by the kinetic diffusion.



For the first identity Lean compares actual representatives of the sum and scalar multiple of the time kernels, then applies \texttt{friction\_\allowbreak{}mul\_\allowbreak{}psi\_\allowbreak{}one}. For the primitive identity, \texttt{Continuous\allowbreak{}Linear\allowbreak{}Map.\allowbreak{}interval\allowbreak{}Integral\_\allowbreak{}comp\_\allowbreak{}comm} transports the deterministic Bochner integral through the isonormal map; its integrability is supplied by the proved L\ensuremath{{}^{2}} time continuity. Finally, ordinary vector-space rearrangement gives the damped Brownian integral equation. No matching-noise law or integral-solution equation is introduced as a premise.

\subsection{\texorpdfstring{Derive the explicit Brownian convolution by an integrating factor}{Derive the explicit Brownian convolution by an integrating factor}}


The proved primitive identity implies \(P'_t=V_t\) for positive time. Combined with \(V_t+\gamma P_t=B_t\), this is the ordinary L\ensuremath{{}^{2}}-valued equation \(P'_t+\gamma P_t=B_t\), with \(P_0=0\). Its integrating factor gives

\[\frac{d}{dt}\bigl(e^{\gamma t}P_t\bigr)=e^{\gamma t}B_t,\qquad P_t=\int_0^t e^{-\gamma(t-s)}B_s\,ds.\]

Substitution into the exact Brownian balance then gives the integration-by-parts representation

\[V_t=B_t-\gamma\int_0^t e^{-\gamma(t-s)}B_s\,ds.\]

Thus the isonormal Wiener integral and the continuous additive-Brownian forcing formulation are proved equal as actual L\ensuremath{{}^{2}} vectors.



Lean first differentiates the indefinite integral defining the primitive. On a neighborhood of a positive time, coercing through \texttt{Real.\allowbreak{}to\allowbreak{}NNReal} does not change the endpoint, so \texttt{Has\allowbreak{}Deriv\allowbreak{}At.\allowbreak{}congr\_\allowbreak{}of\_\allowbreak{}eventually\allowbreak{}Eq} transfers the derivative to \texttt{position\allowbreak{}Lp}. The product rule differentiates the scalar exponential times the L\ensuremath{{}^{2}} vector. The Brownian balance cancels the damping term. The fundamental theorem of calculus on \([0,t]\), requiring derivatives only in its interior, then integrates this derivative. Multiplication by \(e^{-\gamma t}\) cancels the integrating factor. Finally, scalar multiplication commutes with the interval integral, and the exponential addition law converts \(e^{-\gamma t}e^{\gamma s}\) into \(e^{-\gamma(t-s)}\). This proof uses no sample-path differentiability of Brownian motion.

\section{\texorpdfstring{Construct independent correlated innovations on Brownian time intervals}{Construct independent correlated innovations on Brownian time intervals}}
\label{proof-module-KineticLangevin-BrownianSegments-lean}
{\small\color{gray}Source module: \texttt{Kinetic\allowbreak{}Langevin/\allowbreak{}Brownian\allowbreak{}Segments.\allowbreak{}lean}.\par}


The one-step Gaussian covariance is only part of the algorithm: innovations at different iterations must be independent and must come from increments of the same Brownian source used by the exact reference. This module translates deterministic time kernels onto successive intervals of that source. Translation preserves the within-step Gram covariance. Disjoint time supports give zero cross covariances, and joint Gaussianity turns this orthogonality into independence of the entire innovation blocks. Correlations among coordinates within a block are retained.

\subsection{\texorpdfstring{Translate a time kernel without changing its L\ensuremath{{}^{2}} geometry}{Translate a time kernel without changing its L² geometry}}


For an interval beginning at \(a\), translate the original response kernel by \((S_aK)(r)=K(r-a)\). Lebesgue measure is invariant under translation, so \(S_a\) is a real linear isometry of the deterministic time Hilbert space. The translated response is

\[K^{a,t}_k(r)=1_{(a,a+t]}(r)k(a+t-r).\]

Its inner product with another response on the same interval remains \(\int_0^t k(s)\ell(s)\,ds\).



Lean constructs the isometry with \texttt{Lp.\allowbreak{}comp\allowbreak{}Measure\allowbreak{}Preserving\ensuremath{{}_{li}}} and \texttt{measure\allowbreak{}Preserving\_\allowbreak{}sub\_\allowbreak{}right}. The representative theorem combines the almost-everywhere composition formula with the previously proved kernel representative. Elementary inequalities identify \(r-a\in(0,t]\) with \(r\in(a,a+t]\); scalar algebra identifies the two elapsed-time expressions. Inner-product preservation of the linear isometry then reduces the covariance geometry to the unshifted kernel identity.

\subsection{\texorpdfstring{Prove that disjoint time segments are orthogonal}{Prove that disjoint time segments are orthogonal}}


If \(a+t\le b\), the intervals \((a,a+t]\) and \((b,b+u]\) are disjoint. At every time coordinate, at least one of their indicators is zero. Their L\ensuremath{{}^{2}} inner product is therefore zero, regardless of the two continuous response functions.



The Lean proof unfolds the L\ensuremath{{}^{2}} inner product as an integral, substitutes the explicit almost-everywhere segment representatives, and splits on membership in the first interval. Membership implies \(r\le a+t\le b\), which excludes the strict lower bound \(b<r\) of the second interval. Thus the integrand vanishes almost everywhere, and \texttt{integral\_\allowbreak{}eq\_\allowbreak{}zero\_\allowbreak{}of\_\allowbreak{}ae} concludes. The half-open convention makes adjacent segments disjoint exactly, including their common endpoint.

\subsection{\texorpdfstring{Apply the same isonormal map to every time segment}{Apply the same isonormal map to every time segment}}


Define \(Z_k^{a,t}=I(K_k^{a,t})\) using the same isonormal map and Gaussian product sample space as the original Brownian motion. Each response is centered, and every finite family of responses on an interval is jointly Gaussian. The covariance on one interval is the same Gram integral as before. Responses on separated intervals have covariance zero by the deterministic support calculation.



Lean proves these assertions by composing the translation's linear map with the existing Gaussian isonormal map. The finite-family Gaussian theorem applies to that actual composite linear map. For covariances, \texttt{isonormal\_\allowbreak{}covariance} reduces the probability calculation to the deterministic inner product, where the within-interval and disjoint-interval theorems apply. Independence is not assumed at this stage; zero covariance alone would not imply independence without the jointly Gaussian property proved next.

\subsection{\texorpdfstring{Identify the full Gaussian block law on every interval}{Identify the full Gaussian block law on every interval}}


For finitely many response functions \((k_i)\), the Euclidean vector \((Z_{k_i}^{a,t})_i\) has exactly the centered Gaussian law with covariance \(C_{ij}=\int_0^t k_i(s)k_j(s)\,ds\). The starting time \(a\) does not appear in this law, because deterministic translation preserves all inner products.



Lean assembles the actual L\ensuremath{{}^{2}} representatives into a Euclidean random vector and proves its Gaussianity via the continuous linear equivalence from finite function space. Coordinate projections commute with the Bochner integral, giving the zero vector mean. The covariance bilinear form expands into a finite double sum; each entry is rewritten by the proved segment covariance formula. Equality with the Gram Gaussian's bilinear form, together with equality of means, lets \texttt{Is\allowbreak{}Gaussian.\allowbreak{}ext} identify the two probability measures. No distributional equality is supplied as a hypothesis.

\subsection{\texorpdfstring{Derive independence of all successive Brownian innovation blocks}{Derive independence of all successive Brownian innovation blocks}}


On the uniform grid, the \(n\)-th block uses the interval \((nh,(n+1)h]\). \texttt{grid\allowbreak{}Convolution\allowbreak{}Lp} places every response function on that interval. The family indexed by both iteration and response coordinate is a Gaussian process because every member is obtained by applying one Gaussian linear map to a deterministic Hilbert vector.



For \(n<m\), the elementary inequality \((n+1)h\le mh\) proves disjoint support and therefore zero cross covariance between every coordinate of the two blocks. The Gaussian-process independence theorem converts these coordinate covariance equalities into independence of the entire block-valued random variables, over all iteration numbers simultaneously. Packing a finite coordinate family into Euclidean space is a measurable map, so time independence survives that packing.



Lean constructs the jointly Gaussian process with \texttt{is\allowbreak{}Gaussian\allowbreak{}Process\_\allowbreak{}of\_\allowbreak{}linear\allowbreak{}Map\_\allowbreak{}Lp}. It uses \texttt{i\allowbreak{}Indep\allowbreak{}Fun\_\allowbreak{}of\_\allowbreak{}covariance\_\allowbreak{}eq\_\allowbreak{}zero}, proves the disjoint-segment inequality from \texttt{Nat.\allowbreak{}succ\_\allowbreak{}le\_\allowbreak{}of\_\allowbreak{}lt}, and handles the reversed ordering through symmetry of covariance. Finally, \texttt{i\allowbreak{}Indep\allowbreak{}Fun.\allowbreak{}comp} transports independence through \texttt{With\allowbreak{}Lp.\allowbreak{}to\allowbreak{}Lp}. The per-block law is the segment joint-law theorem, so the actual Brownian-derived Euclidean blocks are identically distributed with the required correlated Gram-Gaussian law.

\section{\texorpdfstring{The exact Brownian equations on every numerical time interval}{The exact Brownian equations on every numerical time interval}}
\label{proof-module-KineticLangevin-BrownianSegmentEquations-lean}
{\small\color{gray}Source module: \texttt{Kinetic\allowbreak{}Langevin/\allowbreak{}Brownian\allowbreak{}Segment\allowbreak{}Equations.\allowbreak{}lean}.\par}


The numerical method uses fresh correlated noise on each interval \((a,a+t]\). The preceding module constructs those noises on the same canonical Brownian source and proves independence across disjoint intervals. Here the actual noises also satisfy the integrated kinetic equations, so their relationship to the continuous diffusion is an equality in L\ensuremath{{}^{2}} rather than a matching-law assumption.

\subsection{\texorpdfstring{Identify the Brownian increment on a translated interval}{Identify the Brownian increment on a translated interval}}


Define \(W^a_t=B_{a+t}-B_a\). The unit response kernel on \((a,a+t]\) is exactly the difference \(1_{(0,a+t]}-1_{(0,a]}\). Lean proves the indicator identity almost everywhere using \texttt{Lp.\allowbreak{}ext}, the explicit representatives, and real order arithmetic. Linearity of the isonormal Wiener-integral map then identifies its image with the actual Brownian increment. This fixes the driving noise as a constructed object on the original probability space.

\subsection{\texorpdfstring{Integrate the shifted velocity response}{Integrate the shifted velocity response}}


Write \(V^a_t=\int_a^{a+t}e^{-\gamma(a+t-r)}\,dB_r\) and \(P^a_t=\int_a^{a+t}\psi_1(a+t-r)\,dB_r\). Translation of deterministic time kernels is a continuous linear isometry. Therefore continuity of the unshifted kernels gives continuity of each translated stochastic response in L\ensuremath{{}^{2}}. The already proved deterministic primitive identity \(K_{\psi_1}(t)=\int_0^tK_{\psi_0}(s)ds\) survives translation and application of the isonormal map. The Lean proof composes their continuous linear maps and invokes \texttt{interval\allowbreak{}Integral\_\allowbreak{}comp\_\allowbreak{}comm}, proving \(P^a_t=\int_0^tV^a_sds\) without pointwise evaluation of an L\ensuremath{{}^{2}} class.

\subsection{\texorpdfstring{Prove the actual damped Brownian integral equation}{Prove the actual damped Brownian integral equation}}


The exact coefficient identity \(\psi_0+\gamma\psi_1=1\) gives \(V^a_t+\gamma P^a_t=W^a_t\). Lean first uses injectivity of the isonormal linear isometry to recover the deterministic kernel identity from the unshifted balance, then applies time translation and the same map. Substituting the primitive gives \(V^a_t=W^a_t-\gamma\int_0^tV^a_sds\). This equality supplies the additive forcing needed when a continuous reference trajectory and a numerical step share one Brownian interval.

\subsection{\texorpdfstring{Recover the ordinary Brownian convolution formula}{Recover the ordinary Brownian convolution formula}}


The position response also equals \(P^a_t=\int_0^te^{-\gamma(t-s)}W^a_sds\), hence \(V^a_t=W^a_t-\gamma\int_0^te^{-\gamma(t-s)}W^a_sds\). To prove this for the exact translated construction, Lean pulls the established unshifted formula back through the injective isonormal isometry, then sends the deterministic identity through the translated continuous linear map. Continuity ensures Bochner integrability. Commuting that map with the integral and its scalar factors yields the formula; the velocity version follows by the balance identity. These are exact integration-by-parts identities on the constructed Brownian source.

\section{\texorpdfstring{The canonical vector Brownian innovations and initial Gaussian velocity}{The canonical vector Brownian innovations and initial Gaussian velocity}}
\label{proof-module-KineticLangevin-CanonicalVectorSegments-lean}
{\small\color{gray}Source module: \texttt{Kinetic\allowbreak{}Langevin/\allowbreak{}Canonical\allowbreak{}Vector\allowbreak{}Segments.\allowbreak{}lean}.\par}


The paper draws correlated velocity and position innovations in each spatial coordinate, independently across coordinates and time steps. Here every such draw is realized from the existing product of scalar canonical Brownian sources. The entire constructed array is proved to have exactly the previously defined Gaussian product law; equality of a few moments alone is not used as a substitute for equality of laws.

\subsection{\texorpdfstring{Transport independence along an actual law}{Transport independence along an actual law}}


A measurable map \(Y\) with actual law \(Q\) can pull back any independent measurable family defined on \(Q\). Lean proves this by the characterization of mutual independence through equality with the infinite product of marginal laws. It applies the pushforward composition theorem to the joint family and to each marginal. This elementary tool is used repeatedly to transfer established independence to the concrete Brownian source.

\subsection{\texorpdfstring{Build correlated scalar and spatial blocks}{Build correlated scalar and spatial blocks}}


A scalar block consists of all deterministic response kernels applied to the same Brownian interval. Packing their values with \texttt{With\allowbreak{}Lp.\allowbreak{}to\allowbreak{}Lp} gives a finite Euclidean vector whose law is the Gram Gaussian already proved in \texttt{Brownian\allowbreak{}Segments}. A vector block evaluates that scalar block on each spatial coordinate source. Every map is measurable because each L\ensuremath{{}^{2}} representative is strongly measurable, coordinate evaluation is measurable, and finite or infinite coordinatewise packing preserves measurability.

\subsection{\texorpdfstring{Prove the law of the entire array}{Prove the law of the entire array}}


For one spatial source, disjoint Brownian intervals give independent scalar blocks. Different spatial sources are independent by the explicit product construction. The theorem \texttt{i\allowbreak{}Indep\allowbreak{}Fun\_\allowbreak{}uncurry'} combines these two facts into independence over iteration-coordinate pairs; reindexing swaps their order. Each marginal has the established Gram-Gaussian law, so \texttt{has\allowbreak{}Law\_\allowbreak{}infinite\allowbreak{}Pi} identifies the full array distribution. Currying the array is a measurable equivalence, and the product-measure curry theorem identifies the distribution as an independent stream of complete spatial blocks. Pullback along this actual path law proves independence across time steps.

\subsection{\texorpdfstring{Identify the existing KLMC Gaussian stream exactly}{Identify the existing KLMC Gaussian stream exactly}}


For \(d\) finite spatial coordinates, the pair-indexed array has exactly the existing \texttt{kernel\allowbreak{}Noise\allowbreak{}Path} law. Specializing the response family to \((\psi_0,\psi_1)\) and packing each response into a vector produces the actual unscaled velocity-position innovation. Composition of the stream law with the existing \texttt{klmc\allowbreak{}Innovation} map proves its exact distribution. Independence of the full vector blocks survives this measurable packing. Each vector coordinate lies in L\ensuremath{{}^{2}} by the Brownian isometry; pulling it back along spatial evaluation and using \texttt{Mem\allowbreak{}Lp.\allowbreak{}of\_\allowbreak{}eval\_\allowbreak{}pi\allowbreak{}Lp} proves both phase components have finite second moments.

\subsection{\texorpdfstring{Adjoin the independent initial Gaussian velocity}{Adjoin the independent initial Gaussian velocity}}


The source for a sampler is \(Q=\mathcal N(0,I_d)\otimes P_{\mathrm{Brownian}}\). Its first projection is the actual standard Gaussian initial velocity. Its second projection carries all constructed innovations. Product-measure projection laws prove the velocity law and preserve the innovation law. The product independence theorem proves that initial velocity is independent of the entire future innovation stream, while measure-preserving pullback gives all required L\ensuremath{{}^{2}} bounds. An arbitrary initial-position coupling can subsequently be extended by this same fixed source, independently of that coupling.



Pullback along the second projection also proves mutual independence of complete innovations across all iterations on this enlarged source; the new Gaussian velocity does not alter the noise law.

\section{\texorpdfstring{Continuous vector Brownian responses in L\ensuremath{{}^{2}}}{Continuous vector Brownian responses in L²}}
\label{proof-module-KineticLangevin-VectorSegmentEquations-lean}
{\small\color{gray}Source module: \texttt{Kinetic\allowbreak{}Langevin/\allowbreak{}Vector\allowbreak{}Segment\allowbreak{}Equations.\allowbreak{}lean}.\par}


The scalar Brownian equations must also hold for Euclidean random vectors so they can drive the actual L\ensuremath{{}^{2}} kinetic diffusion used by the contraction proof. A continuous linear spatial assembly map achieves this directly: it preserves sums, scalar multiples, continuity, and Bochner time integrals.

\subsection{\texorpdfstring{Assemble scalar L\ensuremath{{}^{2}} variables through a continuous linear map}{Assemble scalar L² variables through a continuous linear map}}


For a scalar random variable \(Z\), spatial assembly produces the vector \((Z(\omega_i))_i\), evaluating it on independent spatial copies of the canonical source. The Lean definition is a finite sum. Each summand first pulls back along a measure-preserving coordinate evaluation, then inserts the scalar into one Euclidean basis coordinate. Both operations are continuous linear maps on L\ensuremath{{}^{2}}. The representative theorem is proved by the finite-sum almost-everywhere identity and the explicit representatives of these two maps; Euclidean extensionality verifies each coordinate. This construction avoids interchanging pointwise evaluation with an L\ensuremath{{}^{2}} Bochner integral.

\subsection{\texorpdfstring{Define vector responses and prove continuity}{Define vector responses and prove continuity}}


The vector response is the spatial assembly of the actual scalar response on \((a,a+t]\). The vector Brownian input is the assembly of the scalar Brownian increment. Their definitions remain on the existing canonical product source. The response representative theorem gives exactly the expected independent-coordinate vector. Continuity in time follows by composing the continuous scalar response with the continuous linear assembly map.

\subsection{\texorpdfstring{Lift the exact primitive and damped integral equations}{Lift the exact primitive and damped integral equations}}


The position response satisfies \(P^a_t=\int_0^tV^a_sds\), and the velocity response satisfies \(V^a_t=W^a_t-\gamma\int_0^tV^a_sds\). These equalities follow from the already proved scalar equations by applying the spatial assembly map. The proof commutes that continuous linear map with each Bochner integral using \texttt{interval\allowbreak{}Integral\_\allowbreak{}comp\_\allowbreak{}comm}. Thus the exact filtered Brownian input required by the generic mild-equation theorem is constructed for the vector L\ensuremath{{}^{2}} diffusion.

\subsection{\texorpdfstring{Identify the input with canonical vector Brownian increments}{Identify the input with canonical vector Brownian increments}}


The lifted input is proved almost everywhere equal to the vector whose coordinates are \(B_i(a+t)-B_i(a)\) for the existing canonical coordinate process. The scalar subtraction identity is pulled back through each measure-preserving spatial projection and combined with the spatial representative theorem. Continuity follows by identifying the increment with the unit-kernel vector response. This connects the L\ensuremath{{}^{2}} forcing to the actual canonical vector Brownian construction.

\subsection{\texorpdfstring{Identify the OU increment and numerical innovation endpoints}{Identify the OU increment and numerical innovation endpoints}}


The vector velocity response is exactly \(V^a_t=W^a_t-\gamma\int_0^te^{-\gamma(t-s)}W^a_sds\). Lean applies spatial assembly to the established scalar integration-by-parts formula, uses its linearity, and commutes it with the continuous Bochner integrand. This is precisely the OU input increment used by the kinetic splitting limit. At each grid interval, the paired L\ensuremath{{}^{2}} velocity and position response endpoints are proved almost everywhere equal to the pointwise \texttt{canonical\allowbreak{}KLMCInnovation} already equipped with the exact Gaussian law. The endpoint proof combines the two actual representative equalities with \texttt{Prod.\allowbreak{}ext}; it does not replace the constructed inputs by new independent Gaussian variables.

\section{\texorpdfstring{Splitting convergence on the actual constructed vector Brownian source}{Splitting convergence on the actual constructed vector Brownian source}}
\label{proof-module-KineticLangevin-CanonicalSplitting-lean}
{\small\color{gray}Source module: \texttt{Kinetic\allowbreak{}Langevin/\allowbreak{}Canonical\allowbreak{}Splitting.\allowbreak{}lean}.\par}


The deterministic input in the splitting theorem is now identified with the actual canonical vector Brownian motion in L\ensuremath{{}^{2}}. Its filtered increments are exactly the constructed Gaussian segment convolutions, including the fluctuation-dissipation scaling. Consequently the innovation-driven numerical chain and the continuous-input chain coincide on the same source, and the actual L\ensuremath{{}^{2}} chain converges to its kinetic integral solution.

\subsection{\texorpdfstring{Keep scaling and source lifts exact}{Keep scaling and source lifts exact}}


The OU input construction is linear in the additive input. Multiplying \(B\) by a real coefficient \(c\) multiplies both its matching increment and its exponential convolution by \(c\). More generally, a continuous linear map \(L\) commutes with this construction: it commutes with subtraction, scalar multiplication, and the Bochner interval integral of the continuous integrand. Lean proves these exact identities by integral linearity and continuous-linear integral commutation. The latter applies, in particular, to measure-preserving pullbacks that extend the Brownian source by an independent initial state.

\subsection{\texorpdfstring{Construct the actual continuous L\ensuremath{{}^{2}} Brownian input}{Construct the actual continuous L² Brownian input}}


The input is the already constructed vector Brownian increment from time zero, evaluated at \(t_+=\max\{t,0\}\). This supplies a globally defined real-time function, anchored at zero. The proved continuity of vector Brownian increments in L\ensuremath{{}^{2}} and continuity of \(t\mapsto t_+\) prove continuity of this input. At zero the underlying increment is the difference of a Brownian value with itself, so it is exactly zero. Lean unfolds the actual source definitions to check this anchor.

\subsection{\texorpdfstring{Identify actual input increments and OU segment responses}{Identify actual input increments and OU segment responses}}


The difference between the anchored values at \(a+t\) and \(a\) cancels their common time-zero value, leaving the actual canonical Brownian increment on \((a,a+t]\). The independently proved segment convolution identity gives \[Z_{a,t}=\Delta B_{a,t}-\gamma\int_0^t e^{-\gamma(t-s)}\Delta B_{a,s}\,ds.\] This is exactly the deterministic \texttt{ou\allowbreak{}Input\allowbreak{}Increment} used in the splitting theorem. Lean rewrites the endpoint increment, compares the actual interval integrands, and uses nonnegativity of the integration time to remove its real-to-nonnegative-time clamp. Scaling by any \(c\), including \(\sqrt{2\gamma}\), follows from the checked linearity identity. The equality is of actual L\ensuremath{{}^{2}} variables, stronger than merely matching marginal Gaussian laws.

\subsection{\texorpdfstring{Identify the entire canonical innovation chain}{Identify the entire canonical innovation chain}}


On a mesh of nonnegative size \(h\), the continuous-input chain uses the increment \(N_{kh,h}\). The preceding actual identity replaces this by \(cZ_{kh,h}\) for every grid index. The explicit-innovation recursion therefore agrees exactly with the continuous-input recursion. Lean first uses the proved equality between the two chain representations and then substitutes the actual OU increment identity at each nonnegative grid time.

\subsection{\texorpdfstring{Conclude actual L\ensuremath{{}^{2}} convergence on the canonical source}{Conclude actual L² convergence on the canonical source}}


Apply the generic Banach splitting theorem to the continuous canonical input \(cB\), its proved zero initial value, and the actual continuous kinetic integral solution \(Y\) driven by that same input. Rewrite the resulting continuous-input chain by the proved canonical innovation-chain identity. This yields convergence of the explicit chain with actual Gaussian convolution innovations on the mesh \(T/(n+1)\) to \(Y_T\). Lean checks the nonnegative-time casts in the mesh and grid endpoints. The integral-solution premise is the actual equation constructed elsewhere by Picard existence; no stationarity or numerical convergence is supplied as a contract.

\section{\texorpdfstring{Theorem 1 for actual Brownian kinetic diffusions}{Theorem 1 for actual Brownian kinetic diffusions}}
\label{proof-module-KineticLangevin-CanonicalContinuousMixing-lean}
{\small\color{gray}Source module: \texttt{Kinetic\allowbreak{}Langevin/\allowbreak{}Canonical\allowbreak{}Continuous\allowbreak{}Mixing.\allowbreak{}lean}.\par}


The continuous mixing theorem is instantiated with the existing genuine vector Brownian source. Its initial velocity distribution is arbitrary, subject only to a finite second moment, and is independent of the initial positions and Brownian noise. The output measures are actual terminal pushforwards of the selected kinetic integral solution, on any finite coordinate set.

\subsection{\texorpdfstring{Construct arbitrary initial velocity and genuine Brownian forcing}{Construct arbitrary initial velocity and genuine Brownian forcing}}


The source is the independent product of the prescribed velocity distribution and the canonical vector Brownian measure. Its first projection gives the actual L\ensuremath{{}^{2}} initial velocity. The second projection pulls back the canonical Brownian process, multiplied by the fluctuation-dissipation factor \(\sqrt{2\gamma}\). Continuity and the zero initial value follow from the corresponding proved Brownian facts and linear source pullback.

\subsection{\texorpdfstring{Define the actual position laws and prove arbitrary-\allowbreak{}shift mixing}{Define the actual position laws and prove arbitrary-shift mixing}}


For each initial position distribution, the independent-source construction solves the actual kinetic equation with the original gradient and genuine Brownian input. Its terminal position pushforward defines the law. Applying the general product-source coupling theorem proves the complete arbitrary-shift exponential coefficient. The common velocity law needs no Gaussian assumption. The initial transport term is the true Wasserstein distance between the initial position distributions, obtained by infimizing over genuine couplings rather than assuming an optimizer or a prescribed coupling cost.

\section{\texorpdfstring{Exact Gaussian invariance of the velocity thermostat}{Exact Gaussian invariance of the velocity thermostat}}
\label{proof-module-KineticLangevin-Thermostat-lean}
{\small\color{gray}Source module: \texttt{Kinetic\allowbreak{}Langevin/\allowbreak{}Thermostat.\allowbreak{}lean}.\par}


The O substep removes velocity variance through damping and restores exactly that variance through independent Gaussian noise. This module proves an equality of probability measures, not only a matching second moment. Invariance of this substep is one component of kinetic stationarity; the Hamiltonian evolution and the full diffusion require separate arguments.

\subsection{\texorpdfstring{Independent Gaussian linear combinations}{Independent Gaussian linear combinations}}


Let \(G,G\prime\) be independent standard Gaussian vectors in a finite-dimensional real Hilbert space, and suppose \(a^2+b^2=1\). For a continuous linear functional \(L\), the characteristic function of \(aG+bG\prime\) is

\[\exp(-a^2\|L\|^2/2)\exp(-b^2\|L\|^2/2)=\exp(-\|L\|^2/2).\]

Independence is represented by the product of two standard Gaussian measures. Lean constructs the continuous linear map \(A(v,w)=av+bw\), uses \texttt{char\allowbreak{}Fun\allowbreak{}Dual\_\allowbreak{}map} to pull the test functional back through \(A\), and uses \texttt{char\allowbreak{}Fun\allowbreak{}Dual\_\allowbreak{}prod} to factor its characteristic function. The restrictions of \(L\circ A\) to the two factors are \(aL\) and \(bL\), proved by extensionality. The norm of a scalar multiple supplies the squared coefficients. The relation \(a^2+b^2=1\) then simplifies the exponent, and uniqueness of characteristic functions proves the distributional identity.

\subsection{\texorpdfstring{Apply the fluctuation-\allowbreak{}dissipation balance}{Apply the fluctuation-dissipation balance}}


For nonnegative friction and time step, put \(\delta=e^{-\gamma h/2}\) and \(s=\sqrt{1-\delta^2}\). The variance under the square root is nonnegative, so \(s^2=1-\delta^2\). Hence \(\delta^2+s^2=1\), and the preceding theorem gives \(\delta G+sG\prime\sim\mathcal N(0,I)\). Lean applies the general linear-combination theorem, rewrites the square root squared using its proved nonnegativity, unfolds the variance, and finishes with a polynomial identity. This is precisely the standard-temperature O half-step used in OBABO.

\section{\texorpdfstring{The actual OU innovation law and independent time stream}{The actual OU innovation law and independent time stream}}
\label{proof-module-KineticLangevin-OUMarginal-lean}
{\small\color{gray}Source module: \texttt{Kinetic\allowbreak{}Langevin/\allowbreak{}OUMarginal.\allowbreak{}lean}.\par}


The canonical convolution is an actual random variable on the constructed Brownian source. This module proves that its fluctuation-dissipation scaling has exactly the thermal Gaussian law, and that full spatial vectors on distinct grid intervals are independent. The complete stream therefore has the independent Gaussian product law required by the concrete Markov transition.

\subsection{\texorpdfstring{Identify independent spatial copies through their actual scalar law}{Identify independent spatial copies through their actual scalar law}}


A scalar L\ensuremath{{}^{2}} variable \(Z\) is evaluated on each independent coordinate Brownian source. Each copy has the same proved law \(\nu\), and the coordinate product construction proves their independence. The infinite-product law theorem, specialized to finitely many spatial coordinates, gives the finite product measure \(\nu^{\otimes\iota}\). Transporting this law through the actual Euclidean coordinate map identifies the law of \texttt{spatial\allowbreak{}Lift}; its representative theorem supplies the almost-everywhere equality. For \(\nu\) equal to a scaled real standard Gaussian, the product pushforward identity and the definition-independent standard Gaussian theorem identify this law as the same scalar scaling of the Euclidean standard Gaussian.

\subsection{\texorpdfstring{Compute the actual scalar thermal variance}{Compute the actual scalar thermal variance}}


The scalar segment convolution is the constructed isonormal map applied to its actual time kernel \(K\), so it has centered Gaussian variance \(\|K\|^2\). With \(c=\sqrt{2\gamma}\), its scaled variance is \(2\gamma\|K\|^2\). The proved kernel inner-product identity gives \(\|K\|^2=\int_0^h e^{-2\gamma s}\,ds\), and the exact fluctuation-decay balance gives \(2\gamma\|K\|^2=1-e^{-2\gamma h}\). This equals \(\sigma_h^2\) for \(\sigma_h=\sqrt{1-e^{-2\gamma h}}\). Lean derives nonnegativity from the damping bound, uses the actual isonormal Gaussian law, checks the equality of nonnegative-real variances, and rewrites the scaled real Gaussian pushforward. Zero friction is included without division.

\subsection{\texorpdfstring{Identify the entire spatial OU Gaussian law}{Identify the entire spatial OU Gaussian law}}


Apply the independent-spatial-copy result to the proved scalar thermal law. Linearity of \texttt{spatial\allowbreak{}Lift} identifies its scaled scalar input with the actual scaled vector segment convolution. Thus \(\sqrt{2\gamma}\,Z_{a,h}\) has exactly the law of \(\sqrt{1-e^{-2\gamma h}}\,\xi\), where \(\xi\) is the Euclidean standard Gaussian. This proves the full vector distribution rather than only its covariance.

\subsection{\texorpdfstring{Prove independence of actual full-\allowbreak{}vector innovations across time}{Prove independence of actual full-vector innovations across time}}


The canonical KLMC innovation pairs are already proved independent over disjoint time-grid intervals by their actual Brownian kernel supports. Project each pair to its velocity component and apply any fixed scalar scaling. The actual L\ensuremath{{}^{2}} vector-segment representatives agree almost everywhere with these projected canonical innovations. Independence is preserved under these measurable projections and almost-everywhere changes. Lean explicitly combines the representative scalar-multiplication event with the canonical endpoint identity.

\subsection{\texorpdfstring{Identify the full independent Gaussian stream}{Identify the full independent Gaussian stream}}


Every scaled OU grid innovation has the exact thermal Gaussian law, and the preceding theorem proves independence of all grid vectors. The infinite-product law theorem therefore identifies the entire actual innovation stream with the product of this Gaussian law over all iterations. Measurability follows from the strongly measurable L\ensuremath{{}^{2}} representative at each countable grid index. This stream identity is the source-level fact used to transfer the concrete Brownian-driven splitting chain to its iterated Markov laws.

\section{\texorpdfstring{Canonical kinetic innovations on arbitrary finite coordinate sets}{Canonical kinetic innovations on arbitrary finite coordinate sets}}
\label{proof-module-KineticLangevin-GenericCanonicalSource-lean}
{\small\color{gray}Source module: \texttt{Kinetic\allowbreak{}Langevin/\allowbreak{}Generic\allowbreak{}Canonical\allowbreak{}Source.\allowbreak{}lean}.\par}


Particle states are naturally indexed by a finite disjoint union of parameter and particle coordinates. These constructions work directly for any finite index type, while retaining the same actual canonical Brownian source, correlated Gaussian innovations and independent Gaussian velocity.

\subsection{\texorpdfstring{Construct the genuine independent velocity/\allowbreak{}Brownian source}{Construct the genuine independent velocity/Brownian source}}


The source is the explicit product of a standard Euclidean Gaussian velocity and the existing independent-coordinate Brownian source. Product projection laws prove the velocity distribution and its second moment. The phase innovation packs the two actual exponential-response coordinates from every spatial Brownian interval. Measurability, temporal independence and L\ensuremath{{}^{2}} bounds follow from the generic spatial-block construction. Pullback along the source product projection gives the corresponding L\ensuremath{{}^{2}} innovation components on the enlarged source.

\subsection{\texorpdfstring{Identify the actual correlated Gaussian phase law and its complete stream}{Identify the actual correlated Gaussian phase law and its complete stream}}


The scalar block on each coordinate has the two-dimensional Gram-Gaussian law. The already proved spatial/\allowbreak{}time array law identifies each complete spatial block as the product of these scalar laws. Packing its first and second response coordinates gives exactly the correlated phase innovation used by the algorithm. Mutual independence across time identifies the entire implemented innovation stream with the infinite product of that same phase law. The proof retains the correlation between velocity and position in each coordinate.

\subsection{\texorpdfstring{Connect actual L\ensuremath{{}^{2}} endpoints and the exact thermal OU stream}{Connect actual L² endpoints and the exact thermal OU stream}}


The two vector-response representative equalities prove that the L\ensuremath{{}^{2}} endpoint pair is the pointwise canonical innovation. Taking and scaling its velocity component preserves temporal independence. The exact generic vector OU marginal theorem identifies its scale as \(\sqrt{1-\psi_0(h)^2}\). The independent-family law theorem then identifies the entire actual OU stream with the iid thermal Gaussian product required by the kinetic Gibbs stationarity proof.

\section{\texorpdfstring{Stitched exact segments are the actual global kinetic solution}{Stitched exact segments are the actual global kinetic solution}}
\label{proof-module-KineticLangevin-SegmentGlobal-lean}
{\small\color{gray}Source module: \texttt{Kinetic\allowbreak{}Langevin/\allowbreak{}Segment\allowbreak{}Global.\allowbreak{}lean}.\par}


A reference defined one interval at a time must be identified with the actual continuous kinetic process whose stationarity is proved separately. The bridge here uses the true shifted-input equation and uniqueness. It does not assume a Markov transition formula or independence of arbitrarily chosen reference segments.

\subsection{\texorpdfstring{Identify a selected local solution with a shifted global path}{Identify a selected local solution with a shifted global path}}


Start with an actual continuous global solution \(Y\), and suppose the local input is the actual increment \(W_t=B_{a+t}-B_a\). The previously proved time-shift lemma shows \(Y_{a+t}\) satisfies the same kinetic integral equation on the local horizon. The selected local solution starts at \(Y_a\). Lean checks the two initial conditions, the shared additive input, and continuity, then invokes the proved uniqueness theorem for the globally Lipschitz kinetic drift. Their equality follows at every time in the local interval.

\subsection{\texorpdfstring{Identify every stitched grid point with the global solution}{Identify every stitched grid point with the global solution}}


The reference chain starts at \(Y_0\). Inductively assume its state at step \(n\) equals \(Y_{nh}\). The next selected segment begins at that same point, so the local-global uniqueness result identifies its endpoint with \(Y_{nh+h}\). Lean performs induction over the actual chain recursion, checks the finite-horizon inequalities using nonnegative step size, and simplifies the real cast of \(n+1\). The result is the exact identity \(R_n=Y_{nh}\), including zero step size.

\subsection{\texorpdfstring{Identify all intermediate times within a reference segment}{Identify all intermediate times within a reference segment}}


Once the starting grid point is identified, the same uniqueness result proves \(Z_n(t)=Y_{nh+t}\) throughout the whole segment. This lets a stationary global phase law supply both the stationary reference position at grid points and the Gaussian reference velocity at intermediate times required by the KLMC local-error proof. No stationary law is assumed for a newly invented process: the stitched reference is proved to be the actual global solution.

\section{\texorpdfstring{Derive the actual frozen-\allowbreak{}force local defect}{Derive the actual frozen-force local defect}}
\label{proof-module-KineticLangevin-LocalError-lean}
{\small\color{gray}Source module: \texttt{Kinetic\allowbreak{}Langevin/\allowbreak{}Local\allowbreak{}Error.\allowbreak{}lean}.\par}


The velocity defect is an explicit exponentially weighted integral of the changing force; the position defect is its primitive. Bounds follow from a Lipschitz force and a kinematic identity, not from a built-in defect hypothesis. The ambient norm may represent a deterministic Banach path or an L\ensuremath{{}^{2}} random-vector norm. A kinetic stationarity theorem supplying the required velocity-moment bound remains separate.

\subsection{\texorpdfstring{Define the force-\allowbreak{}freezing remainder as an integral}{Define the force-freezing remainder as an integral}}


Set \(D_v(t)=\int_0^te^{-\gamma(t-s)}(g(x_s)-g(x_0))ds\), and \(D_x(t)=\int_0^tD_v(s)ds\). These are exactly the terms left when one compares an exact kinetic equation with a frozen-force interpolation driven by the same additive noise. \texttt{velocity\allowbreak{}Defect} and \texttt{position\allowbreak{}Defect} define those vector-valued interval integrals. The definitions impose no desired error estimate. Interpreting them as a numerical-process difference additionally requires deriving the variation-of-constants identity for that process.

\subsection{\texorpdfstring{Integrating velocity gives a linear position increment}{Integrating velocity gives a linear position increment}}


Assume \(x_t-x_0=\int_0^tv_sds\), with \(\|v_s\|\le C\) on the horizon. The norm-of-integral bound gives \(\|x_t-x_0\|\le C|t|=Ct\) for nonnegative time. Lean rewrites with the actual kinematic identity, applies the constant bound on the interval integral, and translates membership in the unoriented integration interval into \(0\le s\le t\le T\). This is the extra position regularity that kinetic Langevin gains from integrating velocity instead of receiving direct Brownian position noise.

\subsection{\texorpdfstring{Integrate a linear force increment to obtain the quadratic velocity defect}{Integrate a linear force increment to obtain the quadratic velocity defect}}


Lipschitzness gives \(\|g(x_s)-g(x_0)\|\le MCs\). With \(\gamma\ge0\) and \(s\le t\), the exponential kernel is between zero and one, so the norm of the weighted integrand is still at most \(MCs\). The proof derives this pointwise estimate from the preceding increment theorem, explicit exponential positivity and the norm of scalar multiplication. A continuous scalar majorant provides the required interval integrability. Integrating the bound gives \(\|D_v(t)\|\le\int_0^tMCsds=MCt^2/2\); Lean evaluates the integral and normalizes its coefficients with \texttt{ring}.

\subsection{\texorpdfstring{Integrate once more to obtain the cubic position defect}{Integrate once more to obtain the cubic position defect}}


Apply the proved velocity-defect bound at each intermediate time \(s\): \(\|D_v(s)\|\le MCs^2/2\). Its continuous scalar majorant is integrable, so the norm of the position-defect integral is bounded by \(\int_0^tMCs^2/2\,ds=MCt^3/6\). Lean first verifies that every intermediate time stays in the original horizon, rewrites the scalar coefficient into the majorant's form, applies norm-integral comparison, and uses \texttt{integral\_\allowbreak{}pow}. The integration variable inside the bound is \(s\), preserving the correct cubic constant. No stationarity or Gaussian moment is assumed implicitly by this deterministic estimate.

\subsection{\texorpdfstring{Recover the paper’s three-\allowbreak{}quarter transformed local constant}{Recover the paper’s three-quarter transformed local constant}}


Put \(B=MCh^2\). The component bounds give \(\|D_v\|\le B/2\) and \(\|D_x\|\le Bh/6\). If \(\gamma h\le1/4\), then \(\|D_v+\gamma D_x\|\le B/2+B/24=13B/24\). The transformed energy at parameter zero is \(\|D_v+\gamma D_x\|^2+\|D_v\|^2\), which is at most \(((13/24)^2+1/4)B^2\le(3B/4)^2\). Lean keeps all quantities nonnegative, squares each norm bound with \texttt{sq\_\allowbreak{}le\_\allowbreak{}sq\ensuremath{{}_{0}}}, unfolds the transformed energy and verifies the rational comparison by \texttt{nlinarith}. This derives the constant for the explicit defect integrals; identifying them with an actual stationary-chain coupling still requires the kinetic process construction and moment estimates.

\section{\texorpdfstring{Identify the implemented exponential-\allowbreak{}step error with its actual integral defect}{Identify the implemented exponential-step error with its actual integral defect}}
\label{proof-module-KineticLangevin-MildError-lean}
{\small\color{gray}Source module: \texttt{Kinetic\allowbreak{}Langevin/\allowbreak{}Mild\allowbreak{}Error.\allowbreak{}lean}.\par}


The local errors are concrete integrals of changes in the force. This module proves that these integrals really are the errors of the implemented exponential step when compared with a mild exact velocity and its integrated position. The variation-of-constants equation is an explicit input here; constructing its Brownian convolution and deriving it from the original integral equation remain distinct steps.

\subsection{\texorpdfstring{Separate the frozen force from its variation}{Separate the frozen force from its variation}}


Write \(D_v(t)=\int_0^t e^{-\gamma(t-s)}[g(x_s)-g(x_0)]ds\). Linearity and \(\int_0^t e^{-\gamma(t-s)}ds=\psi_1(t)\) give

\[D_v(t)=\int_0^t e^{-\gamma(t-s)}g(x_s)ds-\psi_1(t)g(x_0).\]

Lean proves interval integrability using continuity of the path, drift and exponential kernel; it then distributes scalar multiplication and the integral over subtraction. For a mild exact velocity \(V_t=\psi_0(t)V_0-\int_0^t e^{-\gamma(t-s)}g(x_s)ds+B_t\), subtracting it from the frozen velocity plus the same convolution \(B_t\) gives exactly \(D_v(t)\). The proof is algebra after the integral identity.

\subsection{\texorpdfstring{Integrating the frozen velocity gives the position formula}{Integrating the frozen velocity gives the position formula}}


The frozen position is \(x+\psi_1(t)v-\psi_2(t)g\). Its derivative is the frozen velocity \(\psi_0(t)v-\psi_1(t)g\), and its value at zero is \(x\). The fundamental theorem of calculus gives \(\int_0^t V_{\rm frozen}(s)ds=X_{\rm frozen}(t)-x\). Lean invokes the already proved derivative theorem and proves integrability from continuity of the coefficients. The identity holds also at zero friction because the coefficients were defined by integrals.

\subsection{\texorpdfstring{Both coordinates use the same noise convolution}{Both coordinates use the same noise convolution}}


The position innovation is \(\int_0^h B_sds\), the time integral of the same velocity convolution. If \(x_h-x_0=\int_0^hV_sds\), integrate the velocity error identity to obtain the position defect \(D_x(h)=\int_0^hD_v(s)ds\). Thus

\[\operatorname{exponentialStep}(V_0,x_0;B_h,\int_0^hB_sds)-(V_h,x_h)=(D_v(h),D_x(h)).\]

Lean proves the velocity identity at every point in the interval, changes the integrand using that pointwise equality, distributes the integral over sums and differences, and uses the frozen-velocity integral and the kinematic identity. Extensionality then joins the two coordinate equalities into the phase-space equality. No stochastic error estimate is hidden inside the definitions.

\subsection{\texorpdfstring{Bound the implemented step in the transformed norm}{Bound the implemented step in the transformed norm}}


Apply the exact step-error identity to replace \texttt{exponential\allowbreak{}Step \ensuremath{-} exact endpoint} by the explicit pair \((D_v,D_x)\). The local-error module has already integrated the kinematic velocity bound to prove \(\mathcal E(D_v,D_x)\le(3MC h^2/4)^2\) under \(\gamma h\le1/4\). The norm-square identity and nonnegativity allow \texttt{transform\allowbreak{}Norm\_\allowbreak{}le\_\allowbreak{}of\_\allowbreak{}energy\_\allowbreak{}le} to take square roots, giving \(\|\text{step}-\text{exact}\|_T\le3MC h^2/4\). Every input is visible: the actual mild velocity equation, continuous common convolution noise, the kinematic identity, and the uniform velocity norm bound. No stochastic moment or stationary law is asserted here.

\section{\texorpdfstring{Coordinate Gaussian moments and L\ensuremath{{}^{2}} Brownian continuity}{Coordinate Gaussian moments and L² Brownian continuity}}
\label{proof-module-UnadjustedLangevin-GaussianCoordinates-lean}
{\small\color{gray}Source module: \texttt{Unadjusted\allowbreak{}Langevin/\allowbreak{}Gaussian\allowbreak{}Coordinates.\allowbreak{}lean}.\par}


Coordinatewise Gaussian laws suffice for the second moment \(\mathbb E\|X\|^2=dv\). This module carefully obtains Brownian \(L^2\) continuity and increment norms without assuming independence between spatial coordinates.

\subsection{\texorpdfstring{gaussian coordinates mem\allowbreak{}Lp}{gaussian coordinates memLp}}


Each centered Gaussian coordinate has a finite second moment. In finite Euclidean dimension, \(\|X\|^2=\sum_iX_i^2\), so the finite sum is integrable. Coordinatewise measurability also assembles to vector measurability, giving \(X\in L^2\). Coordinate independence is unnecessary for this moment statement.

\subsection{\texorpdfstring{gaussian coordinates second moment}{gaussian coordinates second moment}}


Use \(\|X\|^2=\sum_iX_i^2\) and finite-sum integral linearity. The coordinate law \(X_i\sim\mathcal N(0,v)\) gives \(\mathbb EX_i^2=v\), hence \(\mathbb E\|X\|^2=|\iota|v\). The dimension factor is exactly the cardinality of the finite coordinate set.

\subsection{\texorpdfstring{gaussian coordinates Lp norm sq}{gaussian coordinates Lp norm sq}}


For the given \(L^2\) class, its squared norm is the integral of the squared norm of its representative. Apply the preceding Gaussian coordinate moment formula to obtain \(\|X\|_2^2=|\iota|v\).

\subsection{\texorpdfstring{gaussian coordinates Lp norm}{gaussian coordinates Lp norm}}


Take nonnegative square roots in \(\|X\|_2^2=|\iota|v\). Norm nonnegativity selects the positive root, giving \(\|X\|_2=\sqrt{|\iota|v}\), including zero variance and zero dimension.

\subsection{\texorpdfstring{gaussian coordinates to\allowbreak{}Lp coordinate\allowbreak{}Law}{gaussian coordinates toLp coordinateLaw}}


The representative of \texttt{Mem\allowbreak{}Lp.\allowbreak{}to\allowbreak{}Lp X} agrees almost everywhere with \(X\). Coordinate evaluation preserves that equality, so each coordinate law transfers to the canonical \(L^2\) representative by almost-everywhere invariance of pushforwards.

\subsection{\texorpdfstring{gaussian coordinates to\allowbreak{}Lp norm}{gaussian coordinates toLp norm}}


Apply the exact \(L^2\) Gaussian norm theorem to \texttt{to\allowbreak{}Lp X}, using its transferred coordinate laws. This gives \(\|[X]\|_2=\sqrt{|\iota|v}\) for the class constructed from the original random vector.

\subsection{\texorpdfstring{gaussian coordinates scaled l2 norm}{gaussian coordinates scaled l2 norm}}


The scalar-norm identity gives \(\|\sqrt2[X]\|_2=\sqrt2\sqrt{|\iota|v}\). Nonnegativity permits the square-root product rules, producing \(\sqrt{2|\iota|}\sqrt v\). This is the precise Langevin noise constant used in the local error estimate.

\subsection{\texorpdfstring{brownian\allowbreak{}Coordinates}{brownianCoordinates}}


Assemble the scalar coordinate processes into \(B_t(\omega)=(W_i(t,\omega))_{i\in\iota}\). Finite Euclidean coordinates determine the vector; no new random source or independence is introduced by this definition.

\subsection{\texorpdfstring{brownian\allowbreak{}Coordinates coordinate\allowbreak{}Law}{brownianCoordinates coordinateLaw}}


At a fixed time \(t\ge0\), each scalar pre-Brownian coordinate has law \(\mathcal N(0,t)\). Unfolding the assembled vector shows that its \(i\)-th coordinate is exactly that scalar variable, giving the required coordinate law.

\subsection{\texorpdfstring{brownian\allowbreak{}Coordinates\allowbreak{}Increment}{brownianCoordinatesIncrement}}


Define the vector increment \(\Delta_aB_t=B_{a+t}-B_a\) coordinate by coordinate. This records both the starting time and elapsed nonnegative time, matching the local noise input on an arbitrary Langevin segment.

\subsection{\texorpdfstring{brownian\allowbreak{}Coordinates\allowbreak{}Increment coordinate\allowbreak{}Law}{brownianCoordinatesIncrement coordinateLaw}}


The scalar Brownian increment law depends only on elapsed time: \(W_i(a+t)-W_i(a)\sim\mathcal N(0,t)\). Simplify \((a+t)-a=t\) in the variance and unfold the vector increment to obtain its coordinate law.

\subsection{\texorpdfstring{brownian\allowbreak{}Coordinates\allowbreak{}Increment mem\allowbreak{}Lp}{brownianCoordinatesIncrement memLp}}


Apply \texttt{gaussian\_\allowbreak{}coordinates\_\allowbreak{}mem\allowbreak{}Lp} to the increment's coordinate laws of variance \(t\). The assembled Brownian vector increment is square integrable even without independence between different coordinate processes.

\subsection{\texorpdfstring{brownian\allowbreak{}Coordinates\allowbreak{}Increment scaled l2 norm}{brownianCoordinatesIncrement scaled l2 norm}}


Apply the exact scaled Gaussian-coordinate norm theorem with variance \(t\). It yields \(\|\sqrt2[\Delta_aB_t]\|_2=\sqrt{2|\iota|}\sqrt t\), establishing the noise radius from Brownian laws alone.

\subsection{\texorpdfstring{brownian\allowbreak{}Coordinates\allowbreak{}Lp}{brownianCoordinatesLp}}


Convert the assembled Brownian vector at time \(t\) to its \(L^2\) class using its Gaussian-coordinate square-integrability proof. This places the time-indexed process in the Hilbert space where the deterministic integral equation is solved.

\subsection{\texorpdfstring{brownian\allowbreak{}Coordinates\allowbreak{}Lp coe}{brownianCoordinatesLp coe}}


The \texttt{to\allowbreak{}Lp} representative theorem gives \([B_t](\omega)=B_t(\omega)\) almost everywhere for each fixed time. This is the bridge from class-valued continuity statements to coordinate laws of the original Brownian process.

\subsection{\texorpdfstring{brownian\allowbreak{}Coordinates\allowbreak{}Lp sub coordinate\allowbreak{}Law}{brownianCoordinatesLp sub coordinateLaw}}


The difference of two selected representatives agrees almost everywhere with \(B_t-B_s\). Brownian increment laws, including the sign reversal when \(t<s\), give a centered Gaussian with variance \(|t-s|\); centered Gaussian symmetry handles the reversed order.

\subsection{\texorpdfstring{brownian\allowbreak{}Coordinates\allowbreak{}Lp norm sub}{brownianCoordinatesLp norm sub}}


Apply the Gaussian \(L^2\) norm formula to the difference class, whose coordinates have variance \(|t-s|\). Thus \(\|[B_t]-[B_s]\|_2=\sqrt{|\iota|\,|t-s|}\). This exact increment norm is stronger than the continuity statement derived from it.

\subsection{\texorpdfstring{brownian\allowbreak{}Coordinates\allowbreak{}Lp continuous}{brownianCoordinatesLp continuous}}


As \(s\to t\), \(|t-s|\to0\), hence the exact norm formula gives \(\|[B_t]-[B_s]\|_2\to0\). Lean uses continuity of the square root and the metric characterization of continuity. This needs finite-dimensional Brownian laws, with no sample-path continuity premise.

\subsection{\texorpdfstring{brownian\allowbreak{}Coordinates\allowbreak{}Lp increment scaled norm}{brownianCoordinatesLp increment scaled norm}}


Set the two times in the difference-norm identity to \(a+t\) and \(a\). Their distance is \(t\), and scalar multiplication by \(\sqrt2\) gives \(\sqrt{2|\iota|}\sqrt t\). This is the class-valued version of the earlier increment estimate.

\subsection{\texorpdfstring{brownian\allowbreak{}Coordinates\allowbreak{}Lp zero}{brownianCoordinatesLp zero}}


At time zero every coordinate has the degenerate centered Gaussian law of variance zero. Equivalently its \(L^2\) norm is zero by the moment formula. Norm-zero equivalence proves that the whole class is zero, irrespective of the chosen pointwise representative.

\subsection{\texorpdfstring{brownian\allowbreak{}Coordinates\allowbreak{}Real\allowbreak{}Lp}{brownianCoordinatesRealLp}}


Extend the nonnegative-time \(L^2\) process to real times by \(t\mapsto[B_{\max(t,0)}]\), implemented with \texttt{Real.\allowbreak{}to\allowbreak{}NNReal}. It is constant zero for negative times and agrees with the original class at nonnegative times.

\subsection{\texorpdfstring{brownian\allowbreak{}Coordinates\allowbreak{}Real\allowbreak{}Lp continuous}{brownianCoordinatesRealLp continuous}}


The clipping map \(t\mapsto\max(t,0)\) into nonnegative reals is continuous. Compose it with the established \(L^2\)-continuous Brownian map to obtain continuity of the real-time extension, including at zero.

\subsection{\texorpdfstring{brownian\allowbreak{}Coordinates\allowbreak{}Real\allowbreak{}Lp zero}{brownianCoordinatesRealLp zero}}


At real time zero the clipping map returns nonnegative zero. The previously proved zero-time \(L^2\) identity therefore simplifies the real-time extension to zero.

\subsection{\texorpdfstring{brownian\allowbreak{}Coordinates\allowbreak{}Real\allowbreak{}Lp increment scaled norm}{brownianCoordinatesRealLp increment scaled norm}}


For \(a,t\ge0\), clipping leaves \(a\) and \(a+t\) unchanged. Rewrite the real-time extension at both endpoints as the nonnegative-time Brownian classes, then apply their scaled increment formula. This supplies \(\sqrt{2d}\sqrt t\) in the real interval-integral formulation.

\section{\texorpdfstring{Gaussian velocity moments and the stochastic local step}{Gaussian velocity moments and the stochastic local step}}
\label{proof-module-KineticLangevin-LpOneStep-lean}
{\small\color{gray}Source module: \texttt{Kinetic\allowbreak{}Langevin/\allowbreak{}Lp\allowbreak{}One\allowbreak{}Step.\allowbreak{}lean}.\par}


This module proves the exact \(\sqrt p\) reference velocity norm from a standard Gaussian marginal, then combines the genuine L\ensuremath{{}^{2}} nonlinear contraction with the actual exponential-integrator error. Its hypotheses display the Gaussian reference marginal and concrete exact integral paths; stationarity proving that marginal is a separate theorem.

\subsection{\texorpdfstring{Obtain the exact velocity moment from its Gaussian law}{Obtain the exact velocity moment from its Gaussian law}}


The standard Gaussian vector is the pushforward of independent scalar \(N(0,1)\) coordinates by \texttt{With\allowbreak{}Lp.\allowbreak{}to\allowbreak{}Lp 2}. Compose that pushforward with coordinate evaluation; the finite-product marginal formula reduces its law to \(N(0,1)\), since every other coordinate measure has mass one. If an L\ensuremath{{}^{2}} velocity random variable has this standard Gaussian vector law, composition of laws gives the scalar coordinate laws. The project's proved Gaussian-coordinate moment identity then yields \(\|V\|_{L^2}=\sqrt{\#\iota}\). Thus the reference moment is derived from the explicit distribution rather than supplied as a numerical norm assumption.

\subsection{\texorpdfstring{Prove the concrete stochastic one-\allowbreak{}step inequality}{Prove the concrete stochastic one-step inequality}}


Let \(A\) be the auxiliary exact L\ensuremath{{}^{2}} path starting at the numerical phase point, and \(R\) a synchronously driven exact reference. Their transformed L\ensuremath{{}^{2}} distance is at most its initial value \(e\), because \(e^{-ms/\gamma}\le1\) for nonnegative time. Velocity projection and the exact Gaussian reference moment therefore give \(\|A_s^v\|_{L^2}\le\sqrt p+e\). Apply the actual mild-equation local-error theorem with this velocity bound: the numerical exponential update differs from \(A_h\) by at most \(3Mh^2(\sqrt p+e)/4\). Insert \(A_h\) with the transformed Hilbert triangle inequality and contract its distance to \(R_h\). The result is \(e_{\mathrm{next}}\le e^{-mh/\gamma}e+\frac34Mh^2(\sqrt p+e)\). Because the ambient space is L\ensuremath{{}^{2}}, the Bochner-integral and Hilbert triangle estimates are precisely the stochastic Minkowski estimates required by the paper.



The proof is factored through a shared theorem with an explicit reference velocity radius \(C\ge0\). A uniform bound \(\|V_t^{\rm ref}\|_{L^2}\le C\) supplies \(C+e_k\) in the local defect and the terminal bias \(Mh\sqrt2C/m\). The standard-Gaussian statement is a thin specialization with the radius \(\sqrt p\) proved from its marginal law. This same shared theorem applies to particle temperature \(1/N\), using the separately proved thermal radius \(\sqrt{(d_\theta+Nd_x)/N}\); there is no duplicate exponential-integrator argument.

\section{\texorpdfstring{The first-\allowbreak{}order scalar accumulation and its exact scope}{The first-order scalar accumulation and its exact scope}}
\label{proof-module-KineticLangevin-Recurrence-lean}
{\small\color{gray}Source module: \texttt{Kinetic\allowbreak{}Langevin/\allowbreak{}Recurrence.\allowbreak{}lean}.\par}


The numerical accumulation part of the KLMC proof begins with an explicit one-step coupling inequality. This module derives the stable contraction coefficient and sums its affine recurrence, retaining the coupling inequality and reference velocity norm as visible hypotheses. It therefore checks the numerical implication without claiming the still-open Gibbs stationarity and actual Gaussian-chain Wasserstein endpoint.

\subsection{\texorpdfstring{Prove the one-\allowbreak{}sided exponential quadratic estimate}{Prove the one-sided exponential quadratic estimate}}


For \(x\ge0\), let \(F(x)=1-x+x^2/2-e^{-x}\). Its derivative is \(-1+x+e^{-x}\), nonnegative because \(e^{-x}\ge1-x\). Therefore \(F\) is nondecreasing and \(F(x)\ge F(0)=0\), proving \(e^{-x}\le1-x+x^2/2\). Lean constructs the derivative through the polynomial and exponential derivative rules, proves its sign with \texttt{Real.\allowbreak{}add\_\allowbreak{}one\_\allowbreak{}le\_\allowbreak{}exp}, applies the derivative monotonicity criterion, and evaluates the zero endpoint. The estimate is proved analytically rather than introduced as an unchecked Taylor remainder.

\subsection{\texorpdfstring{Keep most of the exact-\allowbreak{}flow contraction after the local correction}{Keep most of the exact-flow contraction after the local correction}}


Set \(r=m/\gamma\). Suppose \(r,M,h\ge0\), \(2r^2\le M\), and \(4Mh\le r\). The quadratic exponential estimate gives \(e^{-rh}+3Mh^2/4\le1-rh+r^2h^2/2+3Mh^2/4\). Since \(r^2/2\le M/4\), the positive correction is at most \(Mh^2\le rh/4\). Thus the total is at most \(1-3rh/4\). Lean multiplies the curvature and step inequalities by the nonnegative time powers and combines them with \texttt{nlinarith}. The next theorem derives both scalar conditions from \(m>0\), \(m\le M\), \(\gamma^2\ge m+M\) and \(4\gamma Mh\le m\), dividing only by positive \(\gamma\) or \(\gamma^2\).

\subsection{\texorpdfstring{Sum the affine recurrence with the visible local-\allowbreak{}coupling premise}{Sum the affine recurrence with the visible local-coupling premise}}


Assume nonnegative errors satisfy \(e_{n+1}\le e^{-rh}e_n+3Mh^2(C+e_n)/4\), with \(r,h>0\), \(C\ge0\), and the scalar restrictions above. Put \(q=1-3rh/4\). The coefficient bound and positivity of the exponential prove \(0\le q<1\), and multiplying that bound by \(e_n\ge0\) yields \(e_{n+1}\le qe_n+3Mh^2C/4\). The shared affine-recurrence theorem gives \(e_n\le q^ne_0+(3Mh^2C/4)/(1-q)\). Positive \(r,h\) permit denominator cancellation, leaving \(e_n\le q^ne_0+MhC/r\). Lean uses \texttt{field\_\allowbreak{}simp} and \texttt{ring} only for that final quotient. The one-step coupling inequality is an explicit assumption, so this result must not be presented as the full actual-law KLMC theorem.

\subsection{\texorpdfstring{Restore the paper’s physical parameters and position coefficient}{Restore the paper’s physical parameters and position coefficient}}


Substitute \(r=m/\gamma\) into the scalar recurrence theorem. The derived scalar conditions justify the call, and rational simplification gives \(e_n\le(1-3mh/(4\gamma))^ne_0+\gamma MhC/m\). Suppose in addition that \(d_n\le\sqrt2e_n/\gamma\) and \(e_0\le\gamma d_0\). The nonnegative projection factor and nonnegative power of the proved stable contraction coefficient allow multiplication of these inequalities. Canceling \(\gamma\) yields

\[d_n\le\sqrt2(1-3mh/(4\gamma))^nd_0+Mh\sqrt2C/m.\]

Lean explicitly proves the projection factor and contraction coefficient nonnegative before transporting inequalities, then uses \texttt{field\_\allowbreak{}simp} for the scalar cancellations. Setting \(C=\sqrt p\) matches the paper’s numerical expression. The local coupling, projection and initialization bounds remain explicit hypotheses on sequences; neither declaration identifies \(d_n\) with actual Wasserstein error or constructs a stationary reference.

\section{\texorpdfstring{Iterating concrete stochastic segments to a Wasserstein bound}{Iterating concrete stochastic segments to a Wasserstein bound}}
\label{proof-module-KineticLangevin-LpConvergence-lean}
{\small\color{gray}Source module: \texttt{Kinetic\allowbreak{}Langevin/\allowbreak{}Lp\allowbreak{}Convergence.\allowbreak{}lean}.\par}


The numerical phase knots, auxiliary exact segments, reference segments, common raw inputs and convolution noises are explicit. The local coupling inequality is proved internally from these paths. The remaining law inputs state the actual reference Gaussian velocity and stationary position marginals. They do not assume a convergence conclusion.

\subsection{\texorpdfstring{Accumulate actual update errors and project to laws}{Accumulate actual update errors and project to laws}}


For each step, the auxiliary path begins at the numerical knot, while the reference segment joins consecutive reference knots. The numerical knot update is exactly \texttt{exponential\allowbreak{}Step} with innovation \((B(h),\int_0^hB(s)\,ds)\). Rewriting these endpoint identities in the proved L\ensuremath{{}^{2}} one-step theorem produces the recurrence \(e_{k+1}\le e^{-mh/\gamma}e_k+3Mh^2(\sqrt p+e_k)/4\). The source step restriction \(4\gamma Mh\le m\), together with \(m\le M\), implies \(\gamma h\le1/4\), as needed for the local constant. The scalar recurrence theorem then gives \(e_n\le(1-3mh/(4\gamma))^ne_0+\gamma Mh\sqrt p/m\).



The actual numerical and reference position random variables provide a finite-cost Wasserstein coupling. The reference position law is the stated target \(\pi\), and position projection gives \(W_2\le\sqrt2e_n/\gamma\). Equal initial velocities identify \(e_0=\gamma\|X_0-Y_0\|_{L^2}\). Lean cancels friction and combines \(\sqrt2\sqrt p=\sqrt{2p}\), proving

\[W_2(\mathcal L(X_n),\pi)\le\sqrt2\left(1-\frac{3mh}{4\gamma}\right)^n\|X_0-Y_0\|_{L^2}+\frac{Mh\sqrt{2p}}m.\]

The initial term is the cost of the supplied coupling. Obtaining the infimum over all initial couplings, and discharging reference stationarity from the actual Brownian kinetic process, remain separate probabilistic requirements.



The proof is factored through a shared theorem with an explicit reference velocity radius \(C\ge0\). A uniform bound \(\|V_t^{\rm ref}\|_{L^2}\le C\) supplies \(C+e_k\) in the local defect and the terminal bias \(Mh\sqrt2C/m\). The standard-Gaussian statement is a thin specialization with the radius \(\sqrt p\) proved from its marginal law. This same shared theorem applies to particle temperature \(1/N\), using the separately proved thermal radius \(\sqrt{(d_\theta+Nd_x)/N}\); there is no duplicate exponential-integrator argument.

\section{\texorpdfstring{Detailed balance through accepted and rejected flows}{Detailed balance through accepted and rejected flows}}
\label{proof-module-UnadjustedLangevin-Metropolis-lean}
{\small\color{gray}Source module: \texttt{Unadjusted\allowbreak{}Langevin/\allowbreak{}Metropolis.\allowbreak{}lean}.\par}


The Metropolis correction preserves the initial marginal whenever the accepted proposal flow is invariant under swapping its two states. A positive proposal-pair density makes this symmetry explicit.

\subsection{\texorpdfstring{accepted\allowbreak{}Flow}{acceptedFlow}}


Weight the current/\allowbreak{}proposal joint measure \(\nu(dx,dy)\) by acceptance probability \(a(x,y)\). Thus \(A(dx,dy)=a(x,y)\nu(dx,dy)\), with \texttt{of\allowbreak{}Real} embedding its real nonnegative weight. The bounds \(0\le a\le1\) are supplied in the conservation theorems below.

\subsection{\texorpdfstring{rejected\allowbreak{}Flow}{rejectedFlow}}


The rejected pair measure is \(R(dx,dy)=(1-a(x,y))\nu(dx,dy)\). It retains the proposed state as an auxiliary coordinate; the eventual output maps it to the original state \(x\), preserving all rejection mass.

\subsection{\texorpdfstring{metropolis\allowbreak{}Output}{metropolisOutput}}


Define the output as \(\operatorname{snd}_*A+\operatorname{fst}_*R\). Accepted moves output \(y\), while rejected moves output \(x\). This is the complete transition measure and requires no renormalization of accepted moves.

\subsection{\texorpdfstring{accepted\allowbreak{}Flow add rejected\allowbreak{}Flow}{acceptedFlow add rejectedFlow}}


Since \(0\le a\le1\), the two nonnegative weights add to one after embedding in extended scalars. Additivity of with-density measures then gives \(A+R=\nu\). The proof checks both nonnegativity guards before combining \texttt{of\allowbreak{}Real} values.

\subsection{\texorpdfstring{marginals eq of swap invariant}{marginals eq of swap invariant}}


The identity \(\operatorname{fst}\circ\operatorname{swap}=\operatorname{snd}\) and pushforward composition give \(\operatorname{snd}_*\nu=\operatorname{fst}_*(\operatorname{swap}_*\nu)\). If the pair measure is swap invariant, this is \(\operatorname{fst}_*\nu\), proving equality of its marginals.

\subsection{\texorpdfstring{metropolis\allowbreak{}Output eq initial}{metropolisOutput eq initial}}


Detailed balance makes the two marginals of the accepted flow equal. Replace its second marginal by its first, then combine the accepted and rejected first marginals using \(A+R=\nu\). The output is therefore \(\operatorname{fst}_*\nu\), exactly the original state law.

\subsection{\texorpdfstring{map weighted comp}{map weighted comp}}


For a measurable \(f\) and weight \(w\), prove \(f_*((w\circ f)\mu)=w(f_*\mu)\). On each measurable set, the pushforward definition and nonnegative change-of-variables formula turn both sides into \(\int\mathbf1_A(f(x))w(f(x))d\mu(x)\). This is the weight-transport identity used by the actual MALA proposal map.

\subsection{\texorpdfstring{map weighted comp equiv}{map weighted comp equiv}}


Specialize the preceding pushforward/\allowbreak{}weight identity to a measurable equivalence. Its bundled measurability supplies the required hypothesis. The application below uses the coordinate swap equivalence to prove symmetry of an accepted density.

\subsection{\texorpdfstring{metropolis\allowbreak{}Acceptance}{metropolisAcceptance}}


For positive joint density \(w(x,y)\), define \(a(x,y)=\min(1,w(y,x)/w(x,y))\). The ratio is reverse flow divided by forward flow; positivity ensures both its meaning and nonnegative acceptance.

\subsection{\texorpdfstring{weight mul metropolis\allowbreak{}Acceptance}{weight mul metropolisAcceptance}}


Distribute the positive factor \(w(x,y)\) through the minimum. Cancellation of the nonzero denominator gives \(w(x,y)a(x,y)=\min(w(x,y),w(y,x))\). This algebra is the reason the accepted joint density is symmetric.

\subsection{\texorpdfstring{accepted\allowbreak{}Flow density}{acceptedFlow density}}


Two successive changes of density multiply their weights. Combine the original positive density \(w\) and acceptance \(a\), then apply the weight/\allowbreak{}minimum identity. The accepted flow becomes the base measure weighted by \(\min(w,w\circ\operatorname{swap})\), a visibly symmetric density.

\subsection{\texorpdfstring{metropolis\allowbreak{}Output density eq initial}{metropolisOutput density eq initial}}


The base measure is swap invariant, and the accepted minimum density is unchanged by swapping its arguments. The weight-transport theorem therefore proves accepted-flow detailed balance. The general accepted/\allowbreak{}rejected conservation theorem then identifies the output with the initial marginal.

\section{\texorpdfstring{Exact Gibbs invariance of an explicit corrected Verlet proposal}{Exact Gibbs invariance of an explicit corrected Verlet proposal}}
\label{proof-module-KineticLangevin-HamiltonianMetropolis-lean}
{\small\color{gray}Source module: \texttt{Kinetic\allowbreak{}Langevin/\allowbreak{}Hamiltonian\allowbreak{}Metropolis.\allowbreak{}lean}.\par}


A volume-preserving involution has a symmetric old/\allowbreak{}new-state graph measure. Weighting that graph by a positive target density and applying the ordinary Metropolis ratio makes its accepted flow symmetric. The rejected mass stays at the old state. This module proves exact invariance for the explicit reversed-Verlet proposal. Passing from these corrected proposals to the unadjusted Hamiltonian flow and then to the full kinetic diffusion still requires quantitative limits.

\subsection{\texorpdfstring{Construct the proposal graph and derive its swap symmetry}{Construct the proposal graph and derive its swap symmetry}}


For \(J^2=\mathrm{id}\), the graph map is \(G(z)=(z,Jz)\). Swapping the old and new state gives \(\operatorname{swap}\circ G=G\circ J\). If \(J\) preserves a base measure \(\lambda\), pushforward composition yields \(\operatorname{swap}_\#G_\#\lambda=G_\#J_\#\lambda=G_\#\lambda\). Lean proves the pointwise graph identity using the involution equation, verifies graph measurability, and rewrites the two pushforward compositions. Thus symmetry of the graph measure is derived from the actual proposal rather than introduced as a detailed-balance field.

\subsection{\texorpdfstring{Apply the accepted-\allowbreak{}flow argument to an arbitrary positive density}{Apply the accepted-flow argument to an arbitrary positive density}}


Let the target be \(w\lambda\), with measurable \(w>0\). Its proposal joint law is \(G_\#(w\lambda)\). The pushforward-with-density theorem expresses this as the symmetric base graph weighted by \(w\) at its first coordinate. The Metropolis ratio is \(\min\{1,w(y)/w(x)\}\), so accepted joint density is \(\min\{w(x),w(y)\}\), symmetric under swapping. The checked shared Metropolis theorem proves that accepted new-state mass plus rejected old-state mass has the initial marginal. Lean computes that marginal by projecting the graph back to its first coordinate, giving the exact identity with \(w\lambda\). No invariant measure or acceptance-balance conclusion is assumed.

\subsection{\texorpdfstring{Bound rejection by an actual energy increase}{Bound rejection by an actual energy increase}}


For energy increment \(d\), rejection is \(1-\min\{1,e^{-d}\}\). If \(d\le0\), the proposed energy decreases and rejection is zero. If \(d\ge0\), the exponential tangent inequality \(e^{-d}\ge1-d\) gives rejection at most \(d\). Thus rejection is at most \(\max\{d,0\}\le|d|\). Lean splits the sign cases, rewrites the minimum and maximum with the corresponding exponential bounds, and applies the tangent inequality. This is a scalar probability bound; converting it to an L\ensuremath{{}^{2}} trajectory error also requires the size of the rejected jump and must not silently assume that jump is of order \(h\).

\subsection{\texorpdfstring{Instantiate the exact corrected Verlet Gibbs theorem}{Instantiate the exact corrected Verlet Gibbs theorem}}


Take \(\lambda\) to be product Haar phase measure, \(J_h=R\,V_h\) to be the constructed reversed Verlet proposal, and \(w(z)=e^{-H(z)}\). The Hamiltonian module proves that \(J_h\) is both an involution and measure preserving. Measurability of the potential and gradient proves measurability of the weight and proposal, and the exponential weight is strictly positive. The generic theorem therefore proves exact preservation of the actual Gibbs weighted measure by this Metropolis step. Its accepted state is velocity-reversed Verlet and its rejected state is the old state. Applying a further velocity reversal changes those to ordinary Verlet on acceptance and a velocity flip on rejection; any limiting kinetic argument must account for this distinction explicitly.

\section{\texorpdfstring{Explicit Verlet energy and rejection estimates}{Explicit Verlet energy and rejection estimates}}
\label{proof-module-KineticLangevin-HamiltonianError-lean}
{\small\color{gray}Source module: \texttt{Kinetic\allowbreak{}Langevin/\allowbreak{}Hamiltonian\allowbreak{}Error.\allowbreak{}lean}.\par}


The exact Hamiltonian defect proved earlier becomes quantitative under a globally Lipschitz gradient. The proof uses only first derivatives: two-sided Taylor control bounds the trapezoidal force work, while a norm factorization bounds the force-square correction. These estimates concern the actual algebraic Verlet map and can control the rejection mass of the invariant corrected approximation.

\subsection{\texorpdfstring{Obtain the two-\allowbreak{}sided Taylor remainder}{Obtain the two-sided Taylor remainder}}


The shared smooth upper bound gives \(f(y)-f(x)-\langle\nabla f(x),y-x\rangle\le M\|y-x\|^2/2\). Apply the same theorem to \(-f\). Its gradient is \(-\nabla f\), proved by rewriting the Fréchet derivative and the linear Riesz inverse, and its Lipschitz constant remains \(M\). The second application gives the reverse inequality, so the absolute remainder has the same bound. Lean explicitly rewrites the negative gradient and combines the upper and lower scalar inequalities with \texttt{abs\_\allowbreak{}le}.

\subsection{\texorpdfstring{Average Taylor errors from both endpoints}{Average Taylor errors from both endpoints}}


Write \(d=y-x\), \(r_x=f(y)-f(x)-\langle\nabla f(x),d\rangle\), and \(r_y=f(x)-f(y)+\langle\nabla f(y),d\rangle\). The trapezoidal force-work error is exactly \((r_x-r_y)/2\). Both remainders have absolute value at most \(M\|d\|^2/2\). The triangle inequality therefore gives \[\left|f(y)-f(x)-\tfrac12\langle\nabla f(x)+\nabla f(y),d\rangle\right|\le\tfrac M2\|d\|^2.\] Lean proves the algebraic average identity after expanding the inner product, reverses the displacement in the second Taylor bound, and applies the absolute-value triangle inequality.

\subsection{\texorpdfstring{Control the squared-\allowbreak{}force correction}{Control the squared-force correction}}


Factor \(\|g(y)\|^2-\|g(x)\|^2=(\|g(y)\|-\|g(x)\|)(\|g(y)\|+\|g(x)\|)\). The reverse triangle inequality and Lipschitz continuity bound the first factor by \(M\|y-x\|\). The ordinary triangle inequality bounds the second by \(2\|g(x)\|+M\|y-x\|\). Multiplying these nonnegative estimates proves the displayed force-square bound. Lean derives the endpoint-force growth estimate before applying \texttt{mul\_\allowbreak{}le\_\allowbreak{}mul}, so every sign condition is checked.

\subsection{\texorpdfstring{Bound the complete Verlet energy defect}{Bound the complete Verlet energy defect}}


Set \(w=v-h\nabla f(x)/2\), \(y=x+hw\). The exact identity from the Hamiltonian module writes the defect as a trapezoidal work error plus \(h^2(\|\nabla f(y)\|^2-\|\nabla f(x)\|^2)/8\). Since \(\|y-x\|=|h|\|w\|\), the preceding two lemmas give \[|H(V_hz)-H(z)|\le\frac M2h^2\|w\|^2+\frac{h^2}{8}M|h|\|w\|\bigl(2\|\nabla f(x)\|+M|h|\|w\|\bigr).\] Lean substitutes the actual intermediate stage, changes the inner product to the displacement form, applies the triangle inequality to the two exact defect terms, and simplifies \(|h|^2=h^2\). The leading term is quadratic in the step under only gradient Lipschitz regularity.

\section{\texorpdfstring{The actual accepted-\allowbreak{}Verlet and rejected-\allowbreak{}flip invariant measure}{The actual accepted-Verlet and rejected-flip invariant measure}}
\label{proof-module-KineticLangevin-HamiltonianWeak-lean}
{\small\color{gray}Source module: \texttt{Kinetic\allowbreak{}Langevin/\allowbreak{}Hamiltonian\allowbreak{}Weak.\allowbreak{}lean}.\par}


The reversed-Verlet involution supplied exact Metropolis invariance. Applying one final velocity reversal gives ordinary Verlet on acceptance and a velocity flip on rejection. This module computes the resulting two weighted pushforward measures explicitly and proves that their sum equals the Hamiltonian Gibbs measure. This identity is the concrete starting point for bounded-test estimates and a numerical limit.

\subsection{\texorpdfstring{Compute the two Metropolis graph marginals}{Compute the two Metropolis graph marginals}}


For a deterministic proposal \(J\), the old/\allowbreak{}new graph is \(G(z)=(z,Jz)\). Weighting its pushforward by a measurable acceptance function \(a(z,Jz)\) commutes with mapping the correspondingly weighted initial measure through \(G\). Projecting onto the second coordinate yields \(J_\#(a\pi)\). Projecting rejected mass onto the first coordinate yields \((1-a)\pi\). Lean applies the previously proved pushforward-with-density formula twice and composes each graph with its coordinate projection. This proves an actual equality of measures without an integrability premise.

\subsection{\texorpdfstring{Construct the Hamiltonian Gibbs measure and prove velocity symmetry}{Construct the Hamiltonian Gibbs measure and prove velocity symmetry}}


The measure is \(e^{-H(v,x)}(\mu\otimes\mu)(d(v,x))\), where \(\mu\) is a translation-invariant symmetric Haar factor and \(H(v,x)=f(x)+\|v\|^2/2\). It is initially unnormalized; finiteness and normalization are independent questions. Velocity reversal preserves product Haar measure and preserves \(H\) exactly. The weighted pushforward identity therefore proves that it preserves this Gibbs measure. Lean verifies measurability of the exponential density and rewrites the two exact preservation identities.

\subsection{\texorpdfstring{Use the actual Verlet energy change as acceptance}{Use the actual Verlet energy change as acceptance}}


The acceptance probability is \(a_h(z)=\min\{1,\exp[-(H(V_hz)-H(z))]\}\). Its measurability follows by expanding the explicit kick-drift-kick stages and composing the measurable potential and gradient. Positivity of the exponential gives \(a_h\ge0\), and the minimum gives \(a_h\le1\). These scalar bounds later justify finite weighted measures and nonnegative rejection estimates.

\subsection{\texorpdfstring{Derive exact corrected Verlet invariance}{Derive exact corrected Verlet invariance}}


The proved Metropolis result uses the involution \(J_h=R\circ V_h\). Since \(H\circ R=H\), its ratio is exactly the acceptance \(a_h\) just defined. Compute its accepted and rejected graph marginals, then push the resulting equality through \(R\). The accepted state becomes \(R(J_hz)=V_hz\), whereas the rejected state becomes \(Rz\). Gibbs velocity symmetry fixes the right-hand side, giving \[(V_h)_\#(a_h\pi)+R_\#((1-a_h)\pi)=\pi.\] Lean performs precisely these graph, density, and pushforward rewrites and uses the checked identity \(R^2=\mathrm{id}\). This theorem does not identify rejection with a small spatial jump; the rejected velocity flip is retained explicitly.

\section{\texorpdfstring{Bounded-\allowbreak{}test estimates from the exact corrected measure}{Bounded-test estimates from the exact corrected measure}}
\label{proof-module-KineticLangevin-HamiltonianTests-lean}
{\small\color{gray}Source module: \texttt{Kinetic\allowbreak{}Langevin/\allowbreak{}Hamiltonian\allowbreak{}Tests.\allowbreak{}lean}.\par}


The rejected move can reverse a velocity of order one. A weak approximation therefore uses bounded complex-valued tests rather than assuming the rejected jump is proportional to the step. The exact corrected measure identity implies that one uncorrected proposal has test bias at most twice the test bound times its mean rejection probability.

\subsection{\texorpdfstring{Keep all acceptance and rejection mass}{Keep all acceptance and rejection mass}}


For measurable \(0\le a\le1\), the two weighted measures \(a\pi\) and \((1-a)\pi\) sum exactly to \(\pi\). Lean combines the two densities, uses nonnegativity to rewrite the sum of \texttt{ENNReal.\allowbreak{}of\allowbreak{}Real} values as the value of their sum, and reduces the resulting density to one. In particular both weighted measures are bounded above by the original finite measure.

\subsection{\texorpdfstring{Turn the explicit invariant measure into an integral identity}{Turn the explicit invariant measure into an integral identity}}


Suppose the already constructed corrected measure satisfies \(V_\#(a\pi)+R_\#((1-a)\pi)=\pi\). This identity is proved for the actual Verlet and reversal maps in the preceding module. For bounded measurable complex \(F\), integration gives \[\int F\,d\pi=\int a(z)F(Vz)\,d\pi+\int(1-a(z))F(Rz)\,d\pi.\] Lean first proves finiteness of each weighted measure from the exact mass split. The bound on \(F\) supplies integrability for both pushforward measures. It then applies the Bochner integral formula for a pushforward and the formula for a density-weighted integral; the acceptance and rejection factors are finite nonnegative real numbers.

\subsection{\texorpdfstring{Bound the entire rejected move through the test bound}{Bound the entire rejected move through the test bound}}


Subtract the invariant-test identity from the ordinary proposal expectation. The exact difference is \[\int F(Vz)\,d\pi-\int F\,d\pi=\int(1-a(z))\,[F(Vz)-F(Rz)]\,d\pi.\] If \(\|F\|\le C\), the norm of the bracket is at most \(2C\), irrespective of the size of the rejected physical jump. Since \(1-a\ge0\), the norm of the integral is at most \(2C\int(1-a)\,d\pi\). Lean explicitly checks bounded integrability of the proposal test, accepted test, rejected test, and rejection factor, rewrites the subtraction under the integral, applies the norm-of-integral inequality, and integrates the pointwise triangle bound. Thus the velocity flip is fully accounted for.

\section{\texorpdfstring{Linear accumulation of weak error and exact limiting laws}{Linear accumulation of weak error and exact limiting laws}}
\label{proof-module-KineticLangevin-HamiltonianIteration-lean}
{\small\color{gray}Source module: \texttt{Kinetic\allowbreak{}Langevin/\allowbreak{}Hamiltonian\allowbreak{}Iteration.\allowbreak{}lean}.\par}


Corrected invariance can be used entirely through bounded tests. Composing a test with more deterministic proposal steps preserves its bound. Each telescoping increment is therefore controlled under the same fixed measure, giving a linear accumulation of rejection mass. Pointwise convergence and vanishing accumulated error then identify the limiting law exactly by dual characteristic functions.

\subsection{\texorpdfstring{Telescope tests composed with deterministic iterates}{Telescope tests composed with deterministic iterates}}


Let \(A_k=\int F(V^kz)\,d\pi\). Apply the one-step weak estimate to \(G=F\circ V^k\). Measurability of iterates proves measurability of \(G\), and the same bound \(\|G\|\le C\) holds. Thus \(\|A_{k+1}-A_k\|\le2C\int(1-a)\,d\pi\). The triangle inequality and induction give \[\|A_n-A_0\|\le2Cn\int(1-a)\,d\pi.\] Lean explicitly rewrites the successor iterate as \(V^k\circ V\), applies the checked one-step result, and normalizes the successor natural-number cast. Every increment is integrated against the fixed measure \(\pi\); no law of an unadjusted intermediate chain is assumed invariant.

\subsection{\texorpdfstring{Identify the actual limiting pushforward through characteristic tests}{Identify the actual limiting pushforward through characteristic tests}}


Suppose measurable maps \(Y_n\) converge pointwise to a measurable map \(Y\), and the bounded-test error is eventually at most a real sequence \(\delta_n\to0\). For every continuous real linear functional \(L\), use \(F(z)=\exp(iLz)\), whose norm is one. Dominated convergence under the finite measure gives \(\int F(Y_nz)\,d\pi\to\int F(Yz)\,d\pi\). The error bound gives convergence of the same sequence to \(\int F\,d\pi\). Uniqueness of limits equates the dual characteristic functions of \(Y_\#\pi\) and \(\pi\); the checked uniqueness theorem for finite measures gives \(Y_\#\pi=\pi\). Lean uses dual characteristic functions, so the result applies to Banach phase space with its usual product norm and does not silently require that norm to arise from an inner product.

\section{\texorpdfstring{Quadratic rejection mass under a concrete phase moment}{Quadratic rejection mass under a concrete phase moment}}
\label{proof-module-KineticLangevin-HamiltonianRejection-lean}
{\small\color{gray}Source module: \texttt{Kinetic\allowbreak{}Langevin/\allowbreak{}Hamiltonian\allowbreak{}Rejection.\allowbreak{}lean}.\par}


The exact Verlet energy defect controls its actual Metropolis rejection probability. For steps between zero and one, the state-dependent bound simplifies to a single quadratic phase envelope. Integrating this bound requires precisely the velocity and force second moments displayed in the theorem, and supplies the vanishing error needed for a fixed-horizon weak limit.

\subsection{\texorpdfstring{Apply the scalar rejection inequality to the actual energy defect}{Apply the scalar rejection inequality to the actual energy defect}}


The acceptance is defined by the actual defect \(d=H(V_hz)-H(z)\). The already proved scalar inequality \(1-\min\{1,e^{-d}\}\le |d|\) therefore applies by substitution, with no acceptance approximation.

\subsection{\texorpdfstring{Reduce all intermediate stages to a quadratic phase envelope}{Reduce all intermediate stages to a quadratic phase envelope}}


Put \(a=\|v\|\), \(b=\|\nabla f(x)\|\), \(w=v-h\nabla f(x)/2\), and \(q=a^2+b^2\). For \(0\le h\le1\), the triangle inequality gives \(\|w\|\le a+b\), hence \(\|w\|^2\le2q\). Also \(2\|w\|b\le\|w\|^2+b^2\), so \(\|w\|(2b+Mh\|w\|)\le(2M+3)q\). Substitution in the earlier exact defect estimate yields \[|H(V_hz)-H(z)|\le h^2\frac{M(2M+11)}8\bigl(\|v\|^2+\|\nabla f(x)\|^2\bigr).\] Lean retains the intermediate nonnegative scalar bounds, checks every multiplication by \(h\) and \(M\), and combines them by polynomial arithmetic. The following theorem composes this estimate with the scalar rejection inequality.

\subsection{\texorpdfstring{Integrate the proved rejection envelope}{Integrate the proved rejection envelope}}


For any finite phase measure with integrable \(q(z)=\|v\|^2+\|\nabla f(x)\|^2\), the rejection factor is measurable and lies between zero and one, so it is integrable. The pointwise bound integrates to \[\int(1-a_h)\,d\pi\le h^2\frac{M(2M+11)}8\int q\,d\pi.\] Lean proves bounded integrability of rejection directly from the acceptance bounds and applies monotonicity of the ordinary integral against the explicitly assumed finite quadratic moment. No one-step trajectory or stationary-flow conclusion is supplied as a premise.

\section{\texorpdfstring{Gibbs normalization and second moments}{Gibbs normalization and second moments}}
\label{proof-module-UnadjustedLangevin-Gibbs-lean}
{\small\color{gray}Source module: \texttt{Unadjusted\allowbreak{}Langevin/\allowbreak{}Gibbs.\allowbreak{}lean}.\par}


On a finite-dimensional real Hilbert space, differentiability and \(m\)-strong convexity with \(m>0\) imply that \(\pi(dx)=Z^{-1}e^{-U(x)}dx\) is a probability measure with finite second moment. This establishes existence of the target measure used throughout the coupling proof.

\subsection{\texorpdfstring{gibbs\allowbreak{}Weight}{gibbsWeight}}


The unnormalized density is \(w_U(x)=e^{-U(x)}\), embedded in \([0,\infty]\) with \texttt{ENNReal.\allowbreak{}of\allowbreak{}Real}. This choice lets the partition integral retain the value infinity; positivity of the exponential ensures that embedding does not discard any part of the density.

\subsection{\texorpdfstring{gibbs\allowbreak{}Partition}{gibbsPartition}}


The partition function is \(Z_U=\int e^{-U(x)}\,dx\), formed as a nonnegative extended integral. At this definition no finiteness assumption is made. Later normalization uses separate proofs of \(0<Z_U\) and \(Z_U<\infty\), which are needed to justify cancellation with its inverse.

\subsection{\texorpdfstring{gibbs\allowbreak{}Measure}{gibbsMeasure}}


Define \(\pi_U(dx)=Z_U^{-1}e^{-U(x)}dx\) using \texttt{volume.\allowbreak{}with\allowbreak{}Density}. The inverse is the extended nonnegative inverse. The definition therefore exists for every potential, while the later probability theorem proves that it has total mass one under the stated strong convexity assumptions.

\subsection{\texorpdfstring{integrable exp neg mul norm sq}{integrable exp neg mul norm sq}}


For \(a>0\), \(e^{-a\|x\|^2}\) is Lebesgue integrable in finite dimension. The Lean proof specializes \texttt{Gaussian\allowbreak{}Fourier.\allowbreak{}integrable\_\allowbreak{}cexp\_\allowbreak{}neg\_\allowbreak{}mul\_\allowbreak{}sq\_\allowbreak{}norm\_\allowbreak{}add} to real positive precision \(a\) and zero linear term, then takes the real part. The complex Gaussian integrand is the real embedding of the desired exponential, so the final simplification identifies the functions exactly. This also includes dimension zero.

\subsection{\texorpdfstring{strong\allowbreak{}Convex quadratic lower bound}{strongConvex quadratic lower bound}}


Apply the strong supporting-plane inequality at zero: \(U(x)\ge U(0)+\langle\nabla U(0),x\rangle+\frac m2\|x\|^2\). Cauchy–Schwarz and \(ab\le a^2/m+(m/4)b^2\), valid for \(m>0\), absorb the linear term. The result is \(U(x)\ge U(0)-\|\nabla U(0)\|^2/m+(m/4)\|x\|^2\). Lean isolates this scalar square inequality before its final arithmetic rearrangement.

\subsection{\texorpdfstring{integrable exp neg of strong\allowbreak{}Convex}{integrable exp neg of strongConvex}}


Exponentiating the quadratic lower bound gives \(e^{-U(x)}\le e^{-U(0)+\|\nabla U(0)\|^2/m}e^{-(m/4)\|x\|^2}\). The right side is an integrable Gaussian multiple. Differentiability supplies continuity and measurability of the left side, so \texttt{Integrable.\allowbreak{}mono'} transfers integrability from that explicit dominating function.

\subsection{\texorpdfstring{integrable norm sq mul exp neg mul norm sq}{integrable norm sq mul exp neg mul norm sq}}


For \(a>0\), split the Gaussian decay into two equal factors. The scalar function \(s\mapsto s e^{-(a/2)s}\) is bounded on \([0,\infty)\), so \(\|x\|^2e^{-a\|x\|^2}\le C_a e^{-(a/2)\|x\|^2}\). The remaining Gaussian is integrable; the proof checks measurability and nonnegativity before applying domination.

\subsection{\texorpdfstring{integrable norm sq mul exp neg of strong\allowbreak{}Convex}{integrable norm sq mul exp neg of strongConvex}}


Multiply the previous Gibbs Gaussian domination by \(\|x\|^2\). Its right side has the form \(C\|x\|^2e^{-(m/4)\|x\|^2}\), whose integrability was proved separately. The pointwise inequality uses \(\|x\|^2\ge0\), so multiplying it preserves order; this is the second-moment input for the normalized measure.

\subsection{\texorpdfstring{gibbs\allowbreak{}Partition pos}{gibbsPartition pos}}


The continuous function \(e^{-U}\) is strictly positive everywhere. Its extended integral cannot vanish: a zero integral of a nonnegative measurable function would force it to vanish almost everywhere, contradicting positivity on a nonzero-volume open set. This proves \(0<Z_U\) without requiring an upper bound on \(Z_U\).

\subsection{\texorpdfstring{gibbs\allowbreak{}Partition ne top}{gibbsPartition ne top}}


Strong convexity with \(m>0\) has already given ordinary integrability of \(e^{-U}\). For a nonnegative integrable function its extended integral is finite. Unfolding \texttt{gibbs\allowbreak{}Partition} identifies that integral exactly with \(Z_U\), giving the guard needed for inverse cancellation.

\subsection{\texorpdfstring{gibbs\allowbreak{}Partition eq of\allowbreak{}Real integral}{gibbsPartition eq ofReal integral}}


Nonnegativity and integrability identify the extended integral with the embedded real integral: \(Z_U=\operatorname{ofReal}(\int e^{-U}\,dx)\). The proof uses the genuine integral/\allowbreak{}lintegral comparison theorem, so it cannot identify an infinite partition with the default real value of an undefined integral.

\subsection{\texorpdfstring{integral exp neg pos}{integral exp neg pos}}


Combine \(0<Z_U\) with \(Z_U=\operatorname{ofReal}(\int e^{-U}\,dx)\). Positivity of \texttt{of\allowbreak{}Real} is equivalent to positivity of its real argument, proving that the ordinary normalizing constant is strictly positive and allowing real division by it later.

\subsection{\texorpdfstring{is\allowbreak{}Probability\allowbreak{}Measure gibbs\allowbreak{}Measure}{isProbabilityMeasure gibbsMeasure}}


The mass of a density measure is its extended density integral. Pull the constant \(Z_U^{-1}\) out to get \(\pi_U(H)=Z_U^{-1}Z_U\). The proofs that \(Z_U\ne0\) and \(Z_U\ne\infty\) justify this product being one. This theorem constructs the \texttt{Is\allowbreak{}Probability\allowbreak{}Measure} instance rather than assuming normalization.

\subsection{\texorpdfstring{integrable norm sq gibbs\allowbreak{}Measure}{integrable norm sq gibbsMeasure}}


Changing measure converts integrability under \(\pi_U\) to integrability of \(Z_U^{-1}e^{-U(x)}\|x\|^2\) under volume. The weighted second moment is already integrable, and the real normalizing factor is a finite scalar. The with-density theorem supplies the almost-everywhere measurability and finiteness guards for this conversion.

\subsection{\texorpdfstring{mem\allowbreak{}Lp id gibbs\allowbreak{}Measure}{memLp id gibbsMeasure}}


For the identity random variable, membership in \(L^2(\pi_U;H)\) is equivalent to integrability of \(\|x\|^2\). The identity is measurable, and the preceding weighted-moment theorem supplies the integral. Lean uses the \(p=2\) characterization of \texttt{Mem\allowbreak{}Lp}, retaining the actual function before passing to any quotient class.

\subsection{\texorpdfstring{mem\allowbreak{}Lp sub const gibbs\allowbreak{}Measure}{memLp sub const gibbsMeasure}}


Since \(\pi_U\) is a probability, the constant map \(x\mapsto z\) belongs to \(L^2\). Subtract it from the identity, already in \(L^2(\pi_U)\), to obtain \(x\mapsto x-z\in L^2(\pi_U)\). This is a moment assertion for every chosen center, without requiring that the center minimize the potential.

\section{\texorpdfstring{Gibbs finiteness and the exact phase moment from strong convexity}{Gibbs finiteness and the exact phase moment from strong convexity}}
\label{proof-module-KineticLangevin-HamiltonianMoments-lean}
{\small\color{gray}Source module: \texttt{Kinetic\allowbreak{}Langevin/\allowbreak{}Hamiltonian\allowbreak{}Moments.\allowbreak{}lean}.\par}


The weak Hamiltonian limit needs finite mass and an integrable velocity-plus-force quadratic envelope. This module derives both for the actual exponential Hamiltonian density in finite dimensions. It factors that density into a radial Gaussian velocity weight and the potential Gibbs weight, uses existing direct Gaussian-domination estimates, and transfers position moments to force moments by Lipschitz growth.

\subsection{\texorpdfstring{Factor the actual Hamiltonian density}{Factor the actual Hamiltonian density}}


Because \(H(v,x)=f(x)+\|v\|^2/2\), the positive exponential weight factors as \(e^{-H(v,x)}=e^{-\|v\|^2/2}e^{-f(x)}\). Therefore its product-Haar weighted measure is the product of the separately weighted velocity and position measures. Lean applies the theorem for products of measures with densities, combines the two nonnegative real weights through \texttt{ENNReal.\allowbreak{}of\allowbreak{}Real\_\allowbreak{}mul}, and checks the exponent identity. Normalization has not been assumed or imposed at this step.

\subsection{\texorpdfstring{Derive finite mass of both factors and their product}{Derive finite mass of both factors and their product}}


The velocity weight is the radial Gaussian with precision \(1/2\), whose ordinary integral is finite. The position weight has finite mass by the previously proved strong-convexity Gaussian domination: a supporting quadratic lower bound for \(f\) makes \(e^{-f}\) integrable. Lean converts these finite ordinary integrals to finite extended density integrals, instantiates finite weighted measures, and uses the exact factorization to prove finiteness of the full phase measure.

\subsection{\texorpdfstring{Derive velocity and position second moments before normalization}{Derive velocity and position second moments before normalization}}


A quadratic factor is absorbed by radial Gaussian decay, so \(\int\|v\|^2e^{-\|v\|^2/2}\,dv<\infty\). The same strong-convexity tail estimate gives \(\int\|x\|^2e^{-f(x)}\,dx<\infty\). Lean uses the ordinary-integral characterization of integration under a density, checks that each exponential density is finite and measurable, and invokes the existing Gaussian-tail estimates. These are actual moments of the unnormalized factors, not moment contracts attached to a stationary process.

\subsection{\texorpdfstring{Transfer position moments to force moments}{Transfer position moments to force moments}}


Gradient Lipschitz continuity implies \(\|\nabla f(x)\|\le M\|x\|+\|\nabla f(0)\|\). Squaring and using \((a+b)^2\le2a^2+2b^2\) gives \(\|\nabla f(x)\|^2\le2M^2\|x\|^2+2\|\nabla f(0)\|^2\). The right-hand side is integrable under the finite position Gibbs weight by its second moment. Lean derives the force growth inequality from the Lipschitz difference bound and triangle inequality, then applies domination by this explicit integrable envelope.

\subsection{\texorpdfstring{Assemble the precise phase envelope needed by the weak limit}{Assemble the precise phase envelope needed by the weak limit}}


On the factored finite phase measure, \(\|v\|^2\) is integrable by pulling the velocity moment through the first coordinate, and \(\|\nabla f(x)\|^2\) is integrable by pulling the force moment through the second coordinate. Their sum is exactly the envelope appearing in the rejection theorem. Lean rewrites the phase measure by its proved factorization, supplies finiteness of both factors, and adds these two coordinate integrability results. Thus every tail premise in the later Hamiltonian stationarity limit follows from positive strong convexity and a Lipschitz gradient.

\section{\texorpdfstring{Exact Gibbs invariance of the actual Hamiltonian flow}{Exact Gibbs invariance of the actual Hamiltonian flow}}
\label{proof-module-KineticLangevin-HamiltonianStationarity-lean}
{\small\color{gray}Source module: \texttt{Kinetic\allowbreak{}Langevin/\allowbreak{}Hamiltonian\allowbreak{}Stationarity.\allowbreak{}lean}.\par}


This module closes the deterministic Hamiltonian stationarity argument. Explicit corrected Verlet proposals preserve the Gibbs measure exactly. Their rejection probability is quadratically small, while the uncorrected Verlet iterates converge to the actual selected Hamiltonian integral solution. Bounded tests transfer invariance through this limit. The finite-dimensional strong-convexity theorem derives every required tail premise, so its statement assumes no invariant flow, numerical convergence, or error contract.

\subsection{\texorpdfstring{Make the actual endpoint a measurable function of its initial state}{Make the actual endpoint a measurable function of its initial state}}


The constructed kinetic endpoint is jointly continuous in initialization and continuous forcing. Fix zero forcing, friction zero, inverse mass one, and endpoint time \(T\). Composing this joint map with \(z\mapsto(z,0)\) proves continuity of the selected Hamiltonian endpoint. Lean supplies the compact forcing-space zero explicitly and simplifies the composed definition. The endpoint pushforward is therefore a concrete measurable law.

\subsection{\texorpdfstring{Pass exact corrected invariance to the actual integral solution}{Pass exact corrected invariance to the actual integral solution}}


Let \(\pi\) be the actual Hamiltonian weighted measure, \(B=\int(\|v\|^2+\|\nabla f(x)\|^2)\,d\pi\), and \(K=M(2M+11)/8\). Use \(N=n+1\) Verlet steps of size \(h=T/N\). Exact corrected invariance and the iterated bounded-test estimate imply, for \(\|F\|\le1\), \[\left\|\int F(V_h^Nz)\,d\pi-\int F\,d\pi\right\|\le 2N\int(1-a_h)\,d\pi\le\frac{2T^2KB}{N}.\] The mesh is eventually at most one, so the proved quadratic rejection estimate applies; the displayed bound tends to zero. The independently proved numerical theorem gives \(V_{T/N}^Nz\to\Phi_Tz\), where \(\Phi\) is the actual Picard integral solution. The bounded-test limit theorem identifies \((\Phi_T)_\#\pi=\pi\). Lean assembles these actual results, proves the mesh limit, checks nonnegative multipliers, and simplifies the exact scalar bound. It assumes only the finite measure and precisely displayed quadratic moment required by this intermediate general theorem.

\subsection{\texorpdfstring{Discharge every tail premise from the potential assumptions}{Discharge every tail premise from the potential assumptions}}


In finite dimensions, positive strong convexity and differentiability make the actual Hamiltonian Gibbs measure finite by the factorization and Gaussian domination proved in the moments module. The globally Lipschitz gradient supplies the finite force second moment, and the thermal velocity factor supplies the velocity second moment. Substituting these checked results into the preceding limit theorem gives exact Hamiltonian Gibbs invariance from the original first-order potential assumptions. The theorem applies to every nonnegative finite horizon. It proves the Hamiltonian subflow; the full friction-and-noise kinetic diffusion additionally needs its thermostat splitting limit.

\section{\texorpdfstring{Identify Gaussian proposals with their radial density}{Identify Gaussian proposals with their radial density}}
\label{proof-module-UnadjustedLangevin-GaussianProposalDensity-lean}
{\small\color{gray}Source module: \texttt{Unadjusted\allowbreak{}Langevin/\allowbreak{}Gaussian\allowbreak{}Proposal\allowbreak{}Density.\allowbreak{}lean}.\par}


For \(h>0\), the affine standard Gaussian proposal \(c+\sqrt{2h}\,Z\) has density proportional to \(e^{-\|y-c\|^2/(4h)}\). The proof identifies probability measures by characteristic functions and does not require a preselected coordinate formula for the Gaussian normalizer.

\subsection{\texorpdfstring{radial\allowbreak{}Gaussian\allowbreak{}Integral}{radialGaussianIntegral}}


Set \(J(a)=\int_H e^{-a\|z\|^2}dz\) as an ordinary real integral. The positive-precision results below prove its integrability and positivity before using \(J(a)^{-1}\) as a normalizing constant.

\subsection{\texorpdfstring{radial\allowbreak{}Gaussian\allowbreak{}Normalizer}{radialGaussianNormalizer}}


Define \(C(a)=J(a)^{-1}\). This is an actual integral normalizer on the given finite-dimensional Hilbert space, avoiding a coordinate-dependent choice of a formula for Lebesgue density.

\subsection{\texorpdfstring{radial\allowbreak{}Gaussian\allowbreak{}Integral pos}{radialGaussianIntegral pos}}


For \(a>0\), Gaussian integrability was established in \texttt{integrable\_\allowbreak{}exp\_\allowbreak{}neg\_\allowbreak{}mul\_\allowbreak{}norm\_\allowbreak{}sq}. The integrand is strictly positive, so its ordinary integral is strictly positive by \texttt{integral\_\allowbreak{}exp\_\allowbreak{}pos}.

\subsection{\texorpdfstring{radial\allowbreak{}Gaussian\allowbreak{}Normalizer pos}{radialGaussianNormalizer pos}}


The inverse of a strictly positive real number is strictly positive. Apply this to \(J(a)\), whose positivity was proved immediately above. This also proves the normalizer is nonzero for later cancellation.

\subsection{\texorpdfstring{radial\allowbreak{}Gaussian\allowbreak{}Density}{radialGaussianDensity}}


Define \(q_{a,c}(z)=C(a)e^{-a\|z-c\|^2}\), embedded in extended nonnegative scalars. The center changes the translated argument, while the normalizer depends only on precision and dimension.

\subsection{\texorpdfstring{integrable radial\allowbreak{}Gaussian\allowbreak{}Density real}{integrable radialGaussianDensity real}}


Translation by \(c\) preserves Lebesgue measure, so \(e^{-a\|z-c\|^2}\) is integrable whenever \(a>0\). Multiply by the finite real constant \(C(a)\) to obtain integrability of the displayed real density.

\subsection{\texorpdfstring{integral radial\allowbreak{}Gaussian\allowbreak{}Density real}{integral radialGaussianDensity real}}


Translate \(z\mapsto z-c\) in the integral and pull out \(C(a)\). The result is \(C(a)J(a)=1\), justified by \(J(a)>0\). Thus normalization is independent of the center.

\subsection{\texorpdfstring{radial\allowbreak{}Gaussian\allowbreak{}Density is\allowbreak{}Probability}{radialGaussianDensity isProbability}}


The total mass of the with-density measure is the extended integral of its density. Nonnegativity and the proved real integrability identify it with \texttt{of\allowbreak{}Real} of the real integral, which is one. This constructs probability normalization, including finiteness.

\subsection{\texorpdfstring{integral complex radial\allowbreak{}Gaussian}{integral complex radialGaussian}}


For real \(a\), \(e^{-(a:\mathbb C)\|z\|^2}\) is the complex embedding of the real Gaussian exponential. Rewrite the integrand with \texttt{Complex.\allowbreak{}of\allowbreak{}Real\_\allowbreak{}exp}; the continuous real embedding commutes with integration, giving the complex embedding of \(J(a)\).

\subsection{\texorpdfstring{integral radial\allowbreak{}Gaussian character}{integral radialGaussian character}}


Apply the finite-dimensional complex Gaussian Fourier integral with imaginary linear coefficient. Combining the two exponentials and using \(i^2=-1\) gives \(\int e^{-a\|z\|^2}e^{i\langle z,t\rangle}dz=J(a)e^{-\|t\|^2/(4a)}\). The zero-frequency Gaussian formula identifies the normalization factor with the same real integral \(J(a)\).

\subsection{\texorpdfstring{char\allowbreak{}Fun radial\allowbreak{}Gaussian\allowbreak{}Density}{charFun radialGaussianDensity}}


Change from the density measure to a weighted Lebesgue integral. Translate \(z=u+c\); the character splits into \(e^{i\langle u,t\rangle}e^{i\langle c,t\rangle}\). The Fourier formula gives a factor \(J(a)\), canceled by \(C(a)\), leaving \(e^{-\|t\|^2/(4a)}e^{i\langle c,t\rangle}\).

\subsection{\texorpdfstring{gaussian\allowbreak{}Proposal\allowbreak{}Normalizer}{gaussianProposalNormalizer}}


Specialize the precision to \(a=1/(4h)\). The resulting constant normalizes a Gaussian with covariance \(2hI\). Its independence of the current state is essential when the same factor cancels in the forward and reverse Metropolis ratio.

\subsection{\texorpdfstring{gaussian\allowbreak{}Proposal\allowbreak{}Density}{gaussianProposalDensity}}


Define \(q_h(c,y)=C(1/(4h))e^{-\|y-c\|^2/(4h)}\). For \(h>0\), this is the density of a covariance-\(2hI\) Gaussian centered at \(c\). The density identification is proved later rather than built into the definition.

\subsection{\texorpdfstring{gaussian\allowbreak{}Proposal\allowbreak{}Normalizer pos}{gaussianProposalNormalizer pos}}


For \(h>0\), the precision \(1/(4h)\) is positive. Apply positivity of the radial Gaussian normalizer at that precision. This nonzero constant is used in detailed-balance cancellation.

\subsection{\texorpdfstring{radial\allowbreak{}Gaussian\allowbreak{}Density eq proposal\allowbreak{}Density}{radialGaussianDensity eq proposalDensity}}


Unfold the two density definitions and the proposal normalizer. Their exponents agree by \(-(1/(4h))\|y-c\|^2=-\|y-c\|^2/(4h)\). Function extensionality gives equality of the density functions, with no probabilistic argument needed.

\subsection{\texorpdfstring{gaussian proposal density}{gaussian proposal density}}


Compute characteristic functions on both sides. The affine image \(c+\sqrt{2h}Z\) has character \(e^{-h\|t\|^2}e^{i\langle c,t\rangle}\); the radial-density formula with \(a=1/(4h)\) gives the same expression. Both measures are probabilities, so \texttt{Measure.\allowbreak{}ext\_\allowbreak{}of\_\allowbreak{}char\allowbreak{}Fun} identifies them. The square-root identity uses \(h>0\).

\subsection{\texorpdfstring{gaussian proposal radial density}{gaussian proposal radial density}}


Rewrite the proposal density using \texttt{radial\allowbreak{}Gaussian\allowbreak{}Density\_\allowbreak{}eq\_\allowbreak{}proposal\allowbreak{}Density} in the established affine-Gaussian law equality. This yields the equivalent precision-parameter form used in the detailed-balance proof.

\subsection{\texorpdfstring{ula proposal radial density}{ula proposal radial density}}


Set the affine Gaussian center to \(c=x-h\nabla U(x)\). The fixed-state density theorem then gives \(x-h\nabla U(x)+\sqrt{2h}Z\) its exact radial density. No regularity of \(U\) is required because \(x\) is fixed throughout this statement.

\subsection{\texorpdfstring{ula proposal density}{ula proposal density}}


Expand the radial density at center \(x-h\nabla U(x)\). The exponent becomes \(-\|y-x+h\nabla U(x)\|^2/(4h)\), with the same positive normalizer. This is the explicit ULA transition density used in the Metropolis acceptance ratio.

\section{\texorpdfstring{The actual Gaussian-\allowbreak{}times-\allowbreak{}Gibbs phase target}{The actual Gaussian-times-Gibbs phase target}}
\label{proof-module-KineticLangevin-GibbsTarget-lean}
{\small\color{gray}Source module: \texttt{Kinetic\allowbreak{}Langevin/\allowbreak{}Gibbs\allowbreak{}Target.\allowbreak{}lean}.\par}


The normalized kinetic target is the independent product of a standard Gaussian velocity and the intended Gibbs position measure. This module identifies that product with an explicit scalar normalization of the actual Hamiltonian weighted measure, transfers the proved Hamiltonian invariance, and proves exact Gaussian thermostat invariance on the correct independent-noise product. Its position and velocity marginals are exact pushforward identities. These two invariant substeps are ingredients for kinetic Gibbs stationarity; passage to the continuous kinetic limit is a separate theorem.

\subsection{\texorpdfstring{Identify the Gaussian thermal density by its actual law}{Identify the Gaussian thermal density by its actual law}}


The already checked Gaussian proposal density at step \(h=1/2\) and center zero is exactly the standard Gaussian. Its scalar precision is \(1/(4h)=1/2\), and its affine map is the identity. Consequently its density is \(c_Ge^{-\|v\|^2/2}\), where \(c_G\) is the inverse of the actual finite Gaussian integral. Lean first normalizes the precision and the identity pushforward, then factors the positive normalizer out of the density using \texttt{with\allowbreak{}Density\_\allowbreak{}smul} and \texttt{ENNReal.\allowbreak{}of\allowbreak{}Real\_\allowbreak{}mul}. The remaining real exponential identity follows by scalar algebra. This is an equality with the existing probability measure \texttt{std\allowbreak{}Gaussian}, not an informal naming of a radial weight.

\subsection{\texorpdfstring{Normalize the Hamiltonian weighted measure explicitly}{Normalize the Hamiltonian weighted measure explicitly}}


Define \(\Pi=\varphi\otimes\pi\), velocity first, with \(\varphi\) the actual standard Gaussian and \(\pi\) the normalized Gibbs measure of \(f\). The Gaussian density identity and the existing definition of \(\pi\) give

\[\Pi=(c_G/Z_f)\,e^{-f(x)-\|v\|^2/2}\,dv\,dx.\]

Lean distributes the two scalar normalizers through the product of measures, then rewrites the previously proved factorization of the Hamiltonian Gibbs weight. Strong convexity with \(m>0\) makes the position normalization positive and finite, so the product is a probability measure. Its probability proof invokes the already derived Gibbs probability theorem; normalization is not an assumption about a stationary process.

\subsection{\texorpdfstring{Transfer the proved Hamiltonian invariance to the canonical target}{Transfer the proved Hamiltonian invariance to the canonical target}}


The actual selected Hamiltonian flow already preserves the unnormalized Hamiltonian weighted measure, by the corrected-Verlet approximation and its proved numerical limit. Multiplying a measure by a fixed scalar commutes with a measurable pushforward. The explicit normalization identity therefore gives \((H_t)_\#\Pi=\Pi\). Lean supplies measurability from continuous dependence of the endpoint on its initial phase state, uses \texttt{Measure.\allowbreak{}map\_\allowbreak{}smul}, and applies the proved Hamiltonian invariance theorem. This statement has only the potential assumptions and finite-dimensional setting; no invariant-target hypothesis remains.

\subsection{\texorpdfstring{Prove the thermostat identity with fresh independent Gaussian noise}{Prove the thermostat identity with fresh independent Gaussian noise}}


Let \(X\) have any probability position law, independently of a standard Gaussian velocity \(V\), and let \(\xi\) be a fresh independent standard Gaussian. If \(a^2+b^2=1\), then \((aV+b\xi,X)\) again has the original independent velocity-position law. The input is explicitly \((\varphi\otimes\pi_X)\otimes\varphi\). A measure-preserving reassociation and swap reorder its three coordinates to \((V,\xi),X\). The previously proved Gaussian linear-combination law acts on the first pair; the position coordinate is unchanged. Lean verifies both the coordinate permutation and its product-measure pushforward before applying \texttt{Measure.\allowbreak{}map\_\allowbreak{}prod\_\allowbreak{}map}. Thus independence and the output marginal are derived, not supplied.



Specializing the arbitrary position law to the actual normalized Gibbs target proves thermostat invariance of \(\Pi\). Exact OU evolution chooses \(a=e^{-\gamma h}\) and \(b=\sqrt{1-e^{-2\gamma h}}\); an OBABO half-step uses the corresponding half-interval coefficients. The theorem admits both choices through the explicit variance balance.

\subsection{\texorpdfstring{Recover the two equilibrium marginals exactly}{Recover the two equilibrium marginals exactly}}


Projecting the independent product \(\varphi\otimes\pi\) to position integrates out a unit-mass Gaussian and leaves \(\pi\). Projecting to velocity integrates out the normalized position measure and leaves \(\varphi\). Lean installs the derived Gibbs probability instance and applies the two product-projection identities, simplifying their unit total-mass factors. These equalities supply the target position marginal and Gaussian velocity marginal used by the later L\ensuremath{{}^{2}} coupling theorem once stationarity of the full kinetic dynamics is proved.

\section{\texorpdfstring{Identify recursive sampling laws with iterated Markov transitions}{Identify recursive sampling laws with iterated Markov transitions}}
\label{proof-module-UnadjustedLangevin-MarkovLaw-lean}
{\small\color{gray}Source module: \texttt{Unadjusted\allowbreak{}Langevin/\allowbreak{}Markov\allowbreak{}Law.\allowbreak{}lean}.\par}


The path-space law of an iid-noise recursion agrees with the repeated one-step transition \(\mu\mapsto F_\#(\mu\otimes\eta)\). The proof explicitly establishes independence of the current state and the fresh innovation.

\subsection{\texorpdfstring{finite\allowbreak{}Noise\allowbreak{}Chain}{finiteNoiseChain}}


Evaluate the same recursion using only a finite prefix \((w_0,\ldots,w_{n-1})\). At the successor, recurse on the first \(n\) coordinates and use the final coordinate as the next innovation. This makes finite dependence explicit for the fresh-noise independence proof.

\subsection{\texorpdfstring{finite\allowbreak{}Noise\allowbreak{}Chain eq}{finiteNoiseChain eq}}


Induct on prefix length. Restricting an infinite path to its first \(n\) entries gives the same previous state, and the final finite coordinate is \(w_n\). The successor updates thus agree, identifying finite-prefix and infinite-path implementations pointwise.

\subsection{\texorpdfstring{finite\allowbreak{}Noise\allowbreak{}Chain measurable}{finiteNoiseChain measurable}}


The zero map is the initial projection. At the successor, restriction to the earlier finite coordinates and evaluation of the last coordinate are measurable; combine them with the measurable update and the induction hypothesis. This proves the state depends measurably on a finite prefix.

\subsection{\texorpdfstring{iid prefix independent fresh}{iid prefix independent fresh}}


The index sets \(\{0,\ldots,n-1\}\) and \(\{n\}\) are disjoint. Independence of the coordinate family on the iid product source therefore gives independence of the whole prefix vector and the next coordinate. This is stronger than pairwise independence of individual past entries.

\subsection{\texorpdfstring{markov\allowbreak{}Step}{markovStep}}


Define \(M_\eta(\mu)=F_*(\mu\otimes\eta)\). Its product measure means the innovation is fresh and independent of the current state. The subsequent theorem derives that independence for the canonical recursion, rather than assuming it at each step.

\subsection{\texorpdfstring{markov\allowbreak{}Laws}{markovLaws}}


Set \(\mu_0=\mu\), \(\mu_{n+1}=M_\eta(\mu_n)\). These are iterated independent-noise transition laws. The recursion alone does not yet assert that a previously defined sample-path algorithm realizes them.

\subsection{\texorpdfstring{chain\allowbreak{}Law joint fresh}{chainLaw joint fresh}}


The current state is a measurable function of initialization and the first \(n\) noises. Product-source independence and prefix/\allowbreak{}fresh independence make it independent of \(w_n\). Its marginal is \texttt{chain\allowbreak{}Law}, and the innovation marginal is \(\eta\); their joint law is therefore the product of those actual marginals.

\subsection{\texorpdfstring{chain\allowbreak{}Law succ iid}{chainLaw succ iid}}


The pointwise successor is \(F(X_n,w_n)\). Push the proved product joint law of \((X_n,w_n)\) through \(F\); this yields \(\mu_{n+1}=M_\eta(\mu_n)\) for the canonical recursion laws. No later-time state/\allowbreak{}noise independence is assumed.

\subsection{\texorpdfstring{chain\allowbreak{}Law eq markov\allowbreak{}Laws}{chainLaw eq markovLaws}}


The canonical path law and the iterated Markov law have the same initial value and the same successor equation. Induction on \(n\) identifies them at every index. This closes the connection between sample-path and transition-law definitions.

\subsection{\texorpdfstring{markov\allowbreak{}Step is\allowbreak{}Probability}{markovStep isProbability}}


The independent product of two probability measures has mass one. A measurable pushforward preserves that total mass, so \texttt{markov\allowbreak{}Step} is again a probability measure.

\subsection{\texorpdfstring{markov\allowbreak{}Laws is\allowbreak{}Probability}{markovLaws isProbability}}


Induct from the initial probability and apply \texttt{markov\allowbreak{}Step\_\allowbreak{}is\allowbreak{}Probability} at each successor. This proves normalization of every transition iterate and provides the probability instances used in later law comparisons.

\subsection{\texorpdfstring{noise\allowbreak{}Driven\allowbreak{}Chain law eq markov\allowbreak{}Laws}{noiseDrivenChain law eq markovLaws}}


First identify the independent iid realization with the canonical product-path output law. Then apply \texttt{chain\allowbreak{}Law\_\allowbreak{}eq\_\allowbreak{}markov\allowbreak{}Laws}. Every realization satisfying the original initialization/\allowbreak{}noise assumptions consequently has the same proved Markov iterate law.

\section{\texorpdfstring{Bounded tests of the actual independent-\allowbreak{}noise Markov chain}{Bounded tests of the actual independent-noise Markov chain}}
\label{proof-module-KineticLangevin-MarkovTests-lean}
{\small\color{gray}Source module: \texttt{Kinetic\allowbreak{}Langevin/\allowbreak{}Markov\allowbreak{}Tests.\allowbreak{}lean}.\par}


To prove that a continuous dynamics preserves its Gibbs distribution, approximate it by concrete random steps. A small one-step weak error at the Gibbs measure must be propagated through those actual steps. The dual transition integrates the next test value against the independent fresh noise. Its measurability, uniform bound and exact relation to the forward law are proved here.

\subsection{\texorpdfstring{Integrate the next test against the fresh noise}{Integrate the next test against the fresh noise}}


For a measurable step map \(S(x,\xi)\) and noise probability \(Q\), define \(PF(x)=\int F(S(x,\xi))Q(d\xi)\). The measurable parameter-integral theorem proves that \(PF\) is measurable. If \(\lVert F\rVert\le C\) everywhere, the norm-of-integral inequality gives \(\lVert PF(x)\rVert\le\int C\,dQ=C\). Independence is encoded by the product measure used by the actual Markov update.

\subsection{\texorpdfstring{Prove the duality identity for the implemented recursion}{Prove the duality identity for the implemented recursion}}


The law update is the pushforward of \(\mu\otimes Q\) by \(S\). The pushforward integration formula and Fubini give \(\int F\,d(\mu P)=\int PF\,d\mu\). Boundedness on a probability space supplies integrability. Induction proves that every \(P^nF\) remains measurable and bounded by the original \(C\), and that \(\int F\,d(\mu P^n)=\int P^nF\,d\mu\). Thus the dual operator describes the actual law recursion rather than an assumed sequence of measures.

\subsection{\texorpdfstring{Accumulate one-\allowbreak{}step weak errors without losing the test bound}{Accumulate one-step weak errors without losing the test bound}}


Suppose every measurable test bounded by \(C\) has one-step defect at \(\pi\) at most \(\delta\). Apply this estimate to \(P^nF\), whose bound is still \(C\). Insert and subtract \(\int P^nF\,d\pi\), apply the triangle inequality, and use the induction hypothesis. The resulting estimate is \(\left|\int F\,d(\pi P^n)-\int F\,d\pi\right|\le n\delta\). No exact invariance is assumed. In the splitting application the rejection bound is quadratic in the time step, so \(n\delta\to0\) on a fixed time interval.

\section{\texorpdfstring{Quadratic stationary weak error for the actual kinetic splitting}{Quadratic stationary weak error for the actual kinetic splitting}}
\label{proof-module-KineticLangevin-SplittingTarget-lean}
{\small\color{gray}Source module: \texttt{Kinetic\allowbreak{}Langevin/\allowbreak{}Splitting\allowbreak{}Target.\allowbreak{}lean}.\par}


This module applies the exact corrected Verlet identity and the exact Gaussian thermostat identity to the actual normalized Gaussian-times-Gibbs target. Averaging a bounded test over the thermostat preserves its bound. The uncorrected splitting therefore has a proved quadratic one-step stationary weak defect, and its actual Markov iterate accumulates at most a linear number of these defects.

\subsection{\texorpdfstring{Use the actual thermal Gaussian innovation measure}{Use the actual thermal Gaussian innovation measure}}


The additive innovation measure is the pushforward of the Euclidean standard Gaussian by the coefficient \(\sigma_h=\sqrt{1-e^{-2\gamma h}}\). It is a probability measure because it is a measurable pushforward of a probability measure. The preceding OU marginal theorem identifies this exact law with the actual canonical convolution innovation; here it is the concrete measure used by the Markov transition.

\subsection{\texorpdfstring{Normalize the exact corrected Verlet identity}{Normalize the exact corrected Verlet identity}}


The canonical target is an explicitly proved finite scalar normalization of the actual Hamiltonian weighted measure. A density-weighted measure, its measurable pushforward, and their sum all commute with this scalar multiplication. Scaling the earlier exact identity therefore gives \((V_h)_\#(a_h\pi)+R_\#((1-a_h)\pi)=\pi\) for the actual normalized target. Lean rewrites the explicit normalization, both density scalings, both pushforward scalings, and the corrected Hamiltonian Gibbs identity.

\subsection{\texorpdfstring{Derive the normalized target phase moment}{Derive the normalized target phase moment}}


The unnormalized Hamiltonian measure already has integrable \(\|v\|^2+\|\nabla f(x)\|^2\) from first-order strong convexity and Lipschitz force growth. The explicit normalizing scalar is finite: its Gaussian factor is an ordinary finite real number and its inverse potential partition is finite because the partition is strictly positive. Lean checks these extended-real facts and applies integrability under finite scalar multiplication. Thus the rejection envelope remains an actual finite target moment.

\subsection{\texorpdfstring{Preserve the target with the actual additive OU law}{Preserve the target with the actual additive OU law}}


The damping \(\delta_h=e^{-\gamma h}\) and innovation scale satisfy \(\delta_h^2+\sigma_h^2=1\). The already proved thermostat result preserves the target when driven by a fresh standard Gaussian. Mapping that Gaussian through \(\sigma_h\) changes its law to the explicit additive innovation measure. Lean proves the coefficient identity, maps the independent product source through this noise scaling, composes the two pushforwards, and obtains exact preservation by the additive OU transition.

\subsection{\texorpdfstring{Apply corrected invariance to the thermostat-\allowbreak{}averaged test}{Apply corrected invariance to the thermostat-averaged test}}


For \(\|F\|\le C\), define \(G(z)=\mathbb E_N[F(O_h(z,N))]\). The actual noise integral makes \(G\) measurable with \(\|G\|\le C\). Exact thermostat invariance gives \(\int G\,d\pi=\int F\,d\pi\), while the actual split transition has expectation \(\int G(V_hz)\,d\pi\). Apply the proved corrected-Verlet weak estimate to \(G\), then integrate the actual quadratic rejection envelope. This gives \[\left\|\int F\,d(P_h\pi)-\int F\,d\pi\right\|\le2Ch^2\frac{M(2M+11)}8\int(\|v\|^2+\|\nabla f(x)\|^2)\,d\pi.\] Lean derives both expectation identities from the actual Markov-step integral formula, checks measurability and boundedness of the averaged test, and composes the checked rejection and target moment estimates.

\subsection{\texorpdfstring{Accumulate only a linear number of genuine local defects}{Accumulate only a linear number of genuine local defects}}


The generic Markov test theorem applies the one-step estimate to each iterated dual test, whose bound remains \(C\), and telescopes their expectations under the fixed target. Consequently the actual \(n\)-step splitting law has weak error at most \(n\) times the displayed quadratic defect. Lean verifies nonnegativity of the target moment and all scalar coefficients, supplies the concrete measurable transition, and instantiates the proved generic theorem. With \(h=T/n\), this actual error tends to zero.

\section{\texorpdfstring{Actual phase laws and their L\ensuremath{{}^{2}} limits}{Actual phase laws and their L² limits}}
\label{proof-module-KineticLangevin-LpPhaseLaw-lean}
{\small\color{gray}Source module: \texttt{Kinetic\allowbreak{}Langevin/\allowbreak{}Lp\allowbreak{}Phase\allowbreak{}Law.\allowbreak{}lean}.\par}


The phase variable is a pair of actual L\ensuremath{{}^{2}} classes, and its law is the pushforward of their measurable representatives. Dual characteristic functions depend continuously on this pair. This identifies a limiting stationary phase law from bounded-test estimates without incorrectly treating the usual maximum product norm as a Hilbert norm.

\subsection{\texorpdfstring{Verify measurability of actual phase representatives}{Verify measurability of actual phase representatives}}


Each L\ensuremath{{}^{2}} class has a fixed strongly measurable representative. Taking the pair of its velocity and position representatives gives a measurable phase map. Lean uses the two strong-measurability results and the measurable product constructor, making the subsequent pushforward an actual measure.

\subsection{\texorpdfstring{Reduce phase characteristic functions to one real L\ensuremath{{}^{2}} variable}{Reduce phase characteristic functions to one real L² variable}}


For a continuous real linear functional \(L\) on phase space, let \(L_v(v)=L(v,0)\) and \(L_x(x)=L(0,x)\). The real L\ensuremath{{}^{2}} variable \(R=L_v(V)+L_x(X)\) represents \(L(V(\omega),X(\omega))\) almost everywhere. Its law therefore has scalar characteristic function at frequency one equal to the dual characteristic function of the phase law. The map \((V,X)\mapsto R\) is continuous, since the two continuous linear maps lift continuously to L\ensuremath{{}^{2}} and addition is continuous. The existing scalar characteristic-function continuity theorem proves the desired phase continuity. Lean explicitly intersects the representative events for both lifted linear maps and their sum, proves the coordinate decomposition by linearity, rewrites the actual pushforward through \(L\), and composes the two continuous maps.

\subsection{\texorpdfstring{Identify the law of an actual L\ensuremath{{}^{2}} limit from vanishing weak errors}{Identify the law of an actual L² limit from vanishing weak errors}}


Suppose actual phase pairs \(S_n\) converge in the product of L\ensuremath{{}^{2}} spaces to \(S\), and their bounded complex-test errors relative to a fixed finite measure \(\pi\) are eventually bounded by \(\delta_n\to0\). For each dual functional use the norm-one test \(F(z)=e^{iLz}\). The error squeeze gives convergence of the approximating characteristic functions to that of \(\pi\); the continuity theorem gives convergence of the same sequence to the characteristic function of the actual limiting phase law. Uniqueness of limits equates them, and uniqueness of finite measures from their dual characteristic functions gives \(P\circ S^{-1}=\pi\). Lean rewrites the actual phase expectation by the measurable pushforward formula and invokes the checked scalar squeeze and measure-uniqueness theorems. The weak estimates must be proved by the concrete numerical and probability layers; stationary output is not assumed.

\section{\texorpdfstring{The actual splitting chain has its independent-\allowbreak{}noise Markov law}{The actual splitting chain has its independent-noise Markov law}}
\label{proof-module-KineticLangevin-SplittingLaw-lean}
{\small\color{gray}Source module: \texttt{Kinetic\allowbreak{}Langevin/\allowbreak{}Splitting\allowbreak{}Law.\allowbreak{}lean}.\par}


A law identity must be established for the random variables actually used by the construction. The input here is the complete joint law of the initial phase and the entire innovation stream. Composing this law with the measurable implemented recursion proves the endpoint law. The same identity is then established for the genuine L\ensuremath{{}^{2}} implementation using its proved almost-everywhere representatives.

\subsection{\texorpdfstring{Push the full input law through the actual recursion}{Push the full input law through the actual recursion}}


Let \(C_n(z,w)\) be the explicit recursion run for \(n\) steps. Its joint measurability has already been proved by induction. If \((X,(N_k)_k)\) has the product law \(\pi\otimes Q^{\otimes\mathbb N}\), the pushforward composition rule says its output has law \((\pi\otimes Q^{\otimes\mathbb N})\circ C_n^{-1}\). The existing finite-prefix proof identifies this with the actual iterated Markov law \(\pi P^n\). On an explicit initial-law/\allowbreak{}source product, the initial phase and a source stream are independent by the product pushforward identity. Lean proves both statements directly using \texttt{Has\allowbreak{}Law.\allowbreak{}comp} and \texttt{Measure.\allowbreak{}map\_\allowbreak{}prod\_\allowbreak{}map}.

\subsection{\texorpdfstring{Match the two concrete implementations step for step}{Match the two concrete implementations step for step}}


\texttt{kinetic\allowbreak{}Split\allowbreak{}Noise\allowbreak{}Chain} uses Verlet followed by the exact OU update. Induction shows it is precisely \texttt{noise\allowbreak{}Driven\allowbreak{}Chain} for that step. Continuity of the force makes every kick, drift and thermostat continuous, so their composition is measurable. The law theorem can therefore be applied to the actual splitting map, with no additional transition-measurability premise.

\subsection{\texorpdfstring{Eliminate the choice of L\ensuremath{{}^{2}} representatives}{Eliminate the choice of L² representatives}}


The lifted chain has already been proved equal almost everywhere to the pointwise vector chain with the same initial state and innovations. Replace its pushforward using this equality, then apply the full-input-law theorem. This proves the actual L\ensuremath{{}^{2}} endpoint law is the independent-noise Markov iterate. Quotient representatives cannot create an extra source of randomness or change the chain law.

\section{\texorpdfstring{Exact kinetic Gibbs stationarity from the actual initial/\allowbreak{}noise laws}{Exact kinetic Gibbs stationarity from the actual initial/noise laws}}
\label{proof-module-KineticLangevin-StationarityTransfer-lean}
{\small\color{gray}Source module: \texttt{Kinetic\allowbreak{}Langevin/\allowbreak{}Stationarity\allowbreak{}Transfer.\allowbreak{}lean}.\par}


This theorem assembles the full analytic stationarity limit. It starts with an actual continuous L\ensuremath{{}^{2}} kinetic integral solution. Its explicit input-law premise states the independent target initial phase and complete canonical OU innovation stream; concrete product-source modules derive that premise. Numerical convergence, stationary weak defects, actual output-law transfer, and the final invariant law are all proved inside the theorem.

\subsection{\texorpdfstring{Derive exact stationarity through the actual splitting limit}{Derive exact stationarity through the actual splitting limit}}


Let \(\pi\) be the actual normalized Gaussian-times-Gibbs target, and let \(Y\) solve the actual kinetic integral equation driven by the continuous L\ensuremath{{}^{2}} input \(B\). On a mesh \(h_n=T/(n+1)\), form the concrete Verlet–OU approximation \(S_n\) using the filtered increments of this same \(B\). The proved Banach splitting theorem gives \(S_n\to Y_T\) in the actual product of L\ensuremath{{}^{2}} spaces.



The full initial/\allowbreak{}innovation joint law is the explicit product \(\pi\otimes Q_{h_n}^{\otimes\mathbb N}\). The checked source-law transfer theorem, together with the almost-everywhere vector-chain representation, identifies the actual law of \(S_n\) with the implemented Markov iterate. Its bounded-test error is at most \[2(n+1)h_n^2 K A=\frac{2T^2KA}{n+1},\qquad K=\frac{M(2M+11)}8,\quad A=\int(\|v\|^2+\|\nabla f(x)\|^2)\,d\pi.\] The moment is finite by the proved target tail theorem. The mesh is eventually at most one, so every quadratic rejection estimate applies, and the displayed scalar error tends to zero.



The actual phase-law continuity theorem now identifies the L\ensuremath{{}^{2}} limiting phase law with \(\pi\). Lean constructs the true approximation sequence, applies the proved convergence theorem, rewrites its chain law using the explicit initial/\allowbreak{}noise product distribution, rewrites the bounded test expectation through the actual pushforward, invokes the proved Markov weak-error bound, and checks the scalar \(T^2/(n+1)\) decay. It then applies dual-characteristic uniqueness for finite measures. No stationary solution, local-error conclusion, output-law independence field, or numerical-limit conclusion is assumed. The remaining input-law hypothesis is a concrete source probability identity, supplied by the canonical independent-product construction.

\section{\texorpdfstring{The complete canonical KLMC theorem on arbitrary finite coordinate sets}{The complete canonical KLMC theorem on arbitrary finite coordinate sets}}
\label{proof-module-KineticLangevin-GenericCanonicalKLMC-lean}
{\small\color{gray}Source module: \texttt{Kinetic\allowbreak{}Langevin/\allowbreak{}Generic\allowbreak{}Canonical\allowbreak{}KLMC.\allowbreak{}lean}.\par}


The entire actual-source KLMC proof is stated for a general finite coordinate type \(I\). This includes the particle index \(\{\theta\}\sqcup(\{1,\ldots,N\}\times\{x\})\) without a separate coordinate-isometry transport. All initial couplings, canonical Brownian inputs, selected reference flows, Gibbs stationary laws and genuine Wasserstein infima are the same concrete constructions as in the finite-dimension presentation; the dimension in the final bias is exactly \(|I|\).

\subsection{\texorpdfstring{Specify the concrete probability space and source projection}{Specify the concrete probability space and source projection}}


The phase variable is velocity first. The sample space contains the two coupled initial positions, one shared Gaussian velocity and one vector Brownian source. The measure is the actual initial-coupling product extension; probability follows from its first marginal and from the Gaussian/\allowbreak{}Brownian probability factors. The Brownian-source projection preserves its measure. Pullback along that projection gives a continuous linear map on L\ensuremath{{}^{2}}, and its representative theorem identifies it almost everywhere with evaluation on the actual product sample.

\subsection{\texorpdfstring{Construct the common Brownian input and filtered innovation}{Construct the common Brownian input and filtered innovation}}


The scaled source lift includes the physical amplitude \(\sqrt{2\gamma}\). At iteration \(n\), the input is the actual vector Brownian increment on \((nh,nh+t]\), and the filtered input is the exponential response to that same increment. The unscaled innovation pair consists of the exact velocity and position Wiener responses at the step endpoint. Both input paths are continuous in L\ensuremath{{}^{2}}, and the additive input starts at zero. These follow from the source lift's continuity and linearity and the proved vector-response properties.

\subsection{\texorpdfstring{Identify the exact noise equations and implemented numerical step}{Identify the exact noise equations and implemented numerical step}}


The constructed filtered input satisfies \(B_n(t)=W_n(t)-\gamma\int_0^tB_n(s)ds\). Its endpoint is the scaled velocity innovation, and its actual time primitive is the scaled positional innovation. Lean transports the vector Brownian equations through the continuous linear source lift and commutes that lift with Bochner integration. Induction on the actual recursions then proves that the exponential-input chain is exactly \texttt{klmc\allowbreak{}Chain}, with the canonical correlated innovation pair and the prescribed fluctuation-dissipation scaling.

\subsection{\texorpdfstring{Connect Brownian L\ensuremath{{}^{2}} innovations with the generic output-\allowbreak{}law construction}{Connect Brownian L² innovations with the generic output-law construction}}


The L\ensuremath{{}^{2}} innovation pair has the same representatives as the pointwise canonical Brownian innovation. The proof pulls back the previously established vector-response representative theorem to the coupling product and combines the two component identities. Comparing with the generic product-law layer's representatives gives equality of the actual L\ensuremath{{}^{2}} pairs. Thus the numerical chain analyzed here has the coupling-independent output law already proved for the implemented KLMC chain.

\subsection{\texorpdfstring{Build the actual numerical and stationary-\allowbreak{}reference initial states}{Build the actual numerical and stationary-reference initial states}}


The two initial phases share exactly one Gaussian velocity variable. The numerical phase uses the first coupled position; the reference phase uses the second. The reference representative is proved to be the explicit pair \((V,X_\mathrm{ref})\). Product independence and the two marginal laws prove that its actual law is \(\mathcal N(0,I_{|I|})\otimes\pi\). The finite target second moment and finite coupling cost supply the required L\ensuremath{{}^{2}} initial positions.

\subsection{\texorpdfstring{Construct all three chains and segment families}{Construct all three chains and segment families}}


The numerical chain applies the actual lifted original gradient and the constructed correlated innovations. The reference chain stitches exact selected kinetic segments, each driven by the actual canonical Brownian increment for that iteration. The corresponding reference segments start from these actual reference states. The auxiliary exact segment starts instead from the actual numerical state and receives the identical additive input. The drift is the genuine pointwise gradient lifted to L\ensuremath{{}^{2}}; no substitute potential on the L\ensuremath{{}^{2}} space is assumed.

\subsection{\texorpdfstring{Prove every analytic premise of the KLMC coupling theorem}{Prove every analytic premise of the KLMC coupling theorem}}


The bundled specification supplies continuity of auxiliary paths, reference paths and filtered inputs; exact initial states; reference endpoint matching; the implemented exponential update; both actual common-input integral equations; the exponential mild identity; the kinematic primitive identity; and equal initial velocities. Each field follows from the concrete definitions and earlier proved equations. The mild and kinematic fields are derived internally. The only remaining probabilistic input for a final stationary-error theorem is the stationary law of the actual global reference diffusion, proved separately and connected to these stitched segments by uniqueness.

\subsection{\texorpdfstring{Construct the full canonical input and identify its increments}{Construct the full canonical input and identify its increments}}


The full input is the scalar canonical Brownian value assembled spatially and pulled onto the actual initial-coupling extension with amplitude \(\sqrt{2\gamma}\). It is continuous in L\ensuremath{{}^{2}} because each assembly and pullback is continuous linear. The Brownian value at zero vanishes: Lean verifies that its deterministic time indicator has empty support, then uses linearity of the isonormal map. For nonnegative local time, the constructed segment input equals the full input difference at \(nh+t\) and \(nh\). The proof checks the nonnegative real-to-NNReal casts and uses linearity of spatial assembly and source pullback.

\subsection{\texorpdfstring{Construct the actual full reference solution}{Construct the actual full reference solution}}


The global reference is the selected kinetic solution with the genuine lifted gradient, the canonical reference initial phase and the full canonical Brownian input. Its specification proves continuity, its exact starting phase, and the actual kinetic integral equation over the whole finite horizon. There is no stationary-law requirement in its definition.

\subsection{\texorpdfstring{Identify every stitched reference point by uniqueness}{Identify every stitched reference point by uniqueness}}


The generic stitching theorem now applies because the local inputs are proved to be actual increments of the full input. Consequently \(R_n=Y_{nh}\), and \(Z_n(t)=Y_{nh+t}\) throughout each step whose endpoint lies in the global horizon. Lean instantiates the existing exact uniqueness results with the actual L\ensuremath{{}^{2}} gradient and checks the horizon arithmetic. Thus a stationary law for the full constructed kinetic solution supplies the law of all reference states used by the quantitative KLMC proof.

\subsection{\texorpdfstring{Project the coupling source onto an independent phase/\allowbreak{}Brownian product}{Project the coupling source onto an independent phase/Brownian product}}


The projection sends the full sample to \(((V,X_\mathrm{ref}),\omega_\mathrm{Brownian})\). Its law is \((\mathcal N(0,I_{|I|})\otimes\pi)\otimes P_\mathrm{Brownian}\). Lean first projects the initial coupling onto its second marginal, retaining the independent Gaussian/\allowbreak{}Brownian source. It then applies measure-preserving product associativity and swaps the position and velocity factors. Thus independence of the complete reference initial phase and the complete Brownian source is proved for the actual product construction.

\subsection{\texorpdfstring{Identify the actual full-\allowbreak{}input OU increment}{Identify the actual full-input OU increment}}


The deterministic integration-by-parts input used in the kinetic splitting proof is \(\Delta W-\gamma\int_0^he^{-\gamma(h-s)}\Delta W_sds\). The actual local-input increment identity replaces each difference of the full canonical input by its constructed segment input. Applying the vector exponential-response formula through the continuous linear source lift proves that this OU increment is exactly the filtered velocity response. Its representative is consequently the scaled canonical vector response on the original Brownian source, pulled onto the coupling extension.

\subsection{\texorpdfstring{Derive the full independent thermal Gaussian joint law}{Derive the full independent thermal Gaussian joint law}}


The vector OU stream already has the exact iid Gaussian law with scale \(\sqrt{1-\psi_0(h)^2}\). The product-source joint-law theorem combines this stream with the initial phase target. Composition with the proved phase/\allowbreak{}Brownian projection law transfers that product law to the actual coupling extension. Finally, the representative identities identify its initial phase with the constructed reference initial L\ensuremath{{}^{2}} phase and its stream with the actual full-input OU increments. The resulting joint law is precisely the source obligation of the proved kinetic Gibbs stationarity transfer theorem.

\subsection{\texorpdfstring{Fix the actual numerical output law}{Fix the actual numerical output law}}


The definition uses the original pointwise KLMC method, the standard Gaussian initial velocity and the exact canonical correlated Brownian innovations. The numerical L\ensuremath{{}^{2}} chain is identified with the generic product-law construction by equality of the actual innovation variables. The established output-law theorem then proves that its pushforward is independent of which initial position coupling was chosen.

\subsection{\texorpdfstring{Transfer one actual global stationary phase law to the reference segments}{Transfer one actual global stationary phase law to the reference segments}}


The named stationary property concerns the full constructed kinetic flow: at every time its phase law is \(\mathcal N(0,I_{|I|})\otimes\pi\). It contains no numerical convergence conclusion. Exact chain-to-global identification transfers its position marginal to every reference grid point. Exact segment-to-global identification transfers its Gaussian velocity marginal to every intermediate time in a reference segment. Lean obtains these marginals by composition with the product-measure projections.

\subsection{\texorpdfstring{Apply the proved coupling estimate to one actual initial plan}{Apply the proved coupling estimate to one actual initial plan}}


The initial position L\ensuremath{{}^{2}} discrepancy equals the true cost of the supplied transport coupling; this follows from the concrete extension's representative and norm identities. The segment specification furnishes all analytic local-error premises internally. Together with the derived stationary position and Gaussian velocity marginals, the quantitative L\ensuremath{{}^{2}} KLMC theorem gives \(W_2(\nu_n,\pi)\le\sqrt2(1-3mh/(4\gamma))^n\operatorname{cost}(c)+Mh\sqrt{2|I|}/m\). Output-law equality and initial-cost equality rewrite this into a bound for the actual numerical law.

\subsection{\texorpdfstring{Prove output finiteness and take the transport infimum}{Prove output finiteness and take the transport infimum}}


Finite initial Wasserstein distance supplies an actual finite-cost initial coupling. Its numerical and reference states are genuine L\ensuremath{{}^{2}} variables, so their distance gives a finite transport bound on the output law. This proves output finiteness before real-valued transport arithmetic is used. The scalar recurrence estimate proves the contraction factor is nonnegative under the paper's step restrictions. Lean then applies the proved infimum theorem over all finite-cost couplings, obtaining the same estimate with \(W_2(\mu,\pi)\) in place of the cost of a chosen plan. No existence of an optimal coupling is assumed.

\subsection{\texorpdfstring{Prove Gibbs stationarity of the actual global coupling reference}{Prove Gibbs stationarity of the actual global coupling reference}}


The reference begins with standard Gaussian velocity independent of the target Gibbs position. Its full Brownian forcing and its entire filtered innovation stream are actual constructed objects on the initial-coupling extension. Their exact initial/\allowbreak{}noise product law is already proved. The kinetic stationarity transfer theorem therefore applies directly to the actual integral solution. For a time \(t\) inside a larger horizon, Lean restricts the actual integral equation to \([0,t]\), supplies the continuous input and its zero starting value, and applies the proved splitting-limit stationarity result. This yields the full phase law \(\mathcal N(0,I_{|I|})\otimes\pi\) at every reference time, without a remaining stationarity or Gaussian-marginal hypothesis.

\subsection{\texorpdfstring{Conclude the source-\allowbreak{}paper KLMC Wasserstein estimate}{Conclude the source-paper KLMC Wasserstein estimate}}


Let \(f\) be differentiable, \(m\)-strongly convex with \(m>0\), and have an \(M\)-Lipschitz gradient. If \(\gamma^2\ge m+M\), \(\gamma>0\), and \(0<h\le m/(4\gamma M)\), the actual KLMC position law satisfies \[W_2(\nu_n,\pi)\le\sqrt2\left(1-\frac{3mh}{4\gamma}\right)^nW_2(\mu,\pi)+\frac{Mh\sqrt{2|I|}}m.\] The initial position distribution may be arbitrary with finite initial Wasserstein distance; its velocity is the explicit independent standard Gaussian. Lean first derives the finite Gibbs second moment from strong convexity. It then supplies the just-proved global-reference stationarity for every initial transport coupling to the complete canonical coupling/\allowbreak{}infimum theorem. The conclusion has no conditional local-error, reference-process, Gaussian-law, stationary-law or optimal-coupling premise. The smoothness assumptions are expressed by the actual gradient and strong convexity, so the theorem also covers the paper's Hessian-bounded setting through the existing deterministic geometry results.

\section{\texorpdfstring{The actual continuous Gibbs target and explicit exponential mixing}{The actual continuous Gibbs target and explicit exponential mixing}}
\label{proof-module-KineticLangevin-ContinuousGibbs-lean}
{\small\color{gray}Source module: \texttt{Kinetic\allowbreak{}Langevin/\allowbreak{}Continuous\allowbreak{}Gibbs.\allowbreak{}lean}.\par}


The actual canonical Brownian diffusion starts on an independent position–velocity–noise product. This module proves its exact Gibbs stationarity on that same source, rather than assuming a reference law. Combining it with the arbitrary-shift continuous coupling theorem gives unconditional transport to the true Gibbs target. The simple friction regime yields the explicit factor \(\sqrt{2}e^{-mt/\gamma}\).

\subsection{\texorpdfstring{Reassociate the genuine source and identify each thermal increment}{Reassociate the genuine source and identify each thermal increment}}


The measurable source map sends \((x,(v,B))\) to \(((v,x),B)\). Product associativity and exchange of independent position and velocity prove that this map preserves the exact phase-times-Brownian law. The continuous linear source pullbacks commute with the OU input integral, and the proved Brownian convolution formula identifies each actual filtered increment with the translated exponential kernel on the independent Brownian coordinate.

\subsection{\texorpdfstring{Prove the complete actual initial/\allowbreak{}noise law}{Prove the complete actual initial/noise law}}


Disjoint Brownian intervals give an actual iid stream of thermal Gaussian innovations. On the phase-times-Brownian product this stream is independent of the initial phase. Pulling that law back along the source reassociation and using the L\ensuremath{{}^{2}} representative identities proves the full initial-phase/\allowbreak{}entire-stream product law on the actual position–velocity–noise source.

\subsection{\texorpdfstring{Identify the exact stationary position law}{Identify the exact stationary position law}}


The selected continuous solution satisfies the genuine kinetic integral equation by construction. The actual initial/\allowbreak{}noise product law discharges every source premise of the proved splitting-limit stationarity theorem. This first proves the full actual phase has exactly the independent Gaussian-times-Gibbs target law. Projecting that phase law then gives the Gibbs position distribution at the terminal time. Strong convexity supplies the finite target second moment used in the construction.

\subsection{\texorpdfstring{Derive the full and simple explicit mixing bounds}{Derive the full and simple explicit mixing bounds}}


In the arbitrary-shift two-law theorem, start the reference position from the Gibbs distribution and use its newly proved stationarity. This produces the full explicit shifted coefficient with the true initial Wasserstein distance. For shift zero and \(m+M\le\gamma^2\), the maximum in the exponent equals \(-m\); the prefactor simplifies to \(\sqrt2\). Finite L\ensuremath{{}^{2}} output and Gibbs moments also prove the actual transport distance is finite, ensuring the real-valued endpoint never conceals an infinite distance.

\section{\texorpdfstring{The actual canonical KLMC coupling on each initial transport plan}{The actual canonical KLMC coupling on each initial transport plan}}
\label{proof-module-KineticLangevin-CanonicalCoupledSegments-lean}
{\small\color{gray}Source module: \texttt{Kinetic\allowbreak{}Langevin/\allowbreak{}Canonical\allowbreak{}Coupled\allowbreak{}Segments.\allowbreak{}lean}.\par}


Fix any finite-cost coupling of the initial position law and the target position law. The probability space is its explicit product with an independent standard Gaussian velocity and the canonical vector Brownian source. This module constructs the numerical chain, the exact reference segments and the auxiliary exact segments on that same space, and proves every analytic segment requirement used by the quantitative KLMC theorem.

\subsection{\texorpdfstring{Specify the concrete probability space and source projection}{Specify the concrete probability space and source projection}}


The phase variable is velocity first. The sample space contains the two coupled initial positions, one shared Gaussian velocity and one vector Brownian source. The measure is the actual initial-coupling product extension; probability follows from its first marginal and from the Gaussian/\allowbreak{}Brownian probability factors. The Brownian-source projection preserves its measure. Pullback along that projection gives a continuous linear map on L\ensuremath{{}^{2}}, and its representative theorem identifies it almost everywhere with evaluation on the actual product sample.

\subsection{\texorpdfstring{Construct the common Brownian input and filtered innovation}{Construct the common Brownian input and filtered innovation}}


The scaled source lift includes the physical amplitude \(\sqrt{2\gamma}\). At iteration \(n\), the input is the actual vector Brownian increment on \((nh,nh+t]\), and the filtered input is the exponential response to that same increment. The unscaled innovation pair consists of the exact velocity and position Wiener responses at the step endpoint. Both input paths are continuous in L\ensuremath{{}^{2}}, and the additive input starts at zero. These follow from the source lift's continuity and linearity and the proved vector-response properties.

\subsection{\texorpdfstring{Identify the exact noise equations and implemented numerical step}{Identify the exact noise equations and implemented numerical step}}


The constructed filtered input satisfies \(B_n(t)=W_n(t)-\gamma\int_0^tB_n(s)ds\). Its endpoint is the scaled velocity innovation, and its actual time primitive is the scaled positional innovation. Lean transports the vector Brownian equations through the continuous linear source lift and commutes that lift with Bochner integration. Induction on the actual recursions then proves that the exponential-input chain is exactly \texttt{klmc\allowbreak{}Chain}, with the canonical correlated innovation pair and the prescribed fluctuation-dissipation scaling.

\subsection{\texorpdfstring{Connect Brownian L\ensuremath{{}^{2}} innovations with the generic output-\allowbreak{}law construction}{Connect Brownian L² innovations with the generic output-law construction}}


The L\ensuremath{{}^{2}} innovation pair has the same representatives as the pointwise canonical Brownian innovation. The proof pulls back the previously established vector-response representative theorem to the coupling product and combines the two component identities. Comparing with the generic product-law layer's representatives gives equality of the actual L\ensuremath{{}^{2}} pairs. Thus the numerical chain analyzed here has the coupling-independent output law already proved for the implemented KLMC chain.

\subsection{\texorpdfstring{Build the actual numerical and stationary-\allowbreak{}reference initial states}{Build the actual numerical and stationary-reference initial states}}


The two initial phases share exactly one Gaussian velocity variable. The numerical phase uses the first coupled position; the reference phase uses the second. The reference representative is proved to be the explicit pair \((V,X_\mathrm{ref})\). Product independence and the two marginal laws prove that its actual law is \(\mathcal N(0,I_d)\otimes\pi\). The finite target second moment and finite coupling cost supply the required L\ensuremath{{}^{2}} initial positions.

\subsection{\texorpdfstring{Construct all three chains and segment families}{Construct all three chains and segment families}}


The numerical chain applies the actual lifted original gradient and the constructed correlated innovations. The reference chain stitches exact selected kinetic segments, each driven by the actual canonical Brownian increment for that iteration. The corresponding reference segments start from these actual reference states. The auxiliary exact segment starts instead from the actual numerical state and receives the identical additive input. The drift is the genuine pointwise gradient lifted to L\ensuremath{{}^{2}}; no substitute potential on the L\ensuremath{{}^{2}} space is assumed.

\subsection{\texorpdfstring{Prove every analytic premise of the KLMC coupling theorem}{Prove every analytic premise of the KLMC coupling theorem}}


The bundled specification supplies continuity of auxiliary paths, reference paths and filtered inputs; exact initial states; reference endpoint matching; the implemented exponential update; both actual common-input integral equations; the exponential mild identity; the kinematic primitive identity; and equal initial velocities. Each field follows from the concrete definitions and earlier proved equations. The mild and kinematic fields are derived internally. The only remaining probabilistic input for a final stationary-error theorem is the stationary law of the actual global reference diffusion, proved separately and connected to these stitched segments by uniqueness.

\section{\texorpdfstring{One canonical global kinetic reference behind all coupling segments}{One canonical global kinetic reference behind all coupling segments}}
\label{proof-module-KineticLangevin-CanonicalGlobalSegments-lean}
{\small\color{gray}Source module: \texttt{Kinetic\allowbreak{}Langevin/\allowbreak{}Canonical\allowbreak{}Global\allowbreak{}Segments.\allowbreak{}lean}.\par}


The exact reference segments are constructed on one canonical Brownian source. This module identifies their increment input with the increments of one full L\ensuremath{{}^{2}} Brownian input and proves that every stitched reference state and intermediate segment point belongs to the actual full kinetic solution.

\subsection{\texorpdfstring{Construct the full canonical input and identify its increments}{Construct the full canonical input and identify its increments}}


The full input is the scalar canonical Brownian value assembled spatially and pulled onto the actual initial-coupling extension with amplitude \(\sqrt{2\gamma}\). It is continuous in L\ensuremath{{}^{2}} because each assembly and pullback is continuous linear. The Brownian value at zero vanishes: Lean verifies that its deterministic time indicator has empty support, then uses linearity of the isonormal map. For nonnegative local time, the constructed segment input equals the full input difference at \(nh+t\) and \(nh\). The proof checks the nonnegative real-to-NNReal casts and uses linearity of spatial assembly and source pullback.

\subsection{\texorpdfstring{Construct the actual full reference solution}{Construct the actual full reference solution}}


The global reference is the selected kinetic solution with the genuine lifted gradient, the canonical reference initial phase and the full canonical Brownian input. Its specification proves continuity, its exact starting phase, and the actual kinetic integral equation over the whole finite horizon. There is no stationary-law requirement in its definition.

\subsection{\texorpdfstring{Identify every stitched reference point by uniqueness}{Identify every stitched reference point by uniqueness}}


The generic stitching theorem now applies because the local inputs are proved to be actual increments of the full input. Consequently \(R_n=Y_{nh}\), and \(Z_n(t)=Y_{nh+t}\) throughout each step whose endpoint lies in the global horizon. Lean instantiates the existing exact uniqueness results with the actual L\ensuremath{{}^{2}} gradient and checks the horizon arithmetic. Thus a stationary law for the full constructed kinetic solution supplies the law of all reference states used by the quantitative KLMC proof.

\section{\texorpdfstring{The KLMC law bound with the genuine initial Wasserstein distance}{The KLMC law bound with the genuine initial Wasserstein distance}}
\label{proof-module-KineticLangevin-CanonicalKLMCBound-lean}
{\small\color{gray}Source module: \texttt{Kinetic\allowbreak{}Langevin/\allowbreak{}Canonical\allowbreak{}KLMCBound.\allowbreak{}lean}.\par}


The numerical law is the pushforward of the implemented KLMC chain on the actual independent Gaussian-velocity/\allowbreak{}Brownian product source. Every initial transport coupling yields this same output law. The complete constructed coupling proof gives the paper's first-order bound, and taking the actual transport infimum replaces the supplied coupling cost by the true initial Wasserstein distance. This intermediate module makes the sole stationary global-reference law property explicit; the final Gibbs theorem discharges it from the actual Brownian source.

\subsection{\texorpdfstring{Fix the actual numerical output law}{Fix the actual numerical output law}}


The definition uses the original pointwise KLMC method, the standard Gaussian initial velocity and the exact canonical correlated Brownian innovations. The numerical L\ensuremath{{}^{2}} chain is identified with the generic product-law construction by equality of the actual innovation variables. The established output-law theorem then proves that its pushforward is independent of which initial position coupling was chosen.

\subsection{\texorpdfstring{Transfer one actual global stationary phase law to the reference segments}{Transfer one actual global stationary phase law to the reference segments}}


The named stationary property concerns the full constructed kinetic flow: at every time its phase law is \(\mathcal N(0,I_d)\otimes\pi\). It contains no numerical convergence conclusion. Exact chain-to-global identification transfers its position marginal to every reference grid point. Exact segment-to-global identification transfers its Gaussian velocity marginal to every intermediate time in a reference segment. Lean obtains these marginals by composition with the product-measure projections.

\subsection{\texorpdfstring{Apply the proved coupling estimate to one actual initial plan}{Apply the proved coupling estimate to one actual initial plan}}


The initial position L\ensuremath{{}^{2}} discrepancy equals the true cost of the supplied transport coupling; this follows from the concrete extension's representative and norm identities. The segment specification furnishes all analytic local-error premises internally. Together with the derived stationary position and Gaussian velocity marginals, the quantitative L\ensuremath{{}^{2}} KLMC theorem gives \(W_2(\nu_n,\pi)\le\sqrt2(1-3mh/(4\gamma))^n\operatorname{cost}(c)+Mh\sqrt{2d}/m\). Output-law equality and initial-cost equality rewrite this into a bound for the actual numerical law.

\subsection{\texorpdfstring{Prove output finiteness and take the transport infimum}{Prove output finiteness and take the transport infimum}}


Finite initial Wasserstein distance supplies an actual finite-cost initial coupling. Its numerical and reference states are genuine L\ensuremath{{}^{2}} variables, so their distance gives a finite transport bound on the output law. This proves output finiteness before real-valued transport arithmetic is used. The scalar recurrence estimate proves the contraction factor is nonnegative under the paper's step restrictions. Lean then applies the proved infimum theorem over all finite-cost couplings, obtaining the same estimate with \(W_2(\mu,\pi)\) in place of the cost of a chosen plan. No existence of an optimal coupling is assumed.

\section{\texorpdfstring{A concrete independent initial phase and canonical Brownian source}{A concrete independent initial phase and canonical Brownian source}}
\label{proof-module-KineticLangevin-CanonicalProduct-lean}
{\small\color{gray}Source module: \texttt{Kinetic\allowbreak{}Langevin/\allowbreak{}Canonical\allowbreak{}Product.\allowbreak{}lean}.\par}


The stationarity proof takes place on one explicit probability space: the product of an initial phase law and the constructed vector Brownian source. This module builds the actual initial L\ensuremath{{}^{2}} variables, the actual continuous Brownian forcing and its OU-filtered innovations, then proves the complete initial/\allowbreak{}noise joint law. The resulting splitting law identity has no assumed independence or Gaussian marginal fields.

\subsection{\texorpdfstring{Represent the original initial phase on the product source}{Represent the original initial phase on the product source}}


\texttt{H} is the finite-dimensional Euclidean state, and \texttt{Sample} pairs a phase with a vector Brownian sample. The probability measure is exactly \(\pi\otimes P_B\). Pull the velocity and position projections through the first product projection. Their square integrability follows from that of the original phase law. The two \texttt{to\allowbreak{}Lp} representative identities prove that the resulting phase agrees almost everywhere with the original initial coordinate.

\subsection{\texorpdfstring{Lift the actual continuous Brownian input}{Lift the actual continuous Brownian input}}


The second product projection preserves the Brownian source law. Its pullback is a linear isometry on L\ensuremath{{}^{2}}, hence a continuous linear map. Apply it to \(\sqrt{2\gamma}B_t\), where the constructed Brownian input is continuous in L\ensuremath{{}^{2}} and zero at time zero. Linearity and continuity give the same properties after adjoining the independent initial phase. Thus the forcing used by the deterministic integral solver is an actual Brownian forcing on this very source.

\subsection{\texorpdfstring{Use the exact OU-\allowbreak{}filtered Brownian intervals}{Use the exact OU-filtered Brownian intervals}}


The innovation in interval \([nh,(n+1)h]\) is the pulled-back isonormal convolution
\[\sqrt{2\gamma}\int_{nh}^{(n+1)h}e^{-\gamma((n+1)h-s)}dB_s\]
. Its representative is exactly the original source variable composed with the Brownian projection. Continuous linear maps commute with the deterministic OU filter. Combined with the proved canonical integration-by-parts identity, this shows the continuous-input filter is exactly this innovation, not merely a random variable with a matching covariance.

\subsection{\texorpdfstring{Derive the complete independent Gaussian stream law}{Derive the complete independent Gaussian stream law}}


The canonical OU marginal theorem gives covariance \((1-e^{-2\gamma h})I\); the disjoint-interval theorem gives independence of the full vector innovations across time. Their whole stream therefore has the infinite product of the thermal Gaussian law. The explicit initial-law/\allowbreak{}source product makes the initial phase independent of that entire stream. Finally use the almost-everywhere initial and innovation identities to obtain the complete joint law of the actual L\ensuremath{{}^{2}} variables.

\subsection{\texorpdfstring{Identify both actual splitting endpoint laws}{Identify both actual splitting endpoint laws}}


Apply the proved full-input-law theorem to the actual lifted Verlet–OU chain. Its phase pushforward is exactly the Markov iterate with fresh thermal Gaussian OU innovations and initial law \(\pi\). The exact OU-filter identity also identifies the continuous-input chain with the same noise-driven chain, so its endpoint has that law as well. These identities provide the missing connection between the pathwise convergence theorem and the measure-level weak stationarity calculation.

\section{\texorpdfstring{The actual initial phase and Brownian innovation joint law}{The actual initial phase and Brownian innovation joint law}}
\label{proof-module-KineticLangevin-CanonicalReferenceNoise-lean}
{\small\color{gray}Source module: \texttt{Kinetic\allowbreak{}Langevin/\allowbreak{}Canonical\allowbreak{}Reference\allowbreak{}Noise.\allowbreak{}lean}.\par}


The Gibbs stationarity proof needs the real joint law of the initial phase and all future filtered Brownian increments. This module derives that law on each arbitrary initial-position coupling extension. It combines the explicit product-source projection, the exact integration-by-parts response and the established independent thermal Gaussian stream.

\subsection{\texorpdfstring{Project the coupling source onto an independent phase/\allowbreak{}Brownian product}{Project the coupling source onto an independent phase/Brownian product}}


The projection sends the full sample to \(((V,X_\mathrm{ref}),\omega_\mathrm{Brownian})\). Its law is \((\mathcal N(0,I_d)\otimes\pi)\otimes P_\mathrm{Brownian}\). Lean first projects the initial coupling onto its second marginal, retaining the independent Gaussian/\allowbreak{}Brownian source. It then applies measure-preserving product associativity and swaps the position and velocity factors. Thus independence of the complete reference initial phase and the complete Brownian source is proved for the actual product construction.

\subsection{\texorpdfstring{Identify the actual full-\allowbreak{}input OU increment}{Identify the actual full-input OU increment}}


The deterministic integration-by-parts input used in the kinetic splitting proof is \(\Delta W-\gamma\int_0^he^{-\gamma(h-s)}\Delta W_sds\). The actual local-input increment identity replaces each difference of the full canonical input by its constructed segment input. Applying the vector exponential-response formula through the continuous linear source lift proves that this OU increment is exactly the filtered velocity response. Its representative is consequently the scaled canonical vector response on the original Brownian source, pulled onto the coupling extension.

\subsection{\texorpdfstring{Derive the full independent thermal Gaussian joint law}{Derive the full independent thermal Gaussian joint law}}


The vector OU stream already has the exact iid Gaussian law with scale \(\sqrt{1-\psi_0(h)^2}\). The product-source joint-law theorem combines this stream with the initial phase target. Composition with the proved phase/\allowbreak{}Brownian projection law transfers that product law to the actual coupling extension. Finally, the representative identities identify its initial phase with the constructed reference initial L\ensuremath{{}^{2}} phase and its stream with the actual full-input OU increments. The resulting joint law is precisely the source obligation of the proved kinetic Gibbs stationarity transfer theorem.

\section{\texorpdfstring{The complete first-\allowbreak{}order KLMC theorem for the actual sampler}{The complete first-order KLMC theorem for the actual sampler}}
\label{proof-module-KineticLangevin-CanonicalKLMC-lean}
{\small\color{gray}Source module: \texttt{Kinetic\allowbreak{}Langevin/\allowbreak{}Canonical\allowbreak{}KLMC.\allowbreak{}lean}.\par}


This final theorem concerns the implemented KLMC chain on the explicit independent Gaussian-velocity and canonical Brownian source. It proves the first-order Wasserstein bound of Dalalyan and Riou-Durand with the true initial Wasserstein distance. All continuous reference laws, common-noise identities, local-error estimates, output-law identities and transport-infimum steps are discharged inside the proof chain.

\subsection{\texorpdfstring{Prove Gibbs stationarity of the actual global coupling reference}{Prove Gibbs stationarity of the actual global coupling reference}}


The reference begins with standard Gaussian velocity independent of the target Gibbs position. Its full Brownian forcing and its entire filtered innovation stream are actual constructed objects on the initial-coupling extension. Their exact initial/\allowbreak{}noise product law is already proved. The kinetic stationarity transfer theorem therefore applies directly to the actual integral solution. For a time \(t\) inside a larger horizon, Lean restricts the actual integral equation to \([0,t]\), supplies the continuous input and its zero starting value, and applies the proved splitting-limit stationarity result. This yields the full phase law \(\mathcal N(0,I_d)\otimes\pi\) at every reference time, without a remaining stationarity or Gaussian-marginal hypothesis.

\subsection{\texorpdfstring{Conclude the source-\allowbreak{}paper KLMC Wasserstein estimate}{Conclude the source-paper KLMC Wasserstein estimate}}


Let \(f\) be differentiable, \(m\)-strongly convex with \(m>0\), and have an \(M\)-Lipschitz gradient. If \(\gamma^2\ge m+M\), \(\gamma>0\), and \(0<h\le m/(4\gamma M)\), the actual KLMC position law satisfies \[W_2(\nu_n,\pi)\le\sqrt2\left(1-\frac{3mh}{4\gamma}\right)^nW_2(\mu,\pi)+\frac{Mh\sqrt{2d}}m.\] The initial position distribution may be arbitrary with finite initial Wasserstein distance; its velocity is the explicit independent standard Gaussian. Lean first derives the finite Gibbs second moment from strong convexity. It then supplies the just-proved global-reference stationarity for every initial transport coupling to the complete canonical coupling/\allowbreak{}infimum theorem. The conclusion has no conditional local-error, reference-process, Gaussian-law, stationary-law or optimal-coupling premise. The smoothness assumptions are expressed by the actual gradient and strong convexity, so the theorem also covers the paper's Hessian-bounded setting through the existing deterministic geometry results.

\section{\texorpdfstring{Construct the actual averaged Hessian on a segment}{Construct the actual averaged Hessian on a segment}}
\label{proof-module-KineticLangevin-Hessian-lean}
{\small\color{gray}Source module: \texttt{Kinetic\allowbreak{}Langevin/\allowbreak{}Hessian.\allowbreak{}lean}.\par}


The mean Hessian is a Bochner integral of the continuous derivative of the force along the segment between two positions. The module derives its self-adjointness, lower and upper quadratic-form bounds and exact representation of the force difference. Thus the operator used in the spectral contraction proof is a concrete object, not an assumed matrix already chosen to make the desired equation hold.

\subsection{\texorpdfstring{Define the segment average and commute it with a pairing}{Define the segment average and commute it with a pairing}}


Set \(\bar H(x,y)=\int_0^1H(y+t(x-y))dt\). Continuity of \(H\) makes the operator-valued integrand continuous and integrable on the compact segment. A continuous linear map commutes with a Bochner integral, so applying the operator to \(w\) and pairing with \(z\) gives \(\langle z,\bar H(x,y)w\rangle=\int_0^1\langle z,H(y+t(x-y))w\rangle dt\). Lean first commutes operator evaluation through the interval integral, then commutes the continuous linear inner-product functional. The operator type is complete because the underlying Hilbert space is complete.

\subsection{\texorpdfstring{Average symmetry and both spectral inequalities}{Average symmetry and both spectral inequalities}}


If every \(H(p)\) is self-adjoint, the pairing identity moves both sides into scalar integrals, where their integrands agree pointwise. If \(m\|z\|^2\le\langle z,H(p)z\rangle\le M\|z\|^2\), integral monotonicity bounds the average between the integrals of these two constants. The interval has length one, leaving exactly the original constants. Lean supplies continuity and interval-integrability for each quadratic-form integrand and the constant comparison function, applies \texttt{integral\_\allowbreak{}mono\_\allowbreak{}on}, and simplifies the constant integral. These results provide the genuine symmetry and spectral premises needed by the transformed-energy proof.

\subsection{\texorpdfstring{The fundamental theorem identifies the actual force difference}{The fundamental theorem identifies the actual force difference}}


Assume \(Dg(p)=H(p)\). The affine path \(\ell(t)=y+t(x-y)\) has derivative \(x-y\), so the chain rule gives \((g\circ\ell)'(t)=H(\ell(t))(x-y)\). Integrate this derivative from zero to one. The fundamental theorem yields \(\int_0^1H(\ell(t))(x-y)dt=g(x)-g(y)\). Commuting the operator evaluation with the average integral gives \(\bar H(x,y)(x-y)=g(x)-g(y)\). Lean constructs the affine-path derivative, composes the force derivative at each segment point, verifies integrability from continuity, and simplifies the endpoints. No coordinatewise mean-value theorem or unjustified single Hessian point is used.

\section{\texorpdfstring{The genuine quadratic Taylor remainder}{The genuine quadratic Taylor remainder}}
\label{proof-module-KineticLangevin-HessianRemainder-lean}
{\small\color{gray}Source module: \texttt{Kinetic\allowbreak{}Langevin/\allowbreak{}Hessian\allowbreak{}Remainder.\allowbreak{}lean}.\par}


The second-order integrator linearizes the force at its initial position. Its error requires a quadratic estimate for this actual linearization, obtained from the explicitly integrated Hessian along a segment. The derivative and Hessian Lipschitz assumptions are ordinary analytic hypotheses on the force.

\subsection{\texorpdfstring{Integrate Hessian variation along the segment}{Integrate Hessian variation along the segment}}


The averaged Hessian is the Bochner integral of H(y+t(x\ensuremath{-}y)), with t between zero and one. Subtracting H(y) inside the integral gives an integrand with norm at most M\ensuremath{{}_{2}}t\ensuremath{\Vert}x\ensuremath{-}y\ensuremath{\Vert}. The norm of an integral is at most the integral of this nonnegative envelope, whose value is M\ensuremath{{}_{2}}\ensuremath{\Vert}x\ensuremath{-}y\ensuremath{\Vert}/\allowbreak{}2. Lean checks continuity of the integrand, moves subtraction through the integral, applies the Lipschitz inequality on the segment, and evaluates the elementary integral of t.

\subsection{\texorpdfstring{Apply the averaged operator to the displacement}{Apply the averaged operator to the displacement}}


The fundamental theorem of calculus identifies g(x)\ensuremath{-}g(y) with the averaged Hessian applied to x\ensuremath{-}y. Thus the true Taylor remainder equals (\ensuremath{\bar{H}}\ensuremath{-}H(y))(x\ensuremath{-}y). The operator norm inequality bounds this by \ensuremath{\Vert}\ensuremath{\bar{H}}\ensuremath{-}H(y)\ensuremath{\Vert}\ensuremath{\Vert}x\ensuremath{-}y\ensuremath{\Vert}, yielding M\ensuremath{{}_{2}}\ensuremath{\Vert}x\ensuremath{-}y\ensuremath{\Vert}\ensuremath{{}^{2}}/\allowbreak{}2. This is an actual remainder estimate; the Hessian is never replaced by an assumed scalar or a fictitious common mean-value point.

\section{\texorpdfstring{Bounded random Hessians act on L\ensuremath{{}^{2}}}{Bounded random Hessians act on L²}}
\label{proof-module-KineticLangevin-RandomOperatorLp-lean}
{\small\color{gray}Source module: \texttt{Kinetic\allowbreak{}Langevin/\allowbreak{}Random\allowbreak{}Operator\allowbreak{}Lp.\allowbreak{}lean}.\par}


A Hessian evaluated at a random initial position is itself random. To use it inside Bochner time integrals, we construct its genuine continuous linear action on square-integrable random vectors. The operator may depend on every coupled variable; independence is unnecessary.

\subsection{\texorpdfstring{Construct the action by Hölder multiplication}{Construct the action by Hölder multiplication}}


A measurable operator field A with a deterministic bound \ensuremath{\Vert}A(\ensuremath{\omega})\ensuremath{\Vert}\ensuremath{\le}M belongs to L\ensuremath{\infty}. Evaluation is a continuous bilinear map from operators and vectors to vectors. Mathlib’s Hölder construction therefore sends the actual L\ensuremath{\infty} operator and an actual L\ensuremath{{}^{2}} vector to an L\ensuremath{{}^{2}} vector, with 1/\allowbreak{}\ensuremath{\infty}+1/\allowbreak{}2=1/\allowbreak{}2. Fixing the operator gives the continuous linear map randomOperatorLp.

\subsection{\texorpdfstring{Identify its representatives}{Identify its representatives}}


The Holder construction supplies an almost-everywhere representative identity. The representative of the L\ensuremath{\infty} equivalence class is A(\ensuremath{\omega}) almost everywhere, so the constructed L\ensuremath{{}^{2}} output is exactly A(\ensuremath{\omega})X(\ensuremath{\omega}). This identity preserves the original probability-space coupling and permits pointwise Taylor and truncation inequalities to be applied to the actual output.

\subsection{\texorpdfstring{Retain the deterministic operator bound}{Retain the deterministic operator bound}}


Pointwise, \ensuremath{\Vert}A(\ensuremath{\omega})X(\ensuremath{\omega})\ensuremath{\Vert}\ensuremath{\le}\ensuremath{\Vert}A(\ensuremath{\omega})\ensuremath{\Vert}\ensuremath{\Vert}X(\ensuremath{\omega})\ensuremath{\Vert}\ensuremath{\le}M\ensuremath{\Vert}X(\ensuremath{\omega})\ensuremath{\Vert}. Integrating this norm inequality gives \ensuremath{\Vert}AX\ensuremath{\Vert}\ensuremath{{}_{2}}\ensuremath{\le}M\ensuremath{\Vert}X\ensuremath{\Vert}\ensuremath{{}_{2}}. Taking the supremum over vectors proves the operator norm of the constructed map is at most M. These estimates justify commuting the random Hessian through time integrals and bounding the integrated force correction.

\section{\texorpdfstring{Exact Gaussian fourth moments for the Hessian correction}{Exact Gaussian fourth moments for the Hessian correction}}
\label{proof-module-KineticLangevin-GaussianFourth-lean}
{\small\color{gray}Source module: \texttt{Kinetic\allowbreak{}Langevin/\allowbreak{}Gaussian\allowbreak{}Fourth.\allowbreak{}lean}.\par}


The Hessian-Lipschitz Taylor error is quadratic in the position increment. Its L\ensuremath{{}^{2}} bound therefore needs a fourth moment. The exact standard Gaussian value is \(p^2+2p\); retaining both terms is necessary in dimension one. These identities are proved for the actual Gaussian measures and transferred to any realization by its proved law.

\subsection{\texorpdfstring{Differentiate the actual scalar moment-\allowbreak{}generating function four times}{Differentiate the actual scalar moment-generating function four times}}


A centered Gaussian of variance \(v\) has moment-generating function \(e^{vt^2/2}\). Successive derivatives are \(vt e^{vt^2/2}\), \((v+v^2t^2)e^{vt^2/2}\), and \((3v^2t+v^3t^3)e^{vt^2/2}\). The fourth derivative at zero is \(3v^2\). Lean constructs each derivative from the power, multiplication and exponential rules, then invokes the proved theorem identifying the derivative of a Gaussian moment-generating function with its moment. This also includes variance zero.

\subsection{\texorpdfstring{Separate diagonal fourth moments from independent cross moments}{Separate diagonal fourth moments from independent cross moments}}


For \(G=(G_1,\ldots,G_p)\), expand \(\lVert G\rVert^4=\sum_{i,j}G_i^2G_j^2\). A diagonal term has expectation three by the scalar theorem. Distinct coordinates are independent under the constructed product Gaussian measure, so their squared product has expectation \(1\cdot1=1\). Finite integrability follows from fourth moments and \(G_i^2G_j^2\le G_i^4+G_j^4\). The sum is \(3p+p(p-1)=p^2+2p\). Lean first works on the actual product of scalar Gaussians, then uses its proved Euclidean standard-Gaussian pushforward identity. Empty coordinate sets are included.

\subsection{\texorpdfstring{Transfer the moment to the coupled stationary velocity}{Transfer the moment to the coupled stationary velocity}}


If an actual random vector has the standard Gaussian law, the pushforward integration identity gives the same norm fourth moment on its original probability space. Applying this to each stationary kinetic velocity supplies exactly the \(\sqrt{p^2+2p}\) factor in the second-order Taylor remainder. No independence between this velocity and the numerical position discrepancy is assumed.

\section{\texorpdfstring{The OBABO metric and a rigorous nonlinear average-\allowbreak{}force estimate}{The OBABO metric and a rigorous nonlinear average-force estimate}}
\label{proof-module-KineticLangevin-OBABOEnergy-lean}
{\small\color{gray}Source module: \texttt{Kinetic\allowbreak{}Langevin/\allowbreak{}OBABOEnergy.\allowbreak{}lean}.\par}


The OBABO source metric is \(E(v,x)=\|x\|^2+2\langle x,v\rangle/\gamma+\|v\|^2/L\). Its continuous drift is contractive when \(\gamma^2\ge4L\). OBABO evaluates the force at two different position pairs; the proof below controls their average by an explicit centered-sector error. This avoids selecting a single Hessian point in a vector-valued mean-value theorem. The higher-order discrete remainder still has to be bounded separately.

\subsection{\texorpdfstring{Prove that the mixed energy defines the correct geometry}{Prove that the mixed energy defines the correct geometry}}


Write \(E_{a,b}(v,x)=\|x\|^2+2b\langle x,v\rangle+a\|v\|^2\), and let \(E_{a,0}\) be its diagonal part. Cauchy–Schwarz bounds the mixed inner product in both directions. The nonnegative square \((\|x\|-2b\|v\|)^2\), together with \(4b^2\le a\), gives \(|2b\langle x,v\rangle|\le E_{a,0}/2\). Lean records the two scalar inner-product comparisons, multiplies them by nonnegative \(2b\), and combines them with the squared difference using \texttt{nlinarith}. Thus \(E_{a,0}/2\le E_{a,b}\le3E_{a,0}/2\), and energy is nonnegative. For \(a=1/L,b=1/\gamma\), the positive-denominator division rules show that the coefficient condition is exactly \(\gamma^2\ge4L\).

\subsection{\texorpdfstring{Derive the source drift with an explicit sector error}{Derive the source drift with an explicit sector error}}


Along \(x'=v,v'=-\gamma v-F\), the energy derivative has common-denominator numerator \(-2L\langle x,F\rangle+(2L-2\gamma^2)\|v\|^2-2\gamma\langle v,F\rangle\). Suppose the force obeys \(\langle x,F\rangle\ge m\|x\|^2-e\|x\|\) and \(\|F\|^2\le L\langle x,F\rangle+Le\|x\|+e^2\). The nonnegative square \(\|\gamma v+F\|^2\) controls the velocity-force mixed term. Insert these inequalities, use the metric's upper equivalence, and use \(\gamma^2\ge4L\), \(0\le m\le L\) to discard the remaining nonpositive velocity term. Clearing the positive denominator \(\gamma L\) proves

\[E'\le-\frac{2m}{3\gamma}E+\frac{2e\|x\|}{\gamma}+\frac{e^2}{\gamma L}.\]

The error is explicit and vanishes when the average force satisfies the exact sector at the original position difference.

\subsection{\texorpdfstring{Control the actual average of the two nonlinear force evaluations}{Control the actual average of the two nonlinear force evaluations}}


Let \(F,F'\) satisfy the strong-smooth centered bounds at differences \(x,x'\), and set \(S=(F+F')/2\). With \(\mu=(m+L)/2\) and \(R=(L-m)/2\), decompose \(S-\mu x\) into half of each centered force residual plus \(\mu(x'-x)/2\). The triangle inequality and \(\|x'\|\le\|x\|+\|x'-x\|\) yield \(\|S-\mu x\|\le R\|x\|+(L/2)\|x'-x\|\). Thus the sector error is precisely \(e=(L/2)\|x'-x\|\). A centered residual controls monotonicity by Cauchy–Schwarz. Moving its center from \(\mu x\) to \(Lx/2\) and squaring yields the approximate cocoercivity bound needed above. Lean proves both consequences from the explicit norm inequality.

\subsection{\texorpdfstring{Apply the sector proof to the original potential}{Apply the sector proof to the original potential}}


For two arbitrary position pairs, use the project's proved strong-smooth interpolation and centered-ball theorem to obtain the two force residual bounds directly from differentiability, strong convexity and Lipschitz gradient. Their average is exactly the average of the two gradient differences used by the two OBABO kicks. The preceding average-sector lemma supplies the explicit position-increment error, and the perturbed dissipativity theorem supplies the source drift bound. No Hessian regularity, invented secant matrix, or assumed contraction inequality appears in this nonlinear step. The remaining discrete proof must bound the displayed error and the actual integrator remainder using its step-size restrictions.

\section{\texorpdfstring{OBABO: controlling every discrete remainder in the weighted metric}{OBABO: controlling every discrete remainder in the weighted metric}}
\label{proof-module-KineticLangevin-OBABORemainder-lean}
{\small\color{gray}Source module: \texttt{Kinetic\allowbreak{}Langevin/\allowbreak{}OBABORemainder.\allowbreak{}lean}.\par}


This module proves the numerical heart of the OBABO synchronous contraction. Starting from the explicit two-kick map, it bounds the averaged-force sector error, the Euler drift, and both components of the discrete remainder. These estimates combine to give the source factor \(1-mh/(3\gamma)\) under \(33\gamma^3h\le m\). The actual gradient and exponential damping are instantiated in the next module.

\subsection{\texorpdfstring{Realize the modified energy as an ordinary Hilbert norm}{Realize the modified energy as an ordinary Hilbert norm}}


Complete the square: \(E_{a,b}(v,x)=\|x+bv\|^2+(a-b^2)\|v\|^2\). The map \((v,x)\mapsto(x+bv,\sqrt{a-b^2}v)\) lands in the Euclidean product \texttt{With\allowbreak{}Lp 2}. This explicit choice matters: the ordinary Lean product norm is a maximum norm. \texttt{weighted\allowbreak{}Phase\allowbreak{}Norm\_\allowbreak{}sq} expands the two component norms, uses the square-root identity for nonnegative \(a-b^2\), and normalizes the scalar algebra. The associated bilinear pairing is exactly the Euclidean inner product of these transforms, so its upper bound follows from Cauchy–Schwarz without an extra geometric assumption.

\subsection{\texorpdfstring{Expand the exact energy of an Euler step plus a residual}{Expand the exact energy of an Euler step plus a residual}}


Expand \(E(s+t)=E(s)+2B(s,t)+E(t)\) and \(E(ct)=c^2E(t)\) directly from norms and inner products. Cauchy–Schwarz and \((\|s\|_E-\|t\|_E)^2\ge0\) give \(E(s+t)\le2E(s)+2E(t)\). For \(s_{+}=s+hA+r\), retain the first-order pairing \(2hB(s,A)\), bound the initial-residual pairing by \(2\|s\|_E\|r\|_E\), and apply the two-term estimate only to \(hA+r\). The resulting explicit bound is \(E(s_+)\le E(s)+2hB(s,A)+2\|s\|_E\|r\|_E+2h^2E(A)+2E(r)\). For the kinetic drift \(A=(-\gamma v-F,v)\), Lean checks that \(2B(s,A)\) is precisely the continuous derivative formula in the preceding module.

\subsection{\texorpdfstring{Bound the drift and residual using their actual components}{Bound the drift and residual using their actual components}}


The averaged force obeys the explicit envelope \(\|F\|\le17L\|x\|/16+\gamma\|v\|/8\). Expanding the drift energy and using the nonnegative square \((\gamma\|x\|-2\|v\|)^2\) bounds its mixed term. The friction condition \(4L\le\gamma^2\) and metric equivalence then give \(E(A)\le13\gamma^2E(s)/4\). For the residual, assume only the directly verifiable component bounds \(\|r_x\|\le h^2(\gamma\|v\|+L\|x\|)/2\) and \(\|r_v\|\le h^2(9\gamma L\|x\|/16+5\gamma^2\|v\|/8)\). Squaring each envelope uses \((a-b)^2\ge0\); multiplying by positive coefficients preserves the inequalities. Metric equivalence and \(4L\le\gamma^2\) yield \(E(r)\le3h^4\gamma^4E(s)\). No norm error estimate is supplied as a convergence hypothesis.

\subsection{\texorpdfstring{Pay for the approximate centered sector}{Pay for the approximate centered sector}}


The actual change in position gives \(e\le hL(\|v\|/2+\gamma\|x\|/16)\). The continuous dissipativity therefore contributes the positive correction \(h(2e\|x\|/\gamma+e^2/(\gamma L))\). The proof bounds \(e^2\) with the square of the difference of its two summands, and bounds the position-velocity term using \((\gamma\|x\|-2\|v\|)^2\). The restrictions \(\gamma h\le1\), \(4L\le\gamma^2\), together with the lower metric equivalence, give the explicit budget \(\gamma^2h^2E(s)/2\). Denominators are cleared only with positive factors; the remaining scalar inequalities are checked by Lean arithmetic.

\subsection{\texorpdfstring{Recover the source step-\allowbreak{}size constant from the complete budget}{Recover the source step-size constant from the complete budget}}


The residual energy bound implies \(\|r\|_E\le7h^2\gamma^2\|s\|_E/4\), since \(\sqrt3\le7/4\). Thus the initial-residual cross term costs at most \(7\gamma^2h^2E(s)/2\). The drift term costs \(13\gamma^2h^2E(s)/2\), the sector correction costs \(\gamma^2h^2E(s)/2\), and \(12\gamma^2h^2\le1\) makes twice the residual energy cost at most another \(\gamma^2h^2E(s)/2\). Their sum is exactly \(11\gamma^2h^2E(s)\). Consequently the map contracts by \(1-2mh/(3\gamma)+11\gamma^2h^2\), which is at most \(1-mh/(3\gamma)\) when \(33\gamma^3h\le m\).



The final algebra theorem constructs the residual explicitly: \(r_x=h(\delta-1)v-h^2F_0/2\), \(r_v=(\delta^2-1+\gamma h)v+h(1-\delta)F\). Norm triangle inequalities and the damping envelopes prove the two component bounds internally. The identity between \(s+hA+r\) and the actual difference update is checked by \texttt{module}, so no update-error contract is assumed.

\section{\texorpdfstring{A genuine local state comparison with exact Hamiltonian dynamics}{A genuine local state comparison with exact Hamiltonian dynamics}}
\label{proof-module-KineticLangevin-VerletLocal-lean}
{\small\color{gray}Source module: \texttt{Kinetic\allowbreak{}Langevin/\allowbreak{}Verlet\allowbreak{}Local.\allowbreak{}lean}.\par}


To quantify sampling bias, compare the implemented O–Verlet–O step with a reference that uses exact Hamiltonian motion between the same two thermostat updates. Exact Hamiltonian motion and the thermostat preserve the Gibbs target separately. This module derives the actual local state error from the Hamiltonian integral equations. It works in any real Hilbert space, including L\ensuremath{{}^{2}}, where uniform velocity and force norm bounds are second-moment statements.

\subsection{\texorpdfstring{Integrate the Hamiltonian equations instead of assuming a local error}{Integrate the Hamiltonian equations instead of assuming a local error}}


The Hamiltonian equations are \(v(t)-v(0)=-\int_0^t g(x(s))ds\) and \(x(t)-x(0)=\int_0^t v(s)ds\). If the force norm is at most G, the integral norm inequality gives \(\|v(t)-v(0)\|\le Gt\). Subtract \(tv(0)\) from the position equation; the result is the integral of the velocity increment. Integrating the derived envelope \(Gs\) gives the actual Euler-position defect \(\|x(t)-x(0)-tv(0)\|\le Gt^2/2\). Lean proves the exact integral identities using integrability from continuity before applying either inequality.

\subsection{\texorpdfstring{Compare both Verlet components to the true endpoint}{Compare both Verlet components to the true endpoint}}


The implemented drift-stage position is \(x_V=x_0+hv_0-h^2g(x_0)/2\). Its difference from the Euler position has norm at most \(Gh^2/2\); adding the exact position defect proves \(\|x_V-x(h)\|\le Gh^2\). For velocity, subtract the frozen-force Euler value \(v_0-hg(x_0)\) from both trajectories. The exact difference is the integral of \(g(x(s))-g(x_0)\), bounded by \(MC h^2/2\) using Lipschitzness and the actual kinematics with velocity norm C. The Verlet difference is exactly \(-h(g(x_V)-g(x_0))/2\), whose norm is at most \(Mh^2C/2+MGh^3/4\). Combining the two gives \(\|v_V-v(h)\|\le MC h^2+MGh^3/4\). Every residual is computed from the update and the exact path; none is supplied as an opaque accuracy assumption.

\subsection{\texorpdfstring{Express the state defect in the contracting OBABO metric}{Express the state defect in the contracting OBABO metric}}


The upper metric equivalence and completing-square Hilbert representation imply the convenient component estimate \(\|(v,x)\|_E\le2(\|x\|+\|v\|/\sqrt M)\). The proof squares both nonnegative sides, expands the scalar square, and uses \(\sqrt M^2=M\). Inserting the two genuine component defects yields

\[\|\mathrm{Verlet}_h(v_0,x_0)-(v(h),x(h))\|_E\le2h^2\bigl(G+\sqrt M C+\sqrt M G h/4\bigr).\]

This estimate uses only a Lipschitz force and the explicit moment envelopes; no Hessian regularity is required.

\subsection{\texorpdfstring{Couple the implemented OBABO step with the exact target-\allowbreak{}preserving split}{Couple the implemented OBABO step with the exact target-preserving split}}


Start the Hamiltonian path at the actual first Ornstein–Uhlenbeck half-step. Apply the same final thermostat vector to its endpoint and to the Verlet endpoint. The final additive noise cancels exactly, while the velocity defect is multiplied by \(\delta\). Since \(|\delta|\le1\), its norm is no larger than the Verlet velocity defect, and the position defect is unchanged. Applying the component metric bound therefore gives the same explicit local estimate for the actual O–Verlet–O and O–exact-Hamiltonian–O outputs. In L\ensuremath{{}^{2}}, an invariant Gibbs Hamiltonian reference supplies \(C=\sqrt d\) and \(G\le\sqrt{Md}\) from the Gaussian velocity marginal and Gibbs score moment theorem. The local bound itself keeps the exact path equations and norm envelopes visible so the probabilistic reference construction must discharge them.

\section{\texorpdfstring{Subgaussian Gibbs score and its sharp second moment}{Subgaussian Gibbs score and its sharp second moment}}
\label{proof-module-UnadjustedLangevin-GibbsScore-lean}
{\small\color{gray}Source module: \texttt{Unadjusted\allowbreak{}Langevin/\allowbreak{}Gibbs\allowbreak{}Score.\allowbreak{}lean}.\par}


Differentiability, an \(L\)-Lipschitz gradient, and finite Gibbs normalization suffice for \(\int\|\nabla U\|^2d\pi\le Ld\). Strong convexity is used only by the final wrapper to derive normalization.

\subsection{\texorpdfstring{gibbs\allowbreak{}Weight mul exp score le}{gibbsWeight mul exp score le}}


The descent lemma gives \(U(x+v)\le U(x)+\langle\nabla U(x),v\rangle+(L/2)\|v\|^2\). Negate and exponentiate, then rearrange to \(e^{-U(x)}e^{-\langle\nabla U(x),v\rangle}\le e^{L\|v\|^2/2}e^{-U(x+v)}\). Lean first establishes the real inequality and then embeds it in the extended nonnegative numbers.

\subsection{\texorpdfstring{gibbs score lmgf le}{gibbs score lmgf le}}


Integrate the pointwise exponential-score bound. Translation invariance gives \(\int e^{-U(x+v)}dx=Z_U\); normalization cancels this factor, leaving \(\mathbb E_{\pi_U}e^{-\langle\nabla U,v\rangle}\le e^{L\|v\|^2/2}\). Finiteness of \(Z_U\) is explicit because extended inverse cancellation needs it.

\subsection{\texorpdfstring{gibbs score subgaussian}{gibbs score subgaussian}}


Apply the preceding bound with the vector \(-tv\), for every real \(t\). The exponent becomes \(t\langle\nabla U,v\rangle\) and the right side is \(e^{t^2L\|v\|^2/2}\). This is exactly \texttt{Has\allowbreak{}Subgaussian\allowbreak{}MGF} with variance proxy \(L\|v\|^2\); the proof also establishes exponential integrability required by that structure.

\subsection{\texorpdfstring{sq le exp add exp neg sub two}{sq le exp add exp neg sub two}}


The series \(\cosh a=\sum_{n\ge0}a^{2n}/(2n)!\) has nonnegative terms, so retaining its first two terms gives \(\cosh a\ge1+a^2/2\). Substitute \(2\cosh a=e^a+e^{-a}\) and rearrange to \(a^2\le e^a+e^{-a}-2\). Lean uses \texttt{Real.\allowbreak{}has\allowbreak{}Sum\_\allowbreak{}cosh} with \texttt{sum\_\allowbreak{}le\_\allowbreak{}has\allowbreak{}Sum}, explicitly proving nonnegativity of every discarded even-power term.

\subsection{\texorpdfstring{integral sq le of subgaussian}{integral sq le of subgaussian}}


Integrate \(t^2X^2\le e^{tX}+e^{-tX}-2\), giving \(t^2\mathbb EX^2\le2(e^{ct^2/2}-1)\). Subgaussianity provides the exponential integrability and hence the square integrability needed here. Divide by \(t^2\) for \(t\ne0\), and let \(t\to0\) to obtain \(\mathbb EX^2\le c\). Lean computes the scalar limit before applying closedness of order.

\subsection{\texorpdfstring{integral norm sq gradient gibbs le}{integral norm sq gradient gibbs le}}


Choose an orthonormal basis \((e_i)_{i=1}^d\) and use Parseval: \(\|\nabla U(x)\|^2=\sum_i\langle\nabla U(x),e_i\rangle^2\). Each coordinate is subgaussian with proxy \(L\), hence has second moment at most \(L\). Finite-sum integrability permits integration term by term, yielding \(\mathbb E_{\pi_U}\|\nabla U\|^2\le Ld\).

\subsection{\texorpdfstring{mem\allowbreak{}Lp gradient gibbs\allowbreak{}Measure}{memLp gradient gibbsMeasure}}


The coordinate subgaussian bounds prove integrability of every squared gradient coordinate. Summing over the finite orthonormal basis proves integrability of \(\|\nabla U\|^2\). Lipschitz continuity provides strong measurability, and the \texttt{Mem\allowbreak{}Lp} square criterion packages those facts as \(\nabla U\in L^2(\pi_U)\).

\subsection{\texorpdfstring{integral norm sq gradient gibbs le of strong\allowbreak{}Convex}{integral norm sq gradient gibbs le of strongConvex}}


Strong convexity with \(m>0\) supplies the finite-partition premise of the score theorem. Substitute that proof into the first-order smoothness result to obtain \(\mathbb E_{\pi_U}\|\nabla U\|^2\le Ld\). The strong convexity constant is used for normalization here; it does not appear in the sharp score bound.

\section{\texorpdfstring{Transfer the Gibbs score estimate to an L\ensuremath{{}^{2}} reference}{Transfer the Gibbs score estimate to an L² reference}}
\label{proof-module-UnadjustedLangevin-StationaryMoments-lean}
{\small\color{gray}Source module: \texttt{Unadjusted\allowbreak{}Langevin/\allowbreak{}Stationary\allowbreak{}Moments.\allowbreak{}lean}.\par}


A reference variable \(Y\) whose law is \(\pi\) inherits \(\|\nabla U(Y)\|_{L^2}\le\sqrt{Ld}\). The assertion that its law is Gibbs is a displayed premise here.

\subsection{\texorpdfstring{lp\allowbreak{}Drift norm le of gibbs law}{lpDrift norm le of gibbs law}}


The representative of \texttt{lp\allowbreak{}Drift} equals \(\nabla U(Y)\) almost everywhere. Change variables using \(P\circ Y^{-1}=\pi_U\), so its squared \(L^2(P)\) norm is \(\int\|\nabla U\|^2d\pi_U\le Ld\). Nonnegativity of the norm and of \(Ld\) then gives \(\|\nabla U(Y)\|_2\le\sqrt{Ld}\). The law identity, rather than independence, is the input to this transfer.

\subsection{\texorpdfstring{lp\allowbreak{}Drift norm le of strong\allowbreak{}Convex gibbs law}{lpDrift norm le of strongConvex gibbs law}}


Use strong convexity to prove \(Z_U<\infty\), and apply the previous law-transfer theorem with that proof. The conclusion is the same \(\sqrt{Ld}\) bound for the lifted drift of any \(L^2\) variable with Gibbs law. There is no separate assumed drift-moment bound hidden in this specialization.

\section{\texorpdfstring{The actual L\ensuremath{{}^{2}} Hamiltonian comparator has the exact Gibbs law}{The actual L² Hamiltonian comparator has the exact Gibbs law}}
\label{proof-module-KineticLangevin-LpHamiltonianLaw-lean}
{\small\color{gray}Source module: \texttt{Kinetic\allowbreak{}Langevin/\allowbreak{}Lp\allowbreak{}Hamiltonian\allowbreak{}Law.\allowbreak{}lean}.\par}


This module closes the probabilistic reference obligations for the Hamiltonian part of the OBABO bias argument. It constructs the actual L\ensuremath{{}^{2}} Picard flow, identifies its finite-dimensional output law, derives exact Gibbs stationarity and velocity/\allowbreak{}score moments, and applies the proved local O–Verlet–O comparison. No stationary reference-law or local-error bound is postulated.

\subsection{\texorpdfstring{Identify an L\ensuremath{{}^{2}} limit with a pointwise numerical limit}{Identify an L² limit with a pointwise numerical limit}}


Suppose actual L\ensuremath{{}^{2}} phase variables \(S_n\) converge to S, and their representatives agree almost everywhere with measurable vector approximants \(f_n\) converging almost everywhere to Y. For any continuous linear functional \(\ell\), the characteristic test \(e^{i\ell(z)}\) has norm one. Bounded dominated convergence identifies the characteristic-function limit of the pointwise approximants, while the proved L\ensuremath{{}^{2}} characteristic continuity identifies the limit of the L\ensuremath{{}^{2}} variables. Their equality yields \(\mathcal L(S)=\mathcal L(Y)\) through characteristic uniqueness. The next lemma proves the required almost-everywhere correspondence for every actual Verlet iterate by induction from the lifted kick and drift identities. Thus the two convergence arguments refer to precisely the same implemented approximants.

\subsection{\texorpdfstring{Identify the actual coupled representatives almost everywhere}{Identify the actual coupled representatives almost everywhere}}


Marginal equality alone is insufficient when the local numerical error must be coupled with the same initial state and thermostat inputs. This stronger result identifies the actual representatives almost everywhere. Each component of an L\ensuremath{{}^{2}}-convergent phase sequence converges in measure. The known almost-everywhere equality transfers this convergence to the same pointwise Verlet approximants. Their almost-everywhere convergence also implies convergence in measure, because the source is a probability measure. Uniqueness of the limit in measure therefore gives equality almost everywhere of each component, hence of the phase pair.



Apply this argument to the actual L\ensuremath{{}^{2}} Hamiltonian flow and its actual Verlet approximants. The result states that its representative at time t is almost everywhere the finite-dimensional Hamiltonian flow evaluated at the same representative of the initial state. Thus adjoining any initial state or thermostat variable preserves the coupling used in the local comparison; no equality of marginal laws is substituted for a joint-law argument.

\subsection{\texorpdfstring{Construct the true L\ensuremath{{}^{2}} Hamiltonian integral flow and identify its law}{Construct the true L² Hamiltonian integral flow and identify its law}}


The L\ensuremath{{}^{2}} force is the actual Nemytskii lift \texttt{lp\allowbreak{}Drift}; its global Lipschitz bound permits the existing Picard construction with zero friction and zero noise. Its specification gives a continuous global path, the prescribed initial value, and the actual Hamiltonian integral equation on the finite horizon. Verlet convergence holds in this Banach space. The same vector Verlet maps converge pointwise to the selected finite-dimensional Hamiltonian flow. The proved almost-everywhere pointwise identification and the composition rule for pushforwards give the exact law identity \(\mathcal L(A_t)=(H_t)_\#\mathcal L(A_0)\). This does not infer a pathwise identification from an unproved interchange of representatives and time integrals.

\subsection{\texorpdfstring{Derive stationarity and both moment envelopes from the actual target}{Derive stationarity and both moment envelopes from the actual target}}


If the initial phase law is the canonical independent Gaussian-times-Gibbs target, the exact Hamiltonian invariance theorem and the derived output-law identity prove that the actual L\ensuremath{{}^{2}} flow has that target law at every time. Projecting the actual integral equation with the continuous linear coordinate maps gives the velocity and kinematic component equations; commuting those maps with Bochner integrals is justified by continuity and interval integrability. Projecting the target law gives standard Gaussian velocity and normalized Gibbs position. The exact Gaussian coordinate calculation therefore yields \(\|V_t\|_{L^2}=\sqrt d\), and the already proved Gibbs score estimate gives \(\|\nabla f(X_t)\|_{L^2}\le\sqrt{Md}\). These are conclusions from the actual law, not supplied moment hypotheses.

\subsection{\texorpdfstring{Apply the local OBABO comparison to this constructed stationary reference}{Apply the local OBABO comparison to this constructed stationary reference}}


Start the genuine L\ensuremath{{}^{2}} Hamiltonian path at the actual first thermostat half-step and assume its phase law is the canonical target, as supplied by thermostat invariance. Its continuity, initial value, velocity equation, kinematic equation, and uniform moments have all been proved above. The local O–Verlet–O versus O–Hamiltonian–O theorem therefore applies with \(C=\sqrt d\) and \(G=\sqrt{Md}\). It yields the actual weighted error bound

\[2h^2\left(\sqrt{Md}+\sqrt M\sqrt d+\frac{h}{4}\sqrt M\sqrt{Md}\right).\]

The two actual thermostat noises are retained, their common final contribution cancels, and the damping needs only \(|\delta|\le1\). This provides the numerical local budget for the global sampling-bias estimate.

\section{\texorpdfstring{Fourth moments of actual time integrals}{Fourth moments of actual time integrals}}
\label{proof-module-KineticLangevin-LpMomentIntegral-lean}
{\small\color{gray}Source module: \texttt{Kinetic\allowbreak{}Langevin/\allowbreak{}Lp\allowbreak{}Moment\allowbreak{}Integral.\allowbreak{}lean}.\par}


Kinetic positions are time integrals of velocities in L\ensuremath{{}^{2}}. A fourth-moment estimate must apply to this existing L\ensuremath{{}^{2}} integral, rather than to a separately postulated L4 process. The proof uses closed convex L4 balls inside L\ensuremath{{}^{2}} and the geometric fact that averages remain in a closed convex set.

\subsection{\texorpdfstring{View an L4 ball inside L\ensuremath{{}^{2}}}{View an L4 ball inside L²}}


The set lpFourthBall(R) contains actual L\ensuremath{{}^{2}} equivalence classes whose extended fourth norm is at most R. If a sequence in this set converges in L\ensuremath{{}^{2}}, it converges in measure. Lower semicontinuity of the extended L4 norm under convergence in measure keeps its limit in the same ball. In the metric L\ensuremath{{}^{2}} space, this sequential argument proves closedness.

\subsection{\texorpdfstring{Verify convexity without extra moment assumptions}{Verify convexity without extra moment assumptions}}


For nonnegative a,b with a+b=1, the actual representative of aX+bY equals aX(\ensuremath{\omega})+bY(\ensuremath{\omega}) almost everywhere. The L4 triangle inequality and homogeneity give \ensuremath{\Vert}aX+bY\ensuremath{\Vert}\ensuremath{{}_{4}}\ensuremath{\le}a\ensuremath{\Vert}X\ensuremath{\Vert}\ensuremath{{}_{4}}+b\ensuremath{\Vert}Y\ensuremath{\Vert}\ensuremath{{}_{4}}\ensuremath{\le}R. Extended norms allow this argument even before finite fourth moments are separately established.

\subsection{\texorpdfstring{Average a continuous L\ensuremath{{}^{2}} path}{Average a continuous L² path}}


If every velocity v(t) on [0,T] has fourth norm at most R, its time average lies in the closed convex L4 ball. The average is T\ensuremath{{}^{-1}} times the actual Bochner integral in L\ensuremath{{}^{2}}. Homogeneity then gives \ensuremath{\Vert}\ensuremath{\int}\ensuremath{{}_{0}}\ensuremath{{}^{T}}v(t)dt\ensuremath{\Vert}\ensuremath{{}_{4}}\ensuremath{\le}TR. The zero-duration case is handled separately; for positive duration Lean checks the interval measure is finite and nonzero before applying the convex averaging theorem. This directly supplies displacement fourth moments from Gaussian velocity marginals.

\section{\texorpdfstring{Gaussian fourth moments control the actual Taylor error}{Gaussian fourth moments control the actual Taylor error}}
\label{proof-module-KineticLangevin-LpMomentBounds-lean}
{\small\color{gray}Source module: \texttt{Kinetic\allowbreak{}Langevin/\allowbreak{}Lp\allowbreak{}Moment\allowbreak{}Bounds.\allowbreak{}lean}.\par}


The quadratic Hessian remainder requires a fourth displacement moment, even though the stochastic solution itself lives in L\ensuremath{{}^{2}}. These results connect actual L4 seminorms, exact Gaussian fourth moments, and the existing L\ensuremath{{}^{2}} kinematic equation.

\subsection{\texorpdfstring{Turn a pointwise quadratic remainder into an L\ensuremath{{}^{2}} bound}{Turn a pointwise quadratic remainder into an L² bound}}


Suppose actual random vectors D and X satisfy \ensuremath{\Vert}D(\ensuremath{\omega})\ensuremath{\Vert}\ensuremath{\le}a\ensuremath{\Vert}X(\ensuremath{\omega})\ensuremath{\Vert}\ensuremath{{}^{2}}, with a\ensuremath{\ge}0 and \ensuremath{\Vert}X\ensuremath{\Vert}\ensuremath{{}_{4}}\ensuremath{\le}C. Monotonicity of the extended L\ensuremath{{}^{2}} norm gives \ensuremath{\Vert}D\ensuremath{\Vert}\ensuremath{{}_{2}}\ensuremath{\le}a\ensuremath{\Vert}\ensuremath{\Vert}X\ensuremath{\Vert}\ensuremath{{}^{2}}\ensuremath{\Vert}\ensuremath{{}_{2}}. The power rule identifies the last factor with \ensuremath{\Vert}X\ensuremath{\Vert}\ensuremath{{}_{4}}\ensuremath{{}^{2}}, yielding aC\ensuremath{{}^{2}}. Lean works first with extended nonnegative norms, where the power and monotonicity rules hold without hidden finiteness assumptions, and then converts the finite upper bound back to the actual L\ensuremath{{}^{2}} norm.

\subsection{\texorpdfstring{Compute the actual Gaussian L4 radius}{Compute the actual Gaussian L4 radius}}


A standard Gaussian vector in p dimensions has fourth norm moment p\ensuremath{{}^{2}}+2p. Fernique integrability ensures its fourth seminorm is finite. The integral formula for the L4 seminorm therefore yields (p\ensuremath{{}^{2}}+2p)\textasciicircum{}(1/\allowbreak{}4) on any realization of the Gaussian law. This proof transfers an actual integral identity through HasLaw and keeps the 2p correction, including the one-dimensional case.

\subsection{\texorpdfstring{Apply it to the existing kinetic displacement}{Apply it to the existing kinetic displacement}}


Assume X(t)\ensuremath{-}X(0)=\ensuremath{\int}\ensuremath{{}_{0}}\ensuremath{{}^{t}}V(s)ds in the actual L\ensuremath{{}^{2}} space and every V(s) has the standard Gaussian marginal. Each velocity lies in the same L4 ball of radius (p\ensuremath{{}^{2}}+2p)\textasciicircum{}(1/\allowbreak{}4). The closed-convex averaging theorem for the existing L\ensuremath{{}^{2}} integral gives \ensuremath{\Vert}X(t)\ensuremath{-}X(0)\ensuremath{\Vert}\ensuremath{{}_{4}}\ensuremath{\le}t(p\ensuremath{{}^{2}}+2p)\textasciicircum{}(1/\allowbreak{}4). No independence across times is used or required. This supplies exactly the fourth moment needed by the Hessian Taylor remainder in Proposition 2.

\section{\texorpdfstring{The actual frozen random Hessian and its Taylor remainder}{The actual frozen random Hessian and its Taylor remainder}}
\label{proof-module-KineticLangevin-LpHessian-lean}
{\small\color{gray}Source module: \texttt{Kinetic\allowbreak{}Langevin/\allowbreak{}Lp\allowbreak{}Hessian.\allowbreak{}lean}.\par}


The Hessian in KLMC2 is evaluated at the actual random starting position. It acts on the original coupled L\ensuremath{{}^{2}} variables through the previously constructed random-operator map. Combining its pointwise Taylor remainder with actual fourth displacement moments gives Proposition 2’s quadratic Taylor forcing.

\subsection{\texorpdfstring{Freeze the true Hessian at a random position}{Freeze the true Hessian at a random position}}


The field \ensuremath{\omega}\ensuremath{\mapsto}H(Y(\ensuremath{\omega})) is strongly measurable because H is Lipschitz and Y has a strongly measurable L\ensuremath{{}^{2}} representative. The force Lipschitz bound controls its derivative: \ensuremath{\Vert}H(y)\ensuremath{\Vert}\ensuremath{\le}M follows from HasFDerivAt g (H(y)) y and the global Lipschitz property of g. The actual random-operator construction then produces a continuous linear action on L\ensuremath{{}^{2}}. Its representative is exactly H(Y(\ensuremath{\omega}))X(\ensuremath{\omega}), and its operator norm is at most M.

\subsection{\texorpdfstring{Lift the genuine Taylor estimate}{Lift the genuine Taylor estimate}}


For actual X,Y, expand the representatives of gLp(X)\ensuremath{-}gLp(Y)\ensuremath{-}H(Y)(X\ensuremath{-}Y). Almost everywhere this is the ordinary pointwise Taylor remainder of g. The deterministic Hessian remainder bound is M\ensuremath{{}_{2}}\ensuremath{\Vert}X\ensuremath{-}Y\ensuremath{\Vert}\ensuremath{{}^{2}}/\allowbreak{}2. If the displacement has fourth norm at most C, the quadratic-to-L\ensuremath{{}^{2}} estimate gives a remainder norm at most M\ensuremath{{}_{2}}C\ensuremath{{}^{2}}/\allowbreak{}2. Every variable stays on the same probability space.

\subsection{\texorpdfstring{Obtain the exact dimension factor from stationary velocities}{Obtain the exact dimension factor from stationary velocities}}


The kinematic equation and actual Gaussian velocity marginals give C=t(p\ensuremath{{}^{2}}+2p)\textasciicircum{}(1/\allowbreak{}4). Squaring this radius yields t\ensuremath{{}^{2}}\ensuremath{\surd}(p\ensuremath{{}^{2}}+2p). Consequently the actual L\ensuremath{{}^{2}} Taylor forcing is at most M\ensuremath{{}_{2}}t\ensuremath{{}^{2}}\ensuremath{\surd}(p\ensuremath{{}^{2}}+2p)/\allowbreak{}2. Lean verifies the fractional-power identity using nonnegativity and the real power multiplication rule. No independence of position, Hessian, or velocity is assumed.

\section{\texorpdfstring{The actual random-\allowbreak{}Hessian L\ensuremath{{}^{2}} analysis represents the KLMC2 sampler}{The actual random-Hessian L² analysis represents the KLMC2 sampler}}
\label{proof-module-KineticLangevin-LpKLMC2-lean}
{\small\color{gray}Source module: \texttt{Kinetic\allowbreak{}Langevin/\allowbreak{}Lp\allowbreak{}KLMC2.\allowbreak{}lean}.\par}


The Hessian in the analysis is the actual bounded random operator obtained by evaluating the original Hessian at the random position. This module proves that each L\ensuremath{{}^{2}} update and every L\ensuremath{{}^{2}} iterate agree almost everywhere with the paper’s implemented four-noise vector updates. Thus later stochastic estimates concern the genuine sampler, including both Hessian-times-velocity and Hessian-times-noise corrections.

\subsection{\texorpdfstring{Read each affine Hessian correction pointwise}{Read each affine Hessian correction pointwise}}


Scalar multiplication, subtraction and addition of L\ensuremath{{}^{2}} classes agree with the corresponding operations on their representatives almost everywhere. Intersecting these finitely many identities proves the exact affine expression used by both KLMC2 coordinates has its original pointwise value.

\subsection{\texorpdfstring{Identify the actual step and full chain}{Identify the actual step and full chain}}


The lifted gradient has the original gradient as its representative. The proved random-operator construction similarly represents the actual Hessian evaluated at the random position, applied to velocity or to either Hessian correction noise. Substitution in the two affine coordinates identifies the complete L\ensuremath{{}^{2}} step with the actual vector KLMC2 formula. Induction then proves every numerical iterate agrees with the original chain driven by exactly the same four Brownian responses.

\subsection{\texorpdfstring{Connect to the exact Gaussian Markov chain law}{Connect to the exact Gaussian Markov chain law}}


Continuity of the original force and Hessian makes the full state-and-four-noise update measurable. Its deterministic recursion equals the general random recursion with the same update. Given a proved actual initial/\allowbreak{}entire-noise joint law, the generic Markov law theorem identifies every L\ensuremath{{}^{2}} phase output with the corresponding implemented Markov iterate. No replacement covariance or per-step independence assumption is introduced by the representative correspondence.

\section{\texorpdfstring{One actual KLMC2 output law for every initial coupling}{One actual KLMC2 output law for every initial coupling}}
\label{proof-module-KineticLangevin-KLMC2ProductLaw-lean}
{\small\color{gray}Source module: \texttt{Kinetic\allowbreak{}Langevin/\allowbreak{}KLMC2\allowbreak{}Product\allowbreak{}Law.\allowbreak{}lean}.\par}


The stochastic analysis may choose any initial transport coupling. The implemented sampler has one fixed law, defined independently of this choice. This module proves the two constructions coincide exactly, including the random Hessian and all four source noises. It supplies the law identity needed to pass a uniform coupling estimate to the true Wasserstein infimum.

\subsection{\texorpdfstring{Lift the actual four-\allowbreak{}noise source}{Lift the actual four-noise source}}


The initial coupling and its independent source extension are probability measures. Each square-integrable original noise coordinate is composed with the source projection and represented in L\ensuremath{{}^{2}}. Intersecting the four representative identities proves that these lifted variables retain precisely the original four-coordinate noise value.

\subsection{\texorpdfstring{Define the fixed-\allowbreak{}source actual sampler law}{Define the fixed-source actual sampler law}}


The pointwise original KLMC2 recursion starts at the common source velocity and the independent initial position. Its terminal position defines the output map and its product-source pushforward defines the actual law. The force and Hessian are continuous, so the complete four-noise step is measurable. Induction over the original recursion proves the output map is almost everywhere measurable from the square-integrable source data.

\subsection{\texorpdfstring{Identify the L\ensuremath{{}^{2}} analysis with that single law}{Identify the L² analysis with that single law}}


The pointwise correspondence theorem for the genuine random-Hessian L\ensuremath{{}^{2}} chain identifies every iterate with the original vector recursion. The initial velocity, initial coupling coordinate and all countably many four-noise values have exactly their source representatives. Substituting these identities proves the L\ensuremath{{}^{2}} output equals the original fixed-source output map composed with the initial/\allowbreak{}source projection. That projection has exactly the independent initial-position/\allowbreak{}source product law. Pushforward composition therefore proves the actual output law is independent of the chosen transport coupling. Its measurability also confirms it is a probability measure.

\section{\texorpdfstring{The actual four-\allowbreak{}kernel correlated Gaussian KLMC2 noise}{The actual four-kernel correlated Gaussian KLMC2 noise}}
\label{proof-module-KineticLangevin-CanonicalKLMC2Noise-lean}
{\small\color{gray}Source module: \texttt{Kinetic\allowbreak{}Langevin/\allowbreak{}Canonical\allowbreak{}KLMC2\allowbreak{}Noise.\allowbreak{}lean}.\par}


The second-order method uses four responses of the same Brownian interval: the two exponential-integrator kernels and the two Hessian correction kernels. Their correlations are essential. This module constructs those four responses simultaneously in every spatial coordinate, proves their exact Gram Gaussian law and iid time stream, and lifts the same random variables to arbitrary initial coupling extensions.

\subsection{\texorpdfstring{Pack the exact scalar covariance into spatial vectors}{Pack the exact scalar covariance into spatial vectors}}


The Euclidean space uses an arbitrary finite spatial coordinate set. Packing transposes the array of scalar four-dimensional blocks into four full spatial vectors. Every coordinate projection and Euclidean packing map is measurable. Pushing the product of the actual four-dimensional Gram Gaussians through this pack map defines the exact vector innovation law, which is a probability measure.

\subsection{\texorpdfstring{Construct the actual Brownian innovation stream}{Construct the actual Brownian innovation stream}}


Apply all four deterministic kernels to the same Brownian interval and pack all independent spatial copies. Brownian intervals have disjoint supports, so these full correlated vectors are independent across time. At each step their law is the packed product of the exact scalar Gram Gaussian laws. The independent-family product-law theorem therefore proves the complete implemented stream has the actual iid Gaussian distribution.

\subsection{\texorpdfstring{Identify the L\ensuremath{{}^{2}} variables with the original sampler noise}{Identify the L² variables with the original sampler noise}}


The four L\ensuremath{{}^{2}} innovations are directly the vector Brownian convolutions, so no separate Gaussian realization is substituted. Their representative theorem identifies every coordinate with the corresponding implemented packed response on the same source. Finite second moments follow from the L\ensuremath{{}^{2}} construction and this exact almost-everywhere identity.

\subsection{\texorpdfstring{Keep the actual noise on every initial coupling extension}{Keep the actual noise on every initial coupling extension}}


Pull the same Brownian convolutions onto the initial-position-coupling/\allowbreak{}Gaussian-velocity/\allowbreak{}Brownian product. The Brownian source projection preserves measure. The resulting representatives are exactly the original innovation values at that coordinate. A countable almost-everywhere intersection identifies the entire lifted stream with the original stream, and source projection preserves its exact iid law.

\section{\texorpdfstring{The genuine canonical KLMC2 Gaussian Markov law and coupling bridge}{The genuine canonical KLMC2 Gaussian Markov law and coupling bridge}}
\label{proof-module-KineticLangevin-CanonicalKLMC2Law-lean}
{\small\color{gray}Source module: \texttt{Kinetic\allowbreak{}Langevin/\allowbreak{}Canonical\allowbreak{}KLMC2\allowbreak{}Law.\allowbreak{}lean}.\par}


The second-order sampler is now a single actual Gaussian Markov law, not a law depending on an analysis coupling. Its four-noise stream is built from the genuine canonical Brownian kernels. The module proves the exact complete initial/\allowbreak{}noise product law, identifies the original pointwise output with the correlated-Gaussian Markov iterates, and identifies every actual random-Hessian L\ensuremath{{}^{2}} coupling construction with that same output.

\subsection{\texorpdfstring{Define the independent-\allowbreak{}source actual second-\allowbreak{}order law}{Define the independent-source actual second-order law}}


The source contains independent standard Gaussian velocity and a vector Brownian coordinate. The implemented four-kernel noise reads that Brownian coordinate. Its finite second moments follow from the actual convolution construction and measure-preserving source projection. The pointwise original KLMC2 chain on initial-position times source defines the canonical output law, which is a probability measure by the proved original-step measurability.

\subsection{\texorpdfstring{Prove the complete actual Gaussian Markov identification}{Prove the complete actual Gaussian Markov identification}}


On phase-times-Brownian source, the initial phase is independent of the entire actual iid four-kernel stream. Reassociation of independent initial position and standard Gaussian velocity pulls this exact law back to the sampler’s source. The resulting initial-phase/\allowbreak{}whole-stream product law determines the Markov iterates through the general random-recursion theorem. Projecting to position proves the independently defined canonical Brownian sampler law equals exactly the implemented correlated-Gaussian Markov output.

\subsection{\texorpdfstring{Identify every actual L\ensuremath{{}^{2}} coupling construction}{Identify every actual L² coupling construction}}


The directly lifted Brownian convolutions and the generic independent-source noise extension have the same original source representatives, hence are equal in L\ensuremath{{}^{2}} coordinate by coordinate. The actual numerical chain uses the original gradient lift, the actual bounded random Hessian and those same four convolutions. The general coupling-independent output theorem therefore gives its exact canonical law. The zero and successor identities expose the genuine initial state and recurrence needed by the stochastic convergence proof.

\subsection{\texorpdfstring{Establish finite transport and the exact initial cost}{Establish finite transport and the exact initial cost}}


A finite-cost initial coupling and the finite target second moment construct an actual L\ensuremath{{}^{2}} numerical chain. Its output-law equality transfers this finite second moment to the single canonical sampler law. Existence of a finite-cost initial plan therefore proves the true output Wasserstein distance is finite. Finally, the initial numerical/\allowbreak{}reference position discrepancy is exactly the root-mean-square cost of the supplied initial coupling. These facts permit the final convergence estimate to take the actual infimum over initial couplings without assuming an optimizer or hiding infinite distances through real conversion.

\section{\texorpdfstring{Actual past and future Brownian independence and its L\ensuremath{{}^{2}} closedness}{Actual past and future Brownian independence and its L² closedness}}
\label{proof-module-KineticLangevin-PastFutureIndependence-lean}
{\small\color{gray}Source module: \texttt{Kinetic\allowbreak{}Langevin/\allowbreak{}Past\allowbreak{}Future\allowbreak{}Independence.\allowbreak{}lean}.\par}


The future convolution has support after the terminal time of every past kernel. Their actual isonormal Hilbert inner products vanish. Joint Gaussianity therefore proves independence of the whole past block and the future response. Independent spatial copies retain this property. The final argument proves independence survives genuine L\ensuremath{{}^{2}} limits, which lets numerical approximants establish independence for the actual nonlinear reference solution.

\subsection{\texorpdfstring{Regroup independent arrays without changing their laws}{Regroup independent arrays without changing their laws}}


The exact independent array law is the corresponding infinite product. The measurable currying equivalence transports this law to the product of row laws. Pulling the independent product coordinate projections back along that actual law proves row independence. This handles the change from individual spatial coordinates to entire past and future spatial blocks.

\subsection{\texorpdfstring{Use disjoint deterministic time supports}{Use disjoint deterministic time supports}}


Place all past kernels in one Gaussian block and the future kernel in a second block. The joint family is Gaussian because every coordinate is the same actual isonormal linear map applied to its deterministic time kernel. A past interval ends before the future begins, so every cross covariance is exactly zero. Gaussian block independence follows. Independent coordinate source projections then assemble the full rectangular spatial array, and currying groups it into independent past and future vector blocks.

\subsection{\texorpdfstring{Recover the actual Euclidean convolution variables}{Recover the actual Euclidean convolution variables}}


Pack each spatial row with the actual Euclidean coordinate map and read the future response from any nonempty block coordinate. The spatial-lift representative theorem identifies these packed functions with the genuine vector L\ensuremath{{}^{2}} segment convolutions. A countable almost-everywhere intersection handles the entire past family. Independence is unchanged by these almost-everywhere representative identities.

\subsection{\texorpdfstring{Adjoin initial data and change probability sources}{Adjoin initial data and change probability sources}}


Measure-preserving source composition transports independence by exact pushforward composition. For a separate independent initial coordinate, product mapping first gives the joint initial/\allowbreak{}past/\allowbreak{}future law. Product associativity then places initial data and past together opposite the future coordinate. The two product projections are independent, proving that arbitrary measurable functions of the initial data and past remain independent of the future.

\subsection{\texorpdfstring{Pass independence through actual L\ensuremath{{}^{2}} limits}{Pass independence through actual L² limits}}


A continuous linear functional sends an L\ensuremath{{}^{2}} vector to an L\ensuremath{{}^{2}} scalar. Continuity of the actual scalar characteristic function therefore gives continuity of every dual characteristic function of the vector law. Independence is characterized by factorization of the joint dual characteristic function. For a convergent sequence of past L\ensuremath{{}^{2}} variables and a fixed future variable, both the joint characteristic function and the product of marginal characteristic functions converge. Their equality at every approximation passes to the limit, proving independence of the actual limiting past variable.

\section{\texorpdfstring{The true reference velocity is independent of its following Brownian interval}{The true reference velocity is independent of its following Brownian interval}}
\label{proof-module-KineticLangevin-CanonicalPastFuture-lean}
{\small\color{gray}Source module: \texttt{Kinetic\allowbreak{}Langevin/\allowbreak{}Canonical\allowbreak{}Past\allowbreak{}Future.\allowbreak{}lean}.\par}


This module proves the stochastic source condition needed in the repaired second-order KLMC argument. The actual reference trajectory is constructed from its initial phase and past Brownian input. Its terminal velocity is independent of the genuine future OU convolution. Stationarity then identifies their actual joint law as standard Gaussian velocity times the exact thermal Gaussian innovation.

\subsection{\texorpdfstring{Make dependence on a finite innovation prefix explicit}{Make dependence on a finite innovation prefix explicit}}


A recursion at index n reads only noise entries with index less than n. Induction proves this for the general random recursion and for the deterministic Verlet–OU splitting. These identities permit arbitrary continuation of a finite past noise grid without changing its terminal state.

\subsection{\texorpdfstring{Identify actual source representatives}{Identify actual source representatives}}


The coupling extension has independent Gaussian velocity and a canonical Brownian coordinate. The square-root fluctuation-dissipation factor is already built into its continuous linear noise lift. The representative theorem proves that each lifted segment is the original Brownian random vector, evaluated at that coordinate and multiplied by the exact factor.

\subsection{\texorpdfstring{Prove independence for genuine splitting approximants}{Prove independence for genuine splitting approximants}}


Divide the past horizon into n+1 intervals. Continue its innovation stream by repeating the last past interval, so every deterministic kernel stays entirely before the future interval. The proved Brownian support theorem makes this whole stream independent of the future response. Adjoin the independent initial positions and velocity, and reassociate the actual product source. The finite pointwise splitting recursion is a measurable function of precisely those past inputs. Its L\ensuremath{{}^{2}} representative is the original vector recursion. Prefix dependence identifies its endpoint with the actual continuous-input splitting chain, proving independence of each true approximant.

\subsection{\texorpdfstring{Pass to the actual nonlinear kinetic reference}{Pass to the actual nonlinear kinetic reference}}


The matching-input splitting theorem converges in L\ensuremath{{}^{2}} to the actual continuous kinetic integral solution on any finite horizon. Projecting to velocity preserves this convergence. Independence from the fixed future convolution is closed under L\ensuremath{{}^{2}} convergence, so the limiting reference velocity is independent of that same actual future random vector. The future marginal itself is exactly the translated OU thermal Gaussian law, by the proved Brownian convolution distribution and measure-preserving source projection.

\subsection{\texorpdfstring{Export the exact Gaussian pair laws}{Export the exact Gaussian pair laws}}


The actual reference phase remains Gaussian-times-Gibbs distributed by the proved stationarity theorem. Its velocity marginal is therefore standard Gaussian. Combining this marginal, the exact future thermal Gaussian marginal and the newly proved independence gives their full product joint law. The stitched reference chain equals the same continuous solution at grid points, yielding the identical joint-law statement for every numerical interval.

\section{\texorpdfstring{The free OU velocity has its actual stationary Gaussian marginal}{The free OU velocity has its actual stationary Gaussian marginal}}
\label{proof-module-KineticLangevin-FreeOU-lean}
{\small\color{gray}Source module: \texttt{Kinetic\allowbreak{}Langevin/\allowbreak{}Free\allowbreak{}OU.\allowbreak{}lean}.\par}


The homogeneous OU velocity retains the actual fresh filtered Brownian innovation. Its Gaussian marginal is derived from the actual joint initial-velocity/\allowbreak{}innovation law, and its continuity and primitive are identities in the actual L\ensuremath{{}^{2}} space. The later canonical source theorem proves the needed past/\allowbreak{}future joint law.

\subsection{\texorpdfstring{An independent thermal innovation restores the exact standard Gaussian law}{An independent thermal innovation restores the exact standard Gaussian law}}


If the initial velocity is standard Gaussian and the independent additive innovation has law \(\sqrt{1-a^2}\,\mathcal N(0,I)\), with \(a^2\le1\), their sum \(aV+N\) is exactly standard Gaussian. Lean first rewrites the true product innovation law as the pushforward of two actual independent standard Gaussian measures. It composes the addition map with that product pushforward, then applies the previously proved Gaussian linear-combination theorem using \(a^2+(\sqrt{1-a^2})^2=1\). The conclusion is an exact measure identity, rather than a covariance assumption.

\subsection{\texorpdfstring{Retain the actual free OU path and its actual position primitive}{Retain the actual free OU path and its actual position primitive}}


For an initial L\ensuremath{{}^{2}} velocity \(V\) and the actual scaled filtered innovation path \(N_t\), define \(U_t=\psi_0(\gamma,t)V+N_t\). This is the precise free OU velocity used to combine the initial-velocity and fresh-noise Hessian correction terms. Its L\ensuremath{{}^{2}} continuity follows from continuity of \(\psi_0\) and of the actual innovation path. Its exact primitive is \[\int_0^tU_s\,ds=\psi_1(\gamma,t)V+\int_0^tN_s\,ds.\] Lean splits the actual Bochner integral of the two continuous terms, moves the fixed velocity outside the scalar integral, and unfolds the defining primitive \(\psi_1\). No pointwise regularity or time-integral identity is supplied as an additional premise.

\subsection{\texorpdfstring{Derive the Gaussian marginal from the actual initial/\allowbreak{}fresh-\allowbreak{}noise product law}{Derive the Gaussian marginal from the actual initial/fresh-noise product law}}


For \(\gamma,t\ge0\), the damping lies between zero and one, so \(\psi_0(\gamma,t)^2\le1\). Assume the actual pair \((V,N_t)\) has the explicit product law \[\mathcal N(0,I)\otimes\bigl(\sqrt{1-\psi_0(\gamma,t)^2}\,\mathcal N(0,I)\bigr).\] This premise states the actual initial/\allowbreak{}fresh-noise probability identity; canonical past/\allowbreak{}future independence proves it for the diffusion reference. Lean applies the proved Gaussian addition identity to this actual pair law. It then rewrites the almost-everywhere representatives of the L\ensuremath{{}^{2}} scalar multiplication and sum, obtaining \(\operatorname{Law}(U_t)=\mathcal N(0,I)\) for the actual L\ensuremath{{}^{2}} element. The result can be used in Gaussian truncation even when the initial position and velocity are dependent, since no independence from that position appears in the cutoff theorem.

\section{\texorpdfstring{Actual Wiener response convolution and all \ensuremath{\psi} time primitives}{Actual Wiener response convolution and all ψ time primitives}}
\label{proof-module-KineticLangevin-KernelConvolution-lean}
{\small\color{gray}Source module: \texttt{Kinetic\allowbreak{}Langevin/\allowbreak{}Kernel\allowbreak{}Convolution.\allowbreak{}lean}.\par}


The \ensuremath{\varphi} corrections in KLMC2 are proved to be actual time convolutions of the same Brownian \ensuremath{\psi} responses. The proof establishes an ordinary integrable Fubini identity in the deterministic time-kernel Hilbert space, then applies the constructed isonormal map and spatial lift. Both velocity and position correction noises retain their actual common source.

\subsection{\texorpdfstring{Name and control the ordinary causal exponential convolution}{Name and control the ordinary causal exponential convolution}}


For a continuous scalar response \(k\), define \((\mathcal E_\beta k)(t)=\int_0^t e^{-\beta(t-s)}k(s)\,ds\). Factoring the exponential gives \((\mathcal E_\beta k)(t)=e^{-\beta t}\int_0^t e^{\beta s}k(s)\,ds\). Lean proves this by the exponential addition law inside the exact interval integral. The factored form proves continuity using continuity of an ordinary primitive of a continuous function. Allowing a damping \(\beta\) distinct from the response friction \(\gamma\) makes both the \ensuremath{\varphi} convolution and an ordinary \ensuremath{\psi} primitive instances of the same checked theorem.

\subsection{\texorpdfstring{Describe the triangular time-\allowbreak{}integration region exactly}{Describe the triangular time-integration region exactly}}


A Wiener response at time \(s\) has deterministic kernel \(1_{(0,s]}(r)\psi_k(\gamma,s-r)\). Integrating in \(s\in(0,t]\) therefore uses precisely the triangle \(0<r\le s\le t\). The first helper proves that restricting a rectangular integral by \(r\le s\) gives the inner interval \((0,s]\). The second helper performs the opposite integration order: \[\int_r^t e^{-\beta(t-s)}\psi_k(\gamma,s-r)\,ds=(\mathcal E_\beta\psi_k)(t-r).\] Lean checks the indicator-set intersections, replaces the closed endpoint by the equal Lebesgue integral over the half-open interval, and substitutes \(u=s-r\). These are exact identities on the genuine integration domains.

\subsection{\texorpdfstring{Prove the deterministic Hilbert-\allowbreak{}space Fubini identity}{Prove the deterministic Hilbert-space Fubini identity}}


Let \(K_k(t)(r)=1_{(0,t]}(r)\psi_k(\gamma,t-r)\) be the actual element of L\ensuremath{{}^{2}} time. The key identity is \[K_{\mathcal E_\beta\psi_k}(t)=\int_0^t e^{-\beta(t-s)}K_k(s)\,ds.\] Lean proves equality by pairing with an arbitrary element \(f\) of this actual Hilbert space. On the finite time rectangle, \(f\) is integrable. The continuous scalar response has a finite bound on \([0,t]\), and the damping is at most one on the triangle because \(\beta\ge0\). Thus the actual double integrand is integrable, which justifies ordinary Bochner Fubini. The two triangular-domain helpers identify the two integration orders with the two sides of the claimed Hilbert identity. Commuting the continuous linear inner-product functional with the time integral completes the equality. No stochastic Fubini conclusion is assumed.

\subsection{\texorpdfstring{Recover \ensuremath{\varphi} responses and every \ensuremath{\psi} primitive}{Recover φ responses and every ψ primitive}}


With \(\beta=\gamma\), the causal convolution is exactly \(\varphi_{k+1}\), so the Hilbert kernel for \(\varphi_{k+1}\) is the damped time integral of the kernel for \(\psi_k\). With \(\beta=0\), the causal convolution is the defining ordinary primitive \(\psi_{k+1}\). This gives the exact time-primitive identity for every \ensuremath{\psi} index, including the second primitive needed by the position correction. Lean specializes the already proved general Hilbert identity and rewrites the defining scalar integrals. The ordinary primitive result works for every friction, since its outer damping is zero.

\subsection{\texorpdfstring{Pass the checked identities through the actual Gaussian source and spatial lift}{Pass the checked identities through the actual Gaussian source and spatial lift}}


The actual scalar response on a Brownian interval is the constructed isonormal map applied to the translated deterministic time kernel. Translation and the isonormal map are continuous linear maps, so they commute with the actual Bochner time integrals. The proved Hilbert identities therefore yield \[N_{\varphi_{k+1}}(a,t)=\int_0^t e^{-\gamma(t-s)}N_{\psi_k}(a,s)\,ds,\qquad N_{\psi_{k+1}}(a,t)=\int_0^tN_{\psi_k}(a,s)\,ds.\] Lean proves continuity of each integrand, applies the continuous linear integral-commutation theorem, and rewrites scalar multiplication. It then applies the actual continuous linear spatial lift to obtain the same identities for arbitrary finite Euclidean coordinate sets. All responses use the very same translated Brownian source. In particular, the \ensuremath{\varphi}\ensuremath{{}_{2}} and \ensuremath{\varphi}\ensuremath{{}_{3}} noises in KLMC2 are exactly the damped integrals of the first and second primitives of the fresh OU noise, which is the identity required to retain the omitted stochastic Hessian-difference terms in the repaired comparison proof.

\section{\texorpdfstring{The complete KLMC2 Hessian correction is an actual free-\allowbreak{}OU time integral}{The complete KLMC2 Hessian correction is an actual free-OU time integral}}
\label{proof-module-KineticLangevin-FreeOUCorrection-lean}
{\small\color{gray}Source module: \texttt{Kinetic\allowbreak{}Langevin/\allowbreak{}Free\allowbreak{}OUCorrection.\allowbreak{}lean}.\par}


The fresh stochastic Hessian-difference terms are retained exactly. Their \ensuremath{\varphi}\ensuremath{{}_{2}} and \ensuremath{\varphi}\ensuremath{{}_{3}} responses combine with the initial velocity into the first and second time primitives of one actual free OU velocity. A bounded random Hessian acts through its proved continuous linear L\ensuremath{{}^{2}} operator, so all identities remain valid on the actual coupled source.

\subsection{\texorpdfstring{Keep each deterministic and stochastic response on the same source}{Keep each deterministic and stochastic response on the same source}}


Let \(L\) be the actual continuous linear pullback from the canonical Brownian L\ensuremath{{}^{2}} space to the coupled probability space, including the factor \(\sqrt{2\gamma}\). Define \[R_k(t)=\psi_k(\gamma,t)V_0+L N_{\psi_k}(a,t),\qquad C_k(t)=\varphi_{k+1}(\gamma,t)V_0+L N_{\varphi_{k+1}}(a,t).\] Here every noise response uses exactly the same Brownian interval beginning at \(a\). In particular, \(R_0\) is the actual free OU velocity and \(C_1,C_2\) are the full \ensuremath{\varphi}\ensuremath{{}_{2}} and \ensuremath{\varphi}\ensuremath{{}_{3}} correction bundles. Lean defines these actual L\ensuremath{{}^{2}} elements and proves continuity of each \(R_k\) from the checked response coefficient and canonical convolution continuity.

\subsection{\texorpdfstring{Derive the exact primitive and damped-\allowbreak{}convolution identities}{Derive the exact primitive and damped-convolution identities}}


The deterministic response coefficients and actual stochastic responses satisfy the same checked time-integral identities. Consequently \[R_{k+1}(t)=\int_0^tR_k(s)\,ds,\qquad C_k(t)=\int_0^te^{-\gamma(t-s)}R_k(s)\,ds.\] Lean proves integrability of every continuous summand, splits each Bochner integral, moves the fixed initial velocity outside its scalar integral, and commutes the actual source pullback with the integral. It then substitutes the proved Wiener \ensuremath{\psi} primitive or \ensuremath{\varphi} convolution identity. This establishes the complete correction identity, including the fresh stochastic terms absent from the displayed subtraction in the source paper.

\subsection{\texorpdfstring{Pass a genuine random Hessian operator through the actual integrals}{Pass a genuine random Hessian operator through the actual integrals}}


A uniformly bounded measurable random Hessian difference acts as a proved continuous linear map \(A\) on the actual L\ensuremath{{}^{2}} space. Therefore it commutes with every genuine time primitive and with the damped convolution. Lean proves this for an arbitrary actual continuous linear \(A\); the random-Hessian construction is a valid concrete instantiation. The resulting full correction identities are \[AC_1(t)=\int_0^te^{-\gamma(t-s)}\int_0^sAR_0(r)\,dr\,ds,\] and \[AC_2(t)=\int_0^te^{-\gamma(t-s)}\int_0^s\int_0^rAR_0(u)\,du\,dr\,ds.\] The proof checks continuity of every relevant L\ensuremath{{}^{2}} integrand and applies the continuous linear integral-commutation theorem, then rewrites the proved primitive identities on the nonnegative integration intervals. Thus Gaussian truncation can be applied to \(AR_0(r)\) at each time, with no independence between the random Hessian and the velocity required.

\subsection{\texorpdfstring{Preserve the crucial \ensuremath{\varphi}\ensuremath{{}_{2}}+\ensuremath{\gamma}\ensuremath{\varphi}\ensuremath{{}_{3}} cancellation for the actual noises}{Preserve the crucial φ₂+γφ₃ cancellation for the actual noises}}


The scalar response identity \(\varphi_2+\gamma\varphi_3=\psi_2\) also holds for the actual jointly sampled stochastic responses. Lean proves the corresponding equality of deterministic time-kernel L\ensuremath{{}^{2}} elements by their actual almost-everywhere indicator representatives. It applies the same time translation, isonormal map, and spatial lift, all continuous and linear. This proves \[N_{\varphi_2}(a,t)+\gamma N_{\varphi_3}(a,t)=N_{\psi_2}(a,t).\] Combining this with the identical scalar coefficient relation gives the exact full-bundle identity \[C_1(t)+\gamma C_2(t)=R_2(t).\] This cancellation includes both fresh stochastic corrections. It is what permits the transformed phase estimate to use a true second primitive of the stationary free OU velocity and recover the correct quadratic time coefficient.

\section{\texorpdfstring{The actual fresh free-\allowbreak{}OU velocity has its Gaussian marginal}{The actual fresh free-OU velocity has its Gaussian marginal}}
\label{proof-module-KineticLangevin-CanonicalFreeOU-lean}
{\small\color{gray}Source module: \texttt{Kinetic\allowbreak{}Langevin/\allowbreak{}Canonical\allowbreak{}Free\allowbreak{}OU.\allowbreak{}lean}.\par}


The Gaussian cutoff in the repaired second-order proof applies to a true free-OU path on each actual reference interval. Its marginal law follows from proved stationarity and proved independence of future Brownian increments from the actual past solution.

\subsection{\texorpdfstring{Use the actual reference-\allowbreak{}grid joint law}{Use the actual reference-grid joint law}}


At the grid time \(a=nh\), let \(V_a\) be the velocity of the actual stationary reference chain and let \(N_{\psi_0}(a,r)\) be its actual fresh Brownian response. The checked past/\allowbreak{}future theorem gives the joint law \[\bigl(V_a,\sqrt{2\gamma}N_{\psi_0}(a,r)\bigr)\sim N(0,I)\otimes N\bigl(0,(1-e^{-2\gamma r})I\bigr).\] In particular, the first random variable is independent of this future response even though it depends on the entire past path. Lean specializes this exact joint law to the nonnegative real duration, then applies the proved Gaussian addition identity to \(R_0(r)=e^{-\gamma r}V_a+\sqrt{2\gamma}N_{\psi_0}(a,r)\). Its covariance is \(e^{-2\gamma r}I+(1-e^{-2\gamma r})I=I\), so \(R_0(r)\sim N(0,I)\). All variables are the actual L\ensuremath{{}^{2}} representatives used by the numerical coupling.

\subsection{\texorpdfstring{Apply the same result at arbitrary continuous reference times}{Apply the same result at arbitrary continuous reference times}}


The global reference solution has the same proven past/\allowbreak{}future joint law at every \(a\in[0,T]\). The identical Gaussian calculation therefore proves the free-OU response beginning at an arbitrary continuous reference time is standard Gaussian at every later duration. This supports local cutoff estimates independently of a chosen numerical grid. Neither statement assumes the desired marginal or an adaptedness contract; both invoke the already proved canonical independence and stationarity results.

\section{\texorpdfstring{The actual cubic and quartic errors of the Hessian correction}{The actual cubic and quartic errors of the Hessian correction}}
\label{proof-module-KineticLangevin-SecondOrderLocal-lean}
{\small\color{gray}Source module: \texttt{Kinetic\allowbreak{}Langevin/\allowbreak{}Second\allowbreak{}Order\allowbreak{}Local.\allowbreak{}lean}.\par}


The free-OU path U contains the exact initial velocity and the shared Brownian convolution. The exact velocity is U minus the damped force integral. Correcting the frozen-force approximation by the Hessian acting on the integrated free path leaves two explicit errors: the quadratic force Taylor remainder and a triple force-memory integral. Their orders and constants are proved from these definitions.

\subsection{\texorpdfstring{Build genuine damped integrals and continuous paths}{Build genuine damped integrals and continuous paths}}


dampedIntegral(\ensuremath{\gamma},F,t)=\ensuremath{\int}\ensuremath{{}_{0}}\ensuremath{{}^{t}}exp(\ensuremath{-}\ensuremath{\gamma}(t\ensuremath{-}s))F(s)ds. It is continuous for a continuous F: factor exp(\ensuremath{-}\ensuremath{\gamma}t) outside the primitive of exp(\ensuremath{\gamma}s)F(s), then apply the fundamental theorem of calculus to that primitive. The general primitive continuity lemma also supports all nested integrals used below.

\subsection{\texorpdfstring{Define the corrected approximation and its two actual defects}{Define the corrected approximation and its two actual defects}}


The corrected velocity is U(t)\ensuremath{-}\ensuremath{\psi}\ensuremath{{}_{1}}(t)G(0)\ensuremath{-}H\ensuremath{{}_{0}}\ensuremath{\int}\ensuremath{{}_{0}}\ensuremath{{}^{t}}exp(\ensuremath{-}\ensuremath{\gamma}(t\ensuremath{-}s))\ensuremath{\int}\ensuremath{{}_{0}}\ensuremath{{}^{s}}U(r)dr ds. The Taylor defect integrates G(s)\ensuremath{-}G(0)\ensuremath{-}H\ensuremath{{}_{0}}(X(s)\ensuremath{-}X(0)) against the damping kernel. The memory defect applies H\ensuremath{{}_{0}} to the damped integral of the primitive of the damped force integral. These definitions retain the entire free-OU path, including its stochastic Hessian correction.

\subsection{\texorpdfstring{Derive the local-\allowbreak{}error identity from the exact equations}{Derive the local-error identity from the exact equations}}


From V=U\ensuremath{-}K\ensuremath{\gamma}G and X(s)\ensuremath{-}X(0)=\ensuremath{\int}\ensuremath{{}_{0}}\ensuremath{{}^{s}}V, obtain \ensuremath{\int}\ensuremath{{}_{0}}\ensuremath{{}^{s}}U=X(s)\ensuremath{-}X(0)+\ensuremath{\int}\ensuremath{{}_{0}}\ensuremath{{}^{s}}K\ensuremath{\gamma}G. Substitute this identity into the correction, commute the bounded linear Hessian through the Bochner integral, and separate the frozen constant force using \ensuremath{\int}exp(\ensuremath{-}\ensuremath{\gamma}(t\ensuremath{-}s))ds=\ensuremath{\psi}\ensuremath{{}_{1}}(t). Algebra then proves correctedVelocity\ensuremath{-}V equals the Taylor defect minus the memory defect. The result assumes exact mild and kinematic equations, rather than an error inequality.

\subsection{\texorpdfstring{Integrate polynomial envelopes with their exact constants}{Integrate polynomial envelopes with their exact constants}}


For forward time and \ensuremath{\gamma}\ensuremath{\ge}0, the damping kernel lies in (0,1]. If \ensuremath{\Vert}F(s)\ensuremath{\Vert}\ensuremath{\le}Cs\ensuremath{{}^{n}}, the norm of its damped integral is at most \ensuremath{\int}\ensuremath{{}_{0}}\ensuremath{{}^{t}}Cs\ensuremath{{}^{n}}ds=Ct\ensuremath{{}^{n+1}}/\allowbreak{}(n+1). The undamped case follows at \ensuremath{\gamma}=0. Applying this rule three times to a score bound G\ensuremath{{}_{0}} gives the memory norm at most \ensuremath{\Vert}H\ensuremath{{}_{0}}\ensuremath{\Vert}G\ensuremath{{}_{0}}t\ensuremath{{}^{3}}/\allowbreak{}6.

\subsection{\texorpdfstring{Obtain cubic velocity and quartic position errors}{Obtain cubic velocity and quartic position errors}}


A Taylor forcing bound Rs\ensuremath{{}^{2}} integrates to Rt\ensuremath{{}^{3}}/\allowbreak{}3. Adding the memory estimate yields velocity error at most (2R+\ensuremath{\Vert}H\ensuremath{{}_{0}}\ensuremath{\Vert}G\ensuremath{{}_{0}})t\ensuremath{{}^{3}}/\allowbreak{}6. The corrected position is X(0) plus the integral of the corrected velocity. Its difference from the exact position is therefore the integral of the actual velocity error, giving (2R+\ensuremath{\Vert}H\ensuremath{{}_{0}}\ensuremath{\Vert}G\ensuremath{{}_{0}})t\ensuremath{{}^{4}}/\allowbreak{}24. Continuity of the corrected velocity justifies this subtraction of integrals.

\subsection{\texorpdfstring{Check the paper’s quarter-\allowbreak{}constant in transformed coordinates}{Check the paper’s quarter-constant in transformed coordinates}}


Write B=Ct\ensuremath{{}^{3}}. The component estimates give \ensuremath{\Vert}v\ensuremath{\Vert}\ensuremath{\le}B/\allowbreak{}6 and \ensuremath{\Vert}x\ensuremath{\Vert}\ensuremath{\le}Bt/\allowbreak{}24. Under \ensuremath{\gamma}t\ensuremath{\le}1/\allowbreak{}5, the first transformed coordinate satisfies \ensuremath{\Vert}v+\ensuremath{\gamma}x\ensuremath{\Vert}\ensuremath{\le}7B/\allowbreak{}40; the second has norm \ensuremath{\Vert}v\ensuremath{\Vert}. Squaring and adding gives an energy below (B/\allowbreak{}4)\ensuremath{{}^{2}}, hence transformed norm at most B/\allowbreak{}4. Applying this exact component lemma with C=2R+\ensuremath{\Vert}H\ensuremath{{}_{0}}\ensuremath{\Vert}G\ensuremath{{}_{0}} proves the actual local phase estimate. For stationary kinetic velocities and score moments, C becomes M\ensuremath{{}_{2}}\ensuremath{\surd}(p\ensuremath{{}^{2}}+2p)+M\ensuremath{\surd}(Mp).

\section{\texorpdfstring{The implemented KLMC2 step is the genuine corrected free-\allowbreak{}OU path}{The implemented KLMC2 step is the genuine corrected free-OU path}}
\label{proof-module-KineticLangevin-CorrectedStep-lean}
{\small\color{gray}Source module: \texttt{Kinetic\allowbreak{}Langevin/\allowbreak{}Corrected\allowbreak{}Step.\allowbreak{}lean}.\par}


The local Taylor estimate and the actual four-noise numerical map refer to exactly the same L\ensuremath{{}^{2}} random variables. This module proves the identification, including both stochastic Hessian corrections, by genuine time differentiation and integration.

\subsection{\texorpdfstring{Differentiate the true damped convolution}{Differentiate the true damped convolution}}


For a continuous Banach-valued path \(F\), write \(D_\gamma F(t)=\int_0^te^{-\gamma(t-s)}F(s)\,ds\). Factoring the exponential yields \(D_\gamma F(t)=e^{-\gamma t}\int_0^te^{\gamma s}F(s)\,ds\). Lean applies the fundamental theorem of calculus to the continuous integrand and the product rule, then cancels \(e^{-\gamma t}e^{\gamma t}=1\). The result is the exact equation \((D_\gamma F)\prime=F-\gamma D_\gamma F\), for every real time.

\subsection{\texorpdfstring{Integrate the entire stochastic correction}{Integrate the entire stochastic correction}}


The actual canonical Wiener kernel on an interval of zero duration is zero in L\ensuremath{{}^{2}}; Lean proves this by almost-everywhere equality of its deterministic kernel, then applies the isonormal map and spatial lift. Thus the full correction \(C_k(t)=\varphi_{k+1}(t)V_0+L N_{\varphi_{k+1}}(a,t)\) vanishes initially and is continuous. The checked identity \(C_1+\gamma C_2=R_2\), together with \(C_2=D_\gamma R_2\), gives \(C_2\prime=C_1\) at every positive time. The fundamental theorem then proves \(C_2(t)=\int_0^tC_1(s)\,ds\). Also \(C_1=D_\gamma(\int_0^{\cdot}R_0)\), using the genuine primitive \(R_1=\int R_0\). Both statements include fresh noise terms.

\subsection{\texorpdfstring{Identify the analytic velocity and position with full kernel formulas}{Identify the analytic velocity and position with full kernel formulas}}


The Taylor analysis defines \(\widehat V_t=R_0(t)-\psi_1(t)F_0-A D_\gamma(\int R_0)(t)\). Substituting the proved identity gives \(\widehat V_t=R_0(t)-\psi_1(t)F_0-AC_1(t)\). Integrating this true path gives \(\widehat X_t=x_0+R_1(t)-\psi_2(t)F_0-AC_2(t)\). Lean first proves equality of the velocity integrands on the actual time interval, establishes all continuous summands are integrable, splits their Bochner integrals, and commutes the bounded Hessian operator with the integral. It then uses the checked primitive identities and additive algebra.

\subsection{\texorpdfstring{Match the implemented four-\allowbreak{}noise KLMC2 update exactly}{Match the implemented four-noise KLMC2 update exactly}}


Let \(R\) be the unscaled source pullback and sample the four correlated responses
\[(N_{\psi_0},N_{\psi_1},N_{\varphi_2},N_{\varphi_3})\]
 on the same Brownian interval. The analytic free-OU path uses the thermally scaled map \(L=\sqrt{2\gamma}\,R\). Expanding the proved velocity and position formulas gives precisely the implemented \(\texttt{klmc2Step}\): its deterministic Hessian terms are \(-\varphi_2H_0V_0\) and \(-\varphi_3H_0V_0\), and its fresh terms are \(\sqrt{2\gamma}(N_{\psi_0}-H_0N_{\varphi_2})\) and \(\sqrt{2\gamma}(N_{\psi_1}-H_0N_{\varphi_3})\). Lean reduces the four finite indices, applies linearity of the actual Hessian and source maps, and closes the two component identities by module algebra. No local error or noise-independence conclusion is assumed.

\section{\texorpdfstring{The genuine stationary reference segment supplies every second-\allowbreak{}order analytic hypothesis}{The genuine stationary reference segment supplies every second-order analytic hypothesis}}
\label{proof-module-KineticLangevin-CanonicalSecondOrderSegments-lean}
{\small\color{gray}Source module: \texttt{Kinetic\allowbreak{}Langevin/\allowbreak{}Canonical\allowbreak{}Second\allowbreak{}Order\allowbreak{}Segments.\allowbreak{}lean}.\par}


Each local second-order comparison uses an actual selected kinetic solution on a true canonical Brownian interval. This module exposes its continuous integral equation, exact endpoint stitching, stationary phase law, force-mild formula, and kinematic identity.

\subsection{\texorpdfstring{Obtain continuity and endpoint matching from the actual solution constructor}{Obtain continuity and endpoint matching from the actual solution constructor}}


The reference segment \(Z_n(s)\), \(0\le s\le h\), is the selected continuous solution starting at the true preceding reference state and driven by the actual Brownian increment beginning at \(nh\). The checked constructor proves continuity, \(Z_n(0)=R_n\), \(Z_n(h)=R_{n+1}\), and the genuine kinetic integral equation \[Z_n(t)-Z_n(0)=-\int_0^tb(Z_n(s))\,ds+(\sqrt{2\gamma}\,\Delta W_n(t),0).\] Lean specializes these already constructed solutions with the actual lifted gradient and unit inverse mass. No local path or endpoint-matching assumption is introduced.

\subsection{\texorpdfstring{Derive the true mild and kinematic equations with the same fresh noise}{Derive the true mild and kinematic equations with the same fresh noise}}


Projecting the actual integral equation gives the velocity equation. The proved filtered canonical Brownian response solves \(N(t)=\sqrt{2\gamma}\Delta W_n(t)-\gamma\int_0^tN(s)\,ds\); applying the genuine integrating-factor theorem yields \[V_n(t)=R_0(t)-\int_0^te^{-\gamma(t-s)}\nabla f(X_n(s))\,ds.\] Here \(R_0\) is exactly the actual free-OU response used by the corrected numerical path. The positional projection gives \(X_n(t)-X_n(0)=\int_0^tV_n(s)\,ds\). Lean restricts the checked segment equation to each shorter duration, substitutes the exact filtered-input identity, simplifies unit inverse mass, and reorders the resulting additive terms.

\subsection{\texorpdfstring{Transfer actual global stationarity to every local segment point}{Transfer actual global stationarity to every local segment point}}


The reference segment is already proved equal to the single actual global reference solution at time \(nh+t\). Choose horizon \(T=(n+1)h\); the segment time lies in \([0,T]\). Lean rewrites the local solution using this actual stitching identity and applies the proved unconditional Gibbs stationarity of the global canonical solution. Thus the full phase law at every segment point, including both endpoints, is \(N(0,I)\otimes\pi\). This supplies the genuine stationary moment hypotheses for the local second-order Taylor estimate.

\section{\texorpdfstring{The genuine frozen-\allowbreak{}linear contraction in Proposition 3}{The genuine frozen-linear contraction in Proposition 3}}
\label{proof-module-KineticLangevin-SecondOrderLinear-lean}
{\small\color{gray}Source module: \texttt{Kinetic\allowbreak{}Langevin/\allowbreak{}Second\allowbreak{}Order\allowbreak{}Linear.\allowbreak{}lean}.\par}


The friction can be much larger than the square root of the smoothness constant. A crude Euler remainder proportional to \ensuremath{\gamma}\ensuremath{{}^{2}} would lose the paper’s step-size regime. The identities below instead use the force’s cocoercivity to retain a second-order budget proportional to M, and prove the actual frozen-linear map contracts by 1\ensuremath{-}mt/\allowbreak{}(2\ensuremath{\gamma}) under 5\ensuremath{\gamma}Mt\ensuremath{\le}m.

\subsection{\texorpdfstring{Separate fast damping from the Euler remainder}{Separate fast damping from the Euler remainder}}


For the kinetic difference drift A=(\ensuremath{-}\ensuremath{\gamma}v\ensuremath{-}F,v), its transformed squared norm satisfies the exact identity E(A)=\ensuremath{-}\ensuremath{\gamma}E\ensuremath{\prime}/\allowbreak{}2+2\ensuremath{\Vert}F\ensuremath{\Vert}\ensuremath{{}^{2}}\ensuremath{-}\ensuremath{\gamma}\ensuremath{{}^{2}}\ensuremath{\langle}x,F\ensuremath{\rangle}. The actual Euler phase (v,x)+sA has energy E+sE\ensuremath{\prime}+s\ensuremath{{}^{2}}E(A). Lean expands the two transformed coordinates, uses the real inner-product identities, and checks both equalities by ring normalization.

\subsection{\texorpdfstring{Use cocoercivity before bounding the remainder}{Use cocoercivity before bounding the remainder}}


The force satisfies \ensuremath{\Vert}F\ensuremath{\Vert}\ensuremath{{}^{2}}\ensuremath{\le}M\ensuremath{\langle}x,F\ensuremath{\rangle} and \ensuremath{\langle}x,F\ensuremath{\rangle}\ensuremath{\le}M\ensuremath{\Vert}x\ensuremath{\Vert}\ensuremath{{}^{2}}. Since \ensuremath{\gamma}\ensuremath{{}^{2}}\ensuremath{\ge}m+M, 2M\ensuremath{-}\ensuremath{\gamma}\ensuremath{{}^{2}}\ensuremath{\le}M\ensuremath{-}m. Also \ensuremath{\gamma}\ensuremath{{}^{2}}\ensuremath{\Vert}x\ensuremath{\Vert}\ensuremath{{}^{2}}\ensuremath{\le}2E. Combining these facts gives E(A)\ensuremath{\le}\ensuremath{-}\ensuremath{\gamma}E\ensuremath{\prime}/\allowbreak{}2+2M(M\ensuremath{-}m)E/\allowbreak{}\ensuremath{\gamma}\ensuremath{{}^{2}}. In the Euler expansion, the derivative coefficient becomes s(1\ensuremath{-}\ensuremath{\gamma}s/\allowbreak{}2), which is nonnegative for \ensuremath{\gamma}s\ensuremath{\le}1. Actual force dissipativity E\ensuremath{\prime}\ensuremath{\le}\ensuremath{-}2mE/\allowbreak{}\ensuremath{\gamma} then yields transformed norm at most [1\ensuremath{-}ms/\allowbreak{}\ensuremath{\gamma}+(m/\allowbreak{}2+M(M\ensuremath{-}m)/\allowbreak{}\ensuremath{\gamma}\ensuremath{{}^{2}})s\ensuremath{{}^{2}}] times the initial norm. The final square comparison retains an additional nonnegative square instead of estimating it away.

\subsection{\texorpdfstring{Bound the true frozen-\allowbreak{}Hessian correction}{Bound the true frozen-Hessian correction}}


The remaining phase correction is (\ensuremath{-}\ensuremath{\varphi}\ensuremath{{}_{2}}Hv,\ensuremath{-}\ensuremath{\psi}\ensuremath{{}_{2}}F\ensuremath{-}\ensuremath{\varphi}\ensuremath{{}_{3}}Hv). The exact coefficient identity \ensuremath{\varphi}\ensuremath{{}_{2}}+\ensuremath{\gamma}\ensuremath{\varphi}\ensuremath{{}_{3}}=\ensuremath{\psi}\ensuremath{{}_{2}} transforms it into (\ensuremath{-}\ensuremath{\psi}\ensuremath{{}_{2}}(Hv+\ensuremath{\gamma}F),\ensuremath{\varphi}\ensuremath{{}_{2}}Hv). The bounds \ensuremath{\psi}\ensuremath{{}_{2}},\ensuremath{\varphi}\ensuremath{{}_{2}}\ensuremath{\le}t\ensuremath{{}^{2}}/\allowbreak{}2, \ensuremath{\Vert}H\ensuremath{\Vert}\ensuremath{\le}M, \ensuremath{\Vert}v\ensuremath{\Vert}\ensuremath{\le}N, and \ensuremath{\gamma}\ensuremath{\Vert}x\ensuremath{\Vert}\ensuremath{\le}3N/\allowbreak{}2 give component envelopes 5Mt\ensuremath{{}^{2}}N/\allowbreak{}4 and Mt\ensuremath{{}^{2}}N/\allowbreak{}2. Their squared sum is below (3Mt\ensuremath{{}^{2}}N/\allowbreak{}2)\ensuremath{{}^{2}}. Thus the actual correction costs at most 3Mt\ensuremath{{}^{2}}N/\allowbreak{}2, with no symmetry or common Hessian eigenbasis assumed.

\subsection{\texorpdfstring{Check the scalar coefficient budget with exact \ensuremath{\psi}\ensuremath{{}_{1}}}{Check the scalar coefficient budget with exact ψ₁}}


The friction balance \ensuremath{\gamma}\ensuremath{\psi}\ensuremath{{}_{2}}=t\ensuremath{-}\ensuremath{\psi}\ensuremath{{}_{1}} and \ensuremath{\psi}\ensuremath{{}_{2}}\ensuremath{\le}t\ensuremath{{}^{2}}/\allowbreak{}2 imply \ensuremath{\psi}\ensuremath{{}_{1}}\ensuremath{\ge}t\ensuremath{-}\ensuremath{\gamma}t\ensuremath{{}^{2}}/\allowbreak{}2, while 0\ensuremath{\le}\ensuremath{\psi}\ensuremath{{}_{1}}\ensuremath{\le}t. Moreover m+M(M\ensuremath{-}m)/\allowbreak{}\ensuremath{\gamma}\ensuremath{{}^{2}}\ensuremath{\le}M. These inequalities bound the Euler coefficient plus the frozen correction by 1\ensuremath{-}mt/\allowbreak{}\ensuremath{\gamma}+5Mt\ensuremath{{}^{2}}/\allowbreak{}2. The source step restriction 5\ensuremath{\gamma}Mt\ensuremath{\le}m makes the last term at most mt/\allowbreak{}(2\ensuremath{\gamma}), leaving 1\ensuremath{-}mt/\allowbreak{}(2\ensuremath{\gamma}). All divisions are justified by positive friction.

\subsection{\texorpdfstring{Identify and contract the actual frozen-\allowbreak{}linear map}{Identify and contract the actual frozen-linear map}}


frozenSecondOrderStep is the explicit phase with \ensuremath{\psi}\ensuremath{{}_{0}}v\ensuremath{-}\ensuremath{\psi}\ensuremath{{}_{1}}F\ensuremath{-}\ensuremath{\varphi}\ensuremath{{}_{2}}Hv in velocity and x+\ensuremath{\psi}\ensuremath{{}_{1}}v\ensuremath{-}\ensuremath{\psi}\ensuremath{{}_{2}}F\ensuremath{-}\ensuremath{\varphi}\ensuremath{{}_{3}}Hv in position. The identity \ensuremath{\psi}\ensuremath{{}_{0}}=1\ensuremath{-}\ensuremath{\gamma}\ensuremath{\psi}\ensuremath{{}_{1}} writes this actual map as the Euler phase plus the correction above. Strong-smooth interpolation supplies both cocoercivity and force dissipativity. The transformed triangle inequality and the checked scalar budget therefore prove the source contraction factor for the actual frozen-linear part of KLMC2.

\section{\texorpdfstring{Derive the chi-\allowbreak{}square positive-\allowbreak{}part tail from the Gaussian density}{Derive the chi-square positive-part tail from the Gaussian density}}
\label{proof-module-KineticLangevin-GaussianTail-lean}
{\small\color{gray}Source module: \texttt{Kinetic\allowbreak{}Langevin/\allowbreak{}Gaussian\allowbreak{}Tail.\allowbreak{}lean}.\par}


The Hessian comparison needs a Gaussian tail estimate even when the Gaussian velocity and position discrepancy are dependent. The relevant bound is \(\mathbb E(\lVert G\rVert^2-a)_+\le4e^{-(a-p)/4}\) for \(a\ge5p\). The proof below derives its exponential moment and the tail estimate from actual Gaussian measures.

\subsection{\texorpdfstring{Compute the scalar exponential square moment with its integrability}{Compute the scalar exponential square moment with its integrability}}


Multiply the standard Gaussian density by \(e^{x^2/3}\). Its exponent becomes \(-x^2/2+x^2/3=-x^2/6\). The real Gaussian integral therefore gives \(\mathbb Ee^{G^2/3}=\sqrt3\). The proof establishes integrability of this weighted density before computing its integral, so the Bochner integral cannot conceal divergence. Lean uses the actual Gaussian density, the with-density integrability equivalence, the real Gaussian integral and exact square-root algebra.

\subsection{\texorpdfstring{Multiply the independent coordinate moments}{Multiply the independent coordinate moments}}


Since \(e^{\lVert G\rVert^2/3}=\prod_i e^{G_i^2/3}\), product integrability and finite-product Fubini give \(\mathbb Ee^{\lVert G\rVert^2/3}=(\sqrt3)^p\). The proof works on the constructed product of scalar Gaussian measures and uses its established Euclidean standard-Gaussian pushforward. This handles dimension zero as well.

\subsection{\texorpdfstring{Choose an elementary Chernoff tilt}{Choose an elementary Chernoff tilt}}


The elementary bound \(e^u\ge1+u\) at \(u=1/3\), raised to the fourth power, proves \(\sqrt3\le e^{2/3}\). Also \((y-a)_+\le3e^{(y-a)/3}\), since its positive branch is bounded using the same exponential inequality and its zero branch is nonnegative. These two inequalities give a convenient exponential moment estimate without importing an unproved chi-square concentration result.

\subsection{\texorpdfstring{Integrate the positive-\allowbreak{}part bound and use the threshold}{Integrate the positive-part bound and use the threshold}}


Applying the pointwise bound to \(y=\lVert G\rVert^2\) gives \(\mathbb E(y-a)_+\le3e^{-a/3}(\sqrt3)^p\le3e^{-a/3+2p/3}\). If \(a\ge5p\), the exponent is at most \(-(a-p)/4\). Replacing the factor three by four yields exactly the source bound. Lean first derives integrability of the positive part from its integrable exponential majorant, then applies integral monotonicity, the exact moment and scalar exponent comparison.

\subsection{\texorpdfstring{Use the bound on the original coupled probability space}{Use the bound on the original coupled probability space}}


The pushforward integration identity transfers the estimate to any actual standard-Gaussian velocity. Its joint distribution with a numerical position discrepancy is unrestricted. This is essential for a synchronous sampling proof, where past noise generally makes those two variables dependent.

\section{\texorpdfstring{A checked Gaussian truncation bound with dependent position and velocity}{A checked Gaussian truncation bound with dependent position and velocity}}
\label{proof-module-KineticLangevin-GaussianTruncation-lean}
{\small\color{gray}Source module: \texttt{Kinetic\allowbreak{}Langevin/\allowbreak{}Gaussian\allowbreak{}Truncation.\allowbreak{}lean}.\par}


A Hessian difference can be small through its spatial Lipschitz bound and uniformly bounded through the Hessian norm bound. This module combines those two genuine pointwise controls with the actual Gaussian velocity tail. The coupled position and velocity may be dependent; no independence premise is used.

\subsection{\texorpdfstring{Separate the moderate velocity region from the exact Gaussian excess}{Separate the moderate velocity region from the exact Gaussian excess}}


Suppose \(d\le\min\{B,Lx\}\,v\), with all quantities nonnegative. For a cutoff \(a\ge0\), the exact scalar estimate is \[d^2\le L^2ax^2+B^2(v^2-a)_+.\] Put \(k=\min\{B,Lx\}\). Where \(v^2\le a\), use \(k\le Lx\) and \(d^2\le k^2v^2\le L^2x^2a\). Where \(v^2>a\), write \(k^2v^2=k^2a+k^2(v^2-a)\), use \(k\le Lx\) in the first term and \(k\le B\) in the second. Lean proves the two squared bounds from the nonnegative minimum, splits on the cutoff comparison, rewrites the positive part in each region, and applies these two scalar comparisons. The subtraction of \(a\) in the tail term is retained exactly, which permits the sharp positive-part tail bound.

\subsection{\texorpdfstring{Identify actual Hilbert L\ensuremath{{}^{2}} norms with their actual second moments}{Identify actual Hilbert L² norms with their actual second moments}}


For an actual Hilbert-valued L\ensuremath{{}^{2}} element \(X\), \(\|X\|_{L^2}^2=\int\|X(\omega)\|^2\,dP\). Lean rewrites the squared norm as the real inner product of \(X\) with itself, unfolds the actual L\ensuremath{{}^{2}} inner product as its Bochner integral, and rewrites the pointwise inner products as squared norms. This gives the precise identity needed to integrate the scalar cutoff estimate; it does not replace a possibly divergent moment by a totalized real expression.

\subsection{\texorpdfstring{Integrate the cutoff and use only the actual Gaussian velocity marginal}{Integrate the cutoff and use only the actual Gaussian velocity marginal}}


Let \(D,X,V\) be actual Euclidean L\ensuremath{{}^{2}} random variables on one probability space, and suppose \(V\) has the actual standard Gaussian law in dimension \(p\). Assume \(\|D\|\le\min\{B,L\|X\|\}\|V\|\) almost everywhere. The preceding scalar inequality integrates to \[\|D\|_{L^2}^2\le L^2a\|X\|_{L^2}^2+B^2\mathbb E(\|V\|^2-a)_+.\] Lean proves integrability of the velocity positive part directly from the actual L\ensuremath{{}^{2}} moment, proves integrability of both terms on the right, and applies almost-everywhere monotonicity of the Bochner integral.



The checked Gaussian tail theorem, transferred through the actual velocity law, gives \(\mathbb E(\|V\|^2-a)_+\le4e^{-(a-p)/4}\) for \(a\ge5p\). Hence \[\|D\|_{L^2}\le L\sqrt a\,\|X\|_{L^2}+2B e^{-(a-p)/8}.\] Lean checks the exponential-square identity, compares the two nonnegative squares, and bounds the nonnegative cross term. The only probability fact used is the marginal law of \(V\); arbitrary dependence between \(X\), \(D\), and \(V\) is retained throughout. This is the cutoff estimate required for a random Hessian-difference factor multiplying a stationary Gaussian velocity.

\section{\texorpdfstring{The genuine sharp Hessian-\allowbreak{}difference bound and its Gaussian L\ensuremath{{}^{2}} cutoff}{The genuine sharp Hessian-difference bound and its Gaussian L² cutoff}}
\label{proof-module-KineticLangevin-HessianDifference-lean}
{\small\color{gray}Source module: \texttt{Kinetic\allowbreak{}Langevin/\allowbreak{}Hessian\allowbreak{}Difference.\allowbreak{}lean}.\par}


The actual Hessians of one strongly convex smooth potential differ in operator norm by at most M\ensuremath{-}m. The proof derives this directly from the checked centered gradient sector inequality and genuine derivatives. It then proves the identical bound for the actual random L\ensuremath{{}^{2}} Hessian operators and establishes the precise Gaussian truncation term used in the repaired KLMC2 comparison.

\subsection{\texorpdfstring{Center the gradient and derive its sharp Lipschitz radius}{Center the gradient and derive its sharp Lipschitz radius}}


For an \(m\)-strongly convex potential with \(M\)-Lipschitz gradient, the checked strong-smooth interpolation inequality implies \[\left\|\nabla f(x)-\nabla f(y)-\frac{m+M}{2}(x-y)\right\|\le\frac{M-m}{2}\|x-y\|.\] Thus \(g_c(x)=\nabla f(x)-\frac{m+M}{2}x\) is globally \((M-m)/2\)-Lipschitz. Lean constructs the nonnegative Lipschitz constant, rewrites the difference of the centered gradients as the residual in this previously proved sector estimate, and applies that exact estimate. No Hessian or spectral conclusion is inserted as a premise.

\subsection{\texorpdfstring{Derive the actual Hessian curvature-\allowbreak{}width bound from actual derivatives}{Derive the actual Hessian curvature-width bound from actual derivatives}}


If \(H_x=D\nabla f(x)\), the genuine derivative of the centered gradient is \(H_x-\frac{m+M}{2}I\). The derivative of a globally Lipschitz map has operator norm at most its Lipschitz constant, so \[\left\|H_x-\frac{m+M}{2}I\right\|\le\frac{M-m}{2}.\] Lean differentiates the actual centered gradient by subtracting the scalar identity derivative, then applies the checked derivative/\allowbreak{}Lipschitz operator bound. Subtracting the centered derivatives at two actual points yields \[\|H_x-H_y\|\le\frac{M-m}{2}+\frac{M-m}{2}=M-m.\] This is the sharp Hessian interval-width estimate. It is derived from the original potential assumptions and actual gradient derivatives, without an additional continuous-Hessian, symmetry, or operator-range hypothesis.

\subsection{\texorpdfstring{Identify and bound the actual random Hessian difference on L\ensuremath{{}^{2}}}{Identify and bound the actual random Hessian difference on L²}}


The existing random-Hessian construction is an actual continuous linear map on L\ensuremath{{}^{2}}, with representative \(H(X(\omega))Z(\omega)\). Lean proves that subtracting the two such maps has the exact representative \[(H(X(\omega))-H(Y(\omega)))Z(\omega).\] It rewrites the actual almost-everywhere representatives of the two random Hessians and their L\ensuremath{{}^{2}} difference. The pointwise sharp bound just proved gives \(\|(H(X)-H(Y))Z\|\le(M-m)\|Z\|\), so the true L\ensuremath{{}^{2}} operator norm of the difference is at most \(M-m\). Lean applies the actual L\ensuremath{{}^{2}} norm monotonicity theorem and the continuous linear operator-norm characterization; no operator bound is assumed on an unspecified random map.

\subsection{\texorpdfstring{Prove the exact Hessian-\allowbreak{}difference Gaussian truncation term}{Prove the exact Hessian-difference Gaussian truncation term}}


The original Hessian Lipschitz assumption also gives \(\|H(X)-H(Y)\|\le M_2\|X-Y\|\) pointwise. Combining this with the genuine uniform bound \(M-m\), and retaining an actual standard Gaussian velocity \(V\), gives \[\|(H(X)-H(Y))V\|_{L^2}\le M_2\sqrt\alpha\,\|X-Y\|_{L^2}+2(M-m)e^{-(\alpha-p)/8},\qquad\alpha\ge5p.\] Lean rewrites the exact random-Hessian difference representative and the actual position discrepancy representative, combines the two proved operator bounds by their minimum, and applies the checked Gaussian cutoff theorem. Only the actual marginal Gaussian law of \(V\) is required; the positions and Hessian difference may depend on \(V\). This is the precise random coefficient bound needed at each free-OU time in the corrected KLMC2 argument.

\section{\texorpdfstring{Sharp transformed Gaussian bounds for the repaired random Hessian correction}{Sharp transformed Gaussian bounds for the repaired random Hessian correction}}
\label{proof-module-KineticLangevin-FreeOUCorrectionBound-lean}
{\small\color{gray}Source module: \texttt{Kinetic\allowbreak{}Langevin/\allowbreak{}Free\allowbreak{}OUCorrection\allowbreak{}Bound.\allowbreak{}lean}.\par}


The actual \ensuremath{\varphi}\ensuremath{{}_{2}} and \ensuremath{\varphi}\ensuremath{{}_{3}} stochastic Hessian terms are estimated together. Their exact \ensuremath{\psi}\ensuremath{{}_{2}} cancellation gives the sharp quadratic time coefficient. Gaussian truncation then supplies the source-paper random-Hessian contribution without discarding fresh-noise terms or assuming independence from the Hessian.

\subsection{\texorpdfstring{Prove the sharp first and second primitive estimates}{Prove the sharp first and second primitive estimates}}


If \(\|U_r\|\le C\) on \([0,t]\), then \(\|\int_0^sU_r\,dr\|\le Cs\). Integrating this genuine bound yields \[\left\|\int_0^t\int_0^sU_r\,dr\,ds\right\|\le\frac{Ct^2}{2}.\] The same estimate holds with the outer factor \(e^{-\gamma(t-s)}\) for \(\gamma\ge0\), since that scalar lies between zero and one. Lean first applies the constant integral-norm bound on each actual inner interval, then bounds the outer integral by the integrable scalar function \(Cs\). The exact scalar primitive gives the factor \(1/2\). Every interval endpoint and nonnegative damping condition is checked explicitly.

\subsection{\texorpdfstring{Use the full stochastic cancellation in the actual phase metric}{Use the full stochastic cancellation in the actual phase metric}}


Let \(A\) be the actual continuous linear random-Hessian operator and suppose \(\|AR_0(r)\|_{L^2}\le C\). The complete correction identities give \(AC_1\) as a damped second primitive of \(AR_0\), so \(\|AC_1(t)\|\le Ct^2/2\). The exact identity \(C_1+\gamma C_2=R_2\), including both stochastic responses, gives \(\|AC_1+\gamma AC_2\|=\|AR_2\|\le Ct^2/2\). The transformed phase norm combines exactly these two coordinates, hence \[\|T_{\gamma,0}(AC_1,AC_2)\|\le\frac{\sqrt2}{2}t^2C.\] Lean rewrites the two actual primitive identities, applies the two sharp integral bounds, commutes \(A\) with the exact correction balance, and sums the two squared component bounds. It finally checks the square of the displayed nonnegative bound. This controls the entire implemented correction.

\subsection{\texorpdfstring{Insert the actual Gaussian cutoff at each free-\allowbreak{}OU time}{Insert the actual Gaussian cutoff at each free-OU time}}


If every actual free-OU velocity marginal is standard Gaussian and \(\|AR_0(r)(\omega)\|\le\min\{B,M_2\|X(\omega)\|\}\|R_0(r)(\omega)\|\), the proved truncation theorem gives the same bound at every time: \[\|AR_0(r)\|_{L^2}\le M_2\sqrt\alpha\,\|X\|_{L^2}+2B e^{-(\alpha-p)/8},\qquad\alpha\ge5p.\] Substituting this common \(C\) in the sharp transformed estimate retains both terms with their exact constants. Lean constructs the nonnegative common bound and applies the actual Gaussian truncation theorem at each time. No independence between the Hessian factor, the position discrepancy, and the velocity is required.

\subsection{\texorpdfstring{Verify the cutoff premise for a concrete random operator}{Verify the cutoff premise for a concrete random operator}}


For an actual measurable random operator field \(H(\omega)\), suppose its true operator norm satisfies both \(\|H(\omega)\|\le B\) and \(\|H(\omega)\|\le M_2\|X(\omega)\|\) almost everywhere. The checked random-operator construction gives an actual L\ensuremath{{}^{2}} continuous linear map with representative \(H(\omega)V(\omega)\). Lean rewrites that actual representative, applies the operator norm inequality, combines the two operator bounds by their minimum, and invokes Gaussian truncation for a velocity with the actual standard Gaussian law. Thus the required cutoff estimate follows from the genuine uniform and Lipschitz operator bounds of a Hessian difference.

\section{\texorpdfstring{Proposition 3 with the full stochastic Hessian difference}{Proposition 3 with the full stochastic Hessian difference}}
\label{proof-module-KineticLangevin-SecondOrderComparison-lean}
{\small\color{gray}Source module: \texttt{Kinetic\allowbreak{}Langevin/\allowbreak{}Second\allowbreak{}Order\allowbreak{}Comparison.\allowbreak{}lean}.\par}


The printed comparison omits the difference of the two Hessians multiplying the fresh phi2 and phi3 innovations. The formal comparison retains them. They combine with the initial reference velocity into the complete free-OU response, whose actual Gaussian marginals admit the cutoff estimate without independence of the Hessian.

\subsection{\texorpdfstring{Use the actual transformed phase geometry}{Use the actual transformed phase geometry}}


Negating both phase coordinates negates their transformed coordinates and preserves their Euclidean norm. The centered force sector, when squared and expanded, is exactly the ordinary strong-smooth interpolation inequality. This lets the already proved centered-gradient bound in L2 supply the interpolation hypothesis of the frozen-linear theorem.

\subsection{\texorpdfstring{Project the cutoff displacement with the exact coefficient}{Project the cutoff displacement with the exact coefficient}}


The Hessian cutoff contains M2 sqrt(alpha) times the position difference norm. The phase projection contributes sqrt(2)/\allowbreak{}gamma. The complete free-OU correction carries sqrt(2)t\textasciicircum{}2/\allowbreak{}2. Multiplication cancels sqrt(2) squared divided by two, leaving M2 sqrt(alpha)t\textasciicircum{}2/\allowbreak{}gamma times the initial transformed norm. The Gaussian tail retains sqrt(2)t\textasciicircum{}2 B exp(-(alpha-p)/\allowbreak{}8). This scalar lemma checks the cancellation separately from the probability-space construction.

\subsection{\texorpdfstring{Keep all four innovations on the same source}{Keep all four innovations on the same source}}


klmc2SourceInnovation applies the actual continuous linear Brownian-source lift to the four kernels psi0, psi1, phi2, phi3. klmc2SourceStep writes the actual update using the corresponding free-OU velocity and its complete Hessian responses. With the lift scaled by sqrt(2 gamma), direct linear algebra proves this is exactly Integrator.klmc2Step at the unscaled four-kernel innovations. These are actual coupled Gaussian-response variables.

\subsection{\texorpdfstring{Derive the corrected difference identity}{Derive the corrected difference identity}}


Subtract the two source steps at z and w with common innovations. The common psi noise cancels. Freeze the linear part at the Hessian of the second starting position; its data are the actual velocity, position and force differences. The remainder is minus the Hessian difference applied to the complete phi2 and phi3 free-OU corrections built from the first reference velocity. This exact identity includes both fresh stochastic Hessian terms absent from the printed algebra.

\subsection{\texorpdfstring{Supply the frozen-\allowbreak{}linear estimate from the actual potential}{Supply the frozen-linear estimate from the actual potential}}


The actual L2 gradient retains its Lipschitz constant and centered strong-smooth sector. The actual random Hessian has operator norm at most M. The frozen-linear map therefore contracts by 1-mt/\allowbreak{}(2 gamma) under gamma\textasciicircum{}2 at least m+M and 5 gamma M t at most m. The checked deterministic theorem is applied in the Hilbert space of actual square-integrable random vectors.

\subsection{\texorpdfstring{Combine the actual coupled terms by Minkowski}{Combine the actual coupled terms by Minkowski}}


The generic triangle lemma applies the exact difference identity, invariance of transformed norm under negation, and the Euclidean phase triangle inequality. Its two inputs are the separately proved frozen-linear component and the complete Hessian correction. This elementary composition introduces no stochastic approximation or omitted noise.

\subsection{\texorpdfstring{Prove the repaired source comparison inequality}{Prove the repaired source comparison inequality}}


Let A be the difference of the two actual random Hessian operators. Each actual free-OU velocity is standard Gaussian. The genuine Hessian difference bound M-m and Hessian Lipschitz constant M2 give the cutoff norm M2 sqrt(alpha) times the L2 position difference plus 2(M-m)exp(-(alpha-p)/\allowbreak{}8), for alpha at least 5p. Integrating the complete response yields the sharp quadratic time factor. Combining this with frozen-linear contraction and position projection proves transformed error at most [1-mt/\allowbreak{}(2 gamma)+M2 sqrt(alpha)t\textasciicircum{}2/\allowbreak{}gamma] times the initial transformed error, plus sqrt(2)t\textasciicircum{}2(M-m)exp(-(alpha-p)/\allowbreak{}8). The Gaussian premise is an actual source law and is discharged by the canonical past/\allowbreak{}future construction.

\section{\texorpdfstring{Proposition 2 with actual Gaussian and Gibbs moments}{Proposition 2 with actual Gaussian and Gibbs moments}}
\label{proof-module-KineticLangevin-SecondOrderLpLocal-lean}
{\small\color{gray}Source module: \texttt{Kinetic\allowbreak{}Langevin/\allowbreak{}Second\allowbreak{}Order\allowbreak{}Lp\allowbreak{}Local.\allowbreak{}lean}.\par}


This theorem combines the actual random-Hessian Taylor estimate with the genuine cubic/\allowbreak{}quartic defect decomposition. Its reference is the actual coupled L\ensuremath{{}^{2}} solution: exact mild and kinematic equations, with each time marginal equal to the Gaussian-times-Gibbs target. The moment and local-error inequalities are derived inside the proof.

\subsection{\texorpdfstring{Derive all local constants from target marginals}{Derive all local constants from target marginals}}


Project the actual phase pushforward law to obtain the standard Gaussian velocity marginal and the Gibbs position marginal. The Gaussian fourth moment gives the displacement Taylor forcing bound M\ensuremath{{}_{2}}\ensuremath{\surd}(p\ensuremath{{}^{2}}+2p)t\ensuremath{{}^{2}}/\allowbreak{}2. The Gibbs score theorem gives \ensuremath{\Vert}\ensuremath{\nabla}f(X(t))\ensuremath{\Vert}\ensuremath{{}_{2}}\ensuremath{\le}\ensuremath{\surd}(Mp), while the actual frozen Hessian operator has norm at most M. The generic local decomposition therefore has constant at most M\ensuremath{{}_{2}}\ensuremath{\surd}(p\ensuremath{{}^{2}}+2p)+M\ensuremath{\surd}(Mp). It yields velocity error Ct\ensuremath{{}^{3}}/\allowbreak{}6 and position error Ct\ensuremath{{}^{4}}/\allowbreak{}24. If \ensuremath{\gamma}t\ensuremath{\le}1/\allowbreak{}5, the actual transformed phase error is at most Ct\ensuremath{{}^{3}}/\allowbreak{}4. The conclusions refer to the corrected paths retaining the entire free-OU response; separate kernel identities identify those paths with the implemented KLMC2 formula.

\section{\texorpdfstring{The actual second-\allowbreak{}order one-\allowbreak{}step coupling recurrence}{The actual second-order one-step coupling recurrence}}
\label{proof-module-KineticLangevin-SecondOrderOneStep-lean}
{\small\color{gray}Source module: \texttt{Kinetic\allowbreak{}Langevin/\allowbreak{}Second\allowbreak{}Order\allowbreak{}One\allowbreak{}Step.\allowbreak{}lean}.\par}


These results connect the implemented four-kernel KLMC2 update to the true stationary kinetic segment and combine the genuine local and repaired comparison estimates. Every moment and error constant has already been derived from actual path equations and Gaussian/\allowbreak{}Gibbs marginals.

\subsection{\texorpdfstring{Recover the local friction restriction from the source step size}{Recover the local friction restriction from the source step size}}


The restriction 5 gamma M h at most m, together with m at most M and M positive, implies 5 gamma h at most one. Cancelling the positive M gives gamma h at most one fifth. This is precisely the small-time restriction used to convert the cubic and quartic local component estimates into the transformed quarter-bound.

\subsection{\texorpdfstring{Bound the local error of the actual implemented step}{Bound the local error of the actual implemented step}}


Use the actual corrected-step identity for all four correlated Brownian kernels. The initial velocity and fresh filtered noise define the complete free-OU path. The actual stationary reference has the exact mild and kinematic equations, with Gaussian-times-Gibbs time marginals. The proved local theorem therefore bounds the transformed norm of the implemented step minus the exact endpoint by t\textasciicircum{}3 times [M2 sqrt(p\textasciicircum{}2+2p)+M sqrt(Mp)] divided by four. The exact fourth-moment dimension remains valid in every dimension.

\subsection{\texorpdfstring{Combine local defect and repaired synchronous comparison}{Combine local defect and repaired synchronous comparison}}


Insert the numerical update starting at the stationary reference into the difference between the actual numerical update and exact reference endpoint. The transformed triangle inequality splits this into comparison of two actual numerical steps and the actual reference local defect. Proposition 3 controls the former, including both fresh stochastic Hessian terms. The preceding local theorem controls the latter. The result is the source affine one-step inequality with factor 1-mt/\allowbreak{}(2 gamma)+M2 sqrt(alpha)t\textasciicircum{}2/\allowbreak{}gamma, Gaussian tail sqrt(2)t\textasciicircum{}2(M-m)exp(-(alpha-p)/\allowbreak{}8), and cubic local term. Its remaining path and law premises are supplied by the constructed canonical segment in the final sampler theorem.

\section{\texorpdfstring{The actual canonical KLMC2 one-\allowbreak{}step error recurrence}{The actual canonical KLMC2 one-step error recurrence}}
\label{proof-module-KineticLangevin-CanonicalKLMC2Recurrence-lean}
{\small\color{gray}Source module: \texttt{Kinetic\allowbreak{}Langevin/\allowbreak{}Canonical\allowbreak{}KLMC2\allowbreak{}Recurrence.\allowbreak{}lean}.\par}


The generic second-order one-step theorem is instantiated with the genuine canonical numerical chain and the stationary reference segment driven by the same Brownian interval. Every source, Gaussian, stationarity, mild-equation and kinematic premise is discharged by earlier proved constructions. The resulting recurrence is about the original implemented sampler.

\subsection{\texorpdfstring{Identify the exact four-\allowbreak{}noise input}{Identify the exact four-noise input}}


The canonical innovation is the Brownian source lift of each of the four deterministic convolution kernels. Unfolding this definition and the analytic lifted-noise definition proves exact equality, since converting the nonnegative numerical step to a nonnegative-real duration does not change it.

\subsection{\texorpdfstring{Discharge every premise using actual reference and numerical paths}{Discharge every premise using actual reference and numerical paths}}


The reference segment is continuous, starts at the current reference state and ends at its next state. Its actual Gaussian-times-Gibbs phase law, free-OU mild equation and position primitive are proved internally. Past/\allowbreak{}future independence identifies the free-OU velocity law as standard Gaussian throughout the same fresh interval. Applying the checked generic KLMC2 one-step theorem therefore gives its exact source coefficient, Gaussian cutoff tail and cubic stationary local defect. Equality of the four noises and the numerical successor formula identify its left side with the actual next numerical/\allowbreak{}reference error. Symmetry of the transformed norm reverses the current reference-minus-numerical discrepancy to the chosen numerical-minus-reference convention.

\section{\texorpdfstring{Explicit accuracy budgets and integer iteration counts}{Explicit accuracy budgets and integer iteration counts}}
\label{proof-module-KineticLangevin-Complexity-lean}
{\small\color{gray}Source module: \texttt{Kinetic\allowbreak{}Langevin/\allowbreak{}Complexity.\allowbreak{}lean}.\par}


The paper's step-size and iteration formulas are scalar consequences of its numerical error estimate. This module verifies that implication, retains a stability-compatible bias budget, and rounds the logarithmic threshold to a natural number. It handles zero initial error separately and clips negative logarithmic thresholds at zero. These accuracy corollaries have an explicit error inequality as input; they do not replace the outstanding unconditional stochastic theorem.

\subsection{\texorpdfstring{Replace geometric decay by its exponential envelope}{Replace geometric decay by its exponential envelope}}


For \(\alpha\le1\), the base \(1-\alpha\) is nonnegative. The elementary exponential inequality gives \(1-\alpha\le e^{-\alpha}\). Raising this comparison to a natural power preserves order, so \((1-\alpha)^n\le e^{-\alpha n}\). Lean uses \texttt{Real.\allowbreak{}add\_\allowbreak{}one\_\allowbreak{}le\_\allowbreak{}exp}, \texttt{pow\_\allowbreak{}le\_\allowbreak{}pow\_\allowbreak{}left\ensuremath{{}_{0}}}, and \texttt{Real.\allowbreak{}exp\_\allowbreak{}nat\_\allowbreak{}mul}. The lemma itself needs no positive lower bound on \(\alpha\); that requirement enters when dividing the logarithmic threshold.

\subsection{\texorpdfstring{Round the logarithmic threshold and include small initial errors}{Round the logarithmic threshold and include small initial errors}}


For initial coupling radius \(D\), define

\[n_0=\left\lceil\max\left\{0,\frac{\log(24D/\varepsilon)}\alpha\right\}\right\rceil.\]

This is a natural number even when the logarithm is negative. If \(\alpha>0\) and \(n\ge n_0\), then \(\alpha n\ge\log(24D/\varepsilon)\). Lean combines the ceiling's lower bound, the right-hand inequality for a maximum, and the cast of the natural comparison, before multiplying by the positive \(\alpha\). The explicit maximum and ceiling handle the initialization-dependent small-error regime and integer rounding.

\subsection{\texorpdfstring{Verify the transient budget only when logarithmic cancellation is valid}{Verify the transient budget only when logarithmic cancellation is valid}}


For \(D>0\) and \(\varepsilon>0\), exponentiating the iteration inequality gives \(e^{-\alpha n}\le e^{-\log(24D/\varepsilon)}=\varepsilon/(24D)\). The geometric envelope and multiplication by \(D\) imply \((1-\alpha)^nD\le\varepsilon/24\). Lean proves positivity of the logarithm argument before using \texttt{Real.\allowbreak{}exp\_\allowbreak{}log}, rewrites the negative exponential as an inverse, and cancels the nonzero factors. Zero initial error is handled by the next theorem without a logarithmic cancellation.

\subsection{\texorpdfstring{Combine the manuscript's 0.\allowbreak{}94 bias allocation with the transient}{Combine the manuscript's 0.94 bias allocation with the transient}}


Suppose an explicit estimate gives \(E_n\le\sqrt2(1-\alpha)^nD+b\), with \(0<\alpha\le1\), \(D\ge0\), and \(b\le0.94\varepsilon\). When \(D=0\), the bias already gives the conclusion and there is no iteration requirement. Otherwise use the rounded threshold. The exact rational inequality \(\sqrt2\le36/25\) bounds the remaining contribution by \(0.06\varepsilon\), so the total error is at most \(\varepsilon\).



Lean proves the square-root comparison by squaring, splits on zero versus positive initial radius, derives the logarithmic condition in the positive case, and combines the inequalities with nonnegative scalar multiplication. For KLMC, substitute \(\alpha=3mh/(4\gamma)\) and bias \(Mh\sqrt{2p}/m\), with a step size that is the minimum of the stability threshold and bias allowance. The same scalar argument applies to KIPLMC1 with its thermal radius and a separately budgeted concentration term.

\section{\texorpdfstring{An explicit epsilon guarantee for the actual KLMC sampler}{An explicit epsilon guarantee for the actual KLMC sampler}}
\label{proof-module-KineticLangevin-Accuracy-lean}
{\small\color{gray}Source module: \texttt{Kinetic\allowbreak{}Langevin/\allowbreak{}Accuracy.\allowbreak{}lean}.\par}


The convergence theorem becomes an executable mathematical budget: choose a positive step, an integer iteration count, and obtain a certified error at most \(\varepsilon\). The result concerns the actual constructed sampler. Neither a recurrence nor a local-error estimate is left as a premise.

\subsection{\texorpdfstring{Verify that the geometric factor is admissible}{Verify that the geometric factor is admissible}}


Set \(\alpha=3mh/(4\gamma)\). The curvature, friction and step assumptions imply \(0<\alpha\le1\). The proof reuses the checked exponential-coefficient estimate: \(e^{-mh/\gamma}+(3/4)Mh^2\le1-\alpha\). Its left side is positive, so the right side is nonnegative. This justifies every later geometric and logarithmic estimate.

\subsection{\texorpdfstring{Apply the scalar error allocation to the actual law}{Apply the scalar error allocation to the actual law}}


The actual KLMC theorem gives \(W_2(\nu_n,\pi)\le\sqrt2(1-\alpha)^nD+b\). The checked scalar budget allocates \(b\le0.94\varepsilon\) and \((1-\alpha)^nD\le\varepsilon/24\). The sum is at most \(\varepsilon\). The threshold is the natural ceiling of \(\max\{0,\log(24D/\varepsilon)/\alpha\}\); if \(D=0\), no initial-error iterations are needed. The finite-transport hypothesis prevents the real-valued Wasserstein convention from hiding an infinite distance.

\subsection{\texorpdfstring{Choose the paper step and round the count}{Choose the paper step and round the count}}


For positive dimension choose \(h=\min\{m/(4\gamma M),\,0.94\varepsilon m/(M\sqrt{2p})\}\). Both terms are positive. The first enforces stability and the second enforces the bias budget by multiplication through positive denominators. Construct the corresponding nonnegative-real Lean step, use the rounded logarithmic count, and apply the actual accuracy theorem. At \(\gamma=\sqrt{m+M}\), this is precisely the manuscript’s minimum-step prescription. The explicit reciprocal dependence displays the \(\sqrt p/\varepsilon\) iteration scaling up to the logarithm.

\section{\texorpdfstring{Absorb the Hessian cutoff and accumulate the second-\allowbreak{}order error}{Absorb the Hessian cutoff and accumulate the second-order error}}
\label{proof-module-KineticLangevin-SecondOrderRecurrence-lean}
{\small\color{gray}Source module: \texttt{Kinetic\allowbreak{}Langevin/\allowbreak{}Second\allowbreak{}Order\allowbreak{}Recurrence.\allowbreak{}lean}.\par}


The repaired local comparison has a random-Hessian term proportional to \(M_2\sqrt\alpha h^2\). Choosing the cutoff \(\alpha=(m/(4M_2h))^2\) absorbs it into the contraction gap. This module proves the choice, both Gaussian tails, the affine accumulation, the position projection, and the exact dimension range for the manuscript’s final simplification.

\subsection{\texorpdfstring{Choose the cutoff and retain both tail estimates}{Choose the cutoff and retain both tail estimates}}


For \(m,M_2,h>0\), the positive square root of \(\alpha=(m/(4M_2h))^2\) is \(m/(4M_2h)\). Thus \(M_2\sqrt\alpha h=m/4\). The source dimension restriction \(4\sqrt{5p}M_2h\le m\) implies \(\alpha\ge5p\). It follows that \((\alpha-p)/8\ge p/2\), giving the printed tail. Also \(p\le\alpha/5\) implies \((\alpha-p)/8\ge\alpha/10=m^2/(160M_2^2h^2)\), giving the stronger source footnote. Lean proves the square-root identity, multiplies through positive denominators, compares nonnegative squares, and applies monotonicity of the exponential.

\subsection{\texorpdfstring{Sum the actual one-\allowbreak{}step comparison and project to position}{Sum the actual one-step comparison and project to position}}


The local bound has coefficient \(1-mh/(2\gamma)+M_2\sqrt\alpha h^2/\gamma\). The cutoff makes this at most \(q=1-mh/(4\gamma)\). The additive local terms are \(\sqrt2h^2B e^{-(\alpha-p)/8}+h^3C/4\), where \(B=M-m\) and \(C=M_2\sqrt{p^2+2p}+M\sqrt{Mp}\). The checked affine recurrence divides this budget by \(1-q=mh/(4\gamma)\). Position projection multiplies by \(\sqrt2/\gamma\), and equal initial velocities give \(e_0\le\gamma d_0\). The exact resulting bias is \(\sqrt2h^2C/m+(8B/m)h e^{-(\alpha-p)/8}\), with the desired transient \(\sqrt2q^nd_0\). The recurrence is a reusable scalar lemma; the actual chain endpoint supplies its one-step inequality separately.

\subsection{\texorpdfstring{State the dimension restriction explicitly}{State the dimension restriction explicitly}}


For \(p\ge2\), \(p^2+2p\le2p^2\), hence \(\sqrt2\sqrt{p^2+2p}\le2p\). Distributing this through \(C\) gives the printed Hessian term \(2M_2p\), while \(\sqrt2\sqrt{Mp}=\sqrt{2Mp}\) gives the score term. The supplied proof’s simplification does not justify the same expression at \(p=1\). The general theorem therefore retains the exact fourth-moment expression there; the printed simplification is an explicit \(p\ge2\) corollary. Lean compares nonnegative squares and never uses the invalid inequality in dimension one.

\section{\texorpdfstring{The actual KLMC2 chain: a bound for every initial transport plan}{The actual KLMC2 chain: a bound for every initial transport plan}}
\label{proof-module-KineticLangevin-CanonicalKLMC2Bound-lean}
{\small\color{gray}Source module: \texttt{Kinetic\allowbreak{}Langevin/\allowbreak{}Canonical\allowbreak{}KLMC2\allowbreak{}Bound.\allowbreak{}lean}.\par}


The reference is the constructed kinetic diffusion started from the normalized Gibbs target, with an independent Gaussian velocity and canonical Brownian source. The numerical chain is the implemented four-kernel KLMC2 update on the same source. All local equations, Gaussian marginals and stationary reference laws have been proved before the scalar recurrence is applied.

\subsection{\texorpdfstring{Keep the accumulated contraction factor nonnegative}{Keep the accumulated contraction factor nonnegative}}


The source restrictions imply gamma h at most one fifth and m at most gamma squared. Hence mh/\allowbreak{}(4 gamma) is at most one twentieth, so q=1-mh/\allowbreak{}(4 gamma) is nonnegative. This elementary check is needed when comparing the actual initial error and projecting the iterated recurrence.

\subsection{\texorpdfstring{Iterate the actual canonical coupling error}{Iterate the actual canonical coupling error}}


For every finite-cost initial position coupling, construct the actual numerical and reference chains with shared initial Gaussian velocity and shared canonical Brownian intervals. Let e(k) be the actual transformed L2 phase discrepancy and d(k) the actual L2 position discrepancy. The native canonical one-step theorem supplies the repaired local recurrence, including the fresh stochastic Hessian terms. The cutoff condition M2 sqrt(alpha)h at most m/\allowbreak{}4 absorbs its position-dependent correction into q. The checked affine recurrence and phase projection then give d(n) at most sqrt(2)q\textasciicircum{}n times the initial position cost, plus sqrt(2)h\textasciicircum{}2 C/\allowbreak{}m and 8(M-m)h exp(-(alpha-p)/\allowbreak{}8)/\allowbreak{}m. Here C=M2 sqrt(p\textasciicircum{}2+2p)+M sqrt(Mp). Equal initial velocities identify the initial transformed norm exactly, and the constructed initial-extension variables identify the initial L2 position norm with the supplied transport plan cost. There is no stationary-law, local-error, Brownian-independence or numerical-law premise in this final per-coupling theorem. Taking the infimum over actual initial couplings is handled by the subsequent Wasserstein endpoint.

\section{\texorpdfstring{The actual KLMC2 Wasserstein theorem and its source corrections}{The actual KLMC2 Wasserstein theorem and its source corrections}}
\label{proof-module-KineticLangevin-CanonicalKLMC2-lean}
{\small\color{gray}Source module: \texttt{Kinetic\allowbreak{}Langevin/\allowbreak{}Canonical\allowbreak{}KLMC2.\allowbreak{}lean}.\par}


This is the complete second-order sampling endpoint for the actual implemented four-kernel Gaussian chain. The continuous reference, its true Gibbs law and Gaussian moments, fresh-noise independence, local Taylor and memory error, stochastic Hessian comparison, and output-law identities are proved internally. Taking the true initial-coupling infimum yields the desired initial Wasserstein term.

\subsection{\texorpdfstring{Take the true transport infimum with the exact dimension term}{Take the true transport infimum with the exact dimension term}}


For every finite-cost initial plan, the checked canonical chain theorem bounds its actual L\ensuremath{{}^{2}} position discrepancy from a stationary reference. The output law is independent of that plan. The reference endpoint has the actual Gibbs law. The L\ensuremath{{}^{2}} transport-cost inequality therefore gives the same bound for the actual output Wasserstein distance. A finite-output theorem excludes an infinite distance before passing to real values. The all-couplings affine-infimum lemma then replaces each plan’s cost by the true initial \(W_2\). With \(p=|\iota|\), \(q=1-mh/(4\gamma)\), \(C_p=M_2\sqrt{p^2+2p}+M\sqrt{Mp}\), and \(\alpha=(m/(4M_2h))^2\), the result is \(W_2(\nu_n,\pi)\le\sqrt2q^nW_2(\nu_0,\pi)+\sqrt2h^2C_p/m+8(M-m)h e^{-(\alpha-p)/8}/m\). The exact Gaussian fourth moment remains valid for every finite coordinate set, including dimension one.

\subsection{\texorpdfstring{Retain the source footnote’s sharper tail}{Retain the source footnote’s sharper tail}}


The source dimension restriction implies \(\alpha\ge5p\). Hence \((\alpha-p)/8\ge\alpha/10=m^2/(160M_2^2h^2)\). Multiply this exponential comparison by the nonnegative actual tail coefficient and apply the exact sampler bound. This gives the genuine chain theorem with the sharper inverse-step-square Gaussian tail, not merely a scalar observation detached from the algorithm.

\subsection{\texorpdfstring{Recover the printed bound in its valid dimension range}{Recover the printed bound in its valid dimension range}}


For \(p\ge2\), \(\sqrt2\sqrt{p^2+2p}\le2p\), and \(\sqrt2\sqrt{Mp}=\sqrt{2Mp}\). The local bias becomes \(2M_2ph^2/m+Mh^2\sqrt{2Mp}/m\). Also \(M-m\le M\) and \(e^{-(\alpha-p)/8}\le e^{-p/2}\). These comparisons produce the supplied Theorem 3 formula. Its proof’s hidden dimension restriction is stated explicitly. In dimension one the preceding exact-moment theorem supplies the valid full result. The underlying comparison includes both stochastic Hessian-difference terms omitted by the supplied displayed subtraction; those terms were repaired through the stationary free-OU bundle.

\section{\texorpdfstring{Correcting the Hessian condition-\allowbreak{}number refinement}{Correcting the Hessian condition-number refinement}}
\label{proof-module-KineticLangevin-ConditionNumber-lean}
{\small\color{gray}Source module: \texttt{Kinetic\allowbreak{}Langevin/\allowbreak{}Condition\allowbreak{}Number.\allowbreak{}lean}.\par}


The supplied manuscript informally simplifies its tail-free KLMC2 step/\allowbreak{}count budget using \(\kappa_2=(M_2^{2/3}+Mp^{-1/3})/m\). This module proves a valid refinement with \(\max(\kappa,\kappa_2)\), identifies the extra regime in which the printed shape follows, and constructs an exact natural-dimension family disproving a uniform removal of the ordinary condition number. Every scalar quantity is tied to the actual source definition. These are statements about the explicit budget prescription, with its Gaussian tail intentionally excluded; the separate actual-sampler accuracy theorems include that tail.

\subsection{\texorpdfstring{Keep both condition numbers and normalize the actual source definition}{Keep both condition numbers and normalize the actual source definition}}


The ordinary gradient condition number is \(\kappa=M/m\). The supplied source proposes \(\rho=(M_2^{2/3}+Mp^{-1/3})/m\). Introduce the dimensionless Hessian ratio \(\eta=M_2/(m\sqrt m)\). The first proof establishes the exact identity \(\rho=\eta^{2/3}+\kappa p^{-1/3}\), rather than assuming a convenient relationship between the two quantities. Positivity of \(m\) permits the real-power product and quotient rules: \((m\sqrt m)^{2/3}=m\). Lean rewrites the square root as a half power, combines exponents, distributes the quotient power, and closes the remaining fraction identity with \texttt{ring}.

\subsection{\texorpdfstring{Derive the useful cubic inequalities from the definition}{Derive the useful cubic inequalities from the definition}}


For \(\kappa>0\), \(\eta\ge0\), and \(p>0\), the source quantity \(\rho\) is positive. Each nonnegative summand is at most their sum. Cubing \(\eta^{2/3}\le\rho\) gives \(\eta^2\le\rho^3\). Cubing \(\kappa p^{-1/3}\le\rho\), and using \((p^{-1/3})^3=1/p\), gives \(\kappa^3\le p\rho^3\). These estimates justify the later dimension scaling; they do not imply \(\kappa\le\rho\). The proof uses monotonicity of natural powers on nonnegative numbers, then multiplies through the positive dimension.

\subsection{\texorpdfstring{Record the exact reciprocal-\allowbreak{}step coefficient before simplifying it}{Record the exact reciprocal-step coefficient before simplifying it}}


Choose \(\gamma=\sqrt{m+M}\), put \(s=h\sqrt m\), and target error \(\varepsilon/\sqrt m\). The two stability restrictions and the source quadratic-bias allocation give

\[s^{-1}=\max\left\{5\kappa\sqrt{1+\kappa},\ 4\sqrt{5p}\eta,\ \sqrt{\frac{4\eta p+2\sqrt2\kappa\sqrt\kappa\sqrt p}{\varepsilon}}\right\}.\]

The transient count has coefficient \(4\sqrt{1+\kappa}/s\), before its logarithm and integer rounding. \texttt{second\allowbreak{}Order\allowbreak{}Normalized\allowbreak{}Work} keeps this entire expression, including the constants. The separate \texttt{second\allowbreak{}Order\allowbreak{}Refinement} records the source's proposed shape \(\sqrt p\rho^2+\sqrt{p/\varepsilon}\rho^{5/4}\). This section concerns the source's expressly tail-free comparison; the actual sampler's separate epsilon theorem controls the Gaussian tail as well.

\subsection{\texorpdfstring{Bound the three branches with one explicit universal constant}{Bound the three branches with one explicit universal constant}}


Suppose \(1\le\kappa\le R\), \(p\ge1\), and \(\eta^2\le R^3\). Comparing nonnegative squares gives \(\eta\le R\sqrt R\), while \(\sqrt p\le p\) and \(\kappa\sqrt\kappa\le R\sqrt R\). Thus the numerator of the accuracy branch is at most \(8pR\sqrt R\). Let \(A=\sqrt pR^2\) and \(B=\sqrt{p/\varepsilon}R^{5/4}\). Their exact squares are \(A^2=pR^4\) and \(B^2=(p/\varepsilon)R^2\sqrt R\).



After multiplying by the common count factor, the ordinary stability branch is at most \(40A\). The Hessian stability branch has squared value \(1280p(1+\kappa)\eta^2\le2560pR^4\le(64A)^2\). The accuracy branch has square at most \(256(p/\varepsilon)R^2\sqrt R=(16B)^2\). All quantities are nonnegative, so the square comparisons recover the original inequalities. Taking their maximum gives the checked bound \(64(A+B)\). The Lean proof expands these squares, applies order-preserving multiplication, and uses the real-power identity \((R^{5/4})^2=R^2\sqrt R\).

\subsection{\texorpdfstring{Apply the bound to the source parameters without adding a false comparison}{Apply the bound to the source parameters without adding a false comparison}}


Set \(R=\max\{\kappa,\rho\}\), where \(\rho\) is the actual source quantity. Then \(R\ge\kappa\), and the proved cubic estimate gives \(\eta^2\le\rho^3\le R^3\). The previous theorem therefore yields

\[\mathcal W(\kappa,\eta,p,\varepsilon)\le64\left[\sqrt p\,\max(\kappa,\rho)^2+\sqrt{p/\varepsilon}\,\max(\kappa,\rho)^{5/4}\right].\]

The physical-parameter wrapper substitutes \(\kappa=M/m\), \(\eta=M_2/(m\sqrt m)\), and the literal source definition of \(\rho\). Its assumptions are positive strong convexity, \(m\le M\), a nonnegative Hessian Lipschitz bound, positive accuracy, and \(p\ge1\). No comparison between the condition numbers is hidden in those assumptions.

\subsection{\texorpdfstring{Identify the additional regime in which the printed shape is valid}{Identify the additional regime in which the printed shape is valid}}


If the extra inequality \(\kappa\le\rho\) holds, the maximum in the corrected theorem equals \(\rho\), so the source's proposed shape is a valid upper bound with universal constant 64. The physical condition is exactly

\[M(1-p^{-1/3})\le M_2^{2/3}.\]

The equivalence follows by multiplying the original condition-number inequality by \(m>0\) and rearranging. This is a substantive extra regime, not a consequence of strong convexity or smoothness. The final wrapper makes that condition explicit and specializes the physical max theorem.

\subsection{\texorpdfstring{Construct an exact family exposing the missing ordinary condition number}{Construct an exact family exposing the missing ordinary condition number}}


Take \(t>0\), \(\kappa=t^2\), \(\eta=1\), \(p=t^6\), and \(\varepsilon=1\). The exact fractional-power rules give \((t^6)^{-1/3}=t^{-2}\), so the source condition number is identically \(\rho=2\). Meanwhile the ordinary stability term alone implies \(\mathcal W\ge20\kappa^2=20t^4\). Its proposed simplification is at most \(8t^3\): \(\sqrt{t^6}=t^3\), \(2^{5/4}\le2^2=4\), and the two terms each cost at most \(4t^3\).



The proof therefore compares exact finite expressions, without relying on informal asymptotic notation. This is a counterexample to a scalar step/\allowbreak{}count simplification. It does not establish a lower bound on the possible performance of the stochastic algorithm.

\subsection{\texorpdfstring{Rule out every uniform multiplicative repair of the printed formula}{Rule out every uniform multiplicative repair of the printed formula}}


For any \(C>0\), choose a natural number \(n>C+1\). The preceding family has genuine integer dimension \(n^6\). Its proposed expression times \(C\) is at most \(8Cn^3\), which is strictly smaller than \(20n^4\), a lower bound on the actual reciprocal-step prescription. Hence no dimension-uniform constant \(C\) can turn the printed expression into a bound for this prescription over all admissible parameters.



The last theorem states the same contradiction with the literal physical definition \(\rho=(M_2^{2/3}+Mp^{-1/3})/m\), using \(m=M_2=1\), \(M=n^2\), and \(p=n^6\). The normalization identity identifies this physical expression with the already checked family. The conclusion concerns algebraic certification budgets; no stochastic complexity lower bound is asserted.

\section{\texorpdfstring{Exact tail-\allowbreak{}free KLMC2 step and rounded complexity}{Exact tail-free KLMC2 step and rounded complexity}}
\label{proof-module-KineticLangevin-SecondOrderComplexity-lean}
{\small\color{gray}Source module: \texttt{Kinetic\allowbreak{}Langevin/\allowbreak{}Second\allowbreak{}Order\allowbreak{}Complexity.\allowbreak{}lean}.\par}


The manuscript's final complexity paragraph first drops its exponentially small persistent tail and chooses the remaining quadratic bias at normalized precision \(\varepsilon/\sqrt m\). This chapter formalizes that scalar prescription exactly. It retains every active step restriction as a minimum, computes its reciprocal as a maximum, and rounds the logarithmic time threshold to a natural iteration count. These identities do not replace the actual sampler's full-tail accuracy theorem, which uses \texttt{second\allowbreak{}Order\allowbreak{}Accuracy\allowbreak{}Step} and keeps the proved persistent floor. The separate condition-number chapter analyzes the manuscript's further informal simplification.

\subsection{\texorpdfstring{Choose the exact source minimum at normalized precision}{Choose the exact source minimum at normalized precision}}


Let \(A=2M_2p+M\sqrt{2Mp}\). The two printed quadratic terms add to \(A h^2/m\). To allocate at most half of the normalized target \(\varepsilon/\sqrt m\) to these terms, impose \(h^2\le\varepsilon\sqrt m/(2A)\). The source also requires \(5\gamma Mh\le m\) and \(4\sqrt{5p}M_2h\le m\). Thus define \[h_0=\min\!\left\{\frac{m}{5\gamma M},\;\frac{m}{4\sqrt{5p}M_2},\;\sqrt{\frac{\varepsilon\sqrt m}{4M_2p+2M\sqrt{2Mp}}}\right\}.\] With positive \(m,M,\gamma,M_2,p,\varepsilon\), all three entries are positive, so \(h_0>0\). Lean projects the three inequalities from the minimum, squares the last one using positivity, and multiplies by its positive denominator. It then rewrites the quadratic sum as \((4M_2p+2M\sqrt{2Mp})h_0^2/(2m)\), obtaining exactly \(2h_0^2M_2p/m+h_0^2M\sqrt{2Mp}/m\le\varepsilon/(2\sqrt m)\). This statement concerns the printed quadratic terms and imposes no unproved small-tail conclusion.

\subsection{\texorpdfstring{Invert the minimum without losing a parameter regime}{Invert the minimum without losing a parameter regime}}


For two positive numbers, inversion reverses order, so \((\min(a,b))^{-1}=\max(a^{-1},b^{-1})\). Lean proves this by the two possible order cases, applying the reciprocal order inequality in each case. Apply it twice to the three entries of \(h_0\). The inverse square-root term is simplified with the exact identity \((\sqrt x)^{-1}=\sqrt{x^{-1}}\). The result is \[h_0^{-1}=Q:=\max\!\left\{\frac{5\gamma M}{m},\;\frac{4\sqrt{5p}M_2}{m},\;\sqrt{\frac{4M_2p+2M\sqrt{2Mp}}{\varepsilon\sqrt m}}\right\}.\] Each branch has a distinct role: numerical stability, the Hessian cutoff restriction, or the quadratic bias budget. Keeping the maximum makes the statement valid when any of them dominates.

\subsection{\texorpdfstring{Recover the exact reciprocal-\allowbreak{}step prefactor}{Recover the exact reciprocal-step prefactor}}


Immediately before introducing its refined condition-number shorthand, the manuscript writes the work scale as \(\sqrt M/(m h_\varepsilon)\). The formal equality is \[\frac{\sqrt M}{m h_0}=\frac{\sqrt M}{m}\,Q.\] Its proof substitutes the reciprocal identity and normalizes the field expression. This is an exact identity for the explicit selected step. It keeps all parameters and contains no big-O convention or hidden uniformity assumption. When \(\gamma=\sqrt{m+M}\) and \(m\le M\), the actual geometric iteration prefactor \(4\gamma/(m h_0)\) differs from this displayed scale only by the familiar universal factors relating \(\sqrt{m+M}\) and \(\sqrt M\).

\subsection{\texorpdfstring{Connect the exact step to the dimensionless condition-\allowbreak{}number theorem}{Connect the exact step to the dimensionless condition-number theorem}}


Set \(\kappa=M/m\), \(\eta=M_2/(m\sqrt m)\), and choose the least admissible source friction \(\gamma=\sqrt{m+M}\). The exact normalization is \[\frac{4\sqrt{m+M}}{m}\,Q=\operatorname{Work}(\kappa,\eta,p,\varepsilon),\] where \texttt{second\allowbreak{}Order\allowbreak{}Normalized\allowbreak{}Work} is the concrete maximum analyzed in the condition-number chapter. Lean proves \(\sqrt{m+M}=\sqrt m\sqrt{1+\kappa}\), then factors \(\sqrt m\) from every branch of \(Q\). For the accuracy branch it uses \(\sqrt{2Mp}=\sqrt2\sqrt M\sqrt p\) and \(\sqrt{M/m}=\sqrt M/\sqrt m\); its radicand becomes \(m(4\eta p+2\sqrt2\kappa\sqrt\kappa\sqrt p)/\varepsilon\). Multiplication by the outer rate factor cancels \(m=(\sqrt m)^2\). This is the semantic bridge between the selected scalar step and the corrected condition-number comparison: the two chapters bound the same quantity, without an additional asymptotic premise.

\subsection{\texorpdfstring{Round logarithmic time to an actual integer count}{Round logarithmic time to an actual integer count}}


The existing iteration count is \(N(\alpha,D,\delta)=\lceil\max(0,\log(24D/\delta)/\alpha)\rceil\), with a natural-number ceiling. For \(\alpha>0\), division commutes with the maximum against zero. The ceiling inequality therefore gives \[N(\alpha,D,\delta)<1+\frac{\max\{0,\log(24D/\delta)\}}{\alpha}.\] The maximum is essential when the initial discrepancy is already small; it prevents a negative logarithmic threshold from being interpreted as a negative iteration count. Substitute \(\alpha=mh_0/(4\gamma)\), \(\delta=\varepsilon/\sqrt m\), and \(h_0^{-1}=Q\) to obtain the exact bound \[N\!\left(\frac{mh_0}{4\gamma},D,\frac{\varepsilon}{\sqrt m}\right)<1+\frac{4\gamma}{m}\,Q\,\max\!\left\{0,\log\frac{24D\sqrt m}{\varepsilon}\right\}.\] Lean explicitly rewrites the normalized logarithm argument, applies the positive-rate ceiling estimate, and uses the proved reciprocal identity. In an application \(D\) is the initial Wasserstein discrepancy. This scalar upper bound does not discard any term from the separate full-tail sampler theorem.

\subsection{\texorpdfstring{Check the friction choice directly}{Check the friction choice directly}}


At fixed positive step and nonnegative initial error, increasing positive friction decreases \(mh/(4\gamma)\), hence increases the geometric factor \((1-mh/(4\gamma))^n\). Lean proves the reciprocal inequality for the two friction values, uses the lower-friction stability condition to keep the base nonnegative, and applies monotonicity of a natural power. Thus the transient upper bound is smallest at the smallest admissible friction. For the explicit minimum prescription, increasing friction decreases its stability branch \(m/(5\gamma M)\), while leaving the cutoff and bias branches unchanged; the chosen step can only decrease. This proves the two concrete monotonicity statements behind the manuscript's choice of the least permitted \(\gamma\).

\subsection{\texorpdfstring{Bound the actual rounded count by the corrected physical condition numbers}{Bound the actual rounded count by the corrected physical condition numbers}}


The final corollary applies to the very natural-number count selected by the explicit step \(h_0\), with \(\gamma=\sqrt{m+M}\). Assume \(m>0\), \(m\le M\), \(M_2>0\), \(p\ge1\), and \(\varepsilon>0\). Write \(\kappa=M/m\), \(\kappa_2=(M_2^{2/3}+Mp^{-1/3})/m\), and \(R=\max(\kappa,\kappa_2)\). The proved integer-count bound becomes \[N\!\left(\frac{mh_0}{4\sqrt{m+M}},D,\frac{\varepsilon}{\sqrt m}\right)<1+64\left[\sqrt p\,R^2+\sqrt{p/\varepsilon}\,R^{5/4}\right]\max\!\left\{0,\log\frac{24D\sqrt m}{\varepsilon}\right\}.\] Lean first uses the exact ceiling-count theorem. It rewrites the work prefactor through the normalization identity, then invokes \texttt{second\allowbreak{}Order\allowbreak{}Normalized\allowbreak{}Work\_\allowbreak{}physical\_\allowbreak{}bound}, whose hypotheses follow from these physical parameters. Multiplication by the nonnegative maximum-logarithm preserves that upper bound. This assembles the explicit step, the rounded integer count, and the corrected condition-number refinement into one theorem; it does not assume a comparison between the two condition numbers.

\section{\texorpdfstring{Actual Brownian paths under a kinetic clock change}{Actual Brownian paths under a kinetic clock change}}
\label{proof-module-KineticLangevin-BrownianScaling-lean}
{\small\color{gray}Source module: \texttt{Kinetic\allowbreak{}Langevin/\allowbreak{}Brownian\allowbreak{}Scaling.\allowbreak{}lean}.\par}


Time change is a probability statement as well as a deterministic substitution in the drift. If \(c>0\), the process \(\widehat B_t=\sqrt c B_{t/c}\) has the same Brownian finite-dimensional laws, because its covariance is \(c\min(s/c,t/c)=\min(s,t)\). This module proves that assertion on the already constructed canonical source, retains independence of entire coordinate paths, and calculates the physical noise amplitude.

\subsection{\texorpdfstring{Change the actual coordinate process and prove its Brownian law}{Change the actual coordinate process and prove its Brownian law}}


\texttt{clock\allowbreak{}Brownian} applies the clock directly to each coordinate of the canonical source. Mathlib’s proved Brownian scaling theorem uses \(a^{-1/2}B_{at}\); setting \(a=c^{-1}\) gives the desired formula. Lean simplifies the inverse square root and the double inverse to identify this process exactly. \texttt{Is\allowbreak{}Pre\allowbreak{}Brownian\allowbreak{}Real} supplies all finite-dimensional Gaussian distributions and independent increments, rather than a single matching variance.

\subsection{\texorpdfstring{Preserve independence and recover the kinetic temperature}{Preserve independence and recover the kinetic temperature}}


Each entire changed-clock coordinate path is a measurable function of its own original coordinate path. Independence is therefore preserved by measurable composition. Finally, velocity scaling by \(c^{-1}\) changes the standard thermal input to \(c^{-1}\sqrt{2c\gamma}\,B_{t/c}=\sqrt{2\gamma/c^2}\,\widehat B_t\). The scalar proof uses the square-root product and quotient identities with their nonnegativity hypotheses. For \(c=\sqrt N\), this is exactly the physical particle temperature \(1/N\).

\section{\texorpdfstring{Normalize inverse mass in the actual kinetic equation}{Normalize inverse mass in the actual kinetic equation}}
\label{proof-module-KineticLangevin-MassScaling-lean}
{\small\color{gray}Source module: \texttt{Kinetic\allowbreak{}Langevin/\allowbreak{}Mass\allowbreak{}Scaling.\allowbreak{}lean}.\par}


The manuscript’s Lemma 1 removes inverse mass \(u>0\) by setting \(c=\sqrt u\), changing time to \(t/c\), and dividing velocity by \(c\). This proof works with the complete integral equation in an arbitrary complete normed space, so it applies to the constructed L\ensuremath{{}^{2}} Brownian solution without differentiating its noise.

\subsection{\texorpdfstring{Normalize the true phase drift}{Normalize the true phase drift}}


The continuous linear map \(D_a(v,x)=(av,x)\) changes velocity alone. Expanding both components proves \(c^{-1}D_{c^{-1}}R_{g,\gamma,c^2}(z)=R_{g,\gamma/c,1}(D_{c^{-1}}z)\). In the first component, the time and velocity factors cancel the original \(c^2\) force factor. In the second, the transformed position derivative is the normalized velocity. Lean verifies these as vector equalities using scalar action identities.

\subsection{\texorpdfstring{Transform the integral and its actual Brownian amplitude}{Transform the integral and its actual Brownian amplitude}}


Apply the bounded linear map to the solution identity at \(t/c\), pass it through the Bochner integral, and substitute \(r=s/c\). The transformed equation has friction \(\gamma/c\) and inverse mass one. The additive input becomes \(c^{-1}B_{t/c}\). For the original thermal amplitude \(\sqrt{2\gamma c^2}\), the coefficient identity gives \(c^{-1}\sqrt{2\gamma c^2}=\sqrt{2\gamma/c}\sqrt c\). The separately proved changed-clock Brownian process is \(\sqrt c W_{t/c}\), so the resulting equation has exactly the required fluctuation-dissipation coefficient. Positivity of \(c\) supplies both the interval inclusion and the square-root cancellations.

\section{\texorpdfstring{Optimizing friction and the shifted contraction exponent}{Optimizing friction and the shifted contraction exponent}}
\label{proof-module-KineticLangevin-Friction-lean}
{\small\color{gray}Source module: \texttt{Kinetic\allowbreak{}Langevin/\allowbreak{}Friction.\allowbreak{}lean}.\par}


The general kinetic theorem permits every shift \(0\le a<\gamma/2\). Its exponential coefficient is \(k(a)=\max\{a^2-m,M-(\gamma-a)^2\}/(\gamma-2a)\); a negative coefficient is a contraction rate. This module proves the three scalar optimization regimes in Table 1 of the supplied manuscript. It establishes both the value at the proposed shift and a lower bound valid for every admissible competing shift. Singular square-root and denominator boundaries are explicit. These are algebraic corollaries of the kinetic estimate, separate from construction of the Brownian process and its invariant law.

\subsection{\texorpdfstring{Center the shift interval}{Center the shift interval}}


Write \(d=\gamma-2a>0\). The two branches of the coefficient become

\[k(a)=-\frac\gamma2+\max\left\{\frac{d^2+\gamma^2-4m}{4d},\frac{4M-\gamma^2-d^2}{4d}\right\}.\]

The definition records the original coefficient exactly. For the centered identity, Lean proves each rational identity by expressing its numerator as \((-\gamma/2)d\) plus a quarter of the new numerator, distributing division, and cancelling the nonzero \(d\). Positivity allows the maximum to commute with division. Translating both branches by the same scalar then extracts \(-\gamma/2\). The denominator sign matters; a mere nonzero premise would not justify the maximum identity.

\subsection{\texorpdfstring{Prove optimality of zero shift at low friction}{Prove optimality of zero shift at low friction}}


If \(\gamma^2\le m+M\) and \(m\le M\), the optimum is \(a=0\), with coefficient \((M-\gamma^2)/\gamma\). For any admissible \(a\ge0\), compare only the second branch to this value. After multiplying by the two positive denominators, their difference is \(a(2M-\gamma a)\), which is nonnegative under the regime assumptions. The full maximum is at least its second branch. Lean formalizes those two comparisons with the positive-denominator division lemmas and \texttt{nlinarith}. At zero, \(M-\gamma^2\ge-m\), so the maximum is exactly the second endpoint. In particular, \(\gamma^2\le M\) gives a nonnegative best exponent, explaining the table's absence of a contraction rate in that range.

\subsection{\texorpdfstring{Balance the endpoints in the middle regime}{Balance the endpoints in the middle regime}}


Set \(D=\sqrt{2(m+M)-\gamma^2}\) and \(a=(\gamma-D)/2\). If \(\gamma^2\ge m+M\) and the radicand is strictly positive, then \(0<D\le\gamma\); hence the shift is admissible. Substituting \(D^2=2(m+M)-\gamma^2\) shows that the two curvature branches agree. The resulting coefficient is

\[k(a)=-\frac\gamma2+\frac{M-m}{2D}.\]

Lean proves the root sign and square first, establishes the denominator identity \(\gamma-2a=D\), rewrites the equal branches, and clears the positive denominator. It never evaluates the formula with \(D=0\).

\subsection{\texorpdfstring{Check every competing shift at the middle corner}{Check every competing shift at the middle corner}}


Assume additionally \(\gamma^2\le3m+M\), and let \(d=\gamma-2a>0\) describe any competing shift. If \(d\le D\), compare the second branch with the balanced value. Clearing positive denominators factors its error as a positive multiple of \((D-d)(D^2+Dd+2(M-m))\). If \(d\ge D\), compare the first branch instead; the corresponding factor is \((d-D)(D(d-D)+2D^2-2(M-m))\). The last coefficient is nonnegative because \(\gamma^2\le3m+M\). Lean splits precisely on \(d\le D\), proves these products nonnegative, and uses the square identity for \(D\), also multiplied by \(D\), to close the polynomial comparisons. This proves global optimality, including the nondifferentiable corner of the maximum, without assuming a derivative there.

\subsection{\texorpdfstring{Select the smaller characteristic root at high friction}{Select the smaller characteristic root at high friction}}


For \(\gamma^2\ge3m+M\), let \(D=\sqrt{\gamma^2-4m}>0\) and \(a=(\gamma-D)/2\). Nonnegative \(m\) gives \(D\le\gamma\), so the shift is nonnegative and its coordinate denominator is \(D>0\). The identity \(a(\gamma-a)=m\) makes the first branch equal \(-a\), and the high-friction condition places the second branch below it. Therefore the rate is \(a=(\gamma-\sqrt{\gamma^2-4m})/2\).



For any competing shift, the first branch alone bounds the maximum from below. The inequality that it is at least \(-a\) is exactly the nonnegative square \((\gamma-2a'-D)^2\ge0\) after division by the positive competing denominator. Lean supplies that square to \texttt{nlinarith}, then uses the maximum comparison. The strict discriminant excludes the equal-curvature critical point where the selected transformation becomes singular; a separate limiting argument would be required at that point.

\section{\texorpdfstring{The literal exponential-\allowbreak{}integrator covariance is positive definite}{The literal exponential-integrator covariance is positive definite}}
\label{proof-module-KineticLangevin-GramPositive-lean}
{\small\color{gray}Source module: \texttt{Kinetic\allowbreak{}Langevin/\allowbreak{}Gram\allowbreak{}Positive.\allowbreak{}lean}.\par}


The two-kernel covariance is strictly positive definite on every positive-duration step. This supplies the genuine Cholesky input required by the first particle algorithm, rather than merely positive semidefiniteness of an abstract Gaussian covariance.

\subsection{\texorpdfstring{Use strict positivity of the integrated exponential}{Use strict positivity of the integrated exponential}}


The first integrated response is \(\psi_1(h)=\int_0^he^{-\gamma t}\,dt\). For \(h>0\), the continuous integrand is positive everywhere, hence this integral is strictly positive for every real \(\gamma\). Lean invokes strict positivity of a continuous interval integral, using its positive value at zero. This avoids dividing by friction and includes \(\gamma=0\).

\subsection{\texorpdfstring{Show every nonzero coefficient vector produces a nonzero response on the interval}{Show every nonzero coefficient vector produces a nonzero response on the interval}}


For \(z=(z_0,z_1)\ne0\), consider \(z_0\psi_0(t)+z_1\psi_1(t)\). If \(z_0\ne0\), its value at zero is \(z_0\), because \(\psi_0(0)=1\), \(\psi_1(0)=0\). If \(z_0=0\), then \(z_1\ne0\), and its value at \(h\) is nonzero because \(\psi_1(h)>0\). Lean handles the two finite coordinates explicitly and proves a strictly positive squared response at an actual point of \([0,h]\).

\subsection{\texorpdfstring{Integrate the square to prove strict positive definiteness}{Integrate the square to prove strict positive definiteness}}


The already proved Gram identity is \[z^\top C_hz=\int_0^h\bigl(z_0\psi_0(t)+z_1\psi_1(t)\bigr)^2\,dt.\] The integrand is continuous and nonnegative, and the previous result supplies a point where it is positive. Its integral is therefore strictly positive for every nonzero \(z\). Lean rewrites the matrix quadratic form as the exact scalar integral, proves its continuity from the finite kernel sum, and combines strict integral positivity with the proved covariance symmetry. Thus \(C_h\) is genuinely positive definite, which justifies an invertible triangular Cholesky factor.

\section{\texorpdfstring{The actual Cholesky implementation of the Gaussian innovations}{The actual Cholesky implementation of the Gaussian innovations}}
\label{proof-module-KineticLangevin-GaussianFactor-lean}
{\small\color{gray}Source module: \texttt{Kinetic\allowbreak{}Langevin/\allowbreak{}Gaussian\allowbreak{}Factor.\allowbreak{}lean}.\par}


A Gaussian law specified through its covariance must agree with the paper’s concrete instruction to multiply independent standard normals by a Cholesky factor. This module proves that exact identity, constructs the actual positive-diagonal lower-triangular 2\ensuremath{\times}2 factor, and applies it to the strictly positive exponential Gram covariance.

\subsection{\texorpdfstring{Identify every genuine covariance factor by mean and covariance}{Identify every genuine covariance factor by mean and covariance}}


If \(AA^\mathsf T=S\), then applying \(A\) to a standard Gaussian gives a centered Gaussian with covariance \(S\). Lean uses the matrix-to-continuous-linear-map star-algebra equivalence: transpose becomes adjoint and matrix multiplication becomes operator composition. The mean passes through the continuous linear map and remains zero. The covariance pushforward is \(\langle A^\mathsf Tx,A^\mathsf Ty\rangle=\langle x,AA^\mathsf Ty\rangle\), which is the actual target covariance. Uniqueness of Gaussian measures with equal mean and covariance proves equality of measures, not only equality of moments.

\subsection{\texorpdfstring{Construct and verify the positive-\allowbreak{}diagonal Cholesky factor}{Construct and verify the positive-diagonal Cholesky factor}}


For a positive definite symmetric 2\ensuremath{\times}2 matrix with entries \(a=S_{00}\), \(b=S_{10}\), \(d=S_{11}\), define \(L=\begin{pmatrix}\sqrt a&0\\b/\sqrt a&\sqrt{d-b^2/a}\end{pmatrix}\). Positive definiteness gives \(a>0\) and \(ad-b^2>0\), hence a strictly positive Schur complement. The factor is lower triangular with positive diagonal. Expanding its four matrix entries, using the two square-root-square identities and symmetry, gives \(LL^\mathsf T=S\). The general Gaussian-factor theorem therefore proves that this concrete factor generates exactly the target Gaussian measure.

\subsection{\texorpdfstring{Generate the actual exponential innovation from independent normals}{Generate the actual exponential innovation from independent normals}}


The two response kernels \(\psi_0,\psi_1\) have a strictly positive Gram covariance on every positive-length interval. Apply the explicit Cholesky construction to that actual matrix. The resulting output law is precisely \texttt{klmc\allowbreak{}Gaussian}, already identified with the Brownian convolution pair used by the sampler. Finally, the product of two real standard Gaussian measures pushes forward to the Euclidean standard Gaussian under the coordinate identification. Composing the two exact laws proves the paper’s implementation starting from two independent scalar normal draws. This closes the computational covariance-factor step used by the shared exponential integrator and its particle specialization.

\section{\texorpdfstring{An explicit quadratic bias envelope and second-\allowbreak{}order accuracy step}{An explicit quadratic bias envelope and second-order accuracy step}}
\label{proof-module-KineticLangevin-SecondOrderBudget-lean}
{\small\color{gray}Source module: \texttt{Kinetic\allowbreak{}Langevin/\allowbreak{}Second\allowbreak{}Order\allowbreak{}Budget.\allowbreak{}lean}.\par}


The stronger source Gaussian tail allows a completely explicit second-order accuracy prescription. The exponentially small term is bounded by a quadratic envelope, making the entire persistent floor at most \(B_p h^2\) for \(h\le1\). A positive minimum step then satisfies stability, the Gaussian cutoff dimension restriction, and the epsilon budget simultaneously.

\subsection{\texorpdfstring{Bound the genuine cutoff tail by a reciprocal}{Bound the genuine cutoff tail by a reciprocal}}


For \(x>0\), \(e^x\ge1+x\ge x\), hence \(e^{-x}\le1/x\). Apply this at \(x=m^2/(160M_2^2h^2)\) and use the previously proved sharper Gaussian tail. The result is \(e^{-(\alpha-p)/8}\le160M_2^2h^2/m^2\). This is a global inequality under the source restrictions, with no asymptotic qualification.

\subsection{\texorpdfstring{Collect the complete floor into a quadratic coefficient}{Collect the complete floor into a quadratic coefficient}}


Write \(C_p=M_2\sqrt{p^2+2p}+M\sqrt{Mp}\). The exact floor is \(\sqrt2h^2C_p/m+(8(M-m)/m)h e^{-(\alpha-p)/8}\). Since \(M-m\le M\), the tail envelope bounds the second term by \(1280MM_2^2h^3/m^3\). For \(0<h\le1\), \(h^3\le h^2\). Thus the full floor is at most \(B_ph^2\), where \(B_p=\sqrt2 C_p/m+1280MM_2^2/m^3\). Every coefficient is explicit, and the exact fourth moment remains valid in dimension one.

\subsection{\texorpdfstring{Choose a positive step satisfying every restriction}{Choose a positive step satisfying every restriction}}


For positive dimension and accuracy choose the minimum of \(1\), \(m/(5\gamma M)\), \(m/(4\sqrt{5p}M_2)\), and \(\sqrt{0.94\varepsilon/B_p}\). All four terms are positive. The second and third enforce the two source step restrictions; the last gives \(B_ph^2\le0.94\varepsilon\). Lean proves each minimum bound, compares nonnegative squares for the final term, and multiplies through positive denominators. For fixed curvature constants, \(B_p\) grows linearly with \(p\), so the displayed step and the logarithmic count exhibit the source \(\sqrt{p/\varepsilon}\) iteration dependence.

\section{\texorpdfstring{A finite epsilon guarantee for the actual second-\allowbreak{}order sampler}{A finite epsilon guarantee for the actual second-order sampler}}
\label{proof-module-KineticLangevin-SecondOrderAccuracy-lean}
{\small\color{gray}Source module: \texttt{Kinetic\allowbreak{}Langevin/\allowbreak{}Second\allowbreak{}Order\allowbreak{}Accuracy.\allowbreak{}lean}.\par}


The actual KLMC2 law theorem and the proved quadratic persistent-error envelope combine into explicit accuracy corollaries. The second result chooses a positive step satisfying every source restriction and a rounded natural iteration count, with no remaining local-error, stationarity, law or recurrence premise.

\subsection{\texorpdfstring{Apply the actual sampler theorem to a quadratic error budget}{Apply the actual sampler theorem to a quadratic error budget}}


The exact Gaussian-chain theorem gives \(W_2(\nu_n,\pi)\le\sqrt2q^nD+b(h)\), where \(q=1-mh/(4\gamma)\). The proved source restrictions make \(0\le q<1\). For \(h\le1\), the complete persistent floor satisfies \(b(h)\le B_ph^2\). Allocate \(B_ph^2\le0.94\varepsilon\) and choose the natural ceiling of \(\max\{0,\log(24D/\varepsilon)/(mh/(4\gamma))\}\). The shared scalar accuracy lemma makes the total at most \(\varepsilon\). A zero initial distance needs no transient iterations; finite transport is retained explicitly.

\subsection{\texorpdfstring{Choose the positive step and integer count directly}{Choose the positive step and integer count directly}}


Use the minimum of \(1\), \(m/(5\gamma M)\), \(m/(4\sqrt{5p}M_2)\), and \(\sqrt{0.94\varepsilon/B_p}\). The checked step-specification theorem supplies positivity, both source restrictions and the quadratic bias allocation. Construct the actual nonnegative-real step and the rounded iteration count, then apply the preceding genuine sampler accuracy theorem. For fixed curvature constants, the explicit coefficient is of order \(p\), so the formula displays the \(\sqrt{p/\varepsilon}\) step count up to its visible logarithm. Positive dimension is required for this particular step prescription; the general sampling bound retains all finite dimensions.


\end{document}
